\documentclass[11pt]{article}
\usepackage[a4paper,margin=30mm]{geometry}
\usepackage{amsmath,amssymb,amsthm,mathtools,mathrsfs}
\usepackage{enumitem}
\usepackage[hidelinks]{hyperref}
\setlength{\emergencystretch}{2em}

\newtheorem{theorem}{Theorem}[section]
\newtheorem{proposition}[theorem]{Proposition}
\newtheorem{lemma}[theorem]{Lemma}
\newtheorem{corollary}[theorem]{Corollary}
\newtheorem{definition}[theorem]{Definition}
\newtheorem{remark}[theorem]{Remark}
\newcommand{\cL}{\mathcal L}
\newcommand{\cU}{\mathcal U}
\newcommand{\cD}{\mathcal D}
\newcommand{\cK}{\mathcal K}
\newcommand{\cB}{\mathcal B}
\newcommand{\cQ}{\mathcal Q}
\newcommand{\one}{\mathbf 1}
\newcommand{\TV}{\mathrm{TV}}
\newcommand{\Lip}{\mathrm{Lip}}

\title{The Two-Time Moving-Defect Problem for Dispersing Billiards\\
\large Pair-safe selection lifts, occurrencewise propagated currents,
product-time face control,
global incidence Cauchy,
two-cutoff convergence, and a strongly conditional
pathwise CM2 criterion}
\author{Working note}
\date{13 July 2026 (version 44)}

\begin{document}
\maketitle

\begin{abstract}
We isolate the two-time estimate needed to differentiate an infinite
correlation susceptibility under a genuine motion of dispersing scatterers.
The moving first defect contains regular $M\times M$ terms and singular
one-dimensional physical faces.  Neither a projective-cone trace estimate nor
a nuclear graph-current representation is available.  We avoid both.  In the
central two-time band we first remove contracted coordinates.  The physical
face is then pushed to its two scalar dominant coordinates, smoothed only on
that quotient, and lifted on two same-branch homogeneous sections.  The
branchwise contracted-coordinate estimate intertwines the original graph
pairing with the section pairing without averaging over other symbolic pasts
or singularity cuts.  Thus a two-dimensional Fourier polynomial is never
declared to be a four-dimensional density by a fictitious chart inverse.  At
finite cutoff the lifted low block is exactly a sum of products of measured
standard-curve currents.  Oscillatory section densities are decomposed into
positive regular standard families.  We retain their actual recovery times,
split labels into jointly recovered and one-sided bad tails, and use separate
conditional recovery moments rather than multiplying the two initial boundary
complexities.  Current centering and its test-side adjoint remain
distinct typed operations.  A measure-trivializing transport puts every
moving collision map on one fixed probability space, so both operations,
endpoint unfolding, and standard-family coupling use the same invariant
reference.  At finite difference-quotient level we keep the fixed-gauge
operator target: coordinate commutators are not cancelled against transported
observables, which would change the differentiated object.  Instead the
complete operator quotient is decomposed exactly into regular, physical-face,
product-current, and one-derivative response blocks.  Regular and response
blocks are double-projected individually, whereas every graph face is kept
together with the exact product-current correction generated by its two
marginals.  Three nested distribution spaces separate the two
jobs that a single weak space cannot perform: an intermediate response space
gives the uniform CM2 majorant, while a weaker difference-quotient space gives
operator convergence and hence each fixed two-time limit.  Dominated
convergence then exchanges the limit and the CM2 double sum.  The regular and
physical-face blocks use the same two-level logic: a strong ledger gives the
uniform CM2 majorant, while respectively a weaker Sobolev topology and a
weak oriented-current topology give fixed-term convergence.  Uniform boundary
mass control upgrades the latter convergence to the branchwise tests generated
by a fixed finite singularity refinement.  A separate two-cutoff theorem then
passes from this fixed core to the growing depth $L(s)$: an $s$-uniform
positive coarea/return bound removes the intermediate triangular band, with no
$|s|^{-1}$ loss, before the existing deep remainder is used.
The complementary dominant
modes are first unfolded to the genuine past/future physical endpoints.  Thus
the Fourier coefficients belong to the original smooth observables, while all
time dependence stays in the physical phase; no regularity of an evolved
observable is assumed.  Off the
central band a signed proper-family theorem supplies forward and reverse
coupling.  This removes graph Banach spaces, tensor renewal, and nuclear
decompositions of the regular defect.

The high-frequency face term is controlled by a physical slope energy at the
same prescribed resolving depth used by the phase.  We
use Leclerc's native unstable coordinate and his cutting theorem modulo every
fixed-degree polynomial as a nonvacuity strategy.  The theorem itself assumes
the exact actual-SRB physical polynomial escape estimate that the phase
argument needs.  It does not infer that estimate from a fixed
negative-relative-pressure clean subsystem: conditional normalization
controls bad fractions among survivors but cannot turn exponentially small
survival mass into a logarithmic-time physical cover.  An adaptive
full-mass finite alphabet with fixed rough witnesses is recorded only as a
separately falsifiable route to this primitive input.  No affine jet is absorbed into the polynomial
defining the event, and no infinite fractal conditional is compared with an
absolutely continuous one.  A single-source
conditional identity freezes one oriented endpoint and places the other
endpoint in the context polynomial.  An actual-SRB first-stopping antichain
selects native intervals whose physical diameters are comparable with the
energy scale; its complete positive tree law carries all context and mark moments.
A mass-preserving resolving buffer and Jensen then give the reduced-slope
small-ball estimate.  Before any context, stopping, buffer, gate, or
continuation label is discarded, all branches are placed in one complete
positive energy tree, with a first-failure label for each loss.  Each loss is
then represented by an exact stopped positive source.  Its source density and
remaining propagation norm are dominated by a predictable cemetery envelope;
future factors are set to one only when that remaining propagation is a
positive total-variation contraction.  The survivor weight and the five
cemetery envelopes form one piecewise full event weight, so a tilted failure
probability controls the actual discarded current rather than only an
auxiliary label law.  A full-product $L^{p_W}$ lemma on the complete tree,
including a genuine joint cemetery-envelope moment, gives a nonvacuity route from unweighted tails,
with every global aggregate exponent explicitly degraded by $1/p_W'$; mere
$L^1$ integrability is never used to preserve a tail exponent.  Uniform
parent-context tails instead use a normalized kernel reverse-H\"older
interface on every frozen actual, split, cemetery, and source-incidence
parent.  A global complete-tree moment is explicitly not promoted to this
conditional statement.  The cemetery factor
is never conditioned on a sigma-algebra on which it is already measurable.
Probability clocks and additive currents are not identified.  First-hit
disjointization controls only the bad-event union and survivor mass.  Every
actual additive source occurrence instead lives on one fixed standard-Borel
marked source space.  Two finite positive kernels over each complete-tree
label have distinct duties: one records the exact positive-component
coefficient mass and the
other the source-norm, recovery, opposite-amplitude, and operator envelope.
Every signed source is first measurably decomposed into positive regular
components with occurrence-constant polarity; no arbitrary carrier-dependent
sign is multiplied into a previously costed positive section.  Their
Radon--Nikodym density assembles the exact current against the envelope
kernel and a normalized weakly measurable carrier section.  Global envelope
integrability of this one propagated kernel makes the scalar weak pairing a bounded
functional on one fixed test space.  Thus continuous tree labels are
integrated rather than nakedly summed, while two distinct currents with
identical support are still counted twice by their occurrence marks.  An
exhaustive measurable owner partition and a separate resolution-dependent
main clock-coverage identity prevent unowned or unclocked marked mass.
Tonelli controls the positive envelope before signed recombination.  Source-incidence
complexity, rather than support-event complexity, is subtracted from every
raw source gap.  The central clock system is a clock--passage matrix:
polynomial-in-energy, prescribed PPE depth, and genuine two-time depth each
have separate Jensen and TV columns, while word depth is handled jointly by
the independent face-time, face-depth, and deep ledgers.  The theorem assumes
aggregate summability of both the
probability union and the additive source mass, or derives it from strict
complexity gaps before applying H\"older.  A Jensen cell now carries two
separate quantities: its total base occurrence mass and its fibre energy.
One global Cauchy inequality on the entire finite incidence measure converts
these into a value-source tail.  No square root is taken occurrence by
occurrence.  The main resonant source is removed from the primitive source
clock.  Its energy lift starts from the unselected collision--SRB pair law,
adds a pointwise normalized proposal mark kernel before the main selector,
and places the selector, propagated occurrence weight, reciprocal proposal,
and carrier ratios in one unnormalized source likelihood.  Different
carriers require an actual pushforward reference law, a Radon--Nikodym bound,
and target-scale slope compatibility.  Thus probability PPE is never
conditioned on a main cell for free.  All likelihood and base-mass losses
enter explicit source-energy and value exponents before this one global
Cauchy step.  Only genuine
two-time TV tails enter the central exponent directly; all Jensen rates are
the already-derived base--energy gaps.  The
choice $\sigma=r^{1/(1+\alpha_L)}$ separates Leclerc's cutting exponent from
the billiard $C^{1+\alpha_B}$ remainder.  On the nonresonant set a direct
$C^{1+\alpha}$ fractional nonstationary-phase estimate yields a polynomial
Fourier saving.  The dominant lattice is divided into three ranges: a finite
low block below $K=e^{\eta_{\rm lo}\max\{m,n\}}$, an intermediate phase block
ending at $H=e^{\eta_{\rm hi}\max\{m,n\}}$, and an ultra-high block estimated
only by the Fourier tail of the original observables.  Thus no prescribed
escape estimate is called beyond the finite spacetime resolving depth.  The
raw lattice cutoff is supplemented by an upper-scale grazing ledger: endpoint
derivative growth is charged jointly to the encountered homogeneity indices
and to an initial face/section/chart/coarea scale mark.  The excess tail is
stated as a survivor subprobability estimate under the complete event-weight
tilt, rather than inferred from an unweighted probability tail and mere
integrability of the mark.  No depth-$N$ endpoint observable is read on an
earlier cemetery branch.  Labels with excessive weighted growth are removed
by that tilted survivor tail.  All
buffer, gate, stopping, and continuation losses enter the final phase exponent
explicitly rather than being hidden in an unnamed exponential remainder.

For radial faces we index every connected component and normalize its mass
before forming the boundary functional.  A full tangent one-dimensional area
formula replaces the vertical two-flux formula, which is retained only as a
special case.  It is combined with an oriented component ledger; no component is silently
discarded.  The resulting standard families recover after $O(\log Z)$
iterates.  Qualitative residual transversality is insufficient for the
infinite-depth limit, so the theorem is stated for analytic or circular table
deformations satisfying an explicit pathwise jet-tail condition with a
strict return/grazing margin.  A word depth $L(s)\asymp\log(1/|s|)$ and two
independent time thresholds are used.  Shallow short times use endpoint
Fourier analysis only while all intermediate-frequency resolving depths lie
below $L(s)$; shallow long times are removed directly as finite-$s$
positive-envelope collars.  Their positive source mass is normalized before
the two one-sided recovery estimates and multiplied back afterwards, so an
unrecorded recovery prefactor cannot enter the off-diagonal cones.
Deep short times are controlled by positive source mass, and long-time
central collars by a boundary-crossing estimate for
one positive envelope that simultaneously dominates the actual multi-time law
and every functional-correlation split law, followed by one common
$N$-dependent cylinder approximation;
off-diagonal collars use moving standard-family coupling after the deep
remainder is normalized and shown to have a uniform recovery exponential
moment.  The small total mass is multiplied back only after this normalized
estimate.  This gives an
$o(1)$ tail in the absolute CM2 double sum, not merely pointwise convergence.
The analytic realization requires a fiberwise Remez/valency condition in
normalized Fourier coordinates; finite-branch jet surjectivity alone is not
used as a sublevel theorem.  Circular families retain the quantitative jet
condition.  Under the primitive hyperbolicity, single-source, actual
endpoint--slope, native-stopping, physical-polynomial-escape, exact
unselected product-pair energy, complete-weighted-energy, dual-mass
constant-polarity source-kernel, post-propagation carrier domination,
resolution-exhaustive main-clock
coverage, selection-safe joint occurrence--carrier likelihood,
parent-kernel reverse-H\"older, branch-section,
geometric-flux, endpoint-upper-scale, split-standard-family,
connected-pullback, shallow spacetime refinement, two-cutoff face convergence,
product-time face-envelope, bounded physical-test and physical/tree
matching, normalized shallow/deep-recovery, and pathwise inputs
listed below, including the stopped-source domination and clock--passage
interfaces, the complete first
defect satisfies CM2.  No unconditional generic circular-table claim is made.

\end{abstract}

\section{The exact target}

Throughout, $\{\widehat T_s:|s|<s_0\}$ is a $C^{r+1}$ compact family of
finite-horizon planar dispersing billiard collision maps generated by a
$C^{r+2}$ normal deformation of strictly convex scatterer boundaries, with
collision spaces $M_s$ and invariant collision probabilities
$\widehat\mu_s$.  Write
$M_s=\bigsqcup_{i=1}^{d_0}M_{s,i}$, where
$M_{s,i}=\partial O_{s,i}\times(-\pi/2,\pi/2)$.  A near-identity
diffeomorphism cannot transfer mass between these connected components, so
the following compatibility is a genuine restriction, not a coordinate
construction:
\begin{equation}
\label{eq:constant-component-mass-vector}
 \widehat\mu_s(M_{s,i})
 =\frac{|\partial O_{s,i}|}{\sum_j|\partial O_{s,j}|}
 =\widehat\mu_0(M_{0,i}),\qquad 1\le i\le d_0.
\end{equation}
We work only in deformation charts satisfying
\eqref{eq:constant-component-mass-vector}.  On such a chart we impose the
following measure-trivializing transport interface.  There are $C^{r+1}$ diffeomorphisms
$A_s:M=M_0\to M_s$, $A_0=\mathrm{id}$, such that
\begin{equation}
\label{eq:measure-trivializing-gauge}
 A_s^*\widehat\mu_s=\mu_0,\qquad
 T_s:=A_s^{-1}\widehat T_sA_s .
\end{equation}
Consequently every $T_s$ preserves the \emph{same} probability $\mu_0$.
For the collision measure this is a componentwise smooth mass-coordinate
transport after the component masses have been matched by
\eqref{eq:constant-component-mass-vector}; a transport inside one component
cannot repair a changing component mass.  Existence with the regularity and
chart compatibility used below is part of the interface.  Thus independent
radius changes of several circular scatterers and unconstrained independent
boundary modes are not covered.  Admissible examples include center motions
with fixed shapes and constrained shape/radius paths whose normalized
perimeter vector is constant.  The chosen analytic-boundary or
finite-dimensional circular deformation is subject to
$\mathrm{QT}_{\rm path}$, specified in
Section~\ref{sec:parameter-state}.  The table at $s=0$ contains the
unobstructed clean cycle used in
Proposition~\ref{prop:exact-circular-anosov}.  We consider only the complete
\emph{first} moving-table defect: regular Jacobian/density derivatives and all
same-label physical one-sided faces are assembled before centering.  The tests
$f,g$ belong to $C^r(M)$, with $r$ above the Fourier summability threshold,
and $\mu_0(g)=0$.  Second shape defects and unassembled point atoms are outside
the statement.

On this fixed probability space we use the Perron convention relative to
$\mu_0$:
\begin{equation}
\label{eq:perron-composition-convention}
 P_sh=h\circ T_s^{-1},\qquad P_s^*f=f\circ T_s,\qquad P_s\one=\one.
\end{equation}
Thus $P_s$ is the transfer operator on $L^1(\mu_0)$ and $P_s^*$ is its
$L^\infty(\mu_0)$ adjoint for every $s$.  Let
\[
  \Pi h=\mu_0(h)\one,\qquad Q=P_0-\Pi,
\]
and, for every finite $s$, use the \emph{same} projection
\[
  \Pi_0h=\mu_0(h)\one,\qquad Q_s=P_s-\Pi_0 .
\]
Since $P_s\Pi_0=\Pi_0P_s=\Pi_0$, for every $m\ge0$ and $n\ge1$,
\begin{equation}
\label{eq:centered-endpoint-iterates}
 Q_s^mg=g\circ T_s^{-m}\quad\text{if }\mu_0(g)=0,
 \qquad (Q_s^*)^nf=f\circ T_s^n-\mu_0(f).
\end{equation}
All finite-parameter estimates below are required uniformly on the common
spacetime refinement of the compact table chart.  Let $K_V$ denote the complete first moving-table defect after all regular
terms, same-label physical faces, and grazing blocks have been assembled.  The
two-time estimate needed for a prescribed centered source $g$ is
\begin{equation}
\label{eq:CM2}
 \bigl|\langle Q^nK_VQ^mg,f\rangle\bigr|
 \le C\tau^m\vartheta^n\|g\|_{C^r}\|f\|_{C^r},
 \qquad m,n\ge0,
\end{equation}
for some $\tau,\vartheta<1$.  For the difference quotient it is enough to
have the same estimate on words of depth at most $L(s)$ together with an
$o(1)$ CM2 tail beyond $L(s)$; this is the pathwise diagonal criterion below.

The operator must be centered as a whole:
\begin{equation}
\label{eq:centered-defect}
  K_V=\Pi^\perp\left.\partial_sP_s\right|_{s=0}\Pi^\perp.
\end{equation}
Centering separate faces is neither invariant nor sufficient.  Without
\eqref{eq:centered-defect}, rank-one terms generally give an additive bound
$O(\tau^m+\vartheta^n)$, whose double sum diverges.

Let $\mathfrak R_R:=\mathcal M(M)$ be the strong Banach space of finite
signed Borel measures on one collision coordinate, with total-variation
norm.  Let $\mathfrak R_{\rm BL}$ be its completion for the genuinely weaker
bounded--Lipschitz dual norm
$\|r\|_{\rm BL^*}:=\sup_{\|h\|_{\rm BL}\le1}|r(h)|$.
The inclusion $\mathfrak R_R\hookrightarrow\mathfrak R_{\rm BL}$ is
continuous, $\mu_0\in\mathfrak R_R$, and
\begin{equation}
\label{eq:one-coordinate-response-pairing}
 |r(h)|\le\|r\|_{\mathfrak R_R}\|h\|_\infty,
 \qquad r\in\mathfrak R_R,\quad h\in C^0(M).
\end{equation}
If $J$ is a finite signed graph measure on $M\times M$, its two marginal
maps are the actual pushforwards and obey
\begin{equation}
\label{eq:bounded-face-marginal-maps}
 \|({\rm pr}_\pm)_*J\|_{\mathfrak R_R}
 \le\|J\|_{\rm TV},\qquad |J(1,1)|\le\|J\|_{\rm TV}.
\end{equation}
Thus no independently stronger response norm is hidden in the fourth type.
For later bookkeeping let $\mathfrak P_R$ be the space of finite product
currents
\begin{equation}
\label{eq:product-current-type}
 P=r^-\otimes\mu_0+\mu_0\otimes r^++c\,\mu_0\otimes\mu_0,\qquad
 \|P\|_{\mathfrak P_R}:=\|r^-\|_{\mathfrak R_R}
 +\|r^+\|_{\mathfrak R_R}+|c|,
\end{equation}
The representation is taken modulo the obvious
constant overlap, with the infimum norm.  These three-dimensional mixed
products are neither four-dimensional $W^{R,1}$ densities nor
one-dimensional graph faces.  They therefore form a fourth type even though
their pairing with a product of two individually centered endpoint tests is
zero.

\begin{definition}[Fixed-gauge operator finite-difference interface
$\mathrm{MT}_{\rm DQ}$]
\label{def:measure-gauge-interface}
Let $U_s$ denote the measure-trivializing pullback in
\eqref{eq:measure-trivializing-gauge}.  The differentiated object is the
fixed-gauge operator
\begin{equation}
\label{eq:fixed-gauge-operator-DQ}
 K_s^{\rm DQ}:=\Pi_0^\perp\frac{P_s-P_0}{s}\Pi_0^\perp .
\end{equation}
The tests $f,g$ are fixed in this gauge.  In particular coordinate
commutators are not cancelled against derivatives of transported
observables: such a cancellation would compute a different physical
bilinear response.  Put
$C^r_0(M)=\{h\in C^r(M):\mu_0(h)=0\}$ and let
$\cB_2\hookrightarrow\cB_1\hookrightarrow\cB_0$ be three fixed distribution
spaces, called respectively the strong, response, and difference-quotient
spaces.  The maps $\Pi_0$ and $\Pi_0^\perp$ extend continuously to all three;
a subscript $0$ denotes the corresponding zero-mass subspace.  Condition
$\mathrm{MT}_{\rm DQ}$ requires the following finite-$s$ statements before
any absolute value is taken.
\begin{enumerate}[label=\textup{(MT\arabic*)}]
\item Before projection the quotient has one exact typed decomposition.  A
      raw regular or response block is projected individually:
      \begin{equation}
      \label{eq:typed-block-double-projection}
       \widehat K_{s,j}^{\rm DQ}
       :=\Pi_0^\perp\widetilde K_{s,j}^{\rm DQ}\Pi_0^\perp,
       \qquad j\in\{\mathrm{reg},\mathrm{sw}\}.
      \end{equation}
      In contrast, let $J_{s,\rm face}^{\rm DQ}$ be the assembled raw graph
      current and let $m_{s,\rm face}^{\pm}$ be its two marginals.  Its
      product-current correction and centered cluster are, exactly,
      \begin{align}
       P_{s,\rm prod}^{\rm DQ}
       &:=-m_{s,\rm face}^{-}\otimes\mu_0
          -\mu_0\otimes m_{s,\rm face}^{+}
          +J_{s,\rm face}^{\rm DQ}(1,1)\mu_0\otimes\mu_0,
          \label{eq:face-product-correction}\\
       K_{s,\rm face\text{-}cl}^{\rm DQ}
       &:=J_{s,\rm face}^{\rm DQ}+P_{s,\rm prod}^{\rm DQ}
        =\mathsf D_{\rm cur}J_{s,\rm face}^{\rm DQ}.
          \label{eq:typed-face-product-cluster}
      \end{align}
      Thus the graph and product summands are distinct distributional types
      but one jointly centered cluster.  Neither summand is asserted to have
      zero marginals, and $P_{s,\rm prod}^{\rm DQ}$ is not projected again
      (doing so would kill it).  Marginal and rank-one terms are never
      relabelled as regular densities or graph faces.  In the notation of
      \eqref{eq:product-current-type},
      $r_s^-=-m_{s,\rm face}^-$,
      $r_s^+=-m_{s,\rm face}^+$, and
      $c_s=J_{s,\rm face}^{\rm DQ}(1,1)$.  With this convention
      the complete quotient is
      \[
        K_s^{\rm DQ}=\widehat K_{s,\rm reg}^{\rm DQ}
          +J_{s,\rm face}^{\rm DQ}+P_{s,\rm prod}^{\rm DQ}
          +\widehat K_{s,\rm sw}^{\rm DQ}+E_s,
        \qquad E_s=\Pi_0^\perp E_s\Pi_0^\perp .
      \]
      Here the regular kernel obeys
      \begin{equation}
      \label{eq:uniform-gauge-DQ-regular}
       \sup_{0<|s|<s_0}\sum_\chi
       \|\zeta_\chi k_{s,\rm reg}^{\rm DQ}\|_{W^{R,1}}<\infty ,
      \end{equation}
      while the graph-supported part belongs to the existing physical face
      atlas and obeys the uniform strong face ledger
      \begin{equation}
      \label{eq:uniform-gauge-DQ-face}
       \sup_{0<|s|<s_0}\sum_{\lambda}M_{s,\lambda}
       \|w_{s,\lambda}\|_{\mathcal A^\alpha}
       \mathfrak d_{s,\lambda}^{A_{\rm dom}}
       \operatorname{Mark}(s,\lambda)^2
       \{1+Z(\mathcal G_{s,\lambda}^-)
              +Z(\mathcal G_{s,\lambda}^+)\}<\infty .
      \end{equation}
      This is the face summand of
      \eqref{eq:event-physical-norm}, not a convergence assertion.  Finally,
      \begin{equation}
      \label{eq:uniform-strong-weak-DQ}
       \sup_{0<|s|<s_0}
       \left(\|\widehat K_{s,\rm sw}^{\rm DQ}\|_{\cB_2\to\cB_1}
       +\|E_s\|_{\cB_2\to\cB_1}\right)<\infty.
      \end{equation}
      The product block has a typed marginal representation and satisfies
      \begin{equation}
      \label{eq:uniform-product-current-DQ}
       \sup_{0<|s|<s_0}
       \|P_{s,\rm prod}^{\rm DQ}\|_{\mathfrak P_R}<\infty.
      \end{equation}
      Its two one-coordinate factors and its scalar coefficient are the
      actual marginals and total mass in
      \eqref{eq:face-product-correction}; no unrelated finite measure may be
      substituted.  Their bounds are precisely the contractive marginal
      estimates \eqref{eq:bounded-face-marginal-maps} applied before the
      positive face ledger is summed.
      A graph-supported boundary contribution of the commutator is included
      in $J_{s,\rm face}^{\rm DQ}$; it is never called an $M\times M$
      density.
\item Uniformly for $|s|<s_0$, the centered evolution has the two genuine
      loss-of-memory bounds
      \begin{align}
       \|Q_s^mh\|_{\cB_2}
       &\le C\theta_-^m\|h\|_{C^r},
       &&h\in C^r(M),\ \mu_0(h)=0,
       \label{eq:strong-side-memory}\\
       \|Q_s^nu\|_{(C^r_0)^*}
       &\le C\theta_+^n\|u\|_{\cB_1},
       &&u\in\cB_{1,0},
       \label{eq:response-side-memory}
      \end{align}
      with $0<\theta_\pm<1$.  These are Banach-space estimates, not an
      assumed CM2 estimate for $\widehat K_{s,\rm sw}^{\rm DQ}$.
\item The same typed organization, including the exact face--product cluster
      \eqref{eq:typed-face-product-cluster}, holds for the limiting
      fixed-gauge operator $K_V$, with
      $\|\widehat K_{V,\rm sw}\|_{\cB_2\to\cB_1}\le C$, and
      \begin{equation}
      \label{eq:strong-DQ-convergence}
       \|\widehat K_{s,\rm sw}^{\rm DQ}-\widehat K_{V,\rm sw}\|_{\cB_2\to\cB_0}
       +\|E_s\|_{\cB_2\to\cB_0}\longrightarrow0.
      \end{equation}
      The product-current factors converge in their one-coordinate weak
      topology:
      \begin{equation}
      \label{eq:product-current-DQ-convergence}
       r_{s}^{\pm}\longrightarrow r_V^{\pm}
       \quad\hbox{in }\mathfrak R_{\rm BL},
       \qquad c_s\longrightarrow c_V.
      \end{equation}
      The other two types have their own explicitly weaker fixed-term
      topologies.  The regular kernels satisfy
      \begin{equation}
      \label{eq:regular-DQ-weak-convergence}
       \sum_\chi\|\zeta_\chi
       (k_{s,\rm reg}^{\rm DQ}-k_{V,\rm reg})\|_{W^{R-1,1}}
       \longrightarrow0.
      \end{equation}
      For the face type, first fix $(m,n)$ and a finite depth $L$, and let
      $\mathcal J_{s,L}^{\rm face}$ and $\mathcal J_{V,L}^{\rm face}$ be the
      assembled, oriented graph currents on the common
      $\mathrm{UCR}_{\rm sh}$ atlas.  More precisely, let
      $\mathfrak A_{L,m,n}^{\rm face}$ be the finite disjoint union of the
      component-indexed parameter arcs in that atlas, equipped with its fixed
      chart metric and Borel record map into $M\times M$; currents and tests
      below are pulled back to this common metric space before their
      bounded--Lipschitz norms are taken.  Throughout,
      $\mathcal J_{s,L}^{\rm face}:=\mathcal J_{s,\le L}^{\rm face}$ and
      $\mathcal J_{V,L}^{\rm face}:=\mathcal J_{V,\le L}^{\rm face}$; the
      subscript is a cutoff depth, never an exact-depth label.  With
      $\|\cdot\|_{\rm BL^*}$ denoting
      the bounded--Lipschitz dual norm on that finite atlas, require
      \begin{align}
       \|\mathcal J_{s,L}^{\rm face}
          -\mathcal J_{V,L}^{\rm face}\|_{\rm BL^*}&\longrightarrow0,
       \label{eq:face-DQ-weak-convergence}\\
       \sup_{0<|s|<s_0}|\mathcal J_{s,L}^{\rm face}|
          ([\mathcal S_{L,m,n}]_\rho)
       +|\mathcal J_{V,L}^{\rm face}|([\mathcal S_{L,m,n}]_\rho)
          &\le C_{L,m,n}\rho^{\theta_{L,m,n}},\qquad 0<\rho<1.
       \label{eq:face-DQ-boundary-tightness}
      \end{align}
      Here $\mathcal S_{L,m,n}\subset\mathfrak A_{L,m,n}^{\rm face}$ is the
      finite union of genuine singularity and common-atlas boundaries
      encountered by the depth-$L$ face words
      and by the branchwise tests $Q_s^mg,Q_s^{*n}f$.  In
      particular, bounded--Lipschitz convergence together with
      \eqref{eq:face-DQ-boundary-tightness} implies convergence against every
      bounded test that is branchwise continuous off $\mathcal S_{L,m,n}$.
      The
      chart maps, orientations, amplitudes, and connected-component indices
      used in these currents are those of the common atlas; symbolic
      cylinders without component indices do not define this topology.  For
      each fixed $n$, uniformly near $s=0$,
      \begin{equation}
      \label{eq:DQ-side-fixed-time-continuity}
       \|Q_s^nu\|_{(C^r_0)^*}\le C_n\|u\|_{\cB_0},
       \qquad u\in\cB_{0,0}.
      \end{equation}
      In addition, for every fixed $m,n$,
      \begin{equation}
      \label{eq:fixed-time-evolution-convergence}
       \|Q_s^m-Q_0^m\|_{C^r_0\to\cB_2}\longrightarrow0,
       \qquad
       \|Q_s^n-Q_0^n\|_{\cB_{0,0}\to(C^r_0)^*}\longrightarrow0.
      \end{equation}
      For the regular and face pairings, the corresponding fixed-time tests
      converge branchwise on the same finite atlas away from
      $\mathcal S_{L,m,n}$, with uniform sup norm and the boundary tightness in
      \eqref{eq:face-DQ-boundary-tightness}.  Hence
      \eqref{eq:regular-DQ-weak-convergence} and
      \eqref{eq:face-DQ-weak-convergence} give the regular limit and the
      fixed-core face limit.  Passage from this core to $L(s)$ is not part of
      $\mathrm{MT}_{\rm DQ}$; it is exactly
      $\mathrm{FACE}_{\rm 2CUT}$ below.
      No spectral gap is required on $\cB_0$: it is used only for each fixed
      two-time limit.  All constants are uniform on the shallow spacetime
      cells used below.
      For every fixed $L,m,n$ and $\rho>0$, after identifying the common
      finite atlas, put
      $\mathfrak A_{L,m,n,\rho}^{\rm off}
       :=\mathfrak A_{L,m,n}^{\rm face}
          \setminus[\mathcal S_{L,m,n}]_\rho$.
      The moving face tests satisfy the genuinely uniform clause
      \begin{equation}
      \label{eq:moving-face-test-BL-convergence}
      \begin{aligned}
       \sup_{|s|<s_0}
       \|\Phi_{s,m,n}\|_{{\rm BL}(\mathfrak A_{L,m,n,\rho}^{\rm off})}
       &<\infty,\\
       \|\Phi_{s,m,n}-\Phi_{0,m,n}\|_{{\rm BL}(
          \mathfrak A_{L,m,n,\rho}^{\rm off})}
       &\longrightarrow0.
      \end{aligned}
      \end{equation}
      In particular the family has a common modulus of continuity away from
      the boundary.  Pointwise branchwise convergence with only a uniform
      supremum bound is not accepted.
\end{enumerate}
The interface depends on the chosen measure-trivializing gauge, exactly as
the operator $K_V$ does.  A material or Eulerian physical susceptibility may
be recovered only after adding the corresponding observable-transport terms;
that is a separate theorem and is not called operator CM2 here.
\end{definition}

\begin{definition}[Two-cutoff face difference-quotient interface
$\mathrm{FACE}_{\rm 2CUT}$]
\label{def:face-two-cutoff}
For every fixed $(m,n)$ let $\Phi_{s,m,n}$ be the double-centered
branchwise test induced by $Q_s^mg$ and $Q_s^{*n}f$ on the common oriented
face atlas.  Decompose the coarea current by genuine word depth,
\[
 \mathcal J_{s,\le L}^{\rm face}
   =\sum_{\ell\le L}\mathcal J_{s,\ell}^{\rm face},
 \qquad
 \mathcal J_{V,\le L}^{\rm face}
   =\sum_{\ell\le L}\mathcal J_{V,\ell}^{\rm face}.
\]
Besides the fixed-core convergence and boundary tightness
\eqref{eq:face-DQ-weak-convergence}--
\eqref{eq:face-DQ-boundary-tightness}, require constants
$c_{\rm ret}>A_{\rm face}\ge0$ such that the positive, pre-cancellation
coarea ledger obeys, for every $\ell\ge0$,
\begin{multline}
\label{eq:face-middle-per-depth}
 \sup_{0<|s|<s_0}
 \sum_{\operatorname{depth}(\lambda)=\ell}
 \int_0^1\!\int
 |\partial_sG_\lambda(\theta s,x)|
 \delta(G_\lambda(\theta s,x))
 \mathcal M_{\lambda,m,n}(\theta,s,x)\,dx\,d\theta\\
 \le C_{m,n}e^{-(c_{\rm ret}-A_{\rm face})\ell}.
\end{multline}
The identical bound is required for the limiting positive current at $s=0$.
Here $\mathcal M_{\lambda,m,n}$ is the positive mark dominating the oriented
amplitude, the two fixed-time tests, component multiplicity, and every
actual/split constituent.  It is evaluated in the exact coarea identity
before signed recombination.  It is not a parallel physical-mass ledger:
the positive coarea measure with density $\mathcal M_{\lambda,m,n}$ is
required to be dominated, up to the displayed fixed-time constant
$C_{m,n}$, by the coarea pushforward of the restriction of the propagated
envelope $d\overline M_{s,\mathfrak d}^{q}
=d\Pi_{s,\mathfrak d}q_\lambda$ to
$I_{\rm face/word}\cup I_{\rm main}^{\rm face/word}$.  Hence every operator,
recovery, and opposite-amplitude exponent is already contained in
$A_{\rm face}$ and $c_{\rm face}$.  The domination is uniform in every
suppressed entry of $\mathfrak d$, including the depth-$\ell$ and $(m,n)$
coordinates.  In particular
\eqref{eq:face-middle-per-depth} contains no $|s|^{-1}$ factor.  Put
$c_{\rm face}:=c_{\rm ret}-A_{\rm face}>0$.  Summing the positive
per-depth bound gives the uniform-integrability consequence
\begin{multline}
\label{eq:face-middle-tail}
 \sup_{0<|s|<s_0}
 \left|\left\langle
 \mathcal J_{s,(L_0,L(s)]}^{\rm face},\Phi_{s,m,n}
 \right\rangle\right|
 +\left|\left\langle
 \mathcal J_{V,>L_0}^{\rm face},\Phi_{0,m,n}
 \right\rangle\right|\\
 \le C_{m,n}e^{-c_{\rm face}L_0},\qquad L_0\ge0.
\end{multline}
The terminal part deeper than $L(s)$ is not included in
\eqref{eq:face-middle-tail}; it is the positive-envelope remainder of
$\mathrm{TAIL}_{\rm env}^{+}/\mathrm{NREC}_{\rm deep}$ and is required to
be $o(1)$ in the CM2 norm by
Theorem~\ref{thm:time-depth-diagonal}.  Thus the fixed core, middle band, and
deep remainder are three different estimates.  The no-$|s|^{-1}$ estimate
\eqref{eq:face-middle-per-depth} is an independent positive-current
primitive; neither a common refinement nor $\mathrm{UCR}_{\rm sh}$ implies
it by itself.
\end{definition}

\begin{definition}[Positive-envelope product-time face interface
$\mathrm{FACE\_TIME}_{\rm CM2}$]
\label{def:face-time-cm2}
For the resolution datum belonging to $(s,m,n)$ put
\[
 I_{\rm fw}(s,m,n):=I_{\rm face/word}(s,m,n)
 \sqcup I_{\rm main}^{\rm face/word}(s,m,n).
\]
This is the union of \emph{all} main and nonmain face/word owners in
the static owner partition \eqref{eq:exhaustive-source-owner-partition} and
the main refinement \eqref{eq:main-clock-coverage}; no depth-zero or
fixed-core part is removed.
Condition $\mathrm{FACE\_TIME}_{\rm CM2}$ requires constants
$c_{\rm fw}>0$ and $C<\infty$, uniform in $s,m,n$ and every suppressed
master coordinate, such that the restriction of the same propagated master
envelope satisfies
\begin{equation}
\label{eq:face-time-positive-envelope}
 \overline M_{s,\mathfrak d}^{q}(I_{\rm fw}(s,m,n))
 \le C e^{-c_{\rm fw}\max\{m,n\}}.
\end{equation}
Here $\mathcal E_{\mathfrak d}(f,g)$ denotes the single bounded bilinear physical-test lift
into the master test product; its definition and norm bound are part of this
interface if they have not already been supplied globally.  The lift and
occurrencewise current domination must hold on this restriction with no
time-dependent constant:
\begin{equation}
\label{eq:face-time-test-lift}
 |J_{\rm fw,s,m,n}(\mathcal E_{\mathfrak d}(f,g))|
 \le C\overline M_{s,\mathfrak d}^{q}(I_{\rm fw}(s,m,n))
 \|f\|_{C^r}\|g\|_{C^r}.
\end{equation}
Equivalently, one may replace the right side of
\eqref{eq:face-time-positive-envelope} by a common array $M_{m,n}^{\rm fw}$
with $\sum_{m,n}M_{m,n}^{\rm fw}<\infty$; the exponential form is required
when a numerical CM2 exponent is claimed.  Put
\begin{equation}
\label{eq:face-time-rate}
 c_{\rm face,time}:=c_{\rm fw}/2>0.
\end{equation}
This interface is orthogonal to $\mathrm{FACE}_{\rm 2CUT}$:
$\mathrm{FACE}_{\rm 2CUT}$ controls the word-depth limit for each fixed
$(m,n)$, whereas $\mathrm{FACE\_TIME}_{\rm CM2}$ supplies the missing
product-time majorant.  Neither condition implies the other.
\end{definition}

\begin{lemma}[Face/word product-time majorant]
\label{lem:face-time-majorant}
Under $\mathrm{FACE\_TIME}_{\rm CM2}$, uniformly in $s$ and for the limiting
current,
\begin{equation}
\label{eq:face-time-CM2-majorant}
 |B_{s,m,n}^{\rm fw}|
 \le C e^{-c_{\rm face,time}(m+n)}
 \|f\|_{C^r}\|g\|_{C^r}.
\end{equation}
Hence the face/word channel is absolutely summable in product time.
\end{lemma}

\begin{proof}
Since $\max\{m,n\}\ge(m+n)/2$, combine
\eqref{eq:face-time-positive-envelope} with
\eqref{eq:face-time-test-lift}.  This argument uses positive propagated
mass before cancellation; it does not assume decay of the signed pairing.
\end{proof}

\begin{theorem}[Growing-depth face limit]
\label{thm:face-two-cutoff}
For every fixed $(m,n)$, assume $\mathrm{MT}_{\rm DQ}$ and
$\mathrm{FACE}_{\rm 2CUT}$.  The complete face difference-quotient pairing
converges to the limiting face pairing after fixing the core depth and then
letting $s\to0$.  More precisely,
\begin{equation}
\label{eq:face-growing-depth-convergence}
 \left\langle\mathcal J_{s,\le L(s)}^{\rm face},
       \Phi_{s,m,n}\right\rangle
 \longrightarrow
\left\langle\mathcal J_{V}^{\rm face},
       \Phi_{0,m,n}\right\rangle,
\end{equation}
If, in addition, all hypotheses of
Theorem~\ref{thm:time-depth-diagonal} hold, then adding the part deeper than
$L(s)$ changes the left side by $o(1)$ in the complete CM2 sum.  This last
assertion is not a consequence of $\mathrm{MT}_{\rm DQ}$ and
$\mathrm{FACE}_{\rm 2CUT}$ alone.
\end{theorem}

\begin{proof}
Given $\varepsilon>0$, choose $L_0$ so that the right side of
\eqref{eq:face-middle-tail} is below $\varepsilon$.  With this $L_0$ fixed,
bounded--Lipschitz convergence on the common atlas and
\eqref{eq:face-DQ-boundary-tightness} give convergence of the core pairing
at depth at most $L_0$.  Equation~\eqref{eq:face-middle-tail} removes both
the finite-$s$ triangular band $L_0<\ell\le L(s)$ and the limiting tail
$\ell>L_0$, uniformly in $s$.  Finally
under its additional hypotheses,
Theorem~\ref{thm:time-depth-diagonal} removes the terminal part
$\ell>L(s)$.  The order is therefore $L_0$ first, $s\to0$ second, and,
when invoked, the deep positive-envelope estimate last.  A fixed-depth topology is never used
as a substitute for growing-depth uniform integrability.
\end{proof}

\begin{lemma}[Fixed-gauge centering and typed finite-DQ convergence]
\label{lem:fixed-gauge-centering}
Under $\mathrm{MT}_{\rm DQ}$, the same projection $\Pi_0$ centers the
finite-$s$ operators, endpoint observables, branch-section modes, and both
standard-family couplings.  Every term has precisely one of the regular,
face, product-current, or response types in
Definition~\ref{def:measure-gauge-interface}; the face and product types are
the two inseparable summands of the centered cluster
\eqref{eq:typed-face-product-cluster}.  Moreover
\begin{equation}
\label{eq:derived-strong-weak-CM2}
 |\langle Q_s^n\widehat K_{s,\rm sw}^{\rm DQ}Q_s^mg,f\rangle|
 \le C\theta_-^m\theta_+^n\|g\|_{C^r}\|f\|_{C^r},
\end{equation}
and the same uniform majorant holds for $E_s$.  For every fixed $(m,n)$ the
product-current block converges in
\eqref{eq:product-current-DQ-convergence} and pairs to zero with the two
individually centered endpoint tests.  Hence the centered face--product
cluster has exactly the same pairing as its graph-face summand on those
tests; this does not make either summand separately centered.  The response block converges to
$\widehat K_{V,\rm sw}$ and the $E_s$ pairing
tends to zero.  The regular pairing converges in
\eqref{eq:regular-DQ-weak-convergence}; the face pairing converges on every
fixed core in \eqref{eq:face-DQ-weak-convergence}, and its growing-depth
limit is supplied by Theorem~\ref{thm:face-two-cutoff}.  Consequently, once the
regular and face estimates below provide their summable uniform majorants,
one common dominated-convergence argument exchanges the limit with the
absolute CM2 double sum for all four types.  Explicitly, if
$B_{s,m,n}^{j}$ denotes the pairing of type
$j\in\{\mathrm{reg},\mathrm{face},\mathrm{prod},\mathrm{sw}\}$ and the corresponding
ledger gives
\begin{equation}
\label{eq:typed-common-dominated-convergence}
 |B_{s,m,n}^{j}|\le M_{m,n}^{j},\qquad
 \sum_{j}\sum_{m,n\ge0}M_{m,n}^{j}<\infty,
\end{equation}
then
\[
 \lim_{s\to0}\sum_j\sum_{m,n\ge0}B_{s,m,n}^{j}
 =\sum_j\sum_{m,n\ge0}\lim_{s\to0}B_{s,m,n}^{j}.
\]
\end{lemma}

\begin{proof}
Equation~\eqref{eq:measure-trivializing-gauge} gives
$\mu_0(P_sh)=\mu_0(h)$, hence $P_s\Pi_0=\Pi_0P_s=\Pi_0$ and
$\partial_s\Pi_0=0$.  Apply \eqref{eq:strong-side-memory}, then
\eqref{eq:uniform-strong-weak-DQ}, and finally
\eqref{eq:response-side-memory}; duality gives
\eqref{eq:derived-strong-weak-CM2}.  Replacing
$\widehat K_{s,\rm sw}^{\rm DQ}$ by $E_s$ gives the same uniform majorant.
For fixed $(m,n)$ add and subtract the expressions with $Q_0^m$ and $Q_0^n$.
Equations \eqref{eq:strong-DQ-convergence}--
\eqref{eq:fixed-time-evolution-convergence} show that the response difference
and the $E_s$ pairing tend to zero.  The individual double projection
\eqref{eq:typed-block-double-projection} makes the regular and response
outputs eligible for the zero-mass estimates.  The face and product
summands are instead used only through their exact jointly centered cluster
\eqref{eq:typed-face-product-cluster}; no zero-mass estimate is applied to
either one alone.  Dominated
convergence with \eqref{eq:derived-strong-weak-CM2} completes the claim.  The
regular claim follows from
\eqref{eq:regular-DQ-weak-convergence} and fixed-time branchwise continuity.
The product-current pairing vanishes because each tensor summand contains
$\mu_0$ in one slot while both endpoint tests are centered; its separate
topology prevents this fact from being used to misclassify it as a density
or a graph face, and the equality with the clustered pairing follows before
absolute values are taken.  For the face core, use
\eqref{eq:moving-face-test-BL-convergence} outside
$[\mathcal S_{L,m,n}]_\rho$; then use
\eqref{eq:face-DQ-boundary-tightness} for the two discarded boundary
neighbourhoods and let $\rho\downarrow0$.  Then apply
Theorem~\ref{thm:face-two-cutoff} to pass from the fixed core to the complete
face pairing.  This is the standard
weak-current plus boundary-tightness upgrade; it is not supplied by the
uniform face norm \eqref{eq:uniform-gauge-DQ-face}.  The regular and face
CM2 majorants are proved later from their separate ledgers, after which the
same dominated-convergence theorem applies.
\end{proof}

The divergence formula for smooth Anosov transfer derivatives in
\cite{FroylandPhalempin} maps a strong anisotropic space to an intermediate
response space, whereas the finite difference quotients converge only in a
still weaker operator topology.  The spectral gap is used on the response
space and is not imposed on the weakest space.  This is the three-space type
model behind $\mathrm{MT}_{\rm DQ}$; the derivative is not an $M\times M$
Sobolev density.  The fixed zero-mass organization is
consistent with the measure-valued-section response framework of
\cite{AlvesBahsoun} and the extended-space loss-of-memory framework of
\cite{GalatoloLucarini}.  None of these references supplies the singular
moving-billiard decomposition or its face atlas; those remain explicit
interfaces here.  A pure auxiliary conjugacy provides a decisive type test:
the full physical pairing with transported observables may have zero
derivative while the fixed-gauge commutator is nonzero.  This is why the two
targets are not identified.
The moving-singularity obstruction in \cite{BahsounGalatolo} is a further
reason not to infer this interface from smooth-map response theory.
For the face topology, \cite[Lemma~5.4]{CanestrariHoles} gives the closest
available model: weak-* limits of standard families with uniform boundary
$Z$ may still be tested by bounded functions whose discontinuities lie on a
fixed finite singularity set.  Equations
\eqref{eq:face-DQ-weak-convergence}--
\eqref{eq:face-DQ-boundary-tightness} require the corresponding statement for
the oriented moving-face currents themselves; the cited lemma does not prove
that moving-billiard DQ input.

\section{Why the direct projective-cone bridge fails}

The projective cone for dispersing billiards controls integrals of a density
against positive dynamically H\"older tests on stable curves, short-leaf
mass, and H\"older comparison of such integrals between nearby stable curves.
A moving inverse-branch boundary instead produces, before physical assembly,
a coefficient of the form
\begin{equation}
\label{eq:face-model}
  h\longmapsto w\,(h\vert_{I^+}-h\vert_{I^-})\,\delta_\Gamma,
\end{equation}
where $I^\pm$ are one-sided terminal preimage arcs.  In the radial collision
atlas these arcs land on the grazing seam and are transversal to the stable
test foliation.  Thus \eqref{eq:face-model} asks for a transverse trace, not a
stable-leaf average.

The following elementary local model rules out the missing implication.

\begin{proposition}[Stable-leaf control does not dominate transverse trace]
\label{prop:no-trace}
Let $R=(-2,2)^2$, let the local stable leaves be the horizontal segments
$W_y=(-2,2)\times\{y\}$, and let the transversal face be
$I=\{0\}\times(-1,1)$.  Fix $0<q<1$ and $0<\gamma<1$.  There are smooth
functions $u_\varepsilon$ supported in an $\varepsilon$-tube about $I$ such
that
\[
 u_\varepsilon\vert_I=b\ne0,
\]
while, uniformly over bounded H\"older test functions $\psi$ on stable
subsegments $W\subset W_y$,
\begin{align}
 \left|\int_Wu_\varepsilon\psi\,dm_W\right|
 &\le C\min\{|W|,\varepsilon\}\|\psi\|_\infty,
       \label{eq:leaf-small}\\
 \sup_{0<|W|\le1}|W|^{-q}
 \left|\int_Wu_\varepsilon\psi\,dm_W\right|
 &\le C\varepsilon^{1-q}\|\psi\|_\infty,\label{eq:short-small}
\end{align}
and the H\"older comparison of these integrals between nearby parallel
stable leaves is $O(\varepsilon)$.  Consequently no norm or order estimate
using only the stable-leaf quantities in
\eqref{eq:leaf-small}--\eqref{eq:short-small} and their transverse H\"older
comparisons can bound $\|u\vert_I\|_{C^0}$.
\end{proposition}

\begin{proof}
Choose $\phi\in C_c^\infty((-1,1))$ with $\phi(0)=1$ and
$b\in C_c^\infty((-1,1))$, $b\ne0$, and set
\[
  u_\varepsilon(x,y)=\phi(x/\varepsilon)b(y).
\]
The trace assertion is immediate.  A horizontal subsegment meets the support
in length at most $C\min\{|W|,\varepsilon\}$, proving
\eqref{eq:leaf-small}.  If $|W|\le\varepsilon$, multiplication by
$|W|^{-q}$ gives at most $C\varepsilon^{1-q}$; if
$|W|\ge\varepsilon$, it gives at most
$C\varepsilon|W|^{-q}\le C\varepsilon^{1-q}$, proving
\eqref{eq:short-small}.  Replacing $y$ by $y'$ changes the integral by at most
$C\varepsilon|y-y'|\|b\|_{C^1}\|\psi\|_{C^0}$, with the usual additional
H\"older test-comparison term.  The trace is independent of
$\varepsilon$, so the asserted domination is impossible.
\end{proof}

\begin{corollary}[Failure of the proposed positive cone operator]
The local face map \eqref{eq:face-model} is not known to be positive for the
Demers--Liverani cone order, and its boundedness cannot be inferred from
projective contraction.  In particular, the implication
\[
 \Theta(P_0^mf_+,P_0^mf_-)\le C\chi^m
 \quad\Longrightarrow\quad
 \|K_VP_0^m(f_+-f_-)\|_Y\le C'\chi^m
\]
is invalid without an additional transverse-trace theorem.
\end{corollary}

\begin{proof}
The cone order is tested through stable-leaf integrals.  The face coefficient
contains precisely the trace excluded by Proposition~\ref{prop:no-trace}.
Projective contraction therefore supplies no order comparison for the two
output face measures.
\end{proof}

This is a go/no-go conclusion, not merely a missing estimate: the direct
``past cone $\to$ face $\to$ future family'' proof must be abandoned unless
physical assembly first converts the terminal trace into an integral
functional controlled by the cone.  The existing adjoint dyadic bound does
not perform that conversion.

\section{A sufficient two-sided bridge}

The complete defect is a kernel located at the event time.  It must be
controlled from the past and the future simultaneously.  The following
criterion is an exact sufficient condition, but not a characterization of
CM2.

\begin{definition}[Centered nuclear recovery bridge]
Let $X_-$ be a Banach space of centered past observables and let $Y_+$ be a
Banach space of centered signed measured unstable families.  A centered
operator $K:X_-\to Y_+$ has a nuclear recovery presentation if
\begin{equation}
\label{eq:nuclear}
 K=\sum_{j\ge1}\ell_j\otimes\eta_j,
 \qquad
 \sum_{j\ge1}\|\ell_j\|_{X_-^*}\|\eta_j\|_{Y_+}<\infty,
\end{equation}
where $\ell_j(\one)=0$ and $\eta_j(\one)=0$ after the complete regular and
physical-face assembly.
\end{definition}

\begin{theorem}[Two-sided bridge implies CM2]
\label{thm:bridge}
Assume
\begin{align}
 \|Q^mg\|_{X_-}&\le C_-\tau^m\|g\|_{X_-},
       \label{eq:past}\\
 |\eta(\cU_0^nf)|&\le C_+\vartheta^n
       \|\eta\|_{Y_+}\|f\|_{C^1},\label{eq:future}
\end{align}
for centered $g$ and centered $\eta$.  If $K$ has the presentation
\eqref{eq:nuclear}, then
\[
 |(KQ^mg)(\cU_0^nf)|
 \le C_-C_+\mathfrak N(K)\tau^m\vartheta^n
       \|g\|_{X_-}\|f\|_{C^1},
\]
where
$\mathfrak N(K)=\sum_j\|\ell_j\|_{X_-^*}\|\eta_j\|_{Y_+}$.
\end{theorem}

\begin{proof}
Using \eqref{eq:nuclear}, \eqref{eq:past}, and \eqref{eq:future},
\begin{align*}
 |(KQ^mg)(\cU_0^nf)|
 &\le\sum_j|\ell_j(Q^mg)|\,
              |\eta_j(\cU_0^nf)|\\
 &\le C_-C_+\tau^m\vartheta^n\|g\|_{X_-}\|f\|_{C^1}
       \sum_j\|\ell_j\|_{X_-^*}\|\eta_j\|_{Y_+}.
\end{align*}
\end{proof}

\begin{remark}
Used globally, a one-gap functional-correlation estimate gives only
$O(\tau^m+\vartheta^n)$ and does not by itself prove CM2.  It is nevertheless
exactly sufficient in the central band, where $m\asymp n$: after smoothing
the event kernel, one decoupling gives $e^{-c\max\{m,n\}}$, hence an
$e^{-c'(m+n)}$ bound there.  Forward and reverse standard-family coupling
handle the two off-diagonal regions; Lemma~\ref{lem:three-region} assembles
the product estimate.  Thus no operator-ideal estimate is used in the
principal proof.
\end{remark}

\begin{remark}[The nuclear criterion is generally too strong]
\label{rem:nuclear-too-strong}
A physical face current is supported on a graph in a stable--unstable product
chart.  Its kernel is therefore analogous to a pullback or a Fourier integral
operator, not to a smoothing operator.  Even after double centering, a graph
delta need not be nuclear between natural past and future spaces.  Thus failure
of \eqref{eq:nuclear} or of the projective-tensor condition below does
\emph{not} imply failure of CM2.  Product decay may instead arise dynamically,
from oscillation and transversality after the graph is acted on by both time
directions.
\end{remark}

\section{Recovery-time weighted future space}

For a measured unstable component $(W,\rho)$ let $R(W)$ be the first time at
which its canonical forward subdivision is a proper standard family.  For
$\omega>0$, define schematically
\begin{equation}
\label{eq:Yplus}
 \|\eta\|_{Y_+}
 :=\inf\sum_j|a_j|e^{\omega R(W_j)}
       \bigl(1+\mathrm{Reg}(\rho_j)\bigr),
 \qquad
 \eta=\sum_ja_j(W_j,\rho_j).
\end{equation}
The infimum also allows the positive and negative Jordan parts to be split by
recovery rank.  Standard-family coupling and recovery give
\eqref{eq:future} once the constants in \eqref{eq:Yplus} are chosen within
the recovery tail.

For a radial grazing block $b$, the physical flux factorization gives a mass
bound $M_b\le C2^{-b}$.  If its geometric deterioration is polynomial in
$2^b$, the recovery lemmas give
\[
 R_b\le C_0+C_1b.
\]
Hence
\begin{equation}
\label{eq:recovery-sum}
 \sum_{b\ge0}M_be^{\omega R_b}<\infty
 \quad\text{whenever}\quad \omega C_1<\log2.
\end{equation}
Thus initial finite $Z$ is unnecessary.  What remains is not future recovery,
but the past functional $\ell_j$ in \eqref{eq:nuclear}.

\subsection{Recovery windows and the actual summability target}

For differentiability of the susceptibility, a uniform pointwise CM2 bound
for every grazing block is stronger than necessary.  The required conclusion
is absolute summability of the complete block series.  This distinction
removes an artificial exponential loss in the grazing rank.

\begin{proposition}[Recovery-window summability]
\label{prop:recovery-window-summability}
Let $B_b(m,n)$ be the scalar pairing produced by block $b$ at past and future
times $m,n$.  Suppose that there are masses $M_b$, recovery times
$R_b^-,R_b^+$, constants $C<\infty$, and $0<\tau,\vartheta<1$ such that
\begin{equation}
\label{eq:shifted-block-decay}
 |B_b(m,n)|
 \le CM_b\tau^{(m-R_b^-)_{+}}
             \vartheta^{(n-R_b^+)_{+}},
 \qquad m,n\ge0.
\end{equation}
Then
\begin{equation}
\label{eq:recovery-window-sum}
 \sum_{b,m,n}|B_b(m,n)|
 \le C'
 \sum_bM_b(1+R_b^-)(1+R_b^+).
\end{equation}
In particular, if $M_b\le C2^{-\alpha b}$ for some $\alpha>0$ and
$R_b^\pm\le C_0+C_1b$, then the full threefold series converges.  The same
conclusion holds uniformly for parameter difference quotients if the
constants and block masses are uniform.

The deterministic recovery times may be replaced by a countable refinement
$B_b=\sum_jB_{b,j}$ with masses $M_{b,j}$ and times $R_{b,j}^\pm$, provided
\begin{equation}
\label{eq:bi-recovery-moment}
 \sum_jM_{b,j}(1+R_{b,j}^-)(1+R_{b,j}^+)
 \le CM_b(1+b)^2.
\end{equation}
It is enough for the two-sided recovery tail, after its deterministic
$O(b)$ entrance time, to be exponential.
\end{proposition}

\begin{proof}
For every $R\ge0$ and $0<q<1$,
\[
 \sum_{k\ge0}q^{(k-R)_+}
 \le R+1+\frac1{1-q}.
\]
Sum \eqref{eq:shifted-block-decay} first in $m,n$ and then in $b$.
Exponential block mass dominates the quadratic recovery window.  Applying the
same calculation to every refined piece gives
\eqref{eq:bi-recovery-moment}.  An exponential tail for
$R_{b,j}^-+R_{b,j}^+-O(b)$ has uniformly bounded polynomial moments, proving
the last assertion.
\end{proof}

\begin{lemma}[Two marginal growth tails imply bi-recovery]
\label{lem:marginal-to-bi-recovery}
Let a block of total variation $M_b$ be cut by the common refinement of its
past and future recovery trees, with refined masses $M_{b,j}$ and recovery
times $R_{b,j}^\pm$.  Assume cutting preserves total variation and, for some
$r_b\le C_0+C_1b$, $0<\rho<1$,
\[
 \sum_{j:R_{b,j}^\pm>r_b+t}M_{b,j}
 \le CM_b\rho^t,
 \qquad t\ge0,
\]
for each sign.  Then \eqref{eq:bi-recovery-moment} holds.
\end{lemma}

\begin{proof}
If $R_{b,j}^-+R_{b,j}^+>2r_b+t$, at least one marginal time is larger than
$r_b+t/2$.  The union bound therefore gives an exponential tail for the sum
of the two recovery times under the normalized weights $M_{b,j}/M_b$.
Its second moment is $O((1+r_b)^2)$.  Since
$(1+R^-)(1+R^+)\le(1+R^-+R^+)^2$, the claim follows.
\end{proof}

For a general graph current the relevant geometric input is therefore a
\emph{proper bi-recovery} statement, not an estimate of a Fourier constant on
the original grazing curve.  For the clean first radial current this input is
Corollary~\ref{cor:radial-bi-recovery}.  A central two-time theorem then gives
\eqref{eq:shifted-block-decay}.  The fully unrecovered rectangle is handled by
total variation, the two mixed rectangles by the corresponding recovered
one-sided estimate, and the fully recovered rectangle by the product bound.
The growth-lemma tail sums these windows through
\eqref{eq:bi-recovery-moment}.

Consequently even the generic $M_b=O(2^{-b})$ assembly margin is sufficient
once proper bi-recovery is proved; the stronger first-radial
$O(2^{-(3-\epsilon_*)b})$ margin gives additional room but is no longer
needed to absorb an exponential Fourier constant in $b$.  This decouples the
grazing rank from the return-rank truncation in
Proposition~\ref{prop:countable-truncation}.

\section{Billiard-specific block problem}

For each genuine physical block $b$, the complete difference quotient must be
written, after artificial-face cancellation, as
\begin{equation}
\label{eq:block-kernel}
 K_{V,b}=\sum_{k\ge1}\ell_{b,k}\otimes\eta_{b,k},
\end{equation}
with
\begin{equation}
\label{eq:block-sum}
 \sum_{b,k}\|\ell_{b,k}\|_{X_-^*}
             \|\eta_{b,k}\|_{Y_+}<\infty.
\end{equation}
The desired estimate follows from Theorem~\ref{thm:bridge}.  Three sufficient
mechanisms for proving \eqref{eq:block-sum} are:
\begin{enumerate}[label=\textup{(\roman*)}]
\item a reverse-time conditional-mixing theorem for the terminal slice,
      furnishing $X_-^*$ directly;
\item a two-sided Young-tower lift of the complete centered defect, with an
      event kernel that is nuclear between the past and future tower norms;
\item a smoothing/nuclear decomposition of the terminal restriction, with a
      quantitative bunching inequality strong enough to sum the smoothing
      error over the homogeneity blocks.
\end{enumerate}
The projective-cone contraction alone supplies none of these three inputs.
These mechanisms are not exhaustive because a non-nuclear graph current can
still equidistribute under the two-sided dynamics.

\subsection{A two-sided tower formulation}

The natural extension gives a sharper version of alternative (ii).  Let
$\Delta$ be a two-sided Young tower for the fixed billiard, with past and
future cylinder filtrations $\mathcal F_i^-$ and $\mathcal F_j^+$.  Denote
the associated martingale differences by
\[
 D_i^-=\mathbb E(\,\cdot\mid\mathcal F_i^-)
       -\mathbb E(\,\cdot\mid\mathcal F_{i-1}^-),
 \qquad
 D_j^+=\mathbb E(\,\cdot\mid\mathcal F_j^+)
       -\mathbb E(\,\cdot\mid\mathcal F_{j-1}^+).
\]
A lifted event kernel $\widehat K$ has summable mixed cylinder variation if
\begin{equation}
\label{eq:mixed-cylinder}
 \sum_{i,j\ge1}
 \|D_i^-D_j^+\widehat K\|_{\mathcal B_-^*\widehat\otimes_\pi
                                      \mathcal B_+}<\infty.
\end{equation}
The indices start at one: complete double centering removes the two axis
terms $D_i^-D_0^+\widehat K$ and $D_0^-D_j^+\widehat K$.

\begin{proposition}[Mixed-cylinder criterion]
\label{prop:mixed-cylinder}
Suppose the two-sided tower has exponential tails and the past and future
transfer actions contract the centered martingale blocks at rates
$\tau^m$ and $\vartheta^n$.  If the complete lifted moving defect satisfies
\eqref{eq:mixed-cylinder}, then it has a nuclear recovery presentation
\eqref{eq:nuclear} and hence satisfies CM2.
\end{proposition}

\begin{proof}
Expand
\[
 \widehat K=\sum_{i,j\ge1}D_i^-D_j^+\widehat K
\]
in the projective tensor product.  Each summand is a norm-convergent sum of
rank-one past/future tensors by the definition of
$\widehat\otimes_\pi$.  Summability in \eqref{eq:mixed-cylinder} therefore
gives \eqref{eq:nuclear}.  Shifting $m$ steps into the past and $n$ steps into
the future contracts the two centered tensor factors separately.  Apply
Theorem~\ref{thm:bridge}.
\end{proof}

This criterion identifies one strong geometric issue.  A smooth two-dimensional
density on a tower rectangle has summable mixed cylinder variation.  A face
current is instead supported on a one-dimensional graph in the local
stable--unstable product chart.  Its graph delta is not automatically in the
projective tensor product.  Proving \eqref{eq:mixed-cylinder} is therefore a
genuine conditional-mixing theorem, not routine tower bookkeeping.  The
finite-strip central face supplies uniform transversality and one-sided
recovery, but these facts alone do not prove the mixed tensor estimate.

\section{The intrinsic dynamic graph-current formulation}

Let $\gamma_b:I_b\to M$ parametrize one side of a complete physical face in
block $b$, let $F_b\gamma_b$ denote its assembled terminal image, and let
$w_b$ be the signed coarea weight after artificial-face cancellation.  The
face contribution to CM2 has the model form
\begin{equation}
\label{eq:dynamic-graph}
 B_{m,n}^{(b)}(g,f)
 =\int_{I_b} g(T^{-m}\gamma_b(t))
                  f(T^nF_b\gamma_b(t))w_b(t)\,dt .
\end{equation}
Equivalently, if $J_{V,b}$ is the graph current on $M\times M$ induced by
$t\mapsto(\gamma_b(t),F_b\gamma_b(t))$, then
\begin{equation}
\label{eq:product-current}
 B_{m,n}^{(b)}(g,f)
 =\bigl((T^{-m}\times T^n)_*J_{V,b}\bigr)(g\otimes f).
\end{equation}
Complete double centering removes the two marginals of $J_{V,b}$.  The
intrinsic CM2 statement is therefore two-parameter equidistribution of a
centered graph current under $T^{-1}\times T$ at a product rate.

This formulation explains both the difficulty and the opportunity.  A graph
curve is one-dimensional, whereas the unstable bundle of the product map is
two-dimensional, so ordinary standard-family coupling does not apply
directly.  On the other hand, the current need not be statically nuclear:
hyperbolic stretching sends its Fourier frequencies to exponentially large
scales, and polynomial Fourier decay of the induced slice is already enough
to produce exponential decay in collision time.

\begin{proposition}[Fourier reduction on one induced rectangle]
\label{prop:fourier-reduction}
Assume that an induced return map has a stable--unstable product chart in
which the centered face slice induces a finite signed measure $\nu_b$ on the
contracting coordinate.  Suppose that, uniformly over admissible inverse
words and over the relevant $C^2$ phase maps $F$,
\begin{equation}
\label{eq:fourier-decay}
 |\widehat{F_*\nu_b}(\xi)|\le C_b(1+|\xi|)^{-\eta}
\end{equation}
for some $\eta>0$.  Suppose also that, after excluding phase cancellation,
the expanding frequencies generated by $m$ past returns and $n$ future
returns are bounded below by
\[
 c\max\{\Lambda_-^m,\Lambda_+^n\},
\]
apart from cylinder pieces whose total centered mass is bounded by
$C\tau_0^m\vartheta_0^n$.  Then, for sufficiently regular
centered cylinder observables,
\begin{equation}
 |B_{m,n}^{(b)}(g,f)|
 \le C C_b\,\tau^m\vartheta^n\|g\|_{C^r}\|f\|_{C^r}
\end{equation}
for some $\tau,\vartheta<1$.
\end{proposition}

\begin{proof}
Decompose $g$ and $f$ into Fourier modes in the product chart and split the
sum into low and high frequencies.  Centering and cylinder mixing control the
low-frequency part.  In the high-frequency part, hyperbolicity converts the
two return depths into the lower frequency scale
$c\max\{\Lambda_-^m,\Lambda_+^n\}$.  Since
\[
 \max\{\Lambda_-^m,\Lambda_+^n\}
 \ge (\Lambda_-^m\Lambda_+^n)^{1/2},
\]
\eqref{eq:fourier-decay} gives the product factor
$\Lambda_-^{-\eta m/2}\Lambda_+^{-\eta n/2}$.  Summability of the Fourier
coefficients of $C^r$ tests and the exceptional-cylinder bound complete the
estimate.  The loss of one half in each exponent is harmless.
\end{proof}

This preliminary proposition is not used as an assumption in the final proof.
Uniform \eqref{eq:fourier-decay} is supplied below by the cohomological
local slope-energy theorem, the dominant-coordinate reduction, and the
physical fractional nonstationary-phase estimate.  One-sided mixing alone
would not supply it.

\subsection{A three-region reduction}

Fourier analysis is not needed uniformly over the whole $(m,n)$ quadrant.
One-sided standard-family recovery gives a forward estimate with a past
regularity cost, and reversibility gives its backward analogue.  The following
elementary lemma isolates the remaining diagonal band.

\begin{lemma}[Off-diagonal coupling plus diagonal oscillation]
\label{lem:three-region}
Suppose
\begin{align}
 |B_{m,n}|&\le C e^{\beta_-m-\alpha_+n},\label{eq:forward-cone}\\
 |B_{m,n}|&\le C e^{-\alpha_-m+\beta_+n},\label{eq:backward-cone}
\end{align}
with all four exponents positive.  Choose
$A>\beta_-/\alpha_+$ and $D>\beta_+/\alpha_-$.  Suppose in addition that
\[
 |B_{m,n}|\le Ce^{-\gamma\max\{m,n\}}
\]
whenever $D^{-1}m<n<Am$.  Then there are
$\tau,\vartheta\in(0,1)$ such that
$|B_{m,n}|\le C'\tau^m\vartheta^n$ for all $m,n$.
\end{lemma}

\begin{proof}
If $n\ge Am$, \eqref{eq:forward-cone} is bounded by
$e^{-c_+(m+n)}$ for some $c_+>0$.  If $m\ge Dn$, use
\eqref{eq:backward-cone}.  The remaining indices lie in the stated diagonal
band, where $\max\{m,n\}\ge(m+n)/2$.  Taking the minimum of the three positive
exponents proves the claim.
\end{proof}

For billiards, \eqref{eq:forward-cone} follows from signed
standard-family coupling after cutting the past-weighted density into
homogeneous pieces: the growth lemma gives the coupling factor
$e^{-\alpha_+n}$, while differentiating the $m$-step past weight costs at
most $e^{\beta_-m}$.  The marked-amplitude moment sums the homogeneous
pieces.  Applying the same argument to the reversible collision map gives
\eqref{eq:backward-cone}.  All constants are uniform on the fixed compact
atlas; recovered grazing pieces enter only after their recorded recovery
times; these are the signed standard-family forms of the usual growth and
coupling estimates \cite{Chernov1999,DLreview}.
Consequently the genuinely new slope/Fourier theorem is needed only for
$m\asymp n$.

This also corrects the naive ambient phase test.  The graph
$t\mapsto(T^{-m}\gamma(t),T^nF\gamma(t))$ is one-dimensional in a
four-dimensional product.  Its first two derivatives cannot be uniformly
transverse to every Fourier covector.  In the diagonal band one must first
discard the contracted coordinates and pass to the dominant past-stable and
future-unstable coordinates.  The resulting two-dimensional plane curve is
governed by the stable--unstable temporal-distance/holonomy function.  A
finite-difference slope-value QNL is appropriate for this reduced phase.  If
one insists on working in the full four-dimensional chart, a finite-type
condition involving derivatives through order at least four is required.

\subsection{A verified radial nonlinearity seed}

The first radial face is not an abstract graph.  For two circular scatterers
with radii $R_j,R_i$ and centre distance $d>R_j+R_i$, put the source centre at
the origin, the target centre at $(d,0)$, and parametrize the source boundary
by
\[
 q_j(\alpha)=R_j(\cos\alpha,\sin\alpha).
\]
Let $r(\alpha)=|(d,0)-q_j(\alpha)|$ and
$\beta(\alpha)=\arg((d,0)-q_j(\alpha))$.  The two directions tangent to the
target circle are
\[
 \theta_\sigma(\alpha)
 =\beta(\alpha)+\sigma\arcsin\frac{R_i}{r(\alpha)},
 \qquad \sigma\in\{-1,1\}.
\]
Thus the radial physical face in Birkhoff coordinates is locally
$\varphi_\sigma(\alpha)=\theta_\sigma(\alpha)-\alpha$.

\begin{lemma}[The circular tangency face is genuinely curved]
\label{lem:circular-face-curvature}
At the collinear source point $\alpha=0$, write $r_0=d-R_j$.  Then
\begin{equation}
\label{eq:circular-face-curvature}
 \varphi_\sigma''(0)
 =-\sigma\frac{R_i dR_j}
 {r_0^2\sqrt{r_0^2-R_i^2}}\ne0.
\end{equation}
Moreover the angular coarea derivative of
\[
 G=|((d,0)-q_j)\times v_\theta|^2-R_i^2
\]
is nonzero there.  Hence each tangency choice gives a smooth, non-affine
central face germ with a smooth coarea density.
\end{lemma}

\begin{proof}
Direct differentiation gives
\[
 r'(0)=0,\qquad r''(0)=\frac{dR_j}{r_0},\qquad
 \beta''(0)=0.
\]
Since
\[
 \frac d{dr}\arcsin\frac{R_i}{r}
 =-\frac{R_i}{r\sqrt{r^2-R_i^2}},
\]
equation \eqref{eq:circular-face-curvature} follows.  On the tangency face,
$|(d,0)-q_j|\times v_\theta|=R_i$ while the longitudinal component has
magnitude $\sqrt{r_0^2-R_i^2}$; therefore
$|\partial_\theta G|=2R_i\sqrt{r_0^2-R_i^2}>0$.
\end{proof}

This proves a useful but limited fact: the radial face supplies a concrete
one-step nonlinearity seed and no Baker-type affine resonance is forced by
the circular geometry.  It does not prove uniform quantitative nonlinearity
after arbitrary past and future words.  Near the diagonal where the two
terms in the reduced phase
\begin{equation}
\label{eq:phase}
 \Phi_{m,n,\xi,\zeta}(t)
 =\xi\cdot T^{-m}\gamma(t)+\zeta\cdot T^nF\gamma(t)
\end{equation}
have comparable size, their derivatives can cancel;
this is the only time regime in which a genuine Dolgopyat/sum-product
argument is needed.  Off that diagonal one time direction dominates and
ordinary nonstationary phase applies.

\subsection{A countable-alphabet truncation lemma}

The latest smooth result closest to the required mechanism is Leclerc's
Fourier-decay theorem for equilibrium states of nonlinear, area-preserving
smooth Axiom--A surface diffeomorphisms.  Its quantitative nonlinearity
condition is expressed through a temporal-distance function and is generic,
with an explicit periodic-point cocycle test.  A billiard return scheme is
piecewise smooth with a countable alphabet, so that theorem is not directly
applicable.  The countable alphabet itself, however, is not the fundamental
obstruction.

\begin{proposition}[Exponential-tail reduction to uniform finite truncations]
\label{prop:countable-truncation}
Let an induced symbolic system have a return/singularity rank $R$ with
\[
 \mu\{R>L\}\le C e^{-aL}.
\]
Suppose a Fourier proof at frequency $|\xi|$ uses at most
$N(\xi)\le c_0\log(2+|\xi|)$ successive symbols.  Assume that the subsystem
with all ranks at most $L$ satisfies, uniformly in the truncation,
\[
 \left|\int e^{i\xi\psi}\chi\,d\mu_L\right|
 \le C e^{cL}|\xi|^{-\eta}.
\]
Then the full countable system has polynomial Fourier decay, provided
$a,\eta>0$ and the finite-truncation constants have the displayed
exponential bound.
\end{proposition}

\begin{proof}
The mass of orbit words which see a rank larger than $L$ is at most
$C N(\xi)e^{-aL}$.  On the complement apply the finite-truncation estimate.
Choose $L=\delta\log|\xi|$ with $0<\delta<\eta/c$ (with no upper restriction
if $c=0$).  The good contribution is
$O(|\xi|^{-\eta+c\delta})$ and the bad contribution is
$O(\log|\xi|\,|\xi|^{-a\delta})$.  Both have a positive power saving.
\end{proof}

Thus this fallback would require a \emph{uniform} singular version of
Leclerc's quantitative-nonlinearity estimate on finite homogeneity/return
truncations.  Exponential tower tails could then remove the truncation, but
the truncation-uniform constants are not supplied by that route.  The
direct gate-resolved slope argument below bypasses this fallback entirely.

\subsection{Uniform QNL requires anchored big images}

Inspection of the QNL proof shows that a single nonzero periodic cocycle is
not, by itself, a uniform truncation theorem.  The power exponent also depends
on a cutting constant: at every scale one removes a child interval carrying
at least an $\eta_{\rm cut}$ fraction of its parent's conditional mass.  For
an arbitrary exhaustion by finite subshifts, $\eta_{\rm cut}$ may tend to
zero; then the QNL exponent tends to zero and
Proposition~\ref{prop:countable-truncation} is useless.

The billiard inducing scheme has a possible cure.  A Gibbs--Markov return map
has a finite big-images/preimages core, and on a compact central magnet the
SRB conditional densities on homogeneous unstable leaves are bounded above
and below.  The truncations must therefore be \emph{anchored}: retain a fixed
finite family of full-image connector branches, including the nonlinear
periodic witness, and add all branches up to rank $L$.  The desired uniform
cutting statement is
\begin{equation}
\label{eq:uniform-child-mass}
 \nu_{x,L}(I_{\mathrm{child}})
 \ge\eta_0\,\nu_{x,L}(I_{\mathrm{parent}}),
 \qquad \eta_0>0
\end{equation}
with $\eta_0$ independent of $x,L$, for at least one QNL-separated child at
each scale.  If \eqref{eq:uniform-child-mass} and a uniform UNI constant
$\kappa_0>0$ hold, Leclerc's tree argument gives a QNL exponent
$\Gamma>0$ independent of $L$; only the multiplicative constant is allowed
to grow like $e^{cL}$.

R\"uhr's pressure inequalities quantify weak convergence of Gibbs measures
on countable Markov shifts to equilibrium measures on finite subsystems by
the square root of the pressure gap.  They support the truncation step, but
do not prove \eqref{eq:uniform-child-mass} or uniform UNI.  Strong positive
recurrence/BIP controls the symbolic escape; the billiard geometry must still
provide the nonlinear separated child.

For the collision SRB measure on a compact homogeneous magnet, the unstable
conditional is better than a general fractal Gibbs conditional: in arclength
coordinates its density is bounded above and below.  In that case the
measure-theoretic part of the QNL argument is elementary.

\begin{lemma}[Uniform interval UNI implies power non-concentration]
\label{lem:interval-uni-qnl}
Let $\nu_L$ be measures on an interval whose densities satisfy
$0<c_-\le d\nu_L/dt\le c_+<\infty$, uniformly in $L$.  Let $X_L$ be a family
of functions with the following scale-invariant UNI property: for every
interval $J$ of length $\sigma<\sigma_0$ and every polynomial $P$ in a fixed
finite-dimensional class, there is a subinterval
$J'\subset J$ with $|J'|\ge\kappa|J|$ such that
\[
 |X_L(t)-P(t)|\ge\kappa|J|\quad(t\in J'),
\]
where $\kappa>0$ is independent of $J,P,L$.  Then for every fixed
$\alpha>0$ there are $C,\Gamma>0$, independent of $L$, such that
\[
 \nu_L\{t\in J:|X_L(t)-P(t)|\le |J|^{1+\alpha}\}
 \le C|J|^\Gamma\nu_L(J).
\]
\end{lemma}

\begin{proof}
Partition every surviving interval into a bounded number of comparable
children and discard a child contained in the UNI interval $J'$.  The density
bounds imply that at least a fixed fraction $\eta_0>0$ of the conditional
mass is discarded at each generation.  Iterate until the child scale is
comparable to $|J|^{1+\alpha}$.  The number of generations is
$c_\alpha|\log|J||+O(1)$, so the surviving mass is at most
$(1-\eta_0)^{c_\alpha|\log|J||}\nu_L(J)
\le C|J|^\Gamma\nu_L(J)$.
\end{proof}

Thus for the central homogeneous face the principal task can be reduced
further: establish a uniform geometric UNI statement for the
stable--unstable temporal-distance function.  Finite-subsystem approximation
is then needed for long singular itineraries and difference quotients, but
not to manufacture the QNL exponent on the central magnet.

\subsection{A fixed two-branch UNI certificate}

Bertozzi's 2026 countable-Markov exponential-mixing construction provides a
useful template: after sufficient inducing, two fixed inverse branches are
repeated and the derivative of their Birkhoff-roof difference is bounded away
from zero uniformly in the iteration depth.  The same device can be formulated
for the billiard temporal-distance increment.

Let $h_0,h_1$ be two regular inverse return branches on the central magnet,
with a common contraction bound $\lambda<1$, and let $\tau$ denote the
reduced stable--unstable phase increment (for the geometric application,
$\tau=\log J^uF$ after the stable coboundary normalization).  Define
\begin{equation}
\label{eq:two-branch-temporal-distance}
 \Psi_N(x)=\sum_{k=1}^N
 \bigl[\tau(h_0^kx)-\tau(h_1^kx)\bigr],
 \qquad
 \Psi_\infty=\lim_{N\to\infty}\Psi_N .
\end{equation}

\begin{proposition}[Anchored two-branch certificate]
\label{prop:two-branch-uni}
Suppose the series in \eqref{eq:two-branch-temporal-distance} converges in
$C^1$ and
\[
 \inf_{x\in I}|\Psi_\infty'(x)|\ge c_0>0
\]
on a central unstable interval $I$.  Then, for all sufficiently large $N$,
\[
 \inf_{x\in I}|\Psi_N'(x)|\ge c_0/2.
\]
If the two branches and their fixed connector branches are retained in every
anchored truncation, this UNI constant is independent of the truncation rank.
\end{proposition}

\begin{proof}
$C^1$ convergence makes
$\|\Psi_N'-\Psi_\infty'\|_\infty<c_0/2$ for all large $N$.  The second
assertion follows because the calculation uses only the anchored branches,
not the additional symbols admitted by a larger truncation.
\end{proof}

\begin{corollary}[First-term domination after acceleration]
\label{cor:first-term-uni}
Assume $\|Dh_i\|\le\lambda<1$, $\|D\tau\|_\infty\le M$, and
\[
 \inf_I|D(\tau\circ h_0-\tau\circ h_1)|\ge d_0\lambda .
\]
If
\[
 \frac{2M\lambda^2}{1-\lambda}\le\frac12d_0\lambda,
\]
then $\inf_I|\Psi_\infty'|\ge d_0\lambda/2$.
\end{corollary}

\begin{proof}
The derivative of the sum of all terms with $k\ge2$ is bounded by
$2M\sum_{k\ge2}\lambda^k=2M\lambda^2/(1-\lambda)$.  Subtract this from the
first-term lower bound.
\end{proof}

This changes the next billiard calculation from an infinite statistical
problem into a finite geometric one.  Choose two central return words whose
one-step radial tangency germs have different curvature coefficients in
\eqref{eq:circular-face-curvature}; attach fixed regular connectors to the
magnet; and compute $\Psi_\infty'$ (or a rigorously bounded finite
approximation plus its geometric tail).  A nonzero lower bound is open under
small table perturbations and gives the uniform UNI input for
Lemma~\ref{lem:interval-uni-qnl}.  The non-equality of the two one-step
curvatures is a promising seed, but it is not by itself the required lower
bound because the stable coboundary and all subsequent return terms enter
$\Psi_\infty$.
After accelerating both branches, Corollary~\ref{cor:first-term-uni} shows
that it is enough to retain a lower bound of order $\lambda$ for the
first-term curvature separation; the remaining infinite temporal-distance
tail is only $O(\lambda^2)$.  This remains a useful branchwise fallback.  The
next calculation gives a more direct fixed periodic witness for the smooth
finite truncations.

\subsection{An exact Anosov-cocycle witness from a normal two-cycle}

There is a more direct way to certify the geometric nonlinearity needed by
the smooth finite truncations.  Leclerc's criterion for a smooth
area-preserving Axiom--A surface map says that QNL follows if the stable or
unstable distribution is not $C^2$.  At a hyperbolic fixed point, this can be
checked by the Anosov cocycle.  In determinant-one adapted coordinates, write
the map as $(F,G)$ with derivative
$\operatorname{diag}(\lambda,\mu)$, where $\lambda>1$, $\mu=\lambda^{-1}$.
The relevant scalar is
\begin{equation}
\label{eq:anosov-cocycle}
 \mathfrak A(F,G)
 =\frac{G_{xyy}}{\mu}
 -\frac{F_{xy}G_{xy}}{\lambda-1}
 -\frac{G_{xx}F_{yy}}{1-\mu^3}
 -\frac{G_{xy}G_{yy}}{\mu^2},
\end{equation}
with all derivatives evaluated at the fixed point.

\begin{proposition}[Exact circular normal-cycle witness]
\label{prop:exact-circular-anosov}
Consider two unit circular scatterers whose facing boundary points are
separated by a unit free-flight gap.  Let $\mathcal R=T^2$ be the billiard
return map to the first component near the unobstructed normal period-two
orbit, in Birkhoff coordinates $(s,p)$ with $p=\sin\varphi$.  Then
\[
 D\mathcal R(0)=
 \begin{pmatrix}7&4\\12&7\end{pmatrix}.
\]
For the determinant-one eigenbasis
\[
 S=\begin{pmatrix}
       1&-\sqrt3/6\\[2pt]
       \sqrt3&1/2
    \end{pmatrix},
 \qquad \det S=1,
\]
write $S^{-1}\mathcal R S=(F,G)$.  Then
\begin{align*}
 \lambda&=7+4\sqrt3,&\mu&=7-4\sqrt3,\\
 F_{xy}=F_{yy}=G_{xx}=G_{xy}=G_{yy}&=0,
 &G_{xyy}&=-\mu.
\end{align*}
Consequently
\[
 \mathfrak A(F,G)=-1.
\]
In particular, the local return dynamics has a nonzero QNL witness which is
stable under sufficiently small clean perturbations of the table.
\end{proposition}

\begin{proof}
Place the circle centres at $(0,0)$ and $(3,0)$ and parametrize the first
boundary near the facing point by $q(s)=(\cos s,\sin s)$.  For outgoing
Birkhoff data $(s,p)$, set
\[
 n=(\cos s,\sin s),\qquad t=(-\sin s,\cos s),\qquad
 v=\sqrt{1-p^2}\,n+pt.
\]
For a ray based at $q$ and a target circle with centre $c$ and radius one,
put $w=q-c$, $b=w\cdot v$, and $c_0=|w|^2-1$.  The first collision time and
reflected velocity are
\[
 t_*=-b-\sqrt{b^2-c_0},\qquad
 v^+=v-2(v\cdot n_*)n_*,
 \quad n_*=q+t_*v-c.
\]
Apply this formula first to the second circle and then to the first.  Taylor
expansion through degree three gives the displayed linear matrix.  Reflection
in the line of centres conjugates $(s,p)$ to $(-s,-p)$ and commutes with the
two-collision return, so the return map is odd and all its quadratic jets
vanish.  Conjugating the cubic jet by $S$ gives
$G_{xyy}=-(7-4\sqrt3)$; substitution in
\eqref{eq:anosov-cocycle} gives $-1$.
\end{proof}

The calculation is scale free: if both radii and the free-flight gap are
multiplied by the same positive number, use dimensionless arclength $s/R$ and
obtain the same return jet.  Thus the local configuration can be inserted at
an arbitrarily small geometric scale into a finite-horizon table, provided
the normal two-cycle remains unobstructed.  More generally, on any fixed clean
finite itinerary the return jet and $\mathfrak A$ depend real-analytically on
the circular radii, centres, and collision points.  Proposition
\ref{prop:exact-circular-anosov} therefore implies that the nonzero-cocycle
tables form an open dense subset of every connected analytic parameter chart
containing this sample.

For an anchored finite horseshoe which contains this normal two-cycle and a
regular connector branch, the same fixed point occurs in every larger
anchored truncation.  Hence the geometric nonlinearity witness itself no
longer deteriorates with truncation rank.  What still requires proof is a
uniform smooth extension/Markov realization of those finite billiard cores,
uniform cutting mass for the SRB conditionals, and transfer of Leclerc's QNL
bound from the equilibrium state to the weighted double-centered face slice.
The exact cocycle closes the local geometry, not these singular-exhaustion
steps.

\subsection{Transfer of QNL to a central weighted face slice}

The equilibrium-state/face-current mismatch is harmless on one compact
homogeneous return rectangle.  Let $W^u$ be a reference unstable leaf and let
$H:\Gamma\to W^u$ be stable holonomy from a central smooth face arc.  Write
$m_\Gamma$ for arclength and let $w$ include the physical face velocity,
coarea factor, and the smooth branch Jacobians.

\begin{lemma}[Weighted face-slice transfer]
\label{lem:weighted-face-qnl-transfer}
Assume, uniformly on a compact clean family, that
\begin{enumerate}[label=\textup{(\roman*)}]
\item $H$ is a $C^{1+\alpha}$ diffeomorphism onto its image and
      $0<c_H\le JH\le C_H$;
\item the SRB conditional $\mu^u$ on $W^u$ has arclength density between
      $c_u>0$ and $C_u<\infty$;
\item $w\in C^\alpha(\Gamma)$ and $\|w\|_\infty\le C_w$.
\end{enumerate}
Define the signed face measure
\[
 \nu=H_*(w\,dm_\Gamma).
\]
If, on every interval $J\subset W^u$, the temporal-distance phase $X$
satisfies
\[
 \mu^u\{t\in J:|X(t)-P(t)|\le\epsilon\}
 \le C_0\epsilon^\Gamma\mu^u(J),
\]
then
\[
 |\nu|\{t\in J:|X(t)-P(t)|\le\epsilon\}
 \le C_1\epsilon^\Gamma |J|,
 \qquad
 C_1=C_0\frac{C_wC_u}{c_Hc_u},
\]
after increasing $C_1$ by a uniform chart constant.  The same conclusion
holds for a finite signed sum of physical face pieces, with $C_1$ multiplied
by the sum of their total-variation bounds.
\end{lemma}

\begin{proof}
The change-of-variables formula gives
\[
 \frac{d|\nu|}{dm_{W^u}}(H z)
 =\frac{|w(z)|}{JH(z)}\le\frac{C_w}{c_H}.
\]
The density bounds for $\mu^u$ compare arclength and $\mu^u$ in both
directions.  Apply the assumed small-value estimate and make those two
comparisons.  Total variation and the triangle inequality handle signed
pieces.
\end{proof}

Lemma~\ref{lem:circular-face-curvature} supplies the central coarea
transversality, while standard clean-rectangle holonomy estimates supply
(i)--(ii).  Thus, once the exact cocycle is placed in an anchored smooth
horseshoe, the central weighted face slice inherits the same QNL exponent.
For pieces whose return word or homogeneity rank tends to infinity, this
local transfer is combined below with the central marked moment and the
two-sided recovery sum; Proposition~\ref{prop:global-labelled-ledger} keeps
those constants on one common refinement.

\subsection{A countable-IFS theorem removes the return truncation}
\label{subsec:countable-ifs}

There is a more direct route than passing through smooth finite horseshoes.
Leclerc--Paukkonen--Sahlsten~\cite{LPS} prove power Fourier decay for Gibbs measures of
finite or countable $C^{1+\alpha}$ conformal IFSs, allowing overlaps.  In the
countable case their assumptions are uniform contraction, bounded distortion,
summable variations, and the light-tail condition
\begin{equation}
\label{eq:lps-light-tail}
 \sum_{a}\sup_{x\in I}
 e^{\phi(h_a x)}|h_a'(x)|^{-\gamma_0}<\infty
 \qquad\text{for some }\gamma_0>0.
\end{equation}
Their Theorem~1.1 says that QNL gives power Fourier decay, including for
H\"older amplitudes and nonstationary phases with a small polynomial loss in
their norms.  Sections~6.1--6.2 prove QNL by finite-difference cutting
arguments for finite MNL and for countable UNI systems with a light tail.  No
finite-alphabet exhaustion or smooth extension of the branch maps is used.

This theorem is unusually well matched to the present induced problem.  Let
$F:I\to I$ be the unstable quotient of a full-branch Gibbs--Markov return to
a central billiard magnet and let $\{h_a:I\to I\}_{a\in\mathcal A}$ be its
inverse branches.  After accelerating once if necessary, the desired
dictionary is
\begin{equation}
\label{eq:lps-dictionary}
 \begin{array}{c|c}
 \text{countable-IFS input}&\text{billiard input}\\ \hline
 h_a&\text{homogeneous inverse return branch}\\
 \phi&\text{normalized induced unstable-Jacobian potential}\\
 \text{uniform contraction}&\text{uniform hyperbolicity after inducing}\\
 \text{bounded distortion}&\text{homogeneous return distortion}\\
 \text{light tail}&\text{return tail plus a small grazing moment}\\
 \text{QNL}&\text{a fixed clean finite-difference gate}.
 \end{array}
\end{equation}

There is a small but important point in the first line of this dictionary.
The one-sided moving-face labels must not be inserted by cutting every return
branch into labelled subbranches: those pieces need not have full image and
would destroy the hypothesis to which the countable-IFS theorem is applied.
The correct object is a full-branch dynamical quotient with a marked amplitude
fiber.

For an interval amplitude use the affine scale-normalized norm
\begin{equation}
\label{eq:scaled-amplitude-norm}
 \|w\|_{\mathcal A^\alpha(J)}
 :=\|w\|_{\infty,J}+|J|^\alpha[w]_{C^\alpha(J)}.
\end{equation}
It is unchanged up to a fixed factor by affine normalization and cannot
increase under restriction.  Every connected component is charged
separately; hence endpoint and component complexity is visible in the sum
rather than hidden in a piecewise-H\"older notation.

\begin{proposition}[Marked full-branch quotient]
\label{prop:marked-full-branch-quotient}
Let $\Lambda$ be a Young return rectangle with product structure and let
$\{\Lambda_a\}_{a\in\mathcal A}$ be a return partition such that
$F|_{\Lambda_a}=T^{R_a}$ maps the stable subset $\Lambda_a$ across $\Lambda$
with the usual stable/unstable Markov property.  Fix a reference unstable
leaf $I$ and let $\pi^s:\Lambda\to I$ be stable holonomy.  Then
\[
 \bar F\circ\pi^s=\pi^s\circ F
\]
defines a one-state full-branch expanding map: for
$I_a=\pi^s(\Lambda_a)$, the map $\bar F:I_a\to I$ is onto and has an inverse
branch $h_a:I\to I_a$.

Suppose now that the lift of a complete physical face through $\Lambda_a$ is
refined into one-sided marked pieces $E_{a,\ell}$, after artificial-face
cancellation.  On the stable quotient its action can be written
\begin{equation}
\label{eq:marked-amplitude-disintegration}
 \nu(\varphi)=
 \sum_{a,\ell}\int_{J_{a,\ell}}
     \varphi(h_a x)w_{a,\ell}(x)\,d\mu^u(x),
\end{equation}
where the $J_{a,\ell}\subset I$ are amplitude pieces and the sign is included
in $w_{a,\ell}$.  Thus the quotient dynamical alphabet remains
$\{h_a\}_{a\in\mathcal A}$; the face label $\ell$ is a fiber mark and does
not have to define a full-image branch.  If
\[
 \sum_{a,\ell}\|w_{a,\ell}\|_{\mathcal A^\alpha(J_{a,\ell})}
      \mu^u(J_{a,\ell})<\infty,
\]
then the marked current is a summable piecewise-H\"older amplitude over the
same full-branch Gibbs system.
\end{proposition}

\begin{proof}
The Markov property says that every unstable leaf in $\Lambda_a$ is mapped
across every stable leaf of $\Lambda$.  Quotienting stable leaves therefore
makes $\bar F|_{I_a}$ onto $I$; contraction of inverse branches and bounded
distortion are the usual quotient estimates.  Disintegrate the coarea face
measure first over return atoms and then over its one-sided physical marks.
Stable holonomy transports each piece to the reference leaf and gives
\eqref{eq:marked-amplitude-disintegration}.  No restriction
$h_a|_{J_{a,\ell}}$ is declared to be a dynamical branch.  The last assertion
is the triangle inequality for the H\"older-amplitude norm.  Matched
artificial cuts contribute opposite weights and disappear before this
disintegration.
\end{proof}

Consequently the standard Young quotient and the moving-face bookkeeping are
compatible provided the marked-amplitude sum is finite.  The construction of
the quotient is standard.  Proposition~\ref{prop:central-marked-moment}
supplies this summability on the central return alphabet, while the
bi-recovery refinement supplies the grazing tail.

The SRB potential part of this audit is now supported by a direct billiard
result.  Lemmas~5.2--5.3 of Lima--Obata--Poletti~\cite{LOP} prove that the
collision Lebesgue/SRB measure is an equilibrium state of the geometric
potential
\[
 \phi=-\log\|DT|_{E^u}\|,
 \qquad P(\phi)=0,
\]
and that $\phi$ becomes H\"older on their countable Markov coding for
finite-horizon dispersing billiards.  Hence the induced potential has
summable variations.  This does not by itself give the one-state return tail,
but it removes the former ambiguity between an MME coding and the required
collision-SRB Gibbs weights.

\begin{proposition}[Central amplitude moment]
\label{prop:central-marked-moment}
Assume the Young return scheme has
\[
 \sum_a\sup_I p_a\,e^{\epsilon R_a}<\infty
\]
and restrict first to a compact regular homogeneity core.  For a complete
first moving-face defect, suppose every collision insertion has uniformly
bounded $C^\alpha$ coarea and stable-holonomy factors.  Then the marked
amplitude disintegration \eqref{eq:marked-amplitude-disintegration} satisfies
\[
 \sum_{a,\ell}e^{\epsilon'R_a}
   \|w_{a,\ell}\|_{\mathcal A^\alpha(J_{a,\ell})}
   \mu^u(J_{a,\ell})<\infty
 \qquad(0<\epsilon'<\epsilon).
\]
After high-homogeneity insertions are assigned to the bi-recovery blocks of
Corollary~\ref{cor:radial-bi-recovery}, the same conclusion holds for the
entire assembled first defect.
\end{proposition}

\begin{proof}
A return word of length $R_a$ has at most $R_a$ possible physical insertion
times.  On the compact regular core each insertion contributes at most a
uniform constant times the Gibbs branch weight.  Hence the left-hand side is
bounded by
\[
 C\sum_a R_a e^{\epsilon'R_a}\sup_I p_a,
\]
which is finite because
$R e^{\epsilon'R}\le C_{\epsilon,\epsilon'}e^{\epsilon R}$.
Artificial insertions cancel before taking total variation.  The removed
high-homogeneity pieces satisfy the common two-sided recovery moment and are
summed by Proposition~\ref{prop:recovery-window-summability}.
\end{proof}

The following consequence records the exact countable-IFS interface rather
than hiding it in a finite-horseshoe approximation.

\begin{proposition}[Finite-difference countable-IFS bridge]
\label{prop:countable-ifs-bridge}
Let the central unstable quotient be a one-state full-branch
$C^{1+\alpha}$ IFS with uniform inverse contraction, bounded distortion, a
normalized Gibbs potential of summable variations, and
\eqref{eq:lps-light-tail}.  For words $\mathbf a,\mathbf b$ of length $N$
and $x,y\in I$, put
\[
 \Delta_N((\mathbf a,x),(\mathbf b,y))
 =\{\log|h_{\mathbf a}'(x)|-\log|h_{\mathbf b}'(x)|\}
  -\{\log|h_{\mathbf a}'(y)|-\log|h_{\mathbf b}'(y)|\}.
\]
Assume the LPS quantitative nonlinearity estimate: for some
$C>0$, $\Theta>0$, and $\rho\in(0,1)$,
\begin{equation}
\label{eq:finite-difference-qnl}
 \sum_{\mathbf a,\mathbf b,\mathbf c,\mathbf d\in\mathcal A^N}
 p_{\mathbf a}p_{\mathbf b}p_{\mathbf c}p_{\mathbf d}
 \one_{J}\!\left(
 \Delta_N((\mathbf a,x_{\mathbf c}),
          (\mathbf b,x_{\mathbf d}))\right)
 \le C\{|J|^\Theta+\rho^N\}
\end{equation}
for every interval $J\subset\mathbb R$; Gibbs weights may be evaluated at
arbitrary reference points in their cylinders.  Then the induced SRB Gibbs
measure has the weighted nonstationary-phase Fourier estimate of
Leclerc--Paukkonen--Sahlsten for every phase with the stated
$C^\alpha$ log-derivative bounds.  The conclusion remains valid for the
marked physical face amplitudes of
Lemma~\ref{lem:weighted-face-qnl-transfer}.
\end{proposition}

\begin{proof}
Equation~\eqref{eq:finite-difference-qnl} is precisely the QNL hypothesis in
Theorem~1.1 of \cite{LPS}.  That theorem is formulated for finite or
countable $C^{1+\alpha}$ systems and H\"older amplitudes; it neither
differentiates $\log h_a'$ nor assumes a $C^2$ holonomy.  The light-tail
moment controls the countable alphabet.  The weighted-face conclusion follows
from Lemma~\ref{lem:weighted-face-qnl-transfer} and the summable
scale-normalized amplitude norm \eqref{eq:scaled-amplitude-norm}.
\end{proof}

The former $C^2$ multiplier/UNI route is deliberately not used.  For a
billiard quotient the stable holonomy is generally only $C^{1+\alpha}$, so
the derivative of $\log h_a'$ need not exist.  The correct gate input is the
finite-difference estimate \eqref{eq:finite-difference-qnl}.

\begin{proposition}[Clean homoclinic horseshoe supplies a finite QNL gate]
\label{prop:homoclinic-horseshoe-uni}
Let $p$ be the regular fixed point of the squared billiard map associated with
the normal period-two orbit in Proposition
\ref{prop:exact-circular-anosov}.  If its Anosov cocycle satisfies
$\mathfrak A(p)\ne0$, then there is a compact finite clean horseshoe
$\Omega_p$, contained in regular branches and containing $p$, such that in at
least one time orientation its temporal distance satisfies QNL for every
H\"older equilibrium state.  Thus $\Omega_p$ supplies a finite topological
      QNL seed; Proposition~\ref{prop:context-polynomial-blowup} transfers its
      polynomial escape to the physical collision-SRB bracket.
\end{proposition}

\begin{proof}
Lemma~5.1 of Lima--Obata--Poletti~\cite{LOP} says that finite-horizon
dispersing billiards have a single homoclinic class at regular hyperbolic
points.  Choose a second regular hyperbolic periodic point $q$.  The theorem
gives clean transverse intersections
$W^u(p)\pitchfork W^s(q)$ and $W^u(q)\pitchfork W^s(p)$.  The inclination
lemma then gives a nontrivial clean transverse homoclinic point of $p$; this
avoids the vacuous substitution $x=y=p$.  The Smale--Birkhoff construction, using
a finite segment of this clean homoclinic orbit, gives a compact finite
horseshoe $\Omega_p$ on which the billiard map is smooth.  In symplectic
Birkhoff coordinates the map is area preserving; the local dynamics near
$\Omega_p$ may be viewed as a smooth Axiom--A basic set.

By the contrapositive of the Anosov cocycle obstruction, together with
Theorem~4.3 of Leclerc~\cite{Leclerc2026}, in one of the two time orientations
the temporal distance on $\Omega_p$ satisfies QNL for every H\"older
equilibrium state, in particular for the equilibrium state of the geometric
potential restricted to this finite subsystem.  No comparison of that
horseshoe measure with collision SRB is made here.
\end{proof}

There is a regularity caveat.  The stable quotient of a smooth horseshoe is in
general only $C^{1+\alpha}$, precisely because the stable holonomy need not be
$C^2$.  Therefore the finite QNL conclusion above cannot be fed directly
into the $C^2$ non-linearizability criterion of
Algom--Rodriguez Hertz--Wang.  The route below instead propagates the
physical gate's finite-difference gap to the countable
$C^{1+\alpha}$ inducing scheme.  The LPS theorem is recorded as a background
alternative; it is not invoked in the principal fractional-phase proof.

The light-tail hypothesis has a concrete billiard reduction.  If a branch
$a$ has return time $R_a$ and grazing rank $k_a$, bounded distortion gives,
schematically,
\[
 |h_a'|^{-\gamma_0}\le
 C e^{c\gamma_0R_a}(1+k_a)^{q\gamma_0}.
\]
Thus an exponential return moment and any strictly positive polynomial SRB
tail in $k_a$ imply \eqref{eq:lps-light-tail} after choosing
$\gamma_0>0$ sufficiently small.  This is strictly weaker than uniform
control of every rank-$L$ smooth extension.

\begin{lemma}[Return and grazing moments imply the LPS light tail]
\label{lem:billiard-light-tail}
Assume the normalized branch weights $p_a(x)=e^{\phi(h_ax)}$ satisfy
\[
 \sum_a\sup_x p_a(x)e^{\epsilon R_a}(1+k_a)^\delta<\infty
\]
for some $\epsilon,\delta>0$, and suppose
$|h_a'|^{-1}\le Ce^{cR_a}(1+k_a)^q$.
Then \eqref{eq:lps-light-tail} holds for every
\[
 0<\gamma_0<\min\{\epsilon/c,\delta/q\}.
\]
\end{lemma}

\begin{proof}
Raise the inverse-derivative estimate to $\gamma_0$ and absorb the two
factors into the assumed exponential and polynomial moment.
\end{proof}

There is also a pressure-based way to obtain the required moment.  The
geometric-potential theory for finite-horizon dispersing billiards has a
spectral gap for $-t\log J^uT$ on every compact subinterval of
$(0,t_*)$, where $t_*>1$, and its modified one-step expansion is uniform for
$t\ge t_0>0$~\cite{DLreview}.  Since the SRB value is $t=1$, the LPS factor
$|h_a'|^{-\gamma_0}$ merely changes the geometric exponent from $1$ to
$1-\gamma_0$.

\begin{proposition}[Geometric-pressure route to the light tail]
\label{prop:pressure-light-tail}
Let a homogeneous full-branch inducing scheme have an exponential return
tail at the SRB parameter.  Assume its induced distortion is uniform and
that the geometric transfer estimates hold for
$t\in[1-\varepsilon,1]$.  Then, after decreasing
$\varepsilon>0$ if necessary, its normalized induced SRB potential satisfies
\eqref{eq:lps-light-tail} for every sufficiently small
$\gamma_0>0$.
\end{proposition}

\begin{proof}
Group branches by return time $R_a=n$.  Up to the uniform distortion
coboundary, the summand in \eqref{eq:lps-light-tail} is the inverse unstable
Jacobian to the power $1-\gamma_0$.  The geometric transfer estimate at
$t=1-p\gamma_0$ bounds the corresponding $p$-moment by
$Ce^{c(p\gamma_0)n}$, with $c(p\gamma_0)\to0$ as
$\gamma_0\to0$.  H\"older's inequality combines this with the SRB return
tail $Ce^{-an}$ to give
\[
 \sum_{R_a=n}\sup_x
 e^{\phi(h_ax)}|h_a'(x)|^{-\gamma_0}
 \le C e^{-a n/q+c(p\gamma_0)n/p}.
\]
Choose $p,q>1$ conjugate and then $\gamma_0$ so that the exponent is
negative; summing over $n$ proves the claim.  Homogeneity-strip subdivision
is already included in the geometric one-step expansion.
\end{proof}

\begin{lemma}[Fixed-gate hitting and two-context moments]
\label{lem:fixed-gate-context-moment}
Let $F$ be a one-state full-branch Gibbs--Markov map and let $\mathcal G$ be a
fixed finite gate cylinder of return depth $L$.  Suppose bounded distortion
gives
\begin{equation}
\label{eq:fixed-gate-probability}
 \inf_C\mu^u(\mathcal G\mid C)\ge\eta_G>0
\end{equation}
for every past cylinder $C$, and suppose
\begin{equation}
\label{eq:return-exponential-moment}
 \mathbb E e^{\epsilon R}<\infty
\end{equation}
for some $\epsilon>0$.  If $N_G$
is the first gate-hitting time in $L$-blocks and $C^+$ is the total collision
length of the future entrance context, then there is $\epsilon_G>0$ such that
\[
 \mu^u(N_G>k)\le(1-\eta_G)^k,
 \qquad
 \mathbb E e^{\epsilon_G C^+}<\infty.
\]
The reverse quotient gives the same estimate for $C^-$.  On the natural
extension no independence is needed:
\begin{equation}
\label{eq:two-context-moment}
 \mathbb E e^{\epsilon_G'(C^-+C^+)}<\infty
\end{equation}
for every sufficiently small $\epsilon_G'>0$.
\end{lemma}

\begin{proof}
Condition \eqref{eq:fixed-gate-probability} and the Markov property give
$\mu^u(N_G>k+1\mid N_G>k)\le1-\eta_G$, proving the geometric tail.  Let
$K(\theta)$ be the uniform conditional exponential moment of one return;
$K(0)=1$ and $K(\theta)<\infty$ for small $\theta>0$.  Bounded Gibbs
distortion gives
\[
 \mathbb E\!\left[e^{\theta C^+};N_G=k\right]
 \le C\{(1-\eta_G)K(\theta)^L\}^{k-1}K(\theta)^L.
\]
Choose $\theta$ so that the quantity in braces is smaller than one and sum in
$k$.  Time reversal proves the past estimate.  Finally
\[
 \mathbb P(C^-+C^+>t)
 \le\mathbb P(C^->t/2)+\mathbb P(C^+>t/2),
\]
so the two marginal exponential tails imply
\eqref{eq:two-context-moment} without independence.
\end{proof}

For the marked quotient, choose one fixed finite gate word.  The Gibbs ratio
of that word after an arbitrary prefix differs from its unconditional weight
by at most the uniform distortion factor.  Hence
\eqref{eq:fixed-gate-probability} holds with
$\eta_G=C_D^{-1}\mu^u(\mathcal G)>0$; it is not an additional mixing
assumption.

\begin{lemma}[Joint return--maximal-grazing moment]
\label{lem:homogeneity-grazing-moment}
Let $Y$ be the base of an exponential Young return, $R$ its return time, and
let $B(x)$ be the largest homogeneity rank visited before $R(x)$.  In
$(r,p)=(r,\sin\varphi)$ coordinates suppose
\[
 \mu\{1-|p|\lesssim2^{-2b}\}\le C2^{-2b},
 \qquad \mu_Y(R>L)\le Ce^{-cL}.
\]
Then $B$ has an exponential tail.  More precisely, for every fixed $A>0$,
\begin{equation}
\label{eq:maximal-grazing-tail}
 \mu_Y(B\ge b)
 \le Ce^{-cAb}+CAb\,2^{-2b}.
\end{equation}
Consequently, for sufficiently small $\epsilon,\kappa>0$,
\begin{equation}
\label{eq:small-grazing-moment}
 \mathbb E_Y e^{\epsilon R+\kappa B}<\infty.
\end{equation}
If the inverse-Jacobian and coarea loss of a return word is bounded by
$Ce^{c_0R+qB}$, then \eqref{eq:lps-light-tail} and the marked-amplitude
moment hold after decreasing the exponents.  No independence of $R$ and $B$
is required.
\end{lemma}

\begin{proof}
On $\{R\le Ab\}$ take a union bound over the at most $Ab$ collision times.
Since $\mu_Y(E)\le\mu(E)/\mu(Y)$ for every event $E$ and the collision SRB
measure is invariant, each visit to a rank-$b$ boundary layer costs at most
$C2^{-2b}$.  Adding $\mu_Y(R>Ab)$ proves
\eqref{eq:maximal-grazing-tail}.  Choose $A$ so that both terms decay
exponentially in $b$; hence $\mathbb E e^{q\kappa B}<\infty$ for some
$q>1$ and small $\kappa$.  For conjugate $p,q>1$, H\"older gives
\[
 \mathbb E e^{\epsilon R+\kappa B}
 \le (\mathbb E e^{p\epsilon R})^{1/p}
      (\mathbb E e^{q\kappa B})^{1/q}<\infty.
\]
Absorb $c_0R+qB$ with smaller exponents and apply
Lemma~\ref{lem:billiard-light-tail}.  The same argument applies to the
scale-normalized coarea amplitude.
\end{proof}

Lemma~\ref{lem:homogeneity-grazing-moment} verifies the grazing factor from
the collision density and strip width, while
Proposition~\ref{prop:central-marked-moment} supplies the moving-face
amplitude budget on the same quotient.  Lemma
\ref{lem:fixed-gate-context-moment} supplies the entrance/exit moment.  Thus
the particular full-branch quotient uses one joint return/grazing/context
ledger rather than an unverified appeal to the literature theorem.

This route also exposes a flaw in the finite-horseshoe detour.  A conservative
pasting lemma~\cite{Teixeira} can locally modify or smooth a volume-preserving map, but it
does not automatically extend a finite family of billiard branches to a
global area-preserving Axiom--A diffeomorphism whose equilibrium state has the
required collision-SRB weights.  Even a successful extension can put a
fractal Gibbs conditional on a saddle basic set in place of the absolutely
continuous unstable conditional of the billiard SRB measure.  The direct
countable-IFS quotient preserves the correct measure and should therefore be
the primary route.

Thus light-tail is not expected to be a new principal theorem once the
standard homogeneous inducing scheme is put in the one-state form required
above.  The exact value $\mathfrak A=-1$, the clean homoclinic theorem, and
Proposition~\ref{prop:homoclinic-horseshoe-uni} produce the clean polynomial
escape seed without naming a connector.  Proposition
\ref{prop:physical-legal-child} records the local physical slope transfer, and
the cohomological bracket below identifies its physical collision-SRB object.

\subsection{The cohomological bracket and polynomial context blow-ups}

The quotient regularity is only $C^{1+\alpha}$.  The useful exact object is
therefore a four-point difference of Birkhoff sums, not a derivative of a
holonomy.  Let $f$ denote the billiard map restricted to the finite clean
basic set, let $F$ be its stable factor, and put
$\tau_f=\log J^uf$ and $\tau_F=\log|F'|$.  As in
\cite[Lemma~3.15]{Leclerc2026}, choose a H\"older function $\theta$ such that
\begin{equation}
\label{eq:potential-cohomology}
 \tau_f=\tau_F\circ\pi+\theta\circ f-\theta .
\end{equation}
For $p,q$ in one Markov rectangle define
\begin{equation}
\label{eq:physical-temporal-bracket}
 \Delta(p,q)=\sum_{j\in\mathbb Z}
 \{\tau_f(f^jp)-\tau_f(f^j[p,q])
   -\tau_f(f^j[q,p])+\tau_f(f^jq)\}.
\end{equation}

\begin{lemma}[Cohomological bracket identification]
\label{lem:cohomological-bracket-identification}
Let $\mathbf a,\mathbf b$ be the length-$N$ past words of $p,q$ and let
$x=\pi p$, $y=\pi q$.  With inverse branches $g_{\mathbf a},g_{\mathbf b}$
of $F$, set
\begin{align}
 \mathcal D_N(\mathbf a,\mathbf b;x,y)
  :={}&S_N\tau_F(g_{\mathbf a}x)-S_N\tau_F(g_{\mathbf a}y)\notag\\
      &-S_N\tau_F(g_{\mathbf b}x)+S_N\tau_F(g_{\mathbf b}y).
\label{eq:factor-four-bracket}
\end{align}
Then
\begin{equation}
\label{eq:bracket-tail}
 |\mathcal D_N(\mathbf a,\mathbf b;x,y)-\Delta(p,q)|
 \le C\kappa^{\alpha N}.
\end{equation}
 The constants are uniform on the persistent finite basic set and for
 $|s|<s_0$.  Entrance and exit contexts are not included in this statement;
 their finite jets and holonomy tails are treated explicitly in Proposition
 \ref{prop:context-polynomial-blowup}.
\end{lemma}

\begin{proof}
Insert \eqref{eq:potential-cohomology} in
\eqref{eq:physical-temporal-bracket}.  The four $\theta$-terms telescope.
For $j\ge0$, stable projection gives
$\pi(f^jp)=\pi(f^j[p,q])$ and
$\pi(f^jq)=\pi(f^j[q,p])$, so the positive-time factor terms cancel
exactly.  Writing the negative iterates through the past inverse words gives
\eqref{eq:factor-four-bracket}.  The omitted tail starts at $N+1$; H\"older
regularity of $\tau_F$ and inverse contraction give
\[
 \sum_{j>N}C\kappa^{\alpha j}\le C'\kappa^{\alpha N}.
\]
 This is precisely the proof of \cite[Lemma~3.15]{Leclerc2026} on the clean
 basic set.  Hyperbolic continuation makes the constants uniform in $s$.
\end{proof}

The all-scale escape in
\cite[Propositions~4.31--4.32]{Leclerc2026} supplies an equilibrium/Cantor
cutting template only.  The actual collision--SRB transfer at the prescribed
depth remains the independent interface $\mathrm{PPE}_{N}$ below; it is not a
consequence of the quoted propositions.

The preceding bracket must still be matched to the physical endpoint maps.
For a stable projection $\pi_s$ on a clean rectangle, the stable holonomy
Jacobian satisfies the convergent identity
\begin{equation}
\label{eq:exact-holonomy-series}
 \log J\pi_s(r)=\sum_{j\ge0}
 \{\tau_f(f^jr)-\tau_f(f^j\pi_s r)\}.
\end{equation}
This is the Jacobian formula of \cite[Lemma~C.1]{Leclerc2026}.  In particular,
entrance and exit holonomies are not discarded as unspecified coboundaries:
their four endpoint series are inserted into \eqref{eq:physical-temporal-bracket};
the common terms cancel and the unpaired tail is bounded by
$C\kappa^{\alpha N}$.

\begin{proposition}[Context--polynomial blow-up]
\label{prop:context-polynomial-blowup}
Assume $\mathfrak A(p)=-1$ and let $I_c=t_c+\ell_c[-1,1]$ be any homogeneous
resolving cylinder whose central excursion enters the persistent clean basic
set.  For the full reduced-slope bracket $S_c$, including the two endpoint
holonomy Jacobians, let $N$ be the total central resolving depth and let
$L_c=|c^-|+|c^+|$ be the entrance/exit length.  There are a fixed degree $d$,
a rough template $\mathcal X_{\omega_c}$ from the clean basic set, a
polynomial $P_c$ of degree at most $d$, and a context cost $H_c\ge1$ such that
\begin{equation}
\label{eq:context-polynomial-blowup}
 \ell_c^{-1}\{S_c(t_c+\ell_c u)-S_c(t_c)\}
 =\mathcal X_{\omega_c}(u)+P_c(u)+E_c(u),\qquad
 \|E_c\|_\infty\le C\left(H_c\ell_c^\alpha+
       \frac{\kappa^{\alpha N}}{\ell_c}\right).
\end{equation}
Moreover $H_c\le Ce^{aL_c}$, and for some $p>1$ the context ledger gives
\begin{equation}
\label{eq:context-jet-moment}
 \sup_{|s|<s_0}\sum_c w_c H_c^p<\infty .
\end{equation}
The degree and the escape constants of $\mathcal X_{\omega_c}$ modulo
$\mathbb R_d[X]$ are independent of $c$ and $s$.  In particular, if
\begin{equation}
\label{eq:resolving-depth-scale}
 N\ge q_*|\log\ell_c|,
 \qquad q_*>\frac{1+\alpha}{\alpha|\log\kappa|},
\end{equation}
then the second error in \eqref{eq:context-polynomial-blowup} is
$O(\ell_c^\alpha)$.
\end{proposition}

\begin{proof}
Use \eqref{eq:exact-holonomy-series} for both endpoint maps and split each
series at the entrance and exit times.  The four common tails cancel in the
bracket.  Each remaining finite clean context sum is $C^{1+\alpha}$ on its
homogeneous component.  Its affine Taylor polynomial is linear (and is
regarded as an element of $\mathbb R_d[X]$); the normalized remainder is
bounded by its $C^{1+\alpha}$ norm times $\ell_c^\alpha$.  The chain rule and
uniform one-step distortion give $H_c\le Ce^{aL_c}$.  Choose $p>1$ so that
$pa$ is smaller than the exponent in \eqref{eq:two-context-moment}; this proves
\eqref{eq:context-jet-moment}.  Lemma
\ref{lem:cohomological-bracket-identification} identifies the central part
with the factor four-point cocycle with error $O(\kappa^{\alpha N})$.
Condition \eqref{eq:resolving-depth-scale}, rather than an unrecorded fixed
acceleration, makes its normalized value $O(\ell_c^\alpha)$.  Finally
\cite[Propositions~4.31--4.32]{Leclerc2026}, activated by
$\mathfrak A(p)\ne0$, gives uniform all-location, all-scale escape of the rough
template from every polynomial of degree at most $d$.  No fixed acceleration
is used.
\end{proof}

\subsection{Native-coordinate single-source slope energy}

The following two conditions isolate the genuinely new clean-gate geometry
on the \emph{actual} collision--SRB interval.  They are not inferred from a
smooth extension of a fractal horseshoe conditional.

\begin{definition}[Single-source bracket condition $\mathrm{SS}$]
\label{def:single-source-condition}
Fix one orientation and condition on the opposite symbolic record, the value
$t'$, and the two finite contexts.  For a native interval $Q$ write
$z=\Phi_{\beta,Q}(t)$ and
$\sigma_Q=\operatorname{diam}\Phi_{\beta,Q}(Q)$.  Put
$d_Z=\max\{1,d\}$, where $d$ is
the fixed degree in Proposition~\ref{prop:context-polynomial-blowup}.
Condition $\mathrm{SS}(d_Z)$ means that
\begin{equation}
\label{eq:exact-reduced-slope}
 S_\beta(t)-S_\beta(t')
 =\ell_{\beta,Q}\{X^{[N]}_{\beta,Q}(z)-P_{\beta,Q,t'}(z)\}
  +\mathcal R_{\beta,Q,t'}(z),
\end{equation}
where $P_{\beta,Q,t'}\in\mathbb R_{d_Z}[z]$ and
$X^{[N]}_{\beta,Q}$ is the finite endpoint series on the actual depth-$N$
collision cylinder.  On the Cantor core it is within $C\kappa^N$ of one
native Leclerc rough template, and
\begin{equation}
\label{eq:single-source-coefficient}
 0<c_\ell\le|\ell_{\beta,Q}|\le C_\ell,
 \qquad
 \|\mathcal R_{\beta,Q,t'}\|_\infty
 \le C\{H_\beta\sigma_Q^{1+\alpha_B}+\kappa^N\}.
\end{equation}
The endpoint identity proving that every other $t$-dependent rough term is
absent or belongs to the fixed-degree polynomial is part of the condition;
freezing a symbolic name alone is not enough.
\end{definition}

\begin{definition}[Actual endpoint--slope bridge
$\mathrm{ACTUAL\_SLOPE\_BRIDGE}$]
\label{def:actual-slope-bridge}
On every labelled carrier to which a balanced phase estimate is applied,
let $u$ be the very same carrier coordinate used by
$\mathrm{BR}_{\rm sec}$, and let
$\mathsf X_{\lambda,m,s}(u)$ and
$\mathsf Y_{\lambda,n,s}(u)$ be its two actual scalar endpoint maps.  The
interface requires their derivatives to be nonzero on the labelled
homogeneous component and defines, rather than merely names,
\begin{equation}
\label{eq:actual-derivative-ratio-slope}
 \varepsilon_{{\rm or},\lambda,m,n,s}
 :=\operatorname{sgn}\!\left(
   \frac{\partial_u\mathsf Y_{\lambda,n,s}}
        {\partial_u\mathsf X_{\lambda,m,s}}\right),
 \qquad
 S_{\lambda,m,n,s}(u)
 :=\log\left|
   \frac{\partial_u\mathsf Y_{\lambda,n,s}}
        {\partial_u\mathsf X_{\lambda,m,s}}\right|.
\end{equation}
After the explicitly recorded carrier coordinate changes, this function is
required to be simultaneously the $S_\beta$ in
Definition~\ref{def:single-source-condition}, the $S_{\lambda,a}$ in every
source-clock row, and the $S^0_\eta$ in the unselected reference pair law.
The branch, orientation, time, and carrier labels in these four occurrences
must agree; equality of their distributions is not enough.  In particular,
for the actual endpoint phase
\begin{equation}
\label{eq:actual-balanced-phase-identity}
 \psi_\xi(u)=\xi_1\mathsf X_{\lambda,m,s}(u)
             +\xi_2\mathsf Y_{\lambda,n,s}(u)
\end{equation}
one has the exact typed identity
\begin{equation}
\label{eq:actual-balanced-phase-derivative}
 \partial_u\psi_\xi
 =\partial_u\mathsf X_{\lambda,m,s}
  \bigl(\xi_1+\varepsilon_{{\rm or},\lambda,m,n,s}
                   \xi_2e^{S_{\lambda,m,n,s}}\bigr).
\end{equation}
Condition $\mathrm{SS}(d_Z)$ is imposed on this exact derivative-ratio
slope.  An auxiliary cocycle or an abstract function satisfying
\eqref{eq:exact-reduced-slope} does not satisfy this interface unless the
displayed pointwise identification has been proved.
\end{definition}

\begin{definition}[Actual native stopping antichain
$\mathrm{NST}_{\rm phys}$]
\label{def:actual-native-stopping}
For $0<r<r_0$ put
\(
 \sigma(r)=(r/c_\ell)^{1/(1+\alpha_L)}.
\)
On each actual collision--SRB unstable conditional, follow homogeneous
full-branch descendants and stop at the first descendant $Q$ satisfying
\begin{equation}
\label{eq:native-first-stopping-rule}
 \operatorname{diam}_u(Q)\le\sigma(r)
 <\operatorname{diam}_u(\operatorname{parent}Q).
\end{equation}
Condition $\mathrm{NST}_{\rm phys}$ requires:
\begin{enumerate}[label=\textup{(NST\arabic*)}]
\item the stopped descendants form a measurable antichain with a normalized
      positive refinement kernel; their union carries all actual conditional
      mass except a grazing/return set of mass at most $Cr^{\Theta_{\rm st}}$;
\item for fixed $0<c_{\rm st}<C_{\rm st}<\infty$,
      \begin{equation}
      \label{eq:native-stopping-comparability}
       c_{\rm st}\sigma(r)\le\sigma_Q
       \le C_{\rm st}\sigma(r)
      \end{equation}
      on every stopped interval, with bounded overlap after connected
      component indexing;
\item the homogeneous buffer and every retained continuation are
      mass-preserving, and their total depth is $O(|\log r|)$; along
      $a(\theta s)$ they belong to $\mathrm{UCR}_{\rm sh}$ whenever this
      depth is at most $L(s)$.  There are fixed
      $c_{\rm buf},c_{\rm gate},c_{\rm cont}>0$ such that the discarded
      buffer, gate, and exceptional-continuation masses at resolving depth
      $N$ are bounded respectively by
      $Ce^{-c_{\rm buf}N}$, $Ce^{-c_{\rm gate}N}$, and
      $Ce^{-c_{\rm cont}N}$ under the unweighted complete positive tree
      $\Pi_{s,N}$.  Each discarded mass is retained on its named cemetery
      label.  These raw tails are used in the phase proof only after
      $\mathrm{TENERGY}_{W}$ has placed them under $\Pi_{s,N}^{W}$;
\item the context variable $H_\beta$, the physical Mark--$Z$ factors,
      dominant inverse scale, and $\mathrm{PPE}_{N}$ marks have their stated
      conditional moments under this one complete tree law, including the
      measurable last-mark extension on cemetery labels.  In
      particular, for some $p>1$,
      \begin{equation}
      \label{eq:stopped-H-moment}
       \sup_{r,\beta}
       \mathbb E_{\rm stop}[H_\beta^p\mid\mathcal F_{\rm parent}]<\infty .
      \end{equation}
\end{enumerate}
The scale comparison and the complete-tree moments are part of one interface;
moments from an unrelated native partition are not substituted into the
energy proof.  Leclerc's first-stopping/Vitali construction
\cite[Section~4]{Leclerc2026} gives the comparable-scale model on Markov
rectangles.  Its transfer to actual
collision--SRB conditionals across homogeneity cuts is not automatic.
\end{definition}

\begin{definition}[Prescribed-depth actual physical escape
$\mathrm{PPE}_{N}$]
\label{def:physical-legal-extension}
Assume $\mathrm{NST}_{\rm phys}$ and let $\mu^u_{\beta,Q}$ be the
normalized collision--SRB conditional on one of its stopped native intervals
$Q$.  For an energy scale $0<r<r_0$ set
\begin{equation}
\label{eq:prescribed-energy-depth}
 \sigma(r)=(r/c_\ell)^{1/(1+\alpha_L)},\qquad
 N(r)=\lceil C_{\rm buf}\log(1/r)\rceil ,
 \qquad \kappa^{N(r)}\le r/100 .
\end{equation}
The family satisfies $\mathrm{PPE}_{N}(d_Z)$ if there are
$\Theta_{\rm PPE}>0$ and $C<\infty$ such that the following hold on every
stopped actual conditional and every frozen context used in
$\mathrm{SS}(d_Z)$, all represented in the retained part of the complete
tree $\Pi_{s,N}$ from $\mathrm{NST}_{\rm phys}$.
\begin{enumerate}[label=\textup{(PPE\arabic*)}]
\item For every $P\in\mathbb R_{d_Z}[z]$ the endpoint series at the
      \emph{prescribed depth in \eqref{eq:prescribed-energy-depth}} obeys
      \begin{equation}
      \label{eq:ppe-prescribed-sublevel}
       \mu^u_{\beta,Q}
       \{|X^{[N(r)]}_{\beta,Q}-P|\le C r\}
       \le C r^{\Theta_{\rm PPE}} .
      \end{equation}
      This is an estimate on the full absolutely continuous collision--SRB
      conditional stopped at the native scale, not on a surviving Cantor
      conditional.
\item Multiplication by a retained positive amplitude preserves
      \eqref{eq:ppe-prescribed-sublevel} with its local PPE mark.
      Endpoint/Jacobian derivatives and the physical Lipschitz cost at depth
      $N(r)$ are bounded by the marks already charged in
      \eqref{eq:event-physical-norm}.  This local clause does not control
      context, stopping, buffer, gate, or continuation failures and does not
      by itself imply the full-$\mathscr W_*$ estimate
      \eqref{eq:tilted-PPE-tail}.
\item The depth-$N(r)$ endpoint is defined on retained open physical
      descendants, the resolving buffer has length $O(N(r))$, bounded
      overlap, and retains all continuations.  It is the same buffer as in
      $\mathrm{NST}_{\rm phys}$.  No affine jet is absorbed into $P$ and no
      comparison with an auxiliary depth $J$ is used.
\item At $s=0$ the constants are uniform for all $r$.  Along
      $a(\theta s)$ they are uniform whenever the resolving depth is at most
      $L(s)$ on $\mathrm{UCR}_{\rm sh}$ cells.  Endpoint/section buffers at
      two-time depth $N_{\rm time}$ cost at most
      $C_{\rm ep}N_{\rm time}$.  A phase-mode buffer is added only in the
      bounded intermediate frequency range
      \eqref{eq:intermediate-phase-range}; no such assertion is made for an
      unbounded Fourier complement.
\end{enumerate}
This is the exact physical escape input used by the energy theorem.  In
particular no statement about $X^{[J]}$ is transferred to $X^{[N(r)]}$.
\end{definition}

\begin{lemma}[Negative-pressure trial obstruction]
\label{lem:negative-pressure-trial-obstruction}
Suppose a fixed clean subsystem has relative pressure
$P_{\rm cl}=-\chi<0$, so a depth-$J$ clean trial has physical survival mass
at most $Ce^{-\chi J}$.  Then $O(J)$ such trials expose at most
$O(Je^{-\chi J})$ of the original conditional mass.  In particular they
cannot yield an uncovered gap $O(e^{-cJ})$.
\end{lemma}

\begin{proof}
Sum the survival masses of the trials.  This is the union bound
$CJ e^{-\chi J}=o(1)$, so the unexposed mass tends to one.  In the Bernoulli
model with clean probability $p=e^{-\chi}<1$ this estimate is exact up to a
constant.  Conditioning on survival cancels $p^J$ only inside a bad/good
ratio; it cannot cancel it in the physical cover.
\end{proof}

\begin{remark}[Adaptive full-mass alphabet as a nonvacuity route]
\label{rem:high-pressure-clean-route}
A falsifiable route to $\mathrm{PPE}_{N}$ is to use the actual one-state
quotient rather than a fixed negative-pressure horseshoe.  At depth $N$ take
an alphabet
\[
 \mathcal A_N=\{\text{return/grazing/context mark}\le b_N\}
 \cup\mathcal A_{\rm rough},\qquad b_N=\lambda N,
\]
where the finite set $\mathcal A_{\rm rough}$ contains fixed QNL rough
witnesses.  If the actual quotient tail is $Ce^{-c_b b}$, the mass lost in
the first $N$ steps is at most
$CN e^{-c_b\lambda N}=O(e^{-c'N})$; hence the relative pressure gap tends to
zero and $N\chi_N\to0$.  To obtain \eqref{eq:ppe-prescribed-sublevel} one
must still prove that rough-template, cutting, distortion, component, and
open-descendant constants for $\mathcal A_N$ grow slowly enough to leave a
strict exponent margin against $c_b\lambda$.  This uniform finite-alphabet
theorem is not supplied by the existing literature and is not assumed to
follow from Leclerc's fixed-alphabet theorem.  Failure of that margin leaves
$\mathrm{PPE}_{N}$ primitive.  Open-system normalization in
\cite{CanestrariHoles,DLreview} preserves the survival factor and cannot
evade Lemma~\ref{lem:negative-pressure-trial-obstruction}.
\end{remark}

\begin{proposition}[Single-source physical bracket]
\label{prop:exact-reduced-slope}
Assume
\[
 \mathrm{SS}(d_Z),\qquad
 \mathrm{ACTUAL\_SLOPE\_BRIDGE},
\]
and the dominant clean-gate chart normalization on the same labelled
$\mathrm{BR}_{\rm sec}$ carrier.  Then
\eqref{eq:exact-reduced-slope}--\eqref{eq:single-source-coefficient} hold on
every labelled native interval.  In the dominant coordinates there are
positive labelled scales $a_X,a_Y$ such that
\begin{equation}
\label{eq:dominant-coordinate-normalization}
 c\le |X'|/a_X,|Y'|/a_Y\le C,
 \qquad \Xi=|\xi_1|a_X\asymp|\xi_2|a_Y.
\end{equation}
\end{proposition}

\begin{proof}
The bracket assertion is $\mathrm{SS}(d_Z)$ applied, by
Definition~\ref{def:actual-slope-bridge}, to the actual logarithmic
derivative ratio \eqref{eq:actual-derivative-ratio-slope}.  The last display
is the separately assumed nonvanishing clean-gate chart normalization after
contracted coordinates are removed.  The proposition does not derive this
normalization from $\mathrm{SS}$.  No absolute uniform lower bound is
asserted: the inverse scale cost is the explicit factor
$\mathfrak d_\lambda$ in \eqref{eq:event-physical-norm}.
\end{proof}

\begin{definition}[Frozen-parent kernel reverse-H\"older interface
$\mathrm{KRH}_{p}$]
\label{def:kernel-reverse-holder}
Let $\mathfrak P_{\rm fr}$ be the set of every frozen parent actually used
below: context and orientation parents, frozen opposite endpoints,
actual/split constituents, cemetery continuations, and the normalized
branch--source-incidence parents of
Definition~\ref{def:source-incidence-space}.  For
$\eta\in\mathfrak P_{\rm fr}$ let $K_\eta$ be the normalized positive
descendant kernel on $\Xi_\eta$.  In a branch kernel the descendant is the
complete retained/cemetery label; in a source kernel it is the labelled
source occurrence.  It never includes a main-cell selector, a reciprocal
proposal factor, or the two internal carrier coordinates used by main
energy; that complete change of law belongs exclusively to
Definition~\ref{def:joint-source-likelihood}.  In particular, a parent in
$\mathrm{TENERGY}_{W}$ may freeze physical context, orientation, or the
opposite endpoint, but never a main source cell $A$.  A zero-mass source
kernel is omitted.  Let
$\mathscr W_\eta>0$ be the complete survivor or incidence envelope used on
that kernel and define
\begin{equation}
\label{eq:kernel-normalized-likelihood}
 L_\eta(\xi):=
 \frac{\mathscr W_\eta(\xi)}
      {\int_{\Xi_\eta}\mathscr W_\eta\,dK_\eta},
 \qquad \int L_\eta\,dK_\eta=1.
\end{equation}
The interface $\mathrm{KRH}_{p}$ is the uniform normalized bound
\begin{equation}
\label{eq:kernel-reverse-holder}
 \sup_{\eta\in\mathfrak P_{\rm fr}}
 \int_{\Xi_\eta}L_\eta^p\,dK_\eta\le C_{\rm KRH},
 \qquad p>1.
\end{equation}
All constants are independent of $m,n,s$, scalar mode, polynomial, and final
observable.  A global complete-tree $L^p$ bound is not a substitute for
\eqref{eq:kernel-reverse-holder}; rare parent cells remain separate kernels.
\end{definition}

\begin{lemma}[Parentwise tilted transfer]
\label{lem:parentwise-holder-transfer}
Under $\mathrm{KRH}_{p}$, every $K_\eta$-event $B$ satisfies
\begin{equation}
\label{eq:parentwise-holder-transfer}
 K_\eta^{W}(B):=\int_B L_\eta\,dK_\eta
 \le C_{\rm KRH}^{1/p}K_\eta(B)^{1/p'},
 \qquad p'=\frac p{p-1}.
\end{equation}
Consequently a raw polynomial or exponential tail uniform over the frozen
parents transfers with exponent divided by $p'$.  The statement applies
separately to branch and source-incidence kernels and therefore does not move
a global moment through a parent conditioning.
\end{lemma}

\begin{proof}
Apply H\"older on the single normalized kernel $K_\eta$ and use
\eqref{eq:kernel-reverse-holder}.  No global tree disintegration or signed
cancellation is involved.
\end{proof}

\begin{definition}[Tilted endpoint upper-scale grazing ledger
$\mathrm{USC}_{\rm end}^{W}(q_g^W)$]
\label{def:endpoint-upper-scale}
On a depth-$N$ homogeneous endpoint branch let $k_j^-$ and $k_j^+$ be the
past and future homogeneity-strip indices and put
\begin{equation}
\label{eq:homogeneity-log-marks}
 G_N^-:=\sum_{j<N}\log(1+k_j^-),\qquad
 G_N^+:=\sum_{j<N}\log(1+k_j^+).
\end{equation}
Let $S_\lambda^\pm\ge1$ dominate the positive derivative scale of the initial
physical-face parametrization, same-branch section, scalar endpoint chart,
and coarea/holonomy coordinate change in the two orientations, and put
\begin{equation}
\label{eq:augmented-homogeneity-mark}
 S_\lambda:=1+S_\lambda^-+S_\lambda^+,
 \qquad
 \widetilde G_N^\pm:=G_N^\pm+\tfrac12\log S_\lambda^\pm.
\end{equation}
On a retained label define the physical survivor factor
\begin{equation}
\label{eq:global-event-weight}
 \mathscr W_{\lambda,\rm surv}:=
 e^{\epsilon(|c^-|+|c^+|)}(1+H_\lambda)^{a_H}
 \mathfrak d_\lambda^{A_{\rm dom}}
 \operatorname{Mark}(\lambda)^2
 \{1+Z(\mathcal G_\lambda^-)+Z(\mathcal G_\lambda^+)\}
 S_\lambda^{a_{\rm scl}}.
\end{equation}
Here $a_{\rm scl}>0$ is fixed below the available positive-scale moment.
For a first-failure label of type
$i\in\{\mathrm{ctx},\mathrm{stop},\mathrm{buf},\mathrm{gate},
\mathrm{cont}\}$, let $\mathscr W_{\lambda,i}^{\rm cem}$ be the predictable
stopped-source envelope from Definition~\ref{def:cemetery-source-domination}.
The one weight used everywhere is
\begin{equation}
\label{eq:global-stopped-source-weight}
 \mathscr W_{\lambda,*}:=
 \mathscr W_{\lambda,\rm surv}\one_{\mathsf G}(\lambda)
 +\sum_i\mathscr W_{\lambda,i}^{\rm cem}\one_{B_i}(\lambda).
\end{equation}
The events $\mathsf G,B_i$ are disjoint because the tree records the first
failure.  We retain the superscript $W$ for tilted laws and rates, but it
always means the unique weight $\mathscr W_*$ in
\eqref{eq:global-stopped-source-weight}; there is no separate survivor or
cemetery tilt.
For every frozen parent $\eta\in\mathfrak P_{\rm fr}$ whose descendant
kernel contains a depth-$N$ retained endpoint, the complete actual/split law
extending $\mathrm{NST}_{\rm phys}$ satisfies, with the \emph{same} weight
\eqref{eq:global-stopped-source-weight},
$0<\int\mathscr W_{\eta,*}\,dK_\eta\le C_W$ uniformly in the allowed
parameters and depths, and
\begin{equation}
\label{eq:upper-scale-excess-tail}
 \int_{\Xi_\eta}
  \mathscr W_{\lambda,*}\one_{\mathsf G}
  \one_{\{\widetilde G_N^\pm>\gamma_gN+u\}}\,dK_\eta
 \le C\!\left(\int_{\Xi_\eta}\mathscr W_{\eta,*}\,dK_\eta\right)
 e^{-q_g^Wu},\qquad u\ge0.
\end{equation}
The indicator $\one_{\mathsf G}$ is part of the assertion: the depth-$N$
strip indices, $S_\lambda^\pm$, and $\widetilde G_N^\pm$ are read only on
retained homogeneous endpoint branches.  They are not extended to an early
cemetery branch.  Such branches are charged by $\mathrm{CEM\_DOM}$ and
$\mathrm{CLOCK\_PROB}$ below; additive cemetery currents are charged on the
source-incidence ledger.  The same complete label law is used for the
native stopping, PPE, endpoint marks, and all cemetery failures;
no marginal tail is transferred to a different conditional law, and an
unweighted probability tail is not promoted to a marked tail by mere
$L^1$ integrability.  Equivalently, under every normalized parent likelihood
\eqref{eq:kernel-normalized-likelihood}, the survivor subprobability has
augmented grazing exponent $q_g^W$.  This is the ``normalize the
positive source, estimate,
then multiply its mass back'' organization used in
\cite[Lemma~5.10]{CanestrariHoles}; that lemma is a template and does not
prove the present tilted tail.  A sufficient, but not necessary, route to
\eqref{eq:upper-scale-excess-tail} is the corresponding raw parentwise tail
together with $\mathrm{KRH}_{p_W}$.  Lemma
\ref{lem:parentwise-holder-transfer}, applied separately on every frozen
parent kernel, gives only $q_g^W=q_g/p_W'$, and all ledgers below use this
degraded exponent.  A complete-tree global $L^{p_W}$ estimate gives only a
global aggregate scale tail and is not used here.
\end{definition}

\begin{definition}[Complete positive tilted energy tree
$\Pi_{s,N}^{W}$]
\label{def:complete-tilted-energy-tree}
Fix an energy scale $r$, the prescribed depth $N=N(r)$, one orientation,
and the frozen opposite context.  Before any mass is discarded, compose the
normalized positive refinement kernels for the context gate, actual native
stopping, homogeneous buffer, physical gate, continuation, and retained
PPE descendant.  Whenever a stage fails, send its full positive mass to a
distinct cemetery label
\[
 \partial_{\rm ctx},\quad\partial_{\rm stop},\quad
 \partial_{\rm buf},\quad\partial_{\rm gate},\quad
 \partial_{\rm cont}
\]
and keep the successful labels.  The resulting probability measure on the
complete tree, including the sampled physical point and all actual/split
constituent labels, is denoted $\Pi_{s,N}$.  A cemetery label carries the
last marks measurable before failure.  It is not completed by declaring
unknown future factors equal to one.  Instead, the stopped source and its
remaining propagator are enveloped before the branch enters the cemetery,
as in Definition~\ref{def:cemetery-source-domination}.  Thus the piecewise
weight $\mathscr W_*$ in \eqref{eq:global-stopped-source-weight} is positive
and measurable on the entire tree, not only on the terminal survivor law.
Define
\begin{equation}
\label{eq:complete-energy-tree-tilt}
 d\Pi_{s,N}^{W}:=
 \frac{\mathscr W_*}{\mathbb E_{\Pi_{s,N}}\mathscr W_*}\,d\Pi_{s,N},
 \qquad
 1\le\mathbb E_{\Pi_{s,N}}\mathscr W_*\le C_W .
\end{equation}
All tree laws below, including those for a split constituent, use this same
extension and tilt.  A terminal survivor law which has already forgotten
discarded branches is not a substitute for $\Pi_{s,N}$.
\end{definition}

\begin{definition}[Concrete oriented source and test spaces]
\label{def:concrete-source-spaces}
Fix the common finite physical atlas, the forward/reverse homogeneous cones,
and the separation-time exponent $\alpha_*$.
For $\varepsilon\in\{+,-\}$ let
$\mathfrak S_{\rm sf}^{\varepsilon}$ be the completion of finite signed
combinations of homogeneous measured standard curves
$\Sigma=\sum_j b_j(W_j,\rho_j)$ in the positive-decomposition norm
\begin{equation}
\label{eq:standard-family-source-norm}
 \|\Sigma\|_{\mathfrak S_{\rm sf}^{\varepsilon}}
 :=\inf_{\Sigma=\sum_jb_j(W_j,\rho_j)}
 \sum_j|b_j|M_j
 \{1+\operatorname{Reg}_{\alpha_*}(\rho_j)+Z(W_j)\}.
\end{equation}
Here $M_j=\int_{W_j}\rho_j\,dm_{W_j}$ and $Z(W_j)$ is the usual inverse-length
boundary functional after canonical homogeneous subdivision, and the
infimum is over positive measured-curve decompositions of the Jordan parts.

Let $\mathfrak S_{\rm flux}^{\varepsilon}$ be the completion of finite signed
oriented coarea/face currents
$\Sigma=\sum_jb_j(\iota_j)_*(\rho_j\,ds)$, with each $\iota_j$ a homogeneous
$C^2$ face chart in orientation $\varepsilon$, under
\begin{equation}
\label{eq:flux-source-norm}
 \|\Sigma\|_{\mathfrak S_{\rm flux}^{\varepsilon}}
 :=\inf\sum_j|b_j|M_j
 \{1+\operatorname{Reg}_{\alpha_*}(\rho_j)+Z_{\partial}(\iota_j)
       +\operatorname{Atl}_2(\iota_j)\}.
\end{equation}
$Z_{\partial}$ counts oriented endpoints and homogeneous cuts with their
inverse-length weights, while $\operatorname{Atl}_2$ is the fixed-atlas
$C^2$ chart/coarea mark, including the tangent Jacobian.  Coincident
artificial faces are assembled with orientation before this norm is taken.

A normalized positive standard family is recorded as
$\mathcal G=\{p_j,(W_j,\rho_j)\}$ with $p_j\ge0$, $\sum_jp_j=1$, and each
$\rho_j$ a probability density on $W_j$.  It describes only the carrier
shape.  Its shape cost is
\begin{equation}
\label{eq:normalized-source-shape-cost}
 K_{\rm sf}(\mathcal G):=
  1+\operatorname{Reg}_{\alpha_*}(\mathcal G)+Z(\mathcal G),
 \qquad Z(\mathcal G):=\sum_j\frac{p_j}{|W_j|}.
\end{equation}
For a normalized flux family one analogously defines
$K_{\rm flux}=1+\operatorname{Reg}_{\alpha_*}+Z_\partial+
\operatorname{Atl}_2$.  A positive source of exact mass $m$ is $m\mathcal G$,
and its source-norm cost is $mK(\mathcal G)$; probability normalization never
sets $K(\mathcal G)$ equal to one.  This is the mass--shape organization of
standard families used in \cite{CanestrariHoles}.

The corresponding test spaces are the completions of bounded functions with
\begin{align}
 \|\phi\|_{\mathfrak T_{\rm sf}^{\varepsilon}}
 &:=\|\phi\|_\infty+
   \sup_{W\in\mathcal W^\varepsilon}
   [\phi|_W]_{\alpha_*,s_\varepsilon},
 \label{eq:standard-family-test-norm}\\
 \|\phi\|_{\mathfrak T_{\rm flux}^{\varepsilon}}
 &:=\|\phi\|_{\mathfrak T_{\rm sf}^{\varepsilon}}+
   \sup_{\iota\in\mathcal A_{\rm face}^\varepsilon}
   \|\phi\circ\iota\|_{C^1}.
 \label{eq:flux-test-norm}
\end{align}
The seminorm in the first line is the dynamic H\"older seminorm determined
by separation time.  Forward and reverse transfer act by the canonical
pushforward of measured curves/currents and by pullback on these tests.
Thus $U_s^{k,\varepsilon}$ is an actual operator
$\mathfrak S_{\rm sf/flux}^{\varepsilon}\to
(\mathfrak T_{\rm sf/flux}^{\varepsilon})^*$, not a symbol depending on the
eventual observable.  Any cemetery row not admitting one of these two
representations remains an explicitly named $\mathrm{SOURCE\_PROP}_i$
primitive with its own fully defined Banach source/test pair.
\end{definition}

\begin{lemma}[Source-space separation and operator interface]
\label{lem:source-space-operator-boundedness}
The norms \eqref{eq:standard-family-source-norm} and
\eqref{eq:flux-source-norm} dominate the total variation of every finite
representative.  Hence they separate currents, and their completions are
Banach source spaces.  On every common homogeneous atlas used below,
$\mathrm{SOURCE\_PROP}_i$ includes the quantitative assertion
\begin{equation}
\label{eq:source-operator-boundedness}
 \|U_s^{k,\varepsilon}\Sigma\|_{(\mathfrak T_i)^*}
 \le C_i\gamma_i^{(k-R_i(\Sigma))_+}
       \|\Sigma\|_{\mathfrak S_i^\varepsilon},
\end{equation}
with measurable recovery mark and constants uniform in $s,k$, the marked
occurrence, Fourier mode, polynomial, and final observable.  Thus bounded
pushforward/pullback is an explicit falsifiable part of the interface; it is
not inferred merely from the word ``completion.''
\end{lemma}

\begin{proof}
For every positive decomposition, total variation is at most
$\sum_j|b_j|M_j$, which is bounded by the displayed source cost.  Taking the
infimum proves separation and permits completion.  The operator assertion is
precisely the typed propagation estimate that must be proved for the moving
atlas; once it holds on the dense finite representatives, it extends
uniquely to the completion.
\end{proof}

\begin{definition}[Measurable constant-polarity decomposition interface
$\mathrm{MPD}$]
\label{lem:measurable-positive-decomposition}
Every signed source admitted to a regular source row is required to be
represented on
a countable standard-Borel component extension
$B_{\lambda,a}\subset\mathbb N$ by
\begin{equation}
\label{eq:measurable-positive-decomposition}
 \Sigma_{\lambda,a}
 =\int_{B_{\lambda,a}}s(b)\,m_{\lambda,a}(db)\,
       \nu_{\lambda,a,b},
 \qquad s(b)\in\{+1,-1\},\qquad
 \nu_{\lambda,a,b}\in\mathcal P(X_{\lambda,a,b}).
\end{equation}
Here $m_{\lambda,a}$ is an atomic countable kernel, the component graph and
weights are jointly measurable, each
$\nu_{\lambda,a,b}$ is a positive normalized regular section, and $s(b)$ is
constant on its entire carrier.  The component index $b$ is absorbed into
the fixed countable occurrence/component mark below.  In particular, no
carrier-dependent polar multiplier is applied after the positive section
norm has been evaluated.

Joint Borel measurability of the component graph, atomic kernel, masses,
normalized sections, source norms, recovery marks, and operator marks is an
explicit part of $\mathrm{MPD}$; it is not inferred from fibrewise Jordan
decomposition.  An essentially uncountable decomposition is not admitted
without a separately typed kernel interface.  If only an abstract Jordan
decomposition is available and either positive part falls outside the
declared regular cone, the source is routed to a typed TV or cemetery row.
\end{definition}

\begin{lemma}[Two-component baseline cone split]
\label{lem:baseline-cone-split}
There is a directly checkable two-component construction for a
$C^{\alpha_*}$ signed density $f$ on one homogeneous carrier, with respect
to its normalized arclength.  Choose
\begin{equation}
\label{eq:baseline-cone-split}
 c\ge 2\|f\|_\infty+[f]_{\alpha_*},
 \qquad f=(c+f)-c.
\end{equation}
Both summands are positive, the first is bounded below by $c/2$, and their
normalized densities have a uniform log--H\"older constant.  Their two
positive masses, source norms, chart/boundary costs, recovery times, and
operator costs are charged separately to the envelope kernel.  The same
construction applies chartwise to a regular flux density whenever the
chosen atlas and all normalization data are Borel in the external labels.
\end{lemma}

\begin{proof}
For the displayed construction,
$c+f\ge c/2$ and
\[
 |\log(c+f(x))-\log(c+f(y))|
 \le 2[f]_{\alpha_*}c^{-1}d(x,y)^{\alpha_*}.
\]
The constant component has zero log--H\"older seminorm.  Normalize the two
positive components and multiply their exact masses back.  This proves the
claimed cone control without cutting at the zero set of $f$.  It does not
prove the Borel selection required by $\mathrm{MPD}$ for an arbitrary moving
family.
\end{proof}

\begin{definition}[Cemetery source kernels
$\mathrm{CEM\_DOM}$]
\label{def:cemetery-source-domination}
Let
\[
 \mathcal I_{\rm cem}:=
 \{\mathrm{ctx},\mathrm{stop},\mathrm{buf},
   \mathrm{gate},\mathrm{cont}\}.
\]
For $i\in\mathcal I_{\rm cem}$ let
$\tau_i$ be the first failure stage.  Fix a parameter-independent Banach
source/test pair
$\mathfrak S_i^{\varepsilon_i},\mathfrak T_i$ and a standard-Borel carrier
$X_i$.  On $B_i$ require two measurable scalar masses
$0\le m_i(\lambda)\le q_i(\lambda)<\infty$ and a weakly measurable
probability section $\nu_{i,\lambda}\in\mathcal P(X_i)$.  Here $m_i$ is the
exact positive coefficient/probability mass of that stopped component, while $q_i$ is
only its source-norm envelope.  Put
$K_i(\lambda):=\|\nu_{i,\lambda}\|_{\mathfrak S_i^{\varepsilon_i}}$; in
general $K_i>1$.  By Definition~\ref{lem:measurable-positive-decomposition}, the
label space has already been extended to a positive component and carries a
measurable constant sign $s_i(\lambda)\in\{+1,-1\}$.  Let
$C_{U,i}(\lambda)\ge1$ dominate the remaining typed operator on that
component.  The exact discarded current is
defined by its scalar pairing
\begin{equation}
\label{eq:exact-stopped-source-factorization}
 J_i^{\rm disc}(\Phi):=
 \int_{B_i}s_i(\lambda)m_i(\lambda)
   \langle U_{i,\lambda}\nu_{i,\lambda},\Phi\rangle
   \,d\Pi_{s,N}(\lambda),
 \qquad s_i(\lambda)\in\{+1,-1\}.
\end{equation}
The carrier is integrated inside the typed source/operator pairing before
the label integral is taken.  In particular,
\eqref{eq:exact-stopped-source-factorization} does not assert absolute
integrability of an arbitrary pointwise representative
$U_{i,\lambda}^{*}\Phi$ against $\nu_{i,\lambda}$.
Thus $m_i$ is a density on the tree-label space while $\nu_{i,\lambda}$ is
the normalized shape on its physical carrier; no Radon--Nikodym identity
equates measures living on these different spaces.  Labelwise the honest
source-norm/recovery envelope is
\begin{equation}
\label{eq:cemetery-fixed-source-norm}
 q_i(\lambda)\ge
 m_i(\lambda)K_i(\lambda)\gamma_i^{-R_i(\lambda)}
 \mathscr A_i^{\rm opp}(\lambda)C_{U,i}(\lambda).
\end{equation}

Let $R_i$ be the recovery time and
$\mathscr A_i^{\rm opp}\ge1$ the already frozen opposite-side amplitude.
There must be a predictable cemetery envelope, fixed before $m,n$, the mode,
and the final observable, such that
\begin{equation}
\label{eq:cemetery-kernel-domination}
 q_i(\lambda)\le C\mathscr W_i^{\rm cem}(\lambda)
 \quad\Pi_{s,N}\text{-a.e. on }B_i.
\end{equation}
Weak measurability means that
$\lambda\mapsto\langle U_{i,\lambda}\nu_{i,\lambda},\Phi\rangle$
is measurable for every $\Phi\in\mathfrak T_i$; this is part of the
interface.  The source-operator bound, followed by the scalar triangle
inequality on the label space, gives
\begin{equation}
\label{eq:cemetery-current-weighted-bound}
 |J_i^{\rm disc}(\Phi)|
 \le C\|\Phi\|_{\mathfrak T_i}
       \int_{B_i}\mathscr W_i^{\rm cem}\,d\Pi_{s,N}
 =C\|\Phi\|_{\mathfrak T_i}
   \mathbb E_{\Pi_{s,N}}\mathscr W_*\Pi_{s,N}^{W}(B_i).
\end{equation}
A future factor may be replaced by one only when the remaining propagator is
a positive total-variation contraction.  Otherwise its recovery and
opposite-amplitude costs are frozen in $q_i$ before the cemetery transition.
The fixed owner table is: context and physical-gate failures use oriented
flux carriers; stopping, buffer, and continuation failures use measured
standard-curve carriers.  Each row records its sign, forward/reverse
orientation, $m_i,q_i,\nu_{i,\lambda},K_i,R_i,U_i$, and opposite-side
amplitude.
Any row not admitting a standard-family or flux representation uses a
registered $\mathrm{SOURCE\_PROP}_i$ Banach pair, but it must still admit
the displayed positive CEM factorization.  A row admitting no such
factorization is a no-go, not a separately floating current.
Its contribution is not allowed to live in a parallel ledger:
the Borel embedding into the global cemetery owner, including equality of
the exact mass, envelope, section, polarity, type, recovery, and operator
marks, is required in 
\eqref{eq:cemetery-global-kernel-identification} below.
\end{definition}

\begin{lemma}[Fixed-norm stopped source and summable propagation]
\label{lem:stopped-source-propagation}
For each $i\in\mathcal I_{\rm cem}$ assume that
$U_{i,s}^{k,\varepsilon_i}:\mathfrak S_i^{\varepsilon_i}\to
\mathfrak T_i^*$ is measurable in the preceding sense and, uniformly in the
moving-table parameter, common atlas, time, and orientation,
\begin{equation}
\label{eq:stopped-source-propagation}
 |\langle U_{i,s}^{k,\varepsilon_i}\Sigma,\phi\rangle|
 \le C\|\Sigma\|_{\mathfrak S_i^{\varepsilon_i}}
       \gamma_i^{(k-R_i)_+}\|\phi\|_{\mathfrak T_i},
 \qquad 0<\gamma_i<1.
\end{equation}
Then \eqref{eq:cemetery-fixed-source-norm}--
\eqref{eq:cemetery-kernel-domination} imply
\eqref{eq:cemetery-current-weighted-bound}, uniformly in $m,n,s$, mode, and
test.  Normalize the positive shape, propagate it, and multiply the exact
mass back; domination uses the distinct norm envelope $q_i$.  No signed
cancellation or cross-space Radon--Nikodym formula is
used, and no pointwise carrier Tonelli estimate is claimed.  The results of
\cite{CanestrariPerturb,CanestrariHoles} are templates
for this source-first organization, not proofs for moving billiards.
\end{lemma}

\begin{definition}[Post-propagation carrier domination
$\mathrm{POST\_CARRIER\_DOM}$]
\label{def:post-carrier-domination}
The operator norm in
\eqref{eq:source-operator-boundedness} controls an occurrence after its
carrier has been integrated; it does not by itself license restriction to a
phase-bad subset of that carrier.  Accordingly, every occurrence declared
to have passage type $J$ must carry a jointly measurable, predictable mark
$C_{\rm car}(\lambda,a)\ge1$ and, for every
$\Phi_a\in\mathfrak T_a$, a canonical measurable representative
$u_{\lambda,a,\Phi_a}$, linear in $\Phi_a$, such that
\begin{align}
 \Lambda_{\lambda,a,\Phi_a}(E)
 &:=\int_Eu_{\lambda,a,\Phi_a}(x)\,d\nu_{\lambda,a}(x),
 \label{eq:post-carrier-restriction-identity}\\
 \|u_{\lambda,a,\Phi_a}\|_{L^\infty(\nu_{\lambda,a})}
 +[u_{\lambda,a,\Phi_a}]_{\alpha_*,{\rm sc}}
 &\le C_U(\lambda,a)C_{\rm car}(\lambda,a)K_{\lambda,a}
       \|\Phi_a\|_{\mathfrak T_a}.
 \label{eq:post-carrier-positive-domination}
\end{align}
Here $E\mapsto\Lambda_{\lambda,a,\Phi_a}(E)$ is the declared finite signed
post-propagation carrier measure and
\begin{equation}
\label{eq:post-carrier-full-pairing}
 \Lambda_{\lambda,a,\Phi_a}(X_{\operatorname{type}(a)})
 =\langle U_{\lambda,a}\nu_{\lambda,a},\Phi_a\rangle.
\end{equation}
When $\one_E\nu_{\lambda,a}$ belongs to the original source space, this
measure agrees with
\[
 \langle U_{\lambda,a}(\one_E\nu_{\lambda,a}),\Phi_a\rangle.
\]
For an
arbitrary Borel $E$ the latter is not used as an ill-typed source-space
expression.
Here the scaled seminorm is the one required by the later fractional
nonstationary-phase row; when only measurable restrictions are used its
seminorm term may be omitted explicitly.  In either case
\begin{equation}
 \label{eq:post-carrier-set-domination}
 \int_E|u_{\lambda,a,\Phi_a}|\,d\nu_{\lambda,a}
 \le C_UC_{\rm car}K_{\lambda,a}\nu_{\lambda,a}(E)
       \|\Phi_a\|_{\mathfrak T_a}.
\end{equation}
The factor $C_{\rm car}$ is fixed before $(m,n,\xi,P)$ and the final test
and is absorbed into the propagated $q$ envelope.  A row for which
\eqref{eq:post-carrier-restriction-identity}--
\eqref{eq:post-carrier-positive-domination} is unavailable is assigned
passage type ${\rm TV}$ and is never cut in the carrier coordinate.  Thus
post-propagation cancellation cannot be hidden behind a pre-propagation
constant-polarity decomposition.
\end{definition}

\begin{definition}[Propagated dual-mass marked source kernel]
\label{def:source-incidence-space}
Together with $\mathrm{MPD}$, the interface
$\mathrm{DUAL\_MASS\_KERNEL}$ first fixes the owner and type registries
\[
 \mathcal O_{\rm src}:=
 \{{\rm cem},{\rm main},{\rm exc}_J,{\rm exc}_{\rm TV},
   {\rm face/word},{\rm deep}\},
\]
and
\[
 \mathcal T_{\rm src}:=
 \{{\rm sf}^{+},{\rm sf}^{-},{\rm flux}^{+},{\rm flux}^{-}\}
 \sqcup\{(i,\varepsilon):i\in\mathcal I_{\rm prop}^{\rm extra},
                    \ \varepsilon\in\{+,-\}\}.
\]
Here $\mathcal I_{\rm prop}^{\rm extra}$ is a fixed countable registry of
every additional $\mathrm{SOURCE\_PROP}_i$ Banach source/test pair admitted
by the theorem.  Thus an extra typed row is part of the common kernel rather
than an unnamed functional outside it.  The standard-Borel mark space is
\[
 \begin{aligned}
 Y_{\rm src}:={}&\mathcal O_{\rm src}\times\mathbb N\times
 \{r,\mathrm{PPE},N_{\rm time},\mathrm{word}\}\times
 \mathbb N_0\times\{J,\mathrm{TV}\}\\
 &\times\{+,-\}_{\rm pol}\times\{+,-\}_{\rm or}\times
 \mathcal T_{\rm src}.
 \end{aligned}
\]
The first integer is the occurrence/component index supplied by the common
refinement.  The second is the numerical word depth $\ell_{\rm word}$; it is
set to zero when that coordinate is unused, while the clock field
distinguishes an actual depth-zero word row.  Let $X_{\rm src}$ be the corresponding countable disjoint union
of the standard-family, flux, and registered extra typed carriers.

All source restrictions are taken from one master datum
\[
 \mathfrak d=(N_{\rm tree},r,N_{\rm PPE},N_{\rm time},m,n,L,\xi,P)
 \in\mathfrak D.
\]
On every section used in the product-time CM2 argument the time clock is not
an independent coordinate: throughout this section and every downstream
central-phase call we impose the typed identity
\begin{equation}
\label{eq:time-clock-product-time-identity}
 N_{\rm time}=N_{m,n}:=\max\{m,n\}.
\end{equation}
An auxiliary construction may temporarily carry a free $N_{\rm time}$, but
it is not an admissible central master datum until restricted to
\eqref{eq:time-clock-product-time-identity}.  Thus no time-tail estimate is
silently evaluated at a different depth from the product time.
The interface supplies a jointly measurable family
$\Pi_{s,\mathfrak d}$ whose tree-label marginal is the physical law
$\Pi_{s,N_{\rm tree}}$ and whose compatible coordinate restrictions have
that same marginal.  Put
\[
 d\overline M_{s,\mathfrak d}^{q}(\lambda,a)
 :=d\Pi_{s,\mathfrak d}(\lambda)q_\lambda(da).
\]
Every CLOCK, main, FACE, shallow, and deep cell below is a restriction or
pushforward of this same measure, with all estimates uniform in the
suppressed entries of $\mathfrak d$.  When only a scalar tree or resolving
depth is active, the older notation $\Pi_{s,N}$ and
$\overline M_{s,N}^{q}$ abbreviates this compatible marginal; it never
licenses mixing measures belonging to different data.

More strongly, the raw incidence kernel is projectively fixed before the
forbidden choices.  There is one simultaneous Borel kernel over the
complete tree such that whenever $\mathfrak d$ and $\mathfrak d'$ differ
only in $(m,n,\xi,P)$ or in the final test, their projected
$(\lambda,a)$ laws and the fields
\begin{equation}
\label{eq:source-kernel-projective-consistency}
 (m,q,\nu,U,C_U,C_{\rm car},o,j,\kappa,\pi,s_{\rm pol},
   \ell_{\rm word},\varepsilon_{\rm or},\operatorname{type},R,\gamma,
   \mathscr A^{\rm opp})
\end{equation}
agree identically.  Only the Borel clock, phase, and owner-handler filters
may depend on those choices.  Equivalently, all displayed
$\overline M_{s,\mathfrak d}^{q}$ are restrictions of one simultaneous
positive incidence measure, not unrelated fibrewise envelopes.  A family
of separately chosen kernels indexed by mode or test is not an admissible
$\mathrm{DUAL\_MASS\_KERNEL}$.

There are jointly measurable
weights $w^m(\lambda,a),w^q(\lambda,a)\ge0$ on the fixed marked product,
supported on the same Borel countable-fibre set
$I\subset\Lambda\times Y_{\rm src}$, with $w^m\le w^q$, and
\begin{equation}
\label{eq:marked-weighted-counting-kernel}
 \begin{aligned}
 m_\lambda(da)&:=\sum_{a\in I(\lambda)}w^m(\lambda,a)\delta_a(da),\\
 q_\lambda(da)&:=\sum_{a\in I(\lambda)}w^q(\lambda,a)\delta_a(da),
 \end{aligned}
 \qquad
 \sum_{a\in I(\lambda)}w^q(\lambda,a)<\infty.
\end{equation}
The support $I$ is restricted to the compatibility diagonal on which
$\operatorname{type}(a)=\mathrm{sf}^{\varepsilon_{\rm or}(a)}$,
$\mathrm{flux}^{\varepsilon_{\rm or}(a)}$, or
$(i,\varepsilon_{\rm or}(a))$ for a registered extra type $i$.  Polarity
$s_{\rm pol}(a)$ is an independent field and imposes no orientation change.
Thus $m_\lambda\ll q_\lambda$ are finite positive kernels on the fixed
standard-Borel space.  Write $g_\lambda=dm_\lambda/dq_\lambda$, so
$0\le g_\lambda\le1$.  The kernel $m_\lambda$ records the exact positive
coefficient mass in the chosen constant-polarity decomposition; it is not
the total variation of the recombined signed current.  The kernel
$q_\lambda$ records the larger propagated source-norm envelope.
For every Borel $A$, both kernel evaluations are measurable in $\lambda$.
This is the measurable weighted-counting construction, rather than an
unmeasurable variable-fibre sum; standard kernel composition and Tonelli are
used in the sense of \cite{VakarOng}.  A countable registry may first be
$s$-finite, but every cell used by the central Jensen argument must have
finite envelope mass.  For $q_\lambda$-a.e. mark $a$, let
$\nu_{\lambda,a}\in\mathcal P(X_{\operatorname{type}(a)})$ be a normalized
weakly measurable source section and put
$K_{\lambda,a}:=\|\nu_{\lambda,a}\|_{
\mathfrak S_a^{\varepsilon_{\rm or}(a)}}$.
The section is positive and its polarity is the occurrence-level mark
$s_{\rm pol}(a)\in\{+1,-1\}$ supplied by
Definition~\ref{lem:measurable-positive-decomposition}; $s_{\rm pol}(a)$ is never a
function of the carrier coordinate.
The record
\begin{equation}
\label{eq:source-incidence-record}
\begin{aligned}
 &(j(a),o(a),\operatorname{type}(a),\varepsilon_{\rm or}(a),
   s_{\rm pol}(a),\ell_{\rm word}(a),\kappa(a),\pi(a),\nu_{\lambda,a},\\
 &\qquad U_{\lambda,a},K_{\lambda,a},R_{\lambda,a},\gamma_a,
   C_U(\lambda,a),C_{\rm car}(\lambda,a),
   \mathscr A_{\lambda,a}^{\rm opp})
\end{aligned}
\end{equation}
is encoded by the mark and the measurable operator kernel.  Distinct
physical currents have distinct occurrence indices even when their supports
coincide.

To type the global scalar assembly, use the fixed Banach product
\begin{equation}
\label{eq:global-source-test-space}
 \mathfrak T_{\rm src}:=
 \prod_{a\in Y_{\rm src}}^{\ell^\infty}\mathfrak T_a,
 \qquad
 \|\Phi\|_{\mathfrak T_{\rm src}}
 :=\sup_a\|\Phi_a\|_{\mathfrak T_a}.
\end{equation}
Compatibility with physical observables is quantitative.  The interface
$\mathrm{PHYS\_TEST\_LIFT}$ supplies, before source assembly, one bounded
bilinear map for every admissible master datum (equivalently, one bounded
linear map on the projective tensor product),
\begin{equation}
\label{eq:bounded-physical-test-lift}
 \mathcal E_{\mathfrak d}:C^r(M)\times C^r(M)
       \longrightarrow\mathfrak T_{\rm src},\qquad
 \|\mathcal E_{\mathfrak d}(f,g)\|_{\mathfrak T_{\rm src}}
 \le C\|f\|_{C^r}\|g\|_{C^r},
\end{equation}
uniformly in $(s,m,n,\xi,P)$.  Its coordinates are the actual tests used by
the physical bilinear response term.  Joint linearity on the Cartesian
product is not meant: together with the displayed product bound it would
force the map to vanish.  Any time- or mode-dependent loss in a concrete
realization must be bounded predictably before those choices and absorbed
into $C_U$ (hence into $q$); it cannot be hidden in the lift norm.
The lift also closes the physical/source diagram, not merely the norms:
\begin{equation}
\label{eq:physical-test-lift-commutation}
 B_{\mathfrak d}^{\rm phys}(f,g)
 =\big\langle J_{\mathfrak d}^{\rm src},
       \mathcal E_{\mathfrak d}(f,g)\big\rangle,
\end{equation}
where the left side is the exact physical occurrence sum before any
absolute value and the right side is the grouped source current
\eqref{eq:marked-source-assembled-pairing}.  This identity is required for
every admissible master datum and is compatible with the exact current
clause of $\mathrm{PROP\_Q\_MATCH}$.
Closedness of a putative compatible subspace without
\eqref{eq:bounded-physical-test-lift} is insufficient.  Each
$U_{\lambda,a}^*$ acts on the coordinate
$\Phi_a$; below this restriction is suppressed from the notation.  Hence
heterogeneous standard-family and flux rows are not silently treated as
elements of one undeclared local test space.

The positive envelope measure is
\begin{equation}
\label{eq:source-incidence-space}
 dM_{\rm src}^{q}(\lambda,a,x):=
 d\Pi_{s,\mathfrak d}(\lambda)\,q_\lambda(da)\,d\nu_{\lambda,a}(x).
\end{equation}
This positive measure is used only for carrier-event energies, measurable
coverage, and positive majorants.  It does not define the signed current;
that current is the grouped occurrencewise pairing
\eqref{eq:marked-source-assembled-pairing} below.
The two atomic masses obey the pointwise domination
with $0<\gamma_a<1$, $R_{\lambda,a}\in\mathbb N_0$, and
\begin{equation}
\label{eq:opposite-amplitude-floor}
 \overline{\mathscr A}_{\lambda,a}^{\rm opp}
 :=\max\{1,\mathscr A_{\lambda,a}^{\rm opp}\}:
\end{equation}
\begin{equation}
\label{eq:source-incidence-envelope}
 w^q(\lambda,a)\ge w^m(\lambda,a)K_{\lambda,a}
 \gamma_a^{-R_{\lambda,a}}\overline{\mathscr A}_{\lambda,a}^{\rm opp}
 C_{\rm prop}(\lambda,a),
\end{equation}
for every active occurrence, where
\begin{equation}
\label{eq:propagated-operator-carrier-mark}
 C_{\rm prop}(\lambda,a):=
 C_U(\lambda,a)\max\{1,\one_{\{\pi(a)=J\}}C_{\rm car}(\lambda,a)\},
\end{equation}
$C_U(\lambda,a)\ge1$ is the measurable
operator mark satisfying
$\|U_{\lambda,a}\|_{\mathfrak S_a\to\mathfrak T_a^*}
\le C_U(\lambda,a)$, and every $J$ row satisfies
$\mathrm{POST\_CARRIER\_DOM}$.  Thus operator growth,
post-propagation carrier restriction, recovery, and the frozen
opposite amplitude are absorbed into the same propagated envelope $q$ used
by every source ledger.  For every admissible test define
the exact assembled current by first evaluating each occurrence in its
typed source/test dual and only then integrating the scalar result:
\begin{equation}
\label{eq:marked-source-assembled-pairing}
 \langle J_{\rm src},\Phi\rangle:=
 \int s_{\rm pol}(a)g_\lambda(a)
       \langle U_{\lambda,a}\nu_{\lambda,a},\Phi_a\rangle
       \,d\overline M_{s,\mathfrak d}^{q}(\lambda,a).
\end{equation}
The grouped scalar integrand is required to be jointly measurable.  Formula
\eqref{eq:marked-source-assembled-pairing} is not shorthand for a Tonelli
integral of $[U^*\Phi](x)$ over the carrier.  Carrier-restricted $J$ pieces
are defined by the signed scalar measure
\eqref{eq:post-carrier-restriction-identity}; TV pieces are never cut in the
carrier coordinate.

The local cemetery kernels are identified with, rather than merely
dominated by, the cemetery restriction of this global kernel.  Namely there
are Borel mark maps $a_i:B_i\to Y_{\rm src}$ for which
$\lambda\mapsto(\lambda,a_i(\lambda))$ are disjoint injections into
$I_{\rm cem}$ and whose graph images exhaust $I_{\rm cem}$ modulo
$\overline M_{s,\mathfrak d}^{q}$-null sets and for which
\begin{equation}
\label{eq:cemetery-global-kernel-identification}
\begin{gathered}
 w^m(\lambda,a_i(\lambda))=m_i(\lambda),\qquad
 w^q(\lambda,a_i(\lambda))=q_i(\lambda),\\
 (\nu,s_{\rm pol},U,K,R,\gamma,\mathscr A^{\rm opp},C_U,C_{\rm car})
   (\lambda,a_i(\lambda))
 = (\nu_i,s_i,U_i,K_i,R_i,\gamma_i,\mathscr A_i^{\rm opp},C_{U,i},1)(\lambda).
\end{gathered}
\end{equation}
The type and orientation fields match the declared Banach pair as well.
Consequently $J_{\rm cem}=\sum_iJ_i^{\rm disc}$ as an absolutely convergent
identity in the fixed test dual; there is no second cemetery current.
The global finite-envelope closure is the explicit requirement
\begin{equation}
\label{eq:global-finite-marked-envelope}
 Q_{s,\mathfrak d}^{\rm prop}:=\int q_\lambda(Y_{\rm src})
                  \,d\Pi_{s,\mathfrak d}(\lambda)<\infty.
\end{equation}
The allowed growth of this already operator-weighted quantity with resolving
depth is charged in the BASE, CLOCK, FACE, and deep ledgers.  There is no
second fixed-depth-only operator integral whose exponent can disappear from
the final minimum.  Fibrewise finiteness alone is not enough.  Measurability
of the grouped scalar pairing and
\eqref{eq:global-finite-marked-envelope} make
\eqref{eq:marked-source-assembled-pairing} an absolutely convergent outer
scalar integral defining one bounded functional on the fixed test space;
no carrier-level Gelfand or Bochner integral is asserted.  The outer scalar
triangle inequality gives uniformly on its unit ball
\begin{equation}
\label{eq:source-incidence-TV-bound}
 |\langle J_{\rm src},\Phi\rangle|
 \le Q_{s,\mathfrak d}^{\rm prop}\|\Phi\|_{\mathfrak T_{\rm src}}.
\end{equation}

Finally there is a measurable canonical owner map and, modulo
$M_{\rm src}^{q}$-null sets on the incidence graph
$I\subset\Lambda\times Y_{\rm src}$, the exhaustive identity
\begin{equation}
\label{eq:exhaustive-source-owner-partition}
 I=I_{\rm cem}\sqcup I_{\rm main}\sqcup
 I_{\rm exc,J}\sqcup I_{\rm exc,TV}\sqcup
 I_{\rm face/word}\sqcup I_{\rm deep}.
\end{equation}
The clock $r$, PPE, $N_{\rm time}$, or word is a measurable field/filter on
an owner cell, not a second owner.  This static identity does not by itself
give resolution-dependent coverage inside $I_{\rm main}$; that separate job
is performed by \eqref{eq:main-clock-coverage}.  This construction replaces every
naked sum over continuous labels.  Probability first-hit disjointization is
used only for survivor unions and never deletes a marked occurrence.
\end{definition}

\begin{lemma}[Globally finite marked-current assembly]
\label{lem:globally-finite-marked-current}
Under \eqref{eq:source-incidence-envelope} and
\eqref{eq:global-finite-marked-envelope},
\eqref{eq:marked-source-assembled-pairing} belongs to the fixed test-space
dual and obeys \eqref{eq:source-incidence-TV-bound}.  Moreover,
\eqref{eq:exhaustive-source-owner-partition} gives the absolutely convergent
dual identity
\begin{equation}
\label{eq:assembled-owner-current-identity}
 J_{\rm src}=J_{\rm cem}+J_{\rm main}+J_{\rm exc,J}+J_{\rm exc,TV}
              +J_{\rm face/word}+J_{\rm deep}.
\end{equation}
Thus a local bound on every registered cell cannot leave an unowned current
in the central or final assembly.
\end{lemma}

\begin{proof}
Use the exact identity $g_\lambda q_\lambda=m_\lambda$, evaluate the
carrier inside the grouped dual pairing, use the operator bound on
$m_\lambda\nu_{\lambda,a}$, and apply
\eqref{eq:source-incidence-envelope} to dominate it directly by the
propagated envelope $q_\lambda$.  Absolute scalar integration on the outer
incidence space bounds the result by $Q_{s,\mathfrak d}^{\rm prop}$
uniformly on the test unit ball.
Restrict the same positive envelope measure to the six disjoint owner cells
and apply countable additivity.  Absolute convergence in the dual follows
from the same common majorant, proving
\eqref{eq:assembled-owner-current-identity} without signed cancellation.
\end{proof}

\begin{remark}[Operator absorption and rate bookkeeping]
\label{rem:operator-absorption-rates}
All source rates below are rates of the propagated $q$, not of the bare
physical mass.  Thus a factor $C_U\lesssim r^{-a_U}$ (respectively
$C_U\lesssim e^{a_UN}$) has already enlarged the declared base/likelihood
growth and weakened every direct-TV rate by $a_U$.  In the main selection
route it enters the actual $a_A^L$ and $a_A^B$ in
\eqref{eq:source-value-small-ball}; in exceptional Jensen rows it enters the
displayed BASE--ENERGY gap; and in the face, shallow, and deep routes it is
contained respectively in $A_{\rm face}$, $A_{\rm sh}$, and $A_{\rm loss}$.
It is never subtracted a second time.  Conversely, if an operator mark cannot
be predicted before the time, mode, and final test are chosen, it is not a
legal $q$ mark and must remain a separately typed primitive.  Fixed-depth
finiteness alone supplies no rate.  Every post-absorption gap is required to
be strictly positive.
\end{remark}

\begin{definition}[Complete tilted energy ledger
$\mathrm{TENERGY}_{W}$]
\label{def:tilted-energy-ledger}
Let $B_{\rm ctx}(r)$ be failure of the joint native context gate,
$B_{\rm stop}(r)$ failure of the actual stopping antichain, and let
$B_{\rm buf}(N)$, $B_{\rm gate}(N)$, and $B_{\rm cont}(N)$ be their named
cemetery events.  On a retained good label let $B_{\rm PPE}(r,P)$ be
\[
 |X^{[N(r)]}_{\beta,Q}-P|\le Cr,
 \qquad P\in\mathbb R_{d_Z}[z].
\]
Condition $\mathrm{TENERGY}_{W}$ requires, uniformly in $s$, orientation,
actual/split constituent, parent context, and admissible polynomial,
\begin{align}
 \Pi_{s,N}^{W}(B_{\rm ctx}(r))&\le Cr^{\Theta_{\rm ctx}^{W}},&
 \Pi_{s,N}^{W}(B_{\rm stop}(r))&\le Cr^{\Theta_{\rm st}^{W}},
 \label{eq:tilted-context-stop-tails}\\
 \Pi_{s,N}^{W}(B_i(N))&\le Ce^{-c_i^{W}N},
 &i&\in\{\mathrm{buf},\mathrm{gate},\mathrm{cont}\},
 \label{eq:tilted-auxiliary-tails}\\
 \Pi_{s,N}^{W}(\mathsf G\cap B_{\rm PPE}(r,P))&\le
 Cr^{\Theta_{\rm PPE}^{W}}.
 \label{eq:tilted-PPE-tail}
\end{align}
Every additional exceptional event must also satisfy the aggregate clock
interfaces below.  These estimates are parent-kernel statements: they are
either independent primitives or consequences of uniform raw parent tails
and $\mathrm{KRH}_{p}$.  A global complete-tree moment alone does not imply
them.  They concern the unselected physical law $\mathsf R_0$: ``parent
context'' may freeze physical context, orientation, or the opposite
endpoint, but never the later main selector $A$.  Atomic tail bounds by
themselves are not part of $\mathrm{TENERGY}_{W}$.  Nor does this ledger by
itself identify the two-endpoint product law: an opposite-endpoint-dependent
tilt is only a one-sided conditional estimate.  The exact pair-law content
is separately imposed in $\mathrm{PAIR\_TENERGY}$ below.
\end{definition}

\begin{definition}[Probability clock and survivor ledger
$\mathrm{CLOCK\_PROB}_{W}$]
\label{def:probability-clock}
The five cemetery failures are the named events in
$\mathrm{TENERGY}_{W}$ and are not members of an additional registry.
Every other bad event used only to control a union or survivor mass is given
a first-failure stage and one clock in
$\{r,\mathrm{PPE},N_{\rm time},\mathrm{word}\}$.  Inside a stage fix a
deterministic order and put
\begin{equation}
\label{eq:clock-prob-first-hit}
 \widehat B_{\tau,j}:=B_{\tau,j}\setminus\bigcup_{k<j}B_{\tau,k}.
\end{equation}
This is solely a probability partition; it makes no assertion about how many
additive currents coexist on an atom.  With word-depth currents assigned to
$\mathrm{FACE\_TIME}_{\rm CM2}$ for product time and
$\mathrm{FACE}_{\rm 2CUT}$ for the depth limit, define
\begin{align}
 \mathfrak Z_{r}^{\rm prob}(r)
 &:=\sum_{j\in\mathcal J_r^{\rm prob}}
      \Pi^W(\widehat B_{r,j}(r)),\notag\\
 \mathfrak Z_{\rm PPE}^{\rm prob}(N)
 &:=\sum_{j\in\mathcal J_{\rm PPE}^{\rm prob}}
      \Pi^W(\widehat B_{{\rm PPE},j}(N)),\notag\\
 \mathfrak Z_{\rm time}^{\rm prob}(N)
 &:=\sum_{j\in\mathcal J_{\rm time}^{\rm prob}}
      \Pi^W(\widehat B_{{\rm time},j}(N)),
 \label{eq:clock-prob-partitions}
\end{align}
and require the aggregate tails
\begin{equation}
\label{eq:clock-prob-tails}
 \mathfrak Z_r^{\rm prob}(r)\le Cr^{\Theta_r^{\rm prob}},\quad
 \mathfrak Z_{\rm PPE}^{\rm prob}(N)\le Ce^{-c_{\rm PPE}^{\rm prob}N},
 \quad
 \mathfrak Z_{\rm time}^{\rm prob}(N)\le Ce^{-c_{\rm time}^{\rm prob}N},
\end{equation}
with strictly positive exponents for every nonempty class.  These bounds,
together with the five named $\mathrm{TENERGY}_{W}$ tails, serve only
$\mathrm{SURV\_NORM}$.  Passage through Jensen is irrelevant to this union
ledger.
\end{definition}

\begin{definition}[Exceptional source clock--passage ledger
$\mathrm{CLOCK\_SRC}_{W}$]
\label{def:source-clock-matrix}
Put $M_*:=\mathbb E_{\Pi_{s,N}}\mathscr W_*$.  For
$\kappa\in\{r,\mathrm{PPE},N_{\rm time}\}$ and
$\pi\in\{J,\mathrm{TV}\}$, the interface first gives jointly Borel,
pairwise disjoint non-main filters
$\mathcal A_{\rm nm}^{\kappa,\pi}(\mathfrak d)$ satisfying, modulo one
common $\overline M_{s,\mathfrak d}^{q}$-null set,
\begin{equation}
\label{eq:nonmain-clock-coverage}
 I_{\rm cem}\sqcup I_{{\rm exc},J}\sqcup I_{{\rm exc},\rm TV}
 =\bigsqcup_{\substack{
       \kappa\in\{r,{\rm PPE},N_{\rm time}\}\\
       \pi\in\{J,{\rm TV}\}}}
       \mathcal A_{\rm nm}^{\kappa,\pi}(\mathfrak d).
\end{equation}
Restriction of the one propagated envelope then gives the absolutely
convergent current identity
\begin{equation}
\label{eq:nonmain-clock-current-identity}
 J_{\rm cem}+J_{{\rm exc},J}+J_{{\rm exc},\rm TV}
 =\sum_{\kappa,\pi}J_{\rm nm}^{\kappa,\pi}.
\end{equation}
The filters and bounds are uniform in the full master datum
$\mathfrak d$; hence no non-main clock row can be made vacuous by omitting
its occurrences.

Put $\mathcal A_{\kappa,J}:=\mathcal A_{\rm nm}^{\kappa,J}$.
After the main coverage below is fixed, let
\[
 \mathcal A_{\kappa,\rm TV}:=
 \mathcal A_{\rm nm}^{\kappa,\rm TV}
 \sqcup\mathcal A_{\rm main}^{\kappa,\rm TV}.
\]
Thus the TV cell contains \emph{every} clock-owned occurrence, including a
main-owner occurrence, charged directly in total variation.  Main $J$ cells
are excluded from $\mathcal A_{\kappa,J}$ and are treated by
Theorem~\ref{thm:main-source-energy-lift}; the main TV filters in
\eqref{eq:main-clock-coverage} are exactly the displayed restrictions.

For a Jensen mark let $S_{\lambda,a}:X_{\operatorname{type}(a)}\to\mathbb R$
be its measurable slope/value coordinate and put
$\mu_{\lambda,a}:=(S_{\lambda,a})_*\nu_{\lambda,a}$.
Whenever this row is used for balanced physical phase, this is not a free
measurable coordinate: $\mathrm{ACTUAL\_SLOPE\_BRIDGE}$ identifies it
pointwise with the logarithmic derivative ratio
\eqref{eq:actual-derivative-ratio-slope} on the same labelled carrier.
For every complete Jensen filter used by a phase call, its resonant centre
$z_{\xi,\kappa,\mathfrak d}$ is common to that filter; it does not vary with
$(\lambda,a)$.  Thus the supremum over $z$ below controls the actual lifted
carrier event.  On the carrier lift of every non-main Jensen filter require
the jointly Borel implication
\begin{equation}
\label{eq:nonmain-phase-bad-implication}
 \{(\lambda,a,x)\in\widehat{\mathcal A}_{\rm nm}^{\kappa,J}:
       |\psi'_{\xi,\lambda,a}(x)|\le\Xi r\}
 \subset
 \{|S_{\lambda,a}(x)-z_{\xi,\kappa,\mathfrak d}|
       \le Ch_\kappa(t)\}.
\end{equation}
Here the hat denotes the full carrier lift, as in the main coverage below.
A source filter may carry passage $J$ only when every phase-good carrier is
uniformly admissible for Lemma~\ref{lem:fractional-nonstationary-phase}:
$\psi_\xi\in C^{1+\alpha}$ with the declared derivative H\"older bound, the
positive amplitude has the required scaled $\mathcal A^\alpha$ norm, and
the physical endpoint ledger is the one in
\eqref{eq:phase-endpoint-ledger}.  It must also satisfy
$\mathrm{POST\_CARRIER\_DOM}$, so that both the phase-bad restriction and
the resulting positive carrier majorant are defined after propagation and
their cost is already in $q$.  A cemetery or extra registered type for
which this typed NSP structure is unavailable must be placed in a TV cell.
Thus a merely Borel slope coordinate never licenses the phase-good estimate.
Define the normalized
base, fibre-energy, value, and direct-TV masses
\begin{align}
 \mathcal B_{\kappa,J}(t)&:=M_*^{-1}
  \int_{\mathcal A_{\kappa,J}(t)}q_\lambda(da)\,d\Pi(\lambda),\notag\\
 \mathcal E_{\kappa,J}(t)&:=M_*^{-1}
  \int_{\mathcal A_{\kappa,J}(t)}
  (\mu_{\lambda,a}\otimes\mu_{\lambda,a})
  \{|u-v|\le2h_\kappa(t)\}\,q_\lambda(da)\,d\Pi(\lambda),\notag\\
 \mathcal V_{\kappa,J}(t)&:=M_*^{-1}\sup_z
  \int_{\mathcal A_{\kappa,J}(t)}
  \mu_{\lambda,a}([z-h_\kappa(t),z+h_\kappa(t)])
  \,q_\lambda(da)\,d\Pi(\lambda),\notag\\
 \mathcal D_{\kappa,\mathrm{TV}}(t)&:=M_*^{-1}
  \int_{\mathcal A_{\kappa,\mathrm{TV}}(t)}
       q_\lambda(da)\,d\Pi(\lambda).
 \label{eq:clock-source-partition}
\end{align}
Here $h_r(r)=r$.  The PPE window is tied to the very scale at which the
unselected pair energy is evaluated: with
\begin{equation}
\label{eq:ppe-clock-window-scale}
 r_N:=e^{-N/C_{\rm buf}},\qquad
 0<h_{\rm PPE}(N)\le C_{\rm win}r_N,\qquad
 N_{\rm PPE}(r)=\lceil C_{\rm buf}\log(1/r)\rceil .
\end{equation}
The time window $h_{\rm time}(N)$ is the declared window in the separate
unselected time-pair estimate, and in a central call its depth is
$N=N_{\rm time}=\max\{m,n\}$ by
\eqref{eq:time-clock-product-time-identity}.  Writing a cell as a function
only of $t=r$ or $N$ means its section at fixed suppressed $\mathfrak d$;
the sets, source kernels, slope maps, and all bounds are jointly measurable
and uniform in the suppressed mode, polynomial, time, and depth coordinates.

For nonempty Jensen classes require BASE--ENERGY bounds
\begin{align}
 \mathcal B_{r,J}(r)&\le Cr^{-a_{r,J}^{B}},&
 \mathcal E_{r,J}(r)&\le Cr^{\Theta_{r,J}^{E}},\notag\\
 \mathcal B_{\mathrm{PPE},J}(N)&\le Ce^{a_{\mathrm{PPE},J}^{B}N},&
 \mathcal E_{\mathrm{PPE},J}(N)&\le Ce^{-c_{\mathrm{PPE},J}^{E}N},\notag\\
 \mathcal B_{\mathrm{time},J}(N)&\le Ce^{a_{\mathrm{time},J}^{B}N},&
 \mathcal E_{\mathrm{time},J}(N)&\le Ce^{-c_{\mathrm{time},J}^{E}N}.
 \label{eq:clock-source-matrix-tails}
\end{align}
For the TV column require directly
\begin{equation}
\label{eq:clock-source-TV-tails}
 \mathcal D_{r,\mathrm{TV}}(r)\le Cr^{\Theta_{r,\mathrm{TV}}^{\rm src}},
 \quad
 \mathcal D_{\mathrm{PPE},\mathrm{TV}}(N)
     \le Ce^{-c_{\mathrm{PPE},\mathrm{TV}}^{\rm src}N},
 \quad
 \mathcal D_{\mathrm{time},\mathrm{TV}}(N)
     \le Ce^{-c_{\mathrm{time},\mathrm{TV}}^{\rm src}N}.
\end{equation}
The derived exceptional value rates are
\begin{equation}
\label{eq:exceptional-Jensen-value-rates}
 \Theta_{r,J}^{\rm exc,val}:=
   \frac{\Theta_{r,J}^{E}-a_{r,J}^{B}}2,
 \quad
 c_{\mathrm{PPE},J}^{\rm exc,val}:=
   \frac{c_{\mathrm{PPE},J}^{E}-a_{\mathrm{PPE},J}^{B}}2,
 \quad
 c_{\mathrm{time},J}^{\rm exc,val}:=
   \frac{c_{\mathrm{time},J}^{E}-a_{\mathrm{time},J}^{B}}2.
\end{equation}
Every displayed value gap must be strictly positive.  Equations
\eqref{eq:nonmain-clock-coverage}--
\eqref{eq:nonmain-clock-current-identity}, rather than prose alone, contain
cemetery, direct-TV, and retained exceptional remainders.  Context,
native-stop, buffer, gate, continuation, and adaptive-alphabet sources lie in the
appropriate Jensen cells; retained absolute PPE and contracted-coordinate
remainders in their TV cells.  The main retained Jensen source is not a
primitive row.  Word-depth sources use
$\mathrm{FACE\_TIME}_{\rm CM2}$ for product-time decay,
$\mathrm{FACE}_{\rm 2CUT}$ for growing-depth convergence, and the deep
ledger for the omitted tail.  An empty class has rate
$+\infty$ and disappears from minima.

Source-incidence complexity is the growth of the integrated kernel mass.
For example, an atomic raw energy rate $c$ and occurrence growth $a$ first
give an aggregate energy rate $c-a$; if the base mass also grows like
$e^{aN}$, the Jensen value rate is $(c-2a)/2$.  The TV column needs only
$c>a$.  Any nonpositive gap is a genuine no-go.  A H\"older transfer is
legal only under a normalized incidence-law $L^p$ bound; a complete-tree
branch moment is not a substitute.
\end{definition}

\begin{lemma}[One global incidence Cauchy]
\label{lem:global-incidence-cauchy}
For every finite marked Jensen cell, including the exceptional cells of
Definition~\ref{def:source-clock-matrix} and the main cells of
Definition~\ref{def:main-source-BEV},
\begin{equation}
\label{eq:global-incidence-cauchy}
 \mathcal V_{\kappa,J}(t)^2
 \le \mathcal B_{\kappa,J}(t)\mathcal E_{\kappa,J}(t).
\end{equation}
Consequently the rates in
\eqref{eq:exceptional-Jensen-value-rates} are valid.
\end{lemma}

\begin{proof}
Let $\rho(d\lambda,da)=M_*^{-1}
\one_{\mathcal A_{\kappa,J}(t)}q_\lambda(da)d\Pi(\lambda)$.
For $I=[z-h_\kappa(t),z+h_\kappa(t)]$, Cauchy on the \emph{whole}
incidence measure gives
\[
 \left(\int\mu_{\lambda,a}(I)d\rho\right)^2
 \le \rho(1)\int\mu_{\lambda,a}(I)^2d\rho
 \le \mathcal B_{\kappa,J}(t)\mathcal E_{\kappa,J}(t).
\]
Take the supremum in $z$.  No square root is taken before occurrences are
integrated, so occurrence multiplicity is paid exactly once through the base
and energy masses.
\end{proof}

\begin{remark}[Global probability transfer only]
\label{rem:global-holder-only}
If the unweighted complete-tree probability partitions are already
aggregated and
\begin{equation}
\label{eq:full-product-Lp-energy-route}
 \|\mathscr W_*\|_{L^{p_W}(\Pi_{s,N})}
 \le C\mathbb E_{\Pi_{s,N}}\mathscr W_*,\qquad p_W>1,
\end{equation}
then global H\"older transfers only the global
$\mathrm{CLOCK\_PROB}_{W}$ tails, with exponents divided by $p_W'$.  It
does not imply parentwise $\mathrm{TENERGY}_{W}$ or
$\mathrm{USC}_{\rm end}^{W}$, and it does not control
$\mathrm{CLOCK\_SRC}_{W}$ without the separate incidence-law moment just
described.
\end{remark}

\begin{lemma}[Homogeneity marks control physical endpoint scale]
\label{lem:endpoint-upper-scale}
Assume $\mathrm{USC}_{\rm end}^{W}(q_g^W)$, including the weighted excess
tail \eqref{eq:upper-scale-excess-tail}, on the same complete tree law.
On every retained labelled branch,
\begin{equation}
\label{eq:endpoint-upper-scale-growth}
 \log^+a_X\le C_{\rm init}+C_0m+2G_m^-+\log S_\lambda^-,
 \qquad
 \log^+a_Y\le C_{\rm init}+C_0n+2G_n^++\log S_\lambda^+.
\end{equation}
Fix $\zeta_g>0$ and put
\begin{equation}
\label{eq:scale-good-constants}
 C_{\rm scl}:=C_0+2(\gamma_g+\zeta_g),\qquad
 c_{\rm scl}:=q_g^W\zeta_g.
\end{equation}
For $N=\max\{m,n\}$ call a label scale-good when
\begin{equation}
\label{eq:weighted-scale-good-event}
 \widetilde G_m^-\le\gamma_gm+\zeta_gN,
 \qquad
 \widetilde G_n^+\le\gamma_gn+\zeta_gN.
\end{equation}
The scale-good survivor labels satisfy
$a_X+a_Y\le Ce^{C_{\rm scl}N}$, while the total marked mass of scale-bad
survivor labels, with the full weight \eqref{eq:global-stopped-source-weight}, is at most
$C\mathbb E_{\rm stop}[\mathscr W_{\lambda,*}]e^{-c_{\rm scl}N}$.  The latter
labels may therefore be estimated by weighted total variation and the
absolute Fourier coefficient sum, without invoking $\mathrm{PPE}_{N}$.
\end{lemma}

\begin{proof}
For a sequential billiard, the homogeneous-branch derivative estimate is
\cite{DLcones}
\[
 \frac{\|DT(x)v\|}{\|v\|}\le\frac{C}{\cos\varphi(Tx)},
 \qquad \cos\varphi\asymp k^{-2}\quad\hbox{on }\mathbb H_k.
\]
Multiplication along the dynamic part of the branch gives the $C_0m+2G_m^-$
and $C_0n+2G_n^+$ terms.  The initial physical-face parametrization,
same-branch section, scalar chart, and coarea/holonomy derivatives give the
two explicit $\log S_\lambda^\pm$ terms; they are not consequences of
compactness of the phase-space chart.  Thus
\eqref{eq:endpoint-upper-scale-growth} is equivalent, up to
$C_{\rm init}$, to a bound by $2\widetilde G_N^\pm$.  Apply the
\emph{weighted} tail \eqref{eq:upper-scale-excess-tail} at depths $m,n$ with
$u=\zeta_gN$ to the two orientations and use the union bound.  If the
parentwise sufficient route is used, Lemma
\ref{lem:parentwise-holder-transfer} is applied before this split and the
exponent is $q_g/p_W'$.  This
argument explains the geometric form of the interface; uniformity for the
stopped actual/split weighted laws remains the explicit content of
$\mathrm{USC}_{\rm end}^{W}$.  In particular it is not asserted for
arbitrary billiard invariant measures.
\end{proof}

\begin{lemma}[Tilted stopped gate, buffer, and native scale]
\label{lem:gate-buffer}
Assume $\mathrm{SS}(d_Z)$, $\mathrm{NST}_{\rm phys}$,
$\mathrm{PPE}_{N}(d_Z)$, and $\mathrm{TENERGY}_{W}$; if the displayed
tilted bounds are obtained from raw parentwise bounds rather than assumed
directly, also assume the corresponding $\mathrm{KRH}_{p_W}$ interface.
Let $\alpha_L<\alpha_B$ and put
\begin{equation}
\label{eq:native-energy-scale}
 \sigma=(r/c_\ell)^{1/(1+\alpha_L)},\qquad
 \Theta_{\rm ctx}=\frac{p(\alpha_B-\alpha_L)}{1+\alpha_L}.
\end{equation}
On the antichain $\mathrm{NST}_{\rm phys}(r)$ call a stopped context good
when
\begin{equation}
\label{eq:joint-native-good-gate}
 H_\beta\sigma_Q^{1+\alpha_B}\le r/10.
\end{equation}
Use the prescribed depth $N=N(r)$ of
\eqref{eq:prescribed-energy-depth}.  Then the error in
\eqref{eq:single-source-coefficient} is at most $Cr/10$ on good contexts,
while
\begin{equation}
\label{eq:bad-context-native-scale}
 \Pi_{s,N}^{W}
 \{H_\beta\sigma_Q^{1+\alpha_B}>r/10\}
 \le Cr^{\Theta_{\rm ctx}^{W}}.
\end{equation}
The physical resolving time is $N(r)=O(|\log r|)$ and all buffer
continuations are retained, so the refinement preserves conditional mass.
There are named constants
$c_{\rm buf}^{W},c_{\rm gate}^{W},c_{\rm cont}^{W}>0$ such that the
buffer, gate, and exceptional-continuation cemetery tails under the same
tilted tree are at most
\begin{equation}
\label{eq:named-energy-tails}
 C\left(e^{-c_{\rm buf}^{W}N(r)}+e^{-c_{\rm gate}^{W}N(r)}
 +e^{-c_{\rm cont}^{W}N(r)}\right).
\end{equation}
The exceptional native stopping mass under $\Pi_{s,N}^{W}$ is
$O(r^{\Theta_{\rm st}^{W}})$.
\end{lemma}

\begin{proof}
The first assertion is \eqref{eq:joint-native-good-gate} together with
$\kappa^{N(r)}\le r/100$.  The tilted context, stopping, and auxiliary
claims are precisely
\eqref{eq:tilted-context-stop-tails}--
\eqref{eq:tilted-auxiliary-tails}.  If they are derived from the raw
$\mathrm{NST}_{\rm phys}$ tails, Markov's inequality first gives the
unweighted exponent $\Theta_{\rm ctx}$ displayed above, and
$\mathrm{KRH}_{p_W}$ then degrades it, and every named parentwise
exponential rate, by $1/p_W'$.  The global estimate
\eqref{eq:full-product-Lp-energy-route} is insufficient for this conditional
claim.  No unweighted rate is retained after multiplication by
$\mathscr W_*$.
\end{proof}

\begin{proposition}[Tilted PPE]
\label{prop:physical-legal-child}
Assume $\mathrm{PPE}_{N}(d_Z)$ and $\mathrm{TENERGY}_{W}$.  Uniformly over
admissible context polynomials $P\in\mathbb R_{d_Z}[z]$, the retained part of
the complete tree obeys
\begin{equation}
\label{eq:local-clean-sublevel}
 \Pi_{s,N}^{W}
 \bigl(\mathsf G\cap\{|X^{[N(r)]}_{\beta,Q}-P|\le Cr\}\bigr)
 \le Cr^{\Theta_{\rm PPE}^{W}}.
\end{equation}
Equivalently, before normalization its full positive marked mass is at most
\[
 C\,\mathbb E_{\Pi_{s,N}}\mathscr W_*\,r^{\Theta_{\rm PPE}^{W}}.
\]
\end{proposition}

\begin{proof}
This is \eqref{eq:tilted-PPE-tail}.  The raw physical geometry is
\textup{(PPE1)--(PPE3)}, but the weighted conclusion is supplied by
$\mathrm{TENERGY}_{W}$, either directly or through the parent-kernel
transfer Lemma~\ref{lem:parentwise-holder-transfer}.  It is not a deduction from the
fixed clean subsystem.  Lemma
\ref{lem:negative-pressure-trial-obstruction} explains why the deleted
fixed-pressure trial argument could not provide this estimate.
\end{proof}

\begin{lemma}[Retained-tree normalization]
\label{lem:retained-tree-normalization}
Assume
\[
 \mathrm{TENERGY}_{W},\qquad \mathrm{CLOCK\_PROB}_{W}.
\]
Let the three active resolution coordinates
satisfy $r\downarrow0$, $N_{\rm PPE}=N_{\rm PPE}(r)\to\infty$, and either
the time registry is empty or
\begin{equation}
\label{eq:time-survivor-explicit-smallness}
 \mathfrak Z_{\rm time}^{\rm prob}(N_{\rm time})\le\frac1{12}.
\end{equation}
In a product-time call this is eventually automatic from
$N_{\rm time}=\max\{m,n\}\to\infty$; it is not a consequence of shrinking
$r$ while holding $N_{\rm time}$ fixed.
Let $\mathsf G$ be the retained-label event and set
$M_{\mathsf G}^{W}:=\Pi_{s,N}^{W}(\mathsf G)$.  For sufficiently small $r$,
uniformly in the actual/split parent law,
\begin{equation}
\label{eq:survivor-mass-lower-bound}
\begin{aligned}
 M_{\mathsf G}^{W}&\ge
  1-\Pi^W(B_{\rm ctx})-\Pi^W(B_{\rm stop})
   -\sum_{i\in\{\rm buf,gate,cont\}}\Pi^W(B_i)\\
 &\quad-\mathfrak Z_r^{\rm prob}(r)
   -\mathfrak Z_{\rm PPE}^{\rm prob}(N_{\rm PPE}(r))
   -\mathfrak Z_{\rm time}^{\rm prob}(N_{\rm time})
 \ge\frac12.
\end{aligned}
\end{equation}
Consequently the conditional survivor law is
$\Pi_{s,N}^{W}(\,\cdot\cap\mathsf G)/M_{\mathsf G}^{W}$ and every transfer
from the complete tree to that law costs at most a factor two.  Equivalently,
all energy estimates may be kept in subprobability form, in which case no
division occurs.
\end{lemma}

\begin{proof}
Use \eqref{eq:tilted-context-stop-tails},
\eqref{eq:tilted-auxiliary-tails}, and
\eqref{eq:clock-prob-tails}.  The named failures and the $r$/PPE probability
partition functions tend to zero at their declared scales; the time term is
controlled by the separate hypothesis
\eqref{eq:time-survivor-explicit-smallness}.  Shrink only the coordinates
that actually tend to their limiting regime so the total is at most $1/2$.
No source
incidence mass is subtracted here, and bad context/native stop are not
charged a second time through an additional registry.
\end{proof}

\begin{definition}[Pair-safe reference]
\label{def:pair-safe-reference}
Assume Definition~\ref{def:actual-slope-bridge} and the retained-tree
normalization above.  The law and energy clauses below are denoted,
respectively,
\[
 \mathrm{PAIR\_SAFE\_REFERENCE},\qquad
 \mathrm{PAIR\_TENERGY}.
\]
For each
legitimate frozen physical parent, before
any main-cell selection, let
\begin{equation}
\label{eq:unselected-reference-law}
 \mathsf R_0(d\eta,dx,dx')
 :=K^{\rm prob}(d\eta)\kappa_\eta(dx)\kappa_\eta(dx')
 :=\sum_\beta w_\beta^W\,
   \delta_\beta(d\eta)\mu_\beta^W(dx)\mu_\beta^W(dx')
\end{equation}
be the actual branch-mixture product law of that normalized unselected
survivor.  This identity is $\mathrm{PAIR\_SAFE\_REFERENCE}$; there is no
second reference law.  Its slope is fixed by
\begin{equation}
\label{eq:reference-slope-identification}
 S_\eta^0=S_\beta=S_{\lambda,m,n,s}=S_{\lambda,a}
 \quad \mathsf R_0\text{-a.e. on the corresponding labelled carrier},
\end{equation}
where the equalities include the carrier identifications in
Definition~\ref{def:actual-slope-bridge}.  Put
\begin{equation}
\label{eq:exact-reference-pair-event}
 D_h^0:=\{(\eta,x,x'):
     |S_\eta^0(x)-S_\eta^0(x')|\le2h\}.
\end{equation}
With
\begin{align}
 \Theta_{\rm par}^{W}&:=\min\{\Theta_{\rm PPE}^{W},
        \Theta_{\rm ctx}^{W},\Theta_{\rm st}^{W}\}>0,
 \label{eq:physical-energy-exponent}\\
 \mathcal T_{\rm prob}^{W}(h)&:=
 \sum_{i\in\{\rm buf,gate,cont\}}\Pi^W(B_i(N_{\rm PPE}(h)))
 +\mathfrak Z_r^{\rm prob}(h)
 +\mathfrak Z_{\rm PPE}^{\rm prob}(N_{\rm PPE}(h)),
 \label{eq:complete-energy-tail}
\end{align}
$\mathrm{PAIR\_TENERGY}$ requires the exact, parentwise bound
\begin{equation}
\label{eq:pair-tenergy-bound}
 \mathsf R_0(D_h^0)
 \le Ch^{\Theta_{\rm par}^{W}}+C\mathcal T_{\rm prob}^{W}(h),
 \qquad 0<h<h_0.
\end{equation}
It may be verified by the following factorized sufficient route: after all
retained/cemetery decisions and positive tilts, the conditional carrier law
must be the same $\kappa_\eta$ for both endpoints and, uniformly in
$x'$,
\begin{equation}
\label{eq:pair-tenergy-factorized-route}
 \kappa_\eta\{x:
 |S_\eta^0(x)-S_\eta^0(x')|\le2h\}
 \le Ch^{\Theta_{\rm par}^{W}}+C\mathcal T_{\rm prob}^{W}(h).
\end{equation}
Then Tonelli under the genuine product law proves
\eqref{eq:pair-tenergy-bound}.  A family of frozen-endpoint laws
$\kappa_{\eta,x'}$ depending on $x'$ does not prove this condition and may
not be integrated as if it were
$\kappa_\eta\otimes\kappa_\eta$; in that case
\eqref{eq:pair-tenergy-bound} remains an independent primitive.
\end{definition}

\begin{theorem}[Complete tilted physical single-source energy tree]
\label{thm:direct-slope-small-ball}
Assume $\mathrm{SS}(d_Z)$, $\mathrm{NST}_{\rm phys}$,
$\mathrm{PPE}_{N}(d_Z)$, $\mathrm{TENERGY}_{W}$,
$\mathrm{CLOCK\_PROB}_{W}$,
$\mathrm{ACTUAL\_SLOPE\_BRIDGE}$, the dominant clean-gate chart
normalization of Proposition~\ref{prop:exact-reduced-slope},
Lemma~\ref{lem:retained-tree-normalization}, and
$\mathrm{PAIR\_SAFE\_REFERENCE}$ together with
$\mathrm{PAIR\_TENERGY}$.  Let $w_\beta^{W}$ and $\mu_\beta^{W}$ be the
branch disintegration in \eqref{eq:unselected-reference-law}, and put
$\nu_\beta^{W}=(S_\beta)_*\mu_\beta^{W}$.  Then
one has
\begin{align}
 \sum_\beta w_\beta^{W}\iint
   \mathbf 1_{\{|S_\beta(t)-S_\beta(t')|\le2r\}}
   \,d\mu_\beta^{W}(t)d\mu_\beta^{W}(t')
 &\le Cr^{\Theta_{\rm par}^{W}}+C\mathcal T_{\rm prob}^{W}(r),\label{eq:conditional-energy-bound}\\
 \sup_z\sum_\beta w_\beta^{W}\nu_\beta^{W}([z-r,z+r])
 &\le Cr^{\Theta_{\rm par}^{W}/2}+C\mathcal T_{\rm prob}^{W}(r)^{1/2}.
 \label{eq:slope-value-small-ball}
\end{align}
\end{theorem}

\begin{proof}
The first assertion is exactly \eqref{eq:pair-tenergy-bound} for the law
\eqref{eq:unselected-reference-law}; it is not obtained by recombining
opposite-endpoint-dependent conditional laws.  If the factorized route is
available, integrate \eqref{eq:pair-tenergy-factorized-route} first in
$x'$ against the same $\kappa_\eta$, then in $\eta$.  Finally use
\[
 \nu_\beta^{W}(I)^2\le
 (\nu_\beta^{W}\otimes\nu_\beta^{W})\{|u-v|\le2|I|\}
\]
and Jensen in the positive weights $w_\beta^{W}$ to obtain
\eqref{eq:slope-value-small-ball}.  The conditional normalization and every
probability tail have already been paid in the exact law and
$\mathcal T_{\rm prob}^{W}$; no selected source law enters this theorem.
\end{proof}

\begin{definition}[Main clock coverage and BEV]
\label{def:main-source-BEV}
This definition states the resolution-exhaustive main-clock interface.  Put
\[
 I_{\rm main}:=I\cap\{o(a)={\rm main}\},
 \qquad
 \widehat I_{\rm main}:=
 \{(\lambda,a,x):(\lambda,a)\in I_{\rm main},\
                    x\in X_{\operatorname{type}(a)}\}.
\]
For every admissible resolution/mode datum
$\vartheta=(r,N_{\rm PPE}(r),N_{\rm time},\xi,P)$, viewed as a section of
the master datum $\mathfrak d$, there are jointly Borel,
pairwise disjoint filters for which, modulo one common
$\overline M_{s,\mathfrak d}^{q}$-null set,
\begin{equation}
\label{eq:main-clock-coverage}
\begin{aligned}
 I_{\rm main}
 &=G_{\rm main}(\vartheta)
 \sqcup\!\bigsqcup_{\substack{
      \kappa\in\{r,{\rm PPE},N_{\rm time}\}\\
      \pi\in\{J,{\rm TV}\}}}
      \mathcal A_{\rm main}^{\kappa,\pi}(\vartheta)\\
 &\qquad\sqcup I_{\rm main}^{\rm face/word}(\vartheta)
 \sqcup I_{\rm main}^{\rm deep}(\vartheta).
\end{aligned}
\end{equation}
If several clocks fire, a fixed measurable first-active ordering selects
exactly one filter; this ordering partitions additive occurrences and is not
the probability first-hit partition.  The handler table is part of the
interface: $G_{\rm main}$ is treated by nonstationary phase; every $J$ cell
by the selection-safe energy lift below; every TV cell by the corresponding
row of $\mathrm{CLOCK\_SRC}_{W}$; and the final two cells respectively by
$\mathrm{FACE\_TIME}_{\rm CM2}$ together with
$\mathrm{FACE}_{\rm 2CUT}$, and by
$\mathrm{TAIL}_{\rm env}^{+}/\mathrm{NREC}_{\rm deep}$.  In particular,
membership in $G_{\rm main}$ or in the phase-good complement of a main
$J$ cell additionally requires exactly the typed NSP admissibility imposed
on non-main $J$ rows: on the same carrier
$\psi_\xi\in C^{1+\alpha}$ with the declared H\"older derivative bound, the
positive amplitude has the scaled $\mathcal A^\alpha$ norm used by
Lemma~\ref{lem:fractional-nonstationary-phase}, and the endpoint contribution
is charged in \eqref{eq:phase-endpoint-ledger}.  It also requires
$\mathrm{ACTUAL\_SLOPE\_BRIDGE}$ and
$\mathrm{POST\_CARRIER\_DOM}$ on that carrier.  A main occurrence lacking
any one of these properties must be assigned to the corresponding TV,
face/word, or deep handler; a lower derivative bound by itself does not
license nonstationary phase.  In particular,
write $\widehat A:=\{(\lambda,a,x):(\lambda,a)\in A,\
x\in X_{\operatorname{type}(a)}\}$ for the carrier lift of a mark filter.
The actual phase-bad set is the jointly Borel carrier event
\[
 \widehat B_{\rm phase}(\vartheta):=
 \{(\lambda,a,x)\in\widehat I_{\rm main}:
   |\psi'_{\xi,\lambda,a}(x)|\le \Xi r\}.
\]
On $\widehat G_{\rm main}$ the complementary derivative lower bound holds.
For each Jensen filter the interface supplies a measurable centre
$z_{\xi,\kappa,\vartheta}$ common to that complete filter and the fibrewise
implication
\[
 \widehat B_{\rm phase}(\vartheta)\cap
 \widehat{\mathcal A}_{\rm main}^{\kappa,J}
 \subset
 \{|S_{\lambda,a}(x)-z_{\xi,\kappa,\vartheta}|
          \le C h_\kappa(t)\}.
\]
The remaining phase-bad carrier points obey the typed inclusion
\begin{equation}
\label{eq:main-phase-bad-coverage}
 \widehat B_{\rm phase}(\vartheta)
 \subset
 \bigsqcup_{\kappa}
   \left(\widehat{\mathcal A}_{\rm main}^{\kappa,J}
       \cap\{|S-z_{\xi,\kappa,\vartheta}|\le Ch_\kappa(t)\}\right)
 \sqcup\bigsqcup_{\kappa}
   \widehat{\mathcal A}_{\rm main}^{\kappa,\rm TV}
 \sqcup\widehat I_{\rm main}^{\rm face/word}
 \sqcup\widehat I_{\rm main}^{\rm deep}.
\end{equation}
Restriction of the one propagated positive envelope gives the absolutely
convergent dual identity
\begin{equation}
\label{eq:main-clock-current-identity}
 J_{\rm main}=J_{G_{\rm main}}
 +\sum_{\kappa,\pi}J_{\rm main}^{\kappa,\pi}
 +J_{\rm main}^{\rm face/word}+J_{\rm main}^{\rm deep}.
\end{equation}

For a Jensen filter $A=\mathcal A_{\rm main}^{\kappa,J}(\vartheta)$,
$\kappa\in\{r,{\rm PPE},N_{\rm time}\}$, use the same slope laws
$\mu_{\lambda,a}$ and $M_*=\mathbb E\mathscr W_*$ as above and define
\begin{align}
 \mathcal B_{{\rm main},\kappa}(t)&:=M_*^{-1}
  \int_{\mathcal A_{\rm main}^{\kappa,J}(t)}q_\lambda(da)d\Pi(\lambda),
  \notag\\
 \mathcal E_{{\rm main},\kappa}(t)&:=M_*^{-1}
  \int_{\mathcal A_{\rm main}^{\kappa,J}(t)}
  (\mu_{\lambda,a}\otimes\mu_{\lambda,a})
  \{|u-v|\le2h_\kappa(t)\}\,q_\lambda(da)d\Pi(\lambda),
  \notag\\
 \mathcal V_{{\rm main},\kappa}(t)&:=M_*^{-1}\sup_z
  \int_{\mathcal A_{\rm main}^{\kappa,J}(t)}
  \mu_{\lambda,a}([z-h_\kappa(t),z+h_\kappa(t)])
  \,q_\lambda(da)d\Pi(\lambda).
 \label{eq:main-source-BEV}
\end{align}
These quantities use the operator- and recovery-charged $q$, are finite by
the declared bounds below, and are defined on the incidence graph rather
than by summing currents over labels.  Identity
\eqref{eq:main-clock-coverage}, not the static owner identity alone, is the
coverage statement used by the phase proof.  Suppressing $\xi,P$ or unused
depth entries in $A(t)$ always means a section at frozen $\vartheta$; all
estimates are uniform in those entries.
\end{definition}

\begin{definition}[Selection-safe joint source likelihood
$\mathrm{SEL\_JSL}_{p}$]
\label{def:joint-source-likelihood}
Fix a legitimate frozen physical parent, but do not select a main cell.
Use the exact $K^{\rm prob}$, $\kappa_\eta$, product law $\mathsf R_0$,
slope $S_\eta^0$, and event $D_h^0$ from
\eqref{eq:unselected-reference-law}--\eqref{eq:exact-reference-pair-event}.
They are the objects fixed by $\mathrm{PAIR\_SAFE\_REFERENCE}$ and estimated
by $\mathrm{PAIR\_TENERGY}$, not newly declared look-alikes.  There must exist one probability proposal kernel
$\zeta_\eta(x,x';da)$, measurable with respect to the preselection physical
parent only, with $\zeta_\eta(x,x';Y_{\rm src})=1$ pointwise.  It is fixed
before every resolution/mode section $\vartheta$ and before every cell $A$;
symbolically $\exists\zeta\ \forall(\vartheta,A)$.  Put
\begin{equation}
\label{eq:proposal-lift}
 \overline{\mathsf R}(d\eta,da,dx,dx')
 :=\mathsf R_0(d\eta,dx,dx')\zeta_\eta(x,x';da).
\end{equation}
For the $a$-independent reference resonance event $D_h^0$ in
\eqref{eq:exact-reference-pair-event}, normalization gives the exact
mass-preserving identity
\begin{equation}
\label{eq:proposal-preserves-reference-energy}
\overline{\mathsf R}(D_h^0\times Y_{\rm src})=\mathsf R_0(D_h^0).
\end{equation}
On a same-carrier cell put
$\overline{\mathsf R}^{\rm src}:=\overline{\mathsf R}$; on a
different-carrier cell put
$\overline{\mathsf R}^{\rm src}:=\overline{\mathsf R}^{H}$, where the typed
pushforward law and event inclusion are those in
\eqref{eq:carrier-match-pushforward-law}--\eqref{eq:carrier-match-slope-window}.

For a nonzero main Jensen filter $A$ of
Definition~\ref{def:main-source-BEV}, define the \emph{unnormalized} source
pair measure
\begin{align}
 d\Sigma_A(\eta,a,x,x')
 &:=M_*^{-1}\one_A(\eta,a)
    d\Pi(\eta)q_\eta(da)
    \nu_{\eta,a}(dx)\nu_{\eta,a}(dx').
 \label{eq:unnormalized-selected-source-pair}
\end{align}
Define the actual source pair event by
\begin{equation}
\label{eq:exact-source-pair-event}
 D_h^{\rm src}:=\{(\eta,a,x,x'):
 |S_{\eta,a}(x)-S_{\eta,a}(x')|\le2h\}.
\end{equation}
Thus $\Sigma_A(1)=\mathcal B_A$ and, at the clock scale
$h=h_\kappa(t)$, $\Sigma_A(D_h^{\rm src})=\mathcal E_A(h)$.  In the PPE
row this is exactly the window in \eqref{eq:ppe-clock-window-scale}; it may
not be replaced by a larger independently declared window.  The interface
requires
\begin{equation}
\label{eq:joint-source-likelihood}
\Sigma_A\ll\overline{\mathsf R}^{\rm src},
 \qquad Z_A:=\frac{d\Sigma_A}{d\overline{\mathsf R}^{\rm src}},
 \qquad p\in(1,\infty].
\end{equation}
The proposal support must cover the source support, and the source event must
satisfy the typed measurable compatibility
\[
 \begin{cases}
 D_h^{\rm src}\subset D_{C_0h}^{0}\times Y_{\rm src},
   &\text{on a same-carrier cell},\\
 \mathcal H^{-1}D_h^{\rm src}\subset
   D_{C_0h}^{0}\times Y_{\rm src},
   &\text{on a different-carrier cell}.
 \end{cases}
\]
Both inclusions are understood on the corresponding support of $\Sigma_A$.
For the PPE row they hold with
$h=h_{\rm PPE}(N)\le C_{\rm win}e^{-N/C_{\rm buf}}$ and hence transfer into
$D_{C e^{-N/C_{\rm buf}}}^{0}$, the same scale controlled by
$\mathrm{PAIR\_TENERGY}$.
Uniformly in all
frozen data, the scale-dependent likelihood bounds are
\begin{equation}
\label{eq:selection-safe-likelihood-rates}
 \|Z_{A,r}\|_{L^p(\overline{\mathsf R}^{\rm src})}
 \le Cr^{-a_{{\rm main},r}^{L}},\qquad
 \|Z_{A,\kappa}\|_{L^p(\overline{\mathsf R}^{\rm src})}
 \le Ce^{a_{{\rm main},\kappa}^{L}N},
 \quad \kappa\in\{{\rm PPE},N_{\rm time}\}.
\end{equation}
The selector $\one_A$, the $q$ reweighting, the reciprocal proposal cost,
and all carrier-density costs occur in this single $Z_A$; neither
$\mathsf R_0$ nor $\zeta$ may depend on $A$.  Consequently
\begin{equation}
\label{eq:joint-source-energy-transfer}
 \mathcal E_A(h)=\Sigma_A(D_h^{\rm src})
 \le \|Z_A\|_{L^p(\overline{\mathsf R}^{\rm src})}
       \mathsf R_0(D_{C_0h}^0)^{1/p'},
 \qquad \frac1p+\frac1{p'}=1,
 \quad (p'=1\text{ when }p=\infty).
\end{equation}
Because $\overline{\mathsf R}^{\rm src}$ is a probability,
$\mathcal B_A=\|Z_A\|_1\le\|Z_A\|_p$; the declared likelihood and base
growth rates must be compatible with this elementary check.  A proposal
missing source support or chosen after seeing $A$ is an immediate no-go.
This is the unique change of law at the kernel where selection and carrier
reweighting occur.  Parent $\mathrm{KRH}_p$ contains none of
$\one_A$, $1/\zeta$, or the two internal carrier coordinates and cannot
replace this interface, in agreement with \cite{NiKernel}.
\end{definition}

\begin{lemma}[Selection factor criterion and typed carrier match]
\label{lem:joint-source-factor-criterion}
Suppose first that the proposal is carrier-independent,
$\zeta_\eta(x,x';da)=\zeta_\eta(da)$, and that source and probability
sections use the same carrier.  Define
\begin{equation}
\label{eq:joint-source-likelihood-factorization}
 w_A^{\rm occ}:=
 \frac{d\{M_*^{-1}\one_A(\eta,a)\Pi(d\eta)q_\eta(da)\}}
      {d\{K^{\rm prob}(d\eta)\zeta_\eta(da)\}},
 \qquad
 Z_A=w_A^{\rm occ}(\eta,a)
      \ell_{\eta,a}(x)\ell_{\eta,a}(x'),
 \quad
 \ell_{\eta,a}:=\frac{d\nu_{\eta,a}}{d\kappa_\eta}.
\end{equation}
Thus $w_A^{\rm occ}$ visibly contains the selector, propagated $q$ weight,
and reciprocal proposal cost.  If $w_A^{\rm occ}\in L^{p_o}$ uniformly and
$\sup_{\eta,a}\|\ell_{\eta,a}\|_{L^{p_c}(\kappa_\eta)}<\infty$, then direct
disintegration gives $Z_A\in L^p(\overline{\mathsf R})$ for every
$1<p\le\min\{p_o,p_c\}$.  In particular, when $p_c=\infty$, the endpoint is
available without an arbitrary loss:
\begin{equation}
\label{eq:same-carrier-holder-endpoint}
 p_{\rm joint}=p_o,
 \qquad
 \|Z_A\|_{L^{p_o}(\overline{\mathsf R})}
 \le \|w_A^{\rm occ}\|_{L^{p_o}}
       \sup_{\eta,a}\|\ell_{\eta,a}\|_\infty^2.
\end{equation}
If $\nu_{\eta,a}$ and $\kappa_{\eta}$ are normalized densities in one
uniform log-H\"older cone on the same homogeneous carrier, the cone bounds
\[
 e^{-\varpi}|W|^{-1}\le\rho\le
 C_{\rm cone}e^{\varpi}|W|^{-1}
\]
from the standard-family construction of \cite{CanestrariHoles} give
$p_c=\infty$.  Hence only the occurrence/selection likelihood is paid.  If
that factor is also bounded one may take $p_{\rm joint}=\infty$.

For different carriers, $\mathrm{CARRIER\_MATCH}$ must instead supply a
jointly measurable $H_{\eta,a}:W_{\rm prob}\to W_{\rm src}$ and the actual
pushforward reference law
\begin{equation}
\label{eq:carrier-match-pushforward-law}
\begin{aligned}
 \mathcal H(\eta,a,x,x')&:=
  (\eta,a,H_{\eta,a}x,H_{\eta,a}x'),\\
 \overline{\mathsf R}^{H}:=\mathcal H_*\overline{\mathsf R}
 &=K^{\rm prob}(d\eta)\zeta_\eta(da)
   (H_{\eta,a*}\kappa_\eta)(dy)
   (H_{\eta,a*}\kappa_\eta)(dy').
\end{aligned}
\end{equation}
The entire kernel $H$ is measurable with respect to the same preselection
parent data as $\zeta$ and is fixed before every resolution/mode section and
every cell: symbolically $\exists H\ \forall(\vartheta,A)$.  An
$A$-dependent holonomy is forbidden because it would hide the selected
conditioning inside the reference carrier.
It requires
\begin{equation}
\label{eq:carrier-match-density-ratio}
 \nu_{\eta,a}\ll H_{\eta,a*}\kappa_\eta,
 \qquad \ell_{\eta,a}^{H}:=
 \frac{d\nu_{\eta,a}}{d(H_{\eta,a*}\kappa_\eta)},
 \qquad
 \sup_{\eta,a}\|\ell_{\eta,a}^{H}\|_{
 L^{p_c}(H_{\eta,a*}\kappa_\eta)}\le C,
\end{equation}
and, for a carrier-independent proposal, the exact likelihood
factorization under the pushforward law
\begin{equation}
\label{eq:carrier-match-likelihood-factorization}
 Z_A^{H}(\eta,a,y,y')
 =w_A^{\rm occ}(\eta,a)
   \ell_{\eta,a}^{H}(y)\ell_{\eta,a}^{H}(y').
\end{equation}
On a different-carrier cell the $Z_A$ in
\eqref{eq:joint-source-likelihood} is, by definition, this $Z_A^{H}$.
If $w_A^{\rm occ}\in L^{p_o}$ and the density ratio has the displayed
uniform $L^{p_c}$ bound, then
$Z_A^{H}\in L^p(\overline{\mathsf R}^{H})$ for every
$1<p\le\min\{p_o,p_c\}$; when $p_c=\infty$ the endpoint $p=p_o$ holds with
the analogue of \eqref{eq:same-carrier-holder-endpoint}.
The interface also requires the target-scale compatibility
\begin{equation}
\label{eq:carrier-match-slope-window}
 |S_{\rm src}(\eta,a,H_{\eta,a}x)-S_{\rm prob}(\eta,x)|
 \le C_0h_\kappa(t),
 \quad
 \mathcal H^{-1}D_h^{\rm src}\subset D_{C_1h}^{0},
\end{equation}
with $h_\kappa(t)$ equal to the actual $r$-, PPE-, or time-window scale.
A Jacobian bound relative to the two normalized arclength measures is only a
geometric sufficient ingredient for \eqref{eq:carrier-match-density-ratio};
it does not replace the Radon--Nikodym condition.  If the proposal depends on
$(x,x')$ and destroys the product disintegration, one must verify the full
$Z_A=d\Sigma_A/d\overline{\mathsf R}^{H}$ bound directly.  Without either
the same-carrier or this typed pushforward alternative, the lift is a no-go.
\end{lemma}

\begin{proof}
Integrating $|Z_A|^p$ first in $x,x'$ proves the conditional factor claim;
\eqref{eq:same-carrier-holder-endpoint} is the same computation at
$p=p_o$ with bounded carrier factors.  The cone claim is the ratio of its
uniform upper and lower density bounds.  For different carriers,
\eqref{eq:carrier-match-pushforward-law} defines the law on which the density
ratio in \eqref{eq:carrier-match-density-ratio} lives, and
\eqref{eq:carrier-match-likelihood-factorization} followed by the same
disintegration proves the required selection-safe $L^p$ bound.  Finally,
\eqref{eq:carrier-match-slope-window} transfers the source event to an
enlarged unselected reference event.  These are separate requirements.
\end{proof}

\begin{theorem}[Main-source energy lift
$\mathrm{MAIN\_SOURCE\_ENERGY\_LIFT}$]
\label{thm:main-source-energy-lift}
Assume $\mathrm{ACTUAL\_SLOPE\_BRIDGE}$,
$\mathrm{PAIR\_SAFE\_REFERENCE}$, $\mathrm{PAIR\_TENERGY}$, the parentwise probability energy estimate of
Theorem~\ref{thm:direct-slope-small-ball}, and
$\mathrm{SEL\_JSL}_{p_{\rm joint}}$ on every nonempty main $J$ filter.  Put
$p_{\rm joint}'=p_{\rm joint}/(p_{\rm joint}-1)$, with
$p_{\rm joint}'=1$ in the bounded-likelihood case.  Compress the complete
reference energy theorem into the explicit rates
\begin{align}
 \Theta_{{\rm ref},r}:=\min\{&\Theta_{\rm par}^{W},
  \Theta_r^{\rm prob},C_{\rm buf}c_{\rm PPE}^{\rm prob},
  C_{\rm buf}c_{\rm buf}^{W},C_{\rm buf}c_{\rm gate}^{W},
  C_{\rm buf}c_{\rm cont}^{W}\},\notag\\
 c_{{\rm ref},{\rm PPE}}:=\min\{&\Theta_{\rm par}^{W}/C_{\rm buf},
  \Theta_r^{\rm prob}/C_{\rm buf},c_{\rm PPE}^{\rm prob},
  c_{\rm buf}^{W},c_{\rm gate}^{W},c_{\rm cont}^{W}\}.
 \label{eq:explicit-reference-energy-rates}
\end{align}
The second line is the first line at the declared relation
$r=r_N=e^{-N/C_{\rm buf}}$; by
\eqref{eq:ppe-clock-window-scale} the actual PPE source event transfers into
an enlarged reference event at this very scale.  Every component is already a parentwise
tilted probability rate; any likelihood used to obtain those probability
rates is paid there.  For a nonempty main $N_{\rm time}$ Jensen cell, require
the corresponding \emph{unselected} reference energy rate
\begin{equation}
\label{eq:unselected-time-reference-energy}
 D_{{\rm time},N}^{0}:=D_{h_{\rm time}(N)}^0,
 \qquad \mathsf R_0(D_{{\rm time},N}^{0})
 \le Ce^{-c_{{\rm ref},{\rm time}}N},
 \qquad c_{{\rm ref},{\rm time}}>0.
\end{equation}
This rate is supplied by the actual unselected probability theorem/clock and
is not inferred from the owner or main-clock partition.  The unique source
selection likelihood is paid exactly once through $p_{\rm joint}'$ below.
Suppose
\begin{align}
\label{eq:main-source-base-mass}
 \mathcal B_{{\rm main},r}(r)&\le Cr^{-a_{{\rm main},r}^{B}},\notag\\
 \mathcal B_{{\rm main},{\rm PPE}}(N)
 &\le Ce^{a_{{\rm main},{\rm PPE}}^{B}N},\qquad
 \mathcal B_{{\rm main},{\rm time}}(N)
 \le Ce^{a_{{\rm main},{\rm time}}^{B}N}.
\end{align}
Together with \eqref{eq:selection-safe-likelihood-rates},
\eqref{eq:joint-source-energy-transfer} gives the \emph{unnormalized}
source-energy bounds directly:
\begin{align}
 \mathcal E_{{\rm main},r}(r)&\le
 Cr^{\Theta_{{\rm ref},r}/p_{\rm joint}'-a_{{\rm main},r}^{L}},\notag\\
 \mathcal E_{{\rm main},{\rm PPE}}(N)&\le
 Ce^{-(c_{{\rm ref},{\rm PPE}}/p_{\rm joint}'
       -a_{{\rm main},{\rm PPE}}^{L})N},\notag\\
 \mathcal E_{{\rm main},{\rm time}}(N)&\le
 Ce^{-(c_{{\rm ref},{\rm time}}/p_{\rm joint}'
       -a_{{\rm main},{\rm time}}^{L})N}.
 \label{eq:source-energy-bound}
\end{align}
One application of Lemma~\ref{lem:global-incidence-cauchy} yields
\begin{align}
 \mathcal V_{{\rm main},r}(r)&\le
 Cr^{\Theta_{{\rm main},r}^{\rm val}},&
 \Theta_{{\rm main},r}^{\rm val}&:=
 \frac{\Theta_{{\rm ref},r}/p_{\rm joint}'
             -a_{{\rm main},r}^{L}-a_{{\rm main},r}^{B}}2,\notag\\
 \mathcal V_{{\rm main},{\rm PPE}}(N)&\le
 Ce^{-c_{{\rm main},{\rm PPE}}^{\rm val}N},&
 c_{{\rm main},{\rm PPE}}^{\rm val}&:=
 \frac{c_{{\rm ref},{\rm PPE}}/p_{\rm joint}'
             -a_{{\rm main},{\rm PPE}}^{L}
             -a_{{\rm main},{\rm PPE}}^{B}}2,\notag\\
 \mathcal V_{{\rm main},{\rm time}}(N)&\le
 Ce^{-c_{{\rm main},{\rm time}}^{\rm val}N},&
 c_{{\rm main},{\rm time}}^{\rm val}&:=
 \frac{c_{{\rm ref},{\rm time}}/p_{\rm joint}'
             -a_{{\rm main},{\rm time}}^{L}
             -a_{{\rm main},{\rm time}}^{B}}2.
 \label{eq:source-value-small-ball}
\end{align}
Every nonempty value gap must be positive.  If instead one starts from the
normalized source law $\mathsf S_A=\Sigma_A/\mathcal B_A$ with
$\|d\mathsf S_A/d\overline{\mathsf R}^{\rm src}\|_p\le Ch^{-d_A}$, then
$a_A^{L}=a_A^{B}+d_A$ and the polynomial value rate becomes
\begin{equation}
\label{eq:normalized-selection-rate-special-case}
 \frac12\left(\frac{\Theta_{{\rm ref},r}}{p_{\rm joint}'}
                 -2a_A^{B}-d_A\right).
\end{equation}
Thus the old two-base-cost expression is only the special case $d_A=0$;
it is not inherited after selection.  Define the complete Jensen value rates
\begin{equation}
\label{eq:complete-Jensen-value-rates}
\begin{aligned}
 \Theta_{r,J}^{\rm val}&:=\min\{
   \Theta_{{\rm main},r}^{\rm val},
   \Theta_{r,J}^{\rm exc,val}\},\\
 c_{{\rm PPE},J}^{\rm val}&:=\min\{
   c_{{\rm main},{\rm PPE}}^{\rm val},
   c_{{\rm PPE},J}^{\rm exc,val}\},\\
 c_{{\rm time},J}^{\rm val}&:=\min\{
   c_{{\rm main},{\rm time}}^{\rm val},
   c_{{\rm time},J}^{\rm exc,val}\}.
\end{aligned}
\end{equation}
An absent class has rate $+\infty$.  The main marked value tail is therefore
derived from unselected reference probability energy, the single
selection/source likelihood, finite base mass, and one global incidence
Cauchy step; it is not
assumed in
$\mathrm{CLOCK\_SRC}_{W}$.
\end{theorem}

\begin{proof}
Theorem~\ref{thm:direct-slope-small-ball} on the unselected law
$\mathsf R_0$ gives the $r$/PPE rates
\eqref{eq:explicit-reference-energy-rates}; the separate hypothesis
\eqref{eq:unselected-time-reference-energy} supplies the time rate when that
class is nonempty.  The pointwise normalization of $\zeta$ preserves these reference
energies because $D_h^0$ is $a$-independent, and the source/reference
inclusions use the exact events
\eqref{eq:exact-reference-pair-event} and
\eqref{eq:exact-source-pair-event}.  Apply H\"older once to the unnormalized $\Sigma_A$ against
$\overline{\mathsf R}^{\rm src}$; the single density $Z_A$ simultaneously pays the
selector, propagated occurrence weight, proposal, and both carrier
coordinates, giving \eqref{eq:source-energy-bound}.  Apply
\eqref{eq:global-incidence-cauchy} once to each complete main clock cell.
The likelihood and base costs $a^L+a^B$ are paid before the square root.  No
selected reference law, occurrencewise Jensen, or primitive main source tail
is used.
\end{proof}

\begin{corollary}[Central slope small-ball]
\label{cor:central-g2}
Assume all hypotheses of
Theorem~\ref{thm:direct-slope-small-ball} and of
Theorem~\ref{thm:main-source-energy-lift}.  In particular assume
$\mathrm{PAIR\_SAFE\_REFERENCE}$, $\mathrm{PAIR\_TENERGY}$, and the
$\mathrm{ACTUAL\_SLOPE\_BRIDGE}$.
The complete central collision--SRB quotient, normalized by its positive
event weight, satisfies
\eqref{eq:slope-value-small-ball} and the main marked-source bounds
\eqref{eq:source-value-small-ball} at $s=0$.  At finite $s$ the constants are
uniform precisely for the stopped scales whose total resolving depth is at
most $L(s)$.  Proposition~\ref{prop:central-phase-verification} invokes this
finite-$s$ statement only in its bounded intermediate frequency range.
\end{corollary}

\begin{proof}
Apply Theorem~\ref{thm:direct-slope-small-ball} and
Theorem~\ref{thm:main-source-energy-lift}.  The latter already estimates the
unnormalized selected source measure.  Multiplying the common $M_*$
normalization back gives the complete marked phase mass; no selected
reference law, signed cancellation, occurrencewise square root, or
unweighted Jensen step is used.
\end{proof}

\begin{lemma}[Slope small-ball controls the balanced phase]
\label{lem:slope-good-bad-phase}
Assume
\[
 \mathrm{ACTUAL\_SLOPE\_BRIDGE},\qquad
 \mathrm{CLOCK\_SRC}_{W},\qquad
 \mathrm{MAIN\_CLOCK\_COVERAGE},
\]
Theorem~\ref{thm:main-source-energy-lift}, and the exceptional source rates
\eqref{eq:exceptional-Jensen-value-rates} on all nonempty displayed rows.
For $\psi_\xi=\xi_1X+\xi_2Y$ on an affine-normalized gate chart, let
\begin{equation}
\label{eq:balanced-physical-frequency}
 \Xi=|\xi_1|a_X\asymp|\xi_2|a_Y,
\end{equation}
with $a_X,a_Y$ as in \eqref{eq:dominant-coordinate-normalization}.  On the
restricted $r$/PPE Jensen subcurrent
\[
 J_{r/{\rm PPE},J}^{\rm cen}:=
 \sum_{\kappa\in\{r,{\rm PPE}\}}
 \left(J_{\rm nm}^{\kappa,J}+J_{\rm main}^{\kappa,J}\right),
\]
the value mass of
$\{|\psi_\xi'|\le\Xi r\}$ is at most
$C\{r^{\Theta_{r,J}^{\rm val}}
+e^{-c_{\rm PPE,J}^{\rm val}N_{\rm PPE}(r)}\}$.  This is the additive
source value mass derived by global incidence Cauchy, not the probability of
its support.  Smooth value-space cutoffs on its complement have
$C^\alpha$ norm $O(r^{-1})$ and satisfy
$\inf|\psi_\xi'|\ge c\Xi r$.
\end{lemma}

\begin{proof}
By the exact identity \eqref{eq:actual-balanced-phase-derivative} in
$\mathrm{ACTUAL\_SLOPE\_BRIDGE}$,
\[
 \psi_\xi'=X'(\xi_1+\varepsilon_{\rm or}\xi_2e^S).
\]
The lower and upper bounds in
\eqref{eq:dominant-coordinate-normalization} convert the bracketed factor to
the physical scale \eqref{eq:balanced-physical-frequency}.  Only
opposite-oriented coefficients can resonate, and the resonant set is
contained in one interval $|S-z_\xi|\le Cr$.  Restrict
\eqref{eq:nonmain-clock-current-identity} and
\eqref{eq:main-clock-current-identity} to the two displayed Jensen filters,
and use the carrier-level implications
\eqref{eq:nonmain-phase-bad-implication} and
\eqref{eq:main-phase-bad-coverage}.  Then apply
\eqref{eq:source-value-small-ball}, combine it with the exceptional rates in
\eqref{eq:complete-Jensen-value-rates}, and use
\eqref{eq:marked-source-assembled-pairing}.  Composing a fixed smooth three-term
partition with
$(\xi_1+\varepsilon_{\rm or}\xi_2e^S)/(\Xi r)$ gives the cutoff and norm
bounds.  The $N_{\rm time}$, TV, face/word, and deep restrictions are not
claimed by this lemma; they are split off and handled by their typed rows in
Proposition~\ref{prop:central-phase-verification} and the two cutoff/tail
theorems.
\end{proof}
\begin{lemma}[Typed global double projection before disintegration]
\label{lem:global-label-recentering}
Let $J$ be the complete finite signed response measure on $M\times M$ in
common physical coordinates, let $m_J^\pm$ be its marginals, and let
$\nu^-=\nu^+=\mu_0$ be the fixed invariant reference supplied by
\eqref{eq:measure-trivializing-gauge}.  Define the
current-side projection
\begin{equation}
\label{eq:labelwise-recentering}
 \mathsf D_{\rm cur}J=:J^\circ
 =J-m_J^-\otimes\nu^+-\nu^-\otimes m_J^+
   +J(1,1)\nu^-\otimes\nu^+.
\end{equation}
For a bounded test $F$ define $\mathsf D_{\rm test}F$ by subtracting its two
integrated marginals and adding its product mean.  Then
\[
 \int F\,d(\mathsf D_{\rm cur}J)=
 \int \mathsf D_{\rm test}F\,dJ.
\]
Thus $J^\circ$ is double centered and the current projection is the adjoint,
not the same operator, as test centering.  If the complete response current
already has zero marginals, then
\begin{equation}
\label{eq:global-recentering-sum}
 J^\circ=J.
\end{equation}
Only after this identity is formed in physical coordinates is the singular
part refined by face labels.  If the original physical amplitudes obey
\eqref{eq:global-ledger-sum} and the references have fixed $C^\alpha$
densities, then
\begin{equation}
\label{eq:recenter-norm-ledger}
 \|\mathsf D_{\rm cur}J\|_{\rm TV}
 \le4\|J\|_{\rm TV},\qquad
 \|J\|_{\rm phys,R}
 +\|m_J^-\otimes\mu_0+\mu_0\otimes m_J^+
 -J(1,1)\mu_0\otimes\mu_0\|_{\mathfrak P_R}<\infty .
\end{equation}
\end{lemma}

\begin{proof}
Fubini proves the adjoint identity and each marginal has total variation at
most $\|J\|_{\rm TV}$.  This gives the factor four.  Marginalization is
contractive by \eqref{eq:bounded-face-marginal-maps}, and tensoring with the
fixed probability $\mu_0$ gives exactly the separate product type
\eqref{eq:product-current-type}.  These operations do not preserve dimension
and hence are not assigned the graph-face norm.  The positive face ledger
dominates signed total variation before recombination; absolute summation of
\eqref{eq:global-ledger-sum} and
\eqref{eq:bounded-face-marginal-maps} prove the second assertion.
\end{proof}

\begin{lemma}[Endpoint unfolding and cutoff commutation]
\label{lem:endpoint-unfolding}
Let a localized assembled face at parameter $s$ be
$J_{\lambda,s}=(H_{\lambda,s})_*(a_{\lambda,s}(u)\,du)$ with
$H_{\lambda,s}(u)=(x_{\lambda,s}(u),y_{\lambda,s}(u))$.  Put
\begin{equation}
\label{eq:physical-endpoint-maps}
 X^-_{\lambda,m,s}(u):=T_s^{-m}x_{\lambda,s}(u),\qquad
 Y^+_{\lambda,n,s}(u):=T_s^{n}y_{\lambda,s}(u).
\end{equation}
Let $g^\circ=g-\mu_0(g)$ and $f^\circ=f-\mu_0(f)$.  After the complete
current is assembled and the global current-side double projection is taken,
\begin{equation}
\label{eq:exact-endpoint-unfolding}
 \int (Q_s^mg)(x)(Q_s^{*n}f)(y)\,
       d\mathsf D_{\rm cur}\!\left(\sum_\lambda J_{\lambda,s}\right)
 =\sum_\lambda\int a_{\lambda,s}(u)
   g^\circ(X^-_{\lambda,m,s}(u))
   f^\circ(Y^+_{\lambda,n,s}(u))\,du.
\end{equation}
The two marginal and constant product blocks are therefore assembled before
the face sum is estimated.  This is the exact \emph{physical endpoint}
stage.  It is distinct from the following scalarization stage.

Let $\mathsf X_{\lambda,m,s}(u)$ and
$\mathsf Y_{\lambda,n,s}(u)$ be the scalar base coordinates of the
same-branch sections in $\mathrm{BR}_{\rm sec}$.  Branchwise contraction
replaces the two endpoint factors by
\begin{equation}
\label{eq:endpoint-to-section-stage}
 g^\circ\!\left(\iota^-_\lambda(\mathsf X_{\lambda,m,s}(u))\right),
 \qquad
 f^\circ\!\left(\iota^+_\lambda(\mathsf Y_{\lambda,n,s}(u))\right),
\end{equation}
with the explicit error \eqref{eq:branch-section-contraction}.  Only now do
we expand the localized restrictions of the \emph{original centered}
functions $g^\circ\circ\iota^-_\lambda$ and
$f^\circ\circ\iota^+_\lambda$.  The scaled composition clause in
$\mathrm{BR}_{\rm sec}$ gives their $C^r$ Fourier decay with the recorded
$\mathfrak d_\lambda$ and Mark costs, uniformly in $m,n,s$.  A scalar mode
has phase
\begin{equation}
\label{eq:unfolded-endpoint-phase}
 \psi_{k,l,\lambda,m,n,s}(u)
 =k\mathsf X_{\lambda,m,s}(u)
  +l\mathsf Y_{\lambda,n,s}(u).
\end{equation}
Thus the Fourier coefficients belong to the fixed original observables;
time dependence lies in the scalar endpoint phase only.  Equation
\eqref{eq:exact-endpoint-unfolding} and
\eqref{eq:unfolded-endpoint-phase} are not asserted to be the same identity:
the controlled section-replacement error lies between them.
The homogeneous refinement cuts only the amplitude and endpoint phase.  Its
Jacobian, cutoff derivatives, and genuine singularity endpoints are included
in the physical Mark--$Z$ ledger.  Applying the same scalar multiplier to the
two endpoint coordinates before or after the global double projection gives
the same joint, two axis, and zero blocks.  Hence the low multiplier and its
high complement are disjoint and exhaustive after assembly.
\end{lemma}

\begin{proof}
Equations~\eqref{eq:perron-composition-convention}--
\eqref{eq:centered-endpoint-iterates} and the adjoint identity of
Lemma~\ref{lem:global-label-recentering} give the displayed formula.  The
axis and constant terms cancel only after the complete physical current is
assembled; no face is separately asserted to have zero marginals.  Insert
the fixed physical partition of unity, apply
\eqref{eq:branch-section-contraction}, and only then Fourier-expand the
localized restrictions of the original observables.  No derivative is applied to $Q_s^mg$ or
$Q_s^{*n}f$.  Fubini and the even scalar multiplier show that Fourier cutoff
and the four-term double projection commute.  The last assertion follows
because every boundary created by the homogeneous partition is either paired
with its adjacent cell or is a genuine endpoint recorded in the ledger.
\end{proof}

\begin{definition}[Branch-section interface $\mathrm{BR}_{\rm sec}$]
\label{def:branch-section-interface}
On every clean graph rectangle, the dominant endpoint coordinates admit
branchwise sections
\[
 \iota^-_\lambda:X\mapsto(X,\zeta^-_\lambda(X)),\qquad
 \iota^+_\lambda:Y\mapsto(Y,\zeta^+_\lambda(Y))
\]
inside the same homogeneous inverse branches as the physical graph trace.
Their images, in the appropriate forward and reversed orientations, are
homogeneous measured standard curves of class $C^{r_{\rm obs}}$.  For a
fixed $q_{\rm end}\ge0$, their scalar composition maps satisfy
\begin{align}
 \left|\widehat{\chi^-_\lambda(g\circ\iota^-_\lambda)}(k)\right|
 &\le C\mathfrak d_\lambda^{q_{\rm end}}
       \|g\|_{C^{r_{\rm obs}}}
       (1+|k|)^{-r_{\rm obs}},\notag\\
 \left|\widehat{\chi^+_\lambda(f\circ\iota^+_\lambda)}(l)\right|
 &\le C\mathfrak d_\lambda^{q_{\rm end}}
       \|f\|_{C^{r_{\rm obs}}}
       (1+|l|)^{-r_{\rm obs}}.
 \label{eq:section-original-observable-fourier}
\end{align}
The powers of $\mathfrak d_\lambda$ are charged to $A_{\rm dom}$; cutoff,
Jacobian, and face-amplitude costs already form the single squared mark in
\eqref{eq:event-physical-norm} and are not inserted a second time through the
Fourier coefficients.  Their arclength/coordinate densities,
endpoint $Z$ values, and scaled H\"older marks are bounded by
$\operatorname{Mark}(\lambda)\{1+Z(\mathcal G^-_\lambda)+
Z(\mathcal G^+_\lambda)\}$.  The label includes the endpoint depths and the
parameter.  With the unfolded endpoint maps of
\eqref{eq:physical-endpoint-maps}, branchwise contraction gives
\begin{align}
 \int|g(X^-_{\lambda,M,s}(\bar s))-
 g(\iota^-_\lambda(\mathsf X_{\lambda,M,s}(\bar s)))|\,d|\omega_\lambda|
 &\le Ce^{-c_-M}\|g\|_{C^1}|\omega_\lambda|,\notag\\
 \int|f(Y^+_{\lambda,N,s}(\bar s))-
 f(\iota^+_\lambda(\mathsf Y_{\lambda,N,s}(\bar s)))|\,d|\omega_\lambda|
 &\le Ce^{-c_+N}\|f\|_{C^1}|\omega_\lambda|,
 \label{eq:branch-section-contraction}
\end{align}
where $d\omega_\lambda=M_\lambda w_\lambda d\bar s$.
This is a same-branch condition: no conditional measure is averaged across
different symbolic pasts or singularity cuts.
\end{definition}

\begin{lemma}[Dominant-coordinate branch-section reduction]
\label{lem:dominant-coordinate-reduction}
Assume $\mathrm{BR}_{\rm sec}$.  After assembly and the recentering of Lemma
\ref{lem:global-label-recentering}, retain the rank-one and regular-product
blocks as ordinary product kernels and evolve both factors by their centered
operators.
On a clean graph rectangle in the diagonal band
$D^{-1}M<N<AM$, the graph-current pairing decomposes as
\begin{equation}
\label{eq:dominant-coordinate-decomposition}
 B_{M,N}=B^{\rm dom,dom}_{M,N}
 +B^{\rm dom,rem}_{M,N}+B^{\rm rem,dom}_{M,N}
 +B^{\rm rem,rem}_{M,N}.
\end{equation}
The first term is the scalar quotient integral obtained by evaluating the two
tests on $\iota^-_\lambda(X)$ and $\iota^+_\lambda(Y)$.
Every other term is bounded by
\begin{equation}
\label{eq:contracted-coordinate-error}
 Ce^{-c\max\{M,N\}}\|g\|_{C^1}\|f\|_{C^1}.
\end{equation}
This bound remains valid after multiplication by the complete positive event
weight and integration over $\Pi_{s,N}^{W}$.  The displayed remainder is a
genuine $N_{\rm time}=\max\{M,N\}$ source occurrence in the time--TV cell
of $\mathrm{CLOCK\_SRC}_{W}$.  Any label on which uniform contraction is
unavailable is separately owned according to the stage where failure is
first observed and entered in exactly one clock--passage cell.  It is never
assigned an untyped exponential rate.  Thus no unweighted exceptional tail
or clock conversion is hidden in the contracted-coordinate remainder.
The global joint constant and axis blocks vanish or remain in the explicitly
separated rank-one summands before the graph-rectangle sum is estimated.
Every graph-axis mode created by a boundary cutoff is uniformly
nonstationary and is charged to its cutoff norm.
\end{lemma}

\begin{proof}
In stable--unstable endpoint coordinates, use the same-branch sections of
$\mathrm{BR}_{\rm sec}$ and telescope the two replacements.
Equation~\eqref{eq:branch-section-contraction} gives remainders bounded by
$C\lambda_-^M\|g\|_{C^1}$ and $C\lambda_+^N\|f\|_{C^1}$ after the physical
amplitude is integrated.  Disintegrating first under the positive complete
tilt preserves the same uniform factors; the named exceptional complement is
charged by its corresponding typed clock in
\eqref{eq:clock-source-matrix-tails} or
\eqref{eq:clock-source-TV-tails}.  In the diagonal band these imply
\eqref{eq:contracted-coordinate-error}.  The two frozen functions depend
only on $X,Y$, proving \eqref{eq:dominant-coordinate-decomposition}.  Global
projection separates the constant and two axis blocks before any
label-dependent chart is chosen.  They are estimated by their explicit
product formulas, not redescribed as graph currents; Lemma
\ref{lem:rank-one-product} supplies product decay.  A graph boundary
cutoff can reintroduce an axis term, but it has phase $\xi_1X$ or $\xi_2Y$;
\eqref{eq:dominant-coordinate-normalization} gives the physical derivative
lower bound, and its cutoff ledger is absolutely summable by
\eqref{eq:recenter-norm-ledger}.
\end{proof}

\begin{proposition}[Scaled physical face density]
\label{prop:scaled-face-density}
On a homogeneous physical face cell
\(Q=t_Q+\ell_Q[-1,1]\), the collision--SRB conditional is
\(d\mu_Q^u=\rho_Q(t)\,dt\), where, before normalization,
\begin{equation}
\label{eq:face-conditional-density}
 \rho_Q^{\rm raw}(t)=c_\mu\cos\varphi(t)\,J_{\rm coa}(t)
 J_{\rm hol}(t)J_{\rm chart}(t),
 \qquad J_{\rm coa}=|\nabla G|^{-1}.
\end{equation}
Put
\(\widehat\rho_Q(u)=\ell_Q\rho_Q(t_Q+\ell_Qu)\).  There is a mark
\begin{equation}
\label{eq:density-mark-definition}
 \operatorname{Mark}(Q)=
 1+\mathfrak h_Q+\mathfrak b_Q+
 [\log J_{\rm hol}]_{\alpha,Q}
 +[\log J_{\rm coa}]_{\alpha,Q}
 +[\log J_{\rm chart}]_{\alpha,Q}
 +\|w_Q\|_{\mathcal A^\alpha}
\end{equation}
such that
\begin{equation}
\label{eq:scaled-density-bound}
 C^{-1}\operatorname{Mark}(Q)^{-1}
 \le\widehat\rho_Q\le C\operatorname{Mark}(Q),\qquad
 [\log\widehat\rho_Q]_\alpha
 \le C\operatorname{Mark}(Q).
\end{equation}
For some \(p_d>2\),
\begin{equation}
\label{eq:density-mark-moment}
 \sum_Q M_Q\operatorname{Mark}(Q)^{p_d}<\infty
\end{equation}
uniformly in the small table family.  The same statement holds after the
positive/negative decomposition of a signed face amplitude.
\end{proposition}

\begin{proof}
In collision coordinates the invariant density is
\(c_\mu\cos\varphi\,dr\,d\varphi\).  Coarea along the physical graph, followed
by stable projection to the gate coordinate, gives
\eqref{eq:face-conditional-density}.  On a homogeneous strip,
\(\cos\varphi\), the chart Jacobian, and the holonomy Jacobian have the
standard scaled log-Hölder distortion.  Dividing by
\(\int_Q\rho_Q^{\rm raw}\) gives \eqref{eq:scaled-density-bound}; the factor
\(\ell_Q\) removes the cell scale.

The homogeneity contribution has a fixed polynomial cost, while the strip
mass is exponentially summable.  The holonomy and chart contributions are
summed by the context moment, and the amplitude contribution by the marked
face ledger.  Generalized Hölder with slightly smaller exponents proves
\eqref{eq:density-mark-moment}.  For a signed amplitude \(q\), write
\(q=(2\|q\|_\infty+q)-2\|q\|_\infty\); multiplication by the positive masses
restores the scaled norm and introduces no reciprocal small-mass factor.
\end{proof}

\begin{lemma}[Fractional nonstationary phase]
\label{lem:fractional-nonstationary-phase}
Let \(I\) be an interval, \(a\in C^\alpha(I)\), and
\(\phi\in C^{1+\alpha}(I)\).  On each component of the support of \(a\),
suppose \(\phi'\) has one sign and \(|\phi'|\ge\lambda>0\).  Put
\(\Phi=\sum_C\int_C|\phi'(t)|\,dt\), the total phase length of the support.
For $C=(l_C,r_C)$ define the endpoint functional
$\mathfrak E(a)=\sum_C(|a(l_C+)|+|a(r_C-)|)$.
\begin{equation}
\label{eq:fractional-nsp}
 \left|\int_Ie^{i\phi(t)}a(t)\,dt\right|
 \le C_\alpha\left\{
 \frac{\mathfrak E(a)}{\lambda}
 +\frac{\Phi[a]_\alpha}{\lambda^{1+\alpha}}
 +\frac{\Phi\|a\|_\infty[\phi']_\alpha}
             {\lambda^{2+\alpha}}\right\}.
\end{equation}
The estimate adds over a countable union of sign components without a
component-count loss.
\end{lemma}

\begin{proof}
On one sign component use \(y=\phi(t)\) and write
\(b(y)=a(\phi^{-1}y)/\phi'(\phi^{-1}y)\).  Then
\[
 \|b\|_\infty\le\|a\|_\infty/\lambda,\qquad
 [b]_\alpha\le
 [a]_\alpha\lambda^{-1-\alpha}
 +\|a\|_\infty[\phi']_\alpha\lambda^{-2-\alpha}.
\]
For a shift \(h=\pi\), the identity
\[
 2\int e^{iy}b(y)\,dy
 =\int e^{iy}\{b(y)-b(y+h)\}\,dy+\text{endpoint terms}
\]
gives $C\Phi[b]_\alpha$ plus the two endpoint values of $b$ on that
component.  Since $|b|\le |a|/\lambda$, summing the endpoint values gives
$\mathfrak E(a)/\lambda$, while the sum of the component $y$-lengths is
$\Phi$.  This proves \eqref{eq:fractional-nsp}; no number-of-components
estimate is hidden in the supremum term.
\end{proof}

\begin{lemma}[Rank-one product under double centering]
\label{lem:rank-one-product}
If \(\eta_\pm\) are admissible dual collision currents and
\(R=\eta_-\otimes\eta_+\), then
\[
 R(Q_-^mg\otimes Q_+^nf)
 =\eta_-(Q_-^mg)\eta_+(Q_+^nf).
\]
Thus the two one-sided gaps give \(C\tau^m\vartheta^n\).
\end{lemma}

\begin{proof}
This is the universal property of the projective tensor product followed by
the two admissible dual bounds.  In the principal proof the rank-one
corrections are instead included in the low-frequency functional of
Proposition~\ref{prop:functional-correlation-low-frequency}; this lemma is
only an algebraic consistency check.
\end{proof}

\begin{lemma}[Balanced--unbalanced exhaustion in physical frequency]
\label{lem:balanced-unbalanced-exhaustion}
Assume the clean-gate chart bounds
\begin{equation}
\label{eq:primitive-dominant-chart-bounds}
 c_0a_X\le |X'|\le C_0a_X,\qquad
 c_0a_Y\le |Y'|\le C_0a_Y
\end{equation}
with the fixed orientation signs on one homogeneous face cell.  Put
$C_{\rm bal}=4C_0/c_0$.  Every nonzero scalar mode belongs to exactly one of
the following regions (with the boundary assigned once):
\begin{enumerate}[label=\textup{(\roman*)}]
\item $C_{\rm bal}^{-1}\le
 |\xi_1|a_X/(|\xi_2|a_Y)\le C_{\rm bal}$ (balanced);
\item $|\xi_1|a_X>C_{\rm bal}|\xi_2|a_Y$ (past-dominant);
\item $|\xi_2|a_Y>C_{\rm bal}|\xi_1|a_X$ (future-dominant).
\end{enumerate}
On either unbalanced region the phase
$\psi_\xi=\xi_1X+\xi_2Y$ has one-sign derivative after at most the fixed
orientation cut and
\begin{equation}
\label{eq:unbalanced-derivative-lower-bound}
 |\psi_\xi'|\ge c\max\{|\xi_1|a_X,|\xi_2|a_Y\}.
\end{equation}
Consequently the complete unbalanced contribution is bounded by
\begin{equation}
\label{eq:unbalanced-mode-saving}
 C\operatorname{Mark}(Q)^2\{1+Z(\mathcal G_Q)\}
 \max\{|\xi_1|a_X,|\xi_2|a_Y\}^{-\alpha},
\end{equation}
and its two-dimensional lattice sum converges under the same observable
regularity used for the balanced sum.
\end{lemma}

\begin{proof}
Consider the past-dominant case.  The orientation table and
\eqref{eq:primitive-dominant-chart-bounds} give
\[
 |\psi_\xi'|
 \ge c_0|\xi_1|a_X-C_0|\xi_2|a_Y
 \ge \tfrac12c_0|\xi_1|a_X.
\]
The other case is symmetric.  Cut only at the fixed orientation endpoints;
no frequency-dependent sign partition is introduced.  Apply Lemma
\ref{lem:fractional-nonstationary-phase} with this lower bound and with the
actual amplitude and endpoint ledger.  Its three terms are bounded by the
right side of \eqref{eq:unbalanced-mode-saving}; the exponent $\alpha$ may be
decreased once to absorb fixed chart losses.  By
Lemma~\ref{lem:endpoint-unfolding}, multiplication is by the Fourier
coefficients of the two localized original $C^{r_{\rm obs}}$ observables,
with the scaled restriction cost in
\eqref{eq:section-original-observable-fourier}, not by coefficients of
evolved tests.  Summation over the two
unbalanced cones is finite as soon as $r_{\rm obs}>2$; the stronger balanced
threshold imposed below therefore suffices.  The three regions are disjoint
and exhaustive.
\end{proof}

\begin{proposition}[Central balanced estimate with explicit budget]
\label{prop:central-phase-verification}
Assume
\[
 \begin{gathered}
 \mathrm{ACTUAL\_SLOPE\_BRIDGE},\quad
 \mathrm{PAIR\_SAFE\_REFERENCE},\\
 \mathrm{PAIR\_TENERGY},\quad
 \mathrm{CLOCK\_SRC}_{W},\quad
 \mathrm{MAIN\_CLOCK\_COVERAGE},\\
 \mathrm{SEL\_JSL}_{p_{\rm joint}},\quad
 \mathrm{BR}_{\rm sec},\quad
 \mathrm{USC}_{\rm end}^{W}(q_g^W),\\
 \mathrm{POST\_CARRIER\_DOM},\quad
 \mathrm{PHYS\_TEST\_LIFT},\quad
 \mathrm{PROP\_Q\_MATCH},
 \end{gathered}
\]
the hypotheses of Theorems~\ref{thm:direct-slope-small-ball} and
\ref{thm:main-source-energy-lift}, and the typed NSP admissibility in both
the non-main and main clock definitions,
as well as all hypotheses of
Proposition~\ref{prop:functional-correlation-low-frequency} and
Lemmas~\ref{lem:regular-low-projector} and
\ref{lem:regular-high-tail} for the low block,
with the main-source likelihood and base bounds
\eqref{eq:selection-safe-likelihood-rates} and
\eqref{eq:main-source-base-mass}.  For every nonempty main time--Jensen cell
also assume \eqref{eq:unselected-time-reference-energy} on the same frozen
parent.  Use
the full weighted upper-scale tail \eqref{eq:upper-scale-excess-tail}, the
joint Mark--$Z$ estimate \eqref{eq:joint-mark-Z-ledger}, and the global
physical ledger \eqref{eq:global-ledger-sum} on the same projectively fixed
incidence law.  Use
the complete additive source rates
\eqref{eq:clock-source-matrix-tails},
\eqref{eq:clock-source-TV-tails}, and
\eqref{eq:complete-Jensen-value-rates}.  Set
\[
 A_{\rm frac}=3+\alpha,\qquad
 0<\delta<\frac{\alpha}{2A_{\rm frac}},
\]
\begin{equation}
\label{eq:complete-phase-exponent}
\begin{aligned}
 \beta_{\rm PPE,J}^{\rm val}&:=
   \delta C_{\rm buf}c_{\rm PPE,J}^{\rm val},
 &\beta_{\rm PPE,TV}^{\rm src}&:=
   \delta C_{\rm buf}c_{\rm PPE,TV}^{\rm src},\\
 \beta_{\rm phase}^{W}&:=
 \min\left\{
 \begin{gathered}
  \delta\Theta_{r,J}^{\rm val},
  \delta\Theta_{r,\rm TV}^{\rm src},
  \alpha-A_{\rm frac}\delta,\\
  \beta_{\rm PPE,J}^{\rm val},
  \beta_{\rm PPE,TV}^{\rm src}
 \end{gathered}
 \right\}>0.
\end{aligned}
\end{equation}
Put $N=N_{\rm time}=\max\{m,n\}$ as required by
\eqref{eq:time-clock-product-time-identity}, and choose
\begin{equation}
\label{eq:three-frequency-cutoffs}
 K_N=e^{\eta_{\rm lo}N},\qquad
 H_N=e^{\eta_{\rm hi}N},\qquad
 0<\eta_{\rm lo}<\eta_{\rm hi}.
\end{equation}
For a scalar mode put
$\Lambda=\max\{1,|k|,|l|\}$ and retain
$\Xi=\max\{|k|a_X,|l|a_Y\}$.  Every balanced mode in the intermediate range
\begin{equation}
\label{eq:intermediate-phase-range}
 \Xi>K_N,\qquad \Lambda\le H_N,
\end{equation}
is first split as follows.  Split every Jensen carrier into its phase-bad
fibre restriction and its complement.  Let $J_{\rm ph}^{\rm clock}$ be the
sum of $J_{G_{\rm main}}$, all phase-good Jensen restrictions, and the
$r$/PPE phase-bad Jensen and TV restrictions, with main and non-main parts
both included.  Let $J_{\rm time}^{\rm bad/TV}$ be the $N_{\rm time}$
phase-bad Jensen restriction plus the complete time--TV restriction.  These
two currents, together with the separately owned face/word and deep
currents, give the exact restriction obtained from the global owner identity
and the two clock identities.  Every
balanced mode of $J_{\rm ph}^{\rm clock}$ in
\eqref{eq:intermediate-phase-range} has saving
\(C\Xi^{-\beta_{\rm phase}^{W}}\) times a fixed power of
\(\operatorname{Mark}(Q)\) on scale-good labels.  The scale-bad label mass is
$O(e^{-c_{\rm scl}N})$ and is estimated without phase or PPE.  Modes with
$\Lambda>H_N$ are likewise estimated without
$\mathrm{PPE}_{N}$ and have total bound
\begin{equation}
\label{eq:ultra-high-observable-tail}
 C H_N^{-\rho_\infty}\|g\|_{C^{r_{\rm obs}}}
 \|f\|_{C^{r_{\rm obs}}},\qquad
 \rho_\infty:=r_{\rm obs}-2>0.
\end{equation}
If $r_{\rm obs}>2+\beta_{\rm phase}^{W}$, the complete high lattice sum is
exponentially small in the two induced depths.  Here all modes are those of the original
localized observables furnished by Lemma~\ref{lem:endpoint-unfolding}; the
induced depths occur only in the endpoint phase.  Identity
\eqref{eq:section-quotient-low-high} assigns every localized quotient mode
exactly once.  Low graph modes, rank-one corrections, and
the regular product block are treated together by
Proposition~\ref{prop:functional-correlation-low-frequency}.  If
$a_{\rm ultra}\ge0$ denotes the explicitly recorded time-complexity exponent
outside the original-observable tail (in the present organization
$a_{\rm ultra}=0$), the required ultra-high inequality is
\begin{equation}
\label{eq:ultra-high-exponent-window}
 \eta_{\rm hi}\rho_\infty>a_{\rm ultra}.
\end{equation}
The original-observable lattice sum of the remaining time subcurrent obeys
\begin{equation}
\label{eq:time-clock-complete-lattice-bound}
 \sum_{k,l}|\widehat g_k\widehat f_l|
 \left|\langle J_{\rm time}^{\rm bad/TV},\Phi_{k,l}\rangle\right|
 \le C e^{-\min\{c_{\rm time,J}^{\rm val},
                     c_{\rm time,TV}^{\rm src}\}N}
       \|g\|_{C^{r_{\rm obs}}}\|f\|_{C^{r_{\rm obs}}}.
\end{equation}
This proposition makes no estimate of
$J_{\rm main}^{\rm face/word}+J_{\rm face/word}$ or
$J_{\rm main}^{\rm deep}+J_{\rm deep}$; those exact restrictions are routed
to $\mathrm{FACE\_TIME}_{\rm CM2}/\mathrm{FACE}_{\rm 2CUT}$ and
$\mathrm{TAIL}_{\rm env}^{+}/\mathrm{NREC}_{\rm deep}$.
\end{proposition}

\begin{proof}
For a mode in \eqref{eq:intermediate-phase-range}, take
\(r=\Xi^{-\delta}\).  The energy theorem bounds the resonant mass by
\begin{equation}
\label{eq:phase-energy-all-tails}
 C\Xi^{-\delta\Theta_{r,J}^{\rm val}}
 +C\Xi^{-\delta C_{\rm buf}c_{\rm PPE,J}^{\rm val}}.
\end{equation}
The owner identity \eqref{eq:assembled-owner-current-identity}, followed by
identities \eqref{eq:nonmain-clock-current-identity} and
\eqref{eq:main-clock-current-identity}, together with the lifted inclusion
\eqref{eq:main-phase-bad-coverage}, first give exactly the split stated in
the proposition.  The face/word and deep restrictions are removed before
this phase estimate.  Then
$N_{\rm PPE}(r)=C_{\rm buf}\log(1/r)+O(1)$; no time-depth symbol is
substituted for this resolving depth.  The $r$- and PPE-depth source
occurrences in the TV column are outside Jensen and contribute, by weighted
absolute summation,
\begin{equation}
\label{eq:phase-PPE-TV-tails}
 C\Xi^{-\delta\Theta_{r,\rm TV}^{\rm src}}
 +C\Xi^{-\delta C_{\rm buf}c_{\rm PPE,TV}^{\rm src}};
\end{equation}
this is the origin of $\beta_{\rm PPE,TV}^{\rm src}$ and has no square-root
loss.  The same support may contribute several occurrences, all of which are
present in \eqref{eq:clock-source-partition}.
The $N_{\rm time}$ phase-bad Jensen restriction is bounded by
$c_{\rm time,J}^{\rm val}$ and the time--TV restriction by
$c_{\rm time,TV}^{\rm src}$.  Since the coefficients are those of the
original localized observables, their absolute lattice sum is finite;
weighted absolute summation proves
\eqref{eq:time-clock-complete-lattice-bound}.  On every phase-good Jensen
carrier use two smooth
cutoffs in
\(\psi_\xi'/(\Xi r)\), one for each sign.  They vanish on the resonant
boundaries.  Sum adjacent physical cells before taking absolute values, so
their artificial endpoint terms cancel with opposite orientation.  The
remaining face, singularity, and homogeneity endpoints obey
\begin{equation}
\label{eq:phase-endpoint-ledger}
 \sum_Q\mathfrak E(a_Q)
 \le Cr^{-1}\sum_QM_Q\operatorname{Mark}(Q)
       \{1+Z(\mathcal G_Q)\},
\end{equation}
by Proposition~\ref{prop:radial-direct-Z}.  Proposition
\ref{prop:scaled-face-density} gives
\[
 \|a_\pm\|_\infty+[a_\pm]_\alpha
 \le C\operatorname{Mark}(Q) r^{-1}.
\]
Furthermore
\[
 \lambda=c\Xi r,\qquad
 \Phi\le C\Xi,\qquad
 [\psi_\xi']_\alpha\le C\Xi\operatorname{Mark}(Q).
\]
Lemma~\ref{lem:fractional-nonstationary-phase} therefore gives
\begin{align*}
 |\mathcal I_Q(\xi)|
 &\le C\operatorname{Mark}(Q)^2
 \{(1+Z(\mathcal G_Q))\Xi^{-1}r^{-2}
        +\Xi^{-\alpha}r^{-(3+\alpha)}\}\\
 &\le C\operatorname{Mark}(Q)^2(1+Z(\mathcal G_Q))
       \Xi^{-\beta_{\rm phase}^{W}},
\end{align*}
where \eqref{eq:phase-energy-all-tails} and
\eqref{eq:phase-PPE-TV-tails} are included through the two typed PPE rates
in \eqref{eq:complete-phase-exponent}.
The joint estimate \eqref{eq:joint-mark-Z-ledger} sums the cells without
separating the density and boundary moments.  The two restriction estimates
\eqref{eq:section-original-observable-fourier} add at most
$\mathfrak d_\lambda^{2q_{\rm end}}$, which is absorbed by
$A_{\rm dom}\ge2+2q_{\rm end}$ in the global ledger.

First restrict to the scale-good labels of
Lemma~\ref{lem:endpoint-upper-scale}.  Then
$a_X+a_Y\le Ce^{C_{\rm scl}N}$ and $\Lambda\le H_N$, hence
$\Xi\le Ce^{(C_{\rm scl}+\eta_{\rm hi})N}$.  A mode in
\eqref{eq:intermediate-phase-range} therefore calls the prescribed escape
estimate only to depth
\begin{equation}
\label{eq:intermediate-max-resolving-depth}
 N_{\rm PPE}\le C_{\rm buf}\delta
       \log\!\left(Ce^{(C_{\rm scl}+\eta_{\rm hi})N}\right)
 \le \{C_{\rm buf}\delta(C_{\rm scl}+\eta_{\rm hi})+o(1)\}N.
\end{equation}
Scale-bad labels are never sent through PPE: total variation, the absolute
Fourier coefficient sum, and the tilted estimate
\eqref{eq:upper-scale-excess-tail} applied to the complete weight
$\mathscr W_{\lambda,*}$ give $Ce^{-c_{\rm scl}N}$.  Neither the unweighted
scale-bad probability nor the mere finiteness of
$\sum_\lambda M_\lambda\mathscr W_{\lambda,*}$ is used for this step.
For $\Lambda>H_N$ use no phase cutoff and no physical escape.  Sum the two
original-observable coefficient estimates
\eqref{eq:section-original-observable-fourier} directly.  After decreasing
the exponent once, a dyadic shell of radius $R$ in the two scalar modes has
$O(R^2)$ terms of total size at most $CR^{2-r_{\rm obs}}$.  The shells
$R>H_N$ therefore give $CH_N^{-(r_{\rm obs}-2)}$, proving
\eqref{eq:ultra-high-observable-tail}.  The fixed inverse-scale and cutoff
costs remain in $A_{\rm dom}$; there is no frequency-dependent physical mark
in this tail.

For completeness, put $K=K_N=e^{\eta_{\rm lo}N}$ with
$\eta_{\rm lo}$ in the window
\eqref{eq:delayed-recovery-eta-window}.  After global centering and the
dominant-coordinate reduction, the face-adapted projector
$\mathsf S_K^{\rm dom}$ retains precisely the scalar modes satisfying
$|\xi_1|a_X\le K$ and $|\xi_2|a_Y\le K$.  Proposition
\ref{prop:functional-correlation-low-frequency} treats this finite sum by
delayed recovery and separate label tails.  Every balanced mode of the complementary dominant
sum has $\Xi>K$ by definition, so the preceding saving is
$e^{-\eta_{\rm lo}\beta_{\rm phase}^{W}N}$ in the intermediate range.  The
ultra-high range contributes
$e^{-(\eta_{\rm hi}\rho_\infty-a_{\rm ultra})N}$ by
\eqref{eq:ultra-high-observable-tail}.  Contracted-coordinate terms were removed before
the cutoff and are exponentially small by Lemma
\ref{lem:dominant-coordinate-reduction}.  The difference between the graph
trace and its same-branch section is already part of that contracted error;
no fixed-fiber average is used.  The two unbalanced scalar cones are
estimated, with their full lattice sum, by Lemma
\ref{lem:balanced-unbalanced-exhaustion}.  The two-dimensional dominant
lattice sum of the original $g,f$ coefficients converges because
\(r_{\rm obs}>2+\beta_{\rm phase}^{W}\).  Together with scale-bad labels the
complete central exponent is
\begin{equation}
\label{eq:complete-central-exponent}
 c_{\rm cen}^{W}:=\min\{\eta_{\rm lo}\beta_{\rm phase}^{W},
 c_{\rm time,J}^{\rm val},c_{\rm time,TV}^{\rm src},
 \eta_{\rm hi}\rho_\infty-a_{\rm ultra},c_{\rm scl},
 c_{\rm low}\}>0,
\end{equation}
where $c_{\rm low}$ is supplied by the low recovery theorem and
$c_{\rm time,J}^{\rm val},c_{\rm time,TV}^{\rm src}>0$ are the two genuine
$N_{\rm time}$ value/TV source rates.  The first already contains the global
base--energy Cauchy loss; it is not halved a second time.  The second, which
includes the contracted-coordinate main remainder in the present
organization, is summed by TV.  All
$N_{\rm PPE}$ and $r$-polynomial exceptional events have already entered
$\beta_{\rm phase}^{W}$ and do not appear a second time.
Lemma~\ref{lem:dominant-coordinate-reduction} supplies the actual typed
remainder, while Lemma~\ref{lem:gate-buffer} supplies the
resolving depth only in the intermediate range, where
\eqref{eq:intermediate-max-resolving-depth} is finite.  Global marginal corrections are inserted through the typed
current/test adjoint identity before any absolute value is taken.  The
regular $M\times M$ density is not placed in the face lattice:
Corollary~\ref{cor:regular-functional-correlation} uses its independent
$W^{R,1}$ smoothing tail.
\end{proof}

\subsection{Low frequencies after dominant-coordinate reduction}

The cutoff is imposed only after the graph term has been reduced to its two
scalar dominant coordinates.  This avoids the false implication that a large
ambient Fourier vector must have a large restriction to a one-dimensional
face.  Contracted coordinates are estimated before the cutoff.

\begin{definition}[Dominant event norm and separate dynamic-H\"older norm]
\label{def:event-current-norm}
Fix a finite physical product atlas and assemble coincident artificial faces
before inserting its partition of unity.  Use the already fixed references
$\nu^-=\nu^+=\mu_0$ and set
$dm_{M\times M}=d\nu^-\otimes d\nu^+$.  For a face label $\lambda$ let
$a_{X,\lambda},a_{Y,\lambda}>0$ be the two scales in
\eqref{eq:dominant-coordinate-normalization} and put
\[
 \mathfrak d_\lambda=1+a_{X,\lambda}^{-1}+a_{Y,\lambda}^{-1}.
\]
Write
\[
 \mathcal K_{\rm ev}=k_{\rm reg}\,dm_{M\times M}
 +\sum_\lambda M_\lambda(H_\lambda)_*(w_\lambda\,d\bar s)
\].
Fix a dominant-scale ledger exponent
$A_{\rm dom}\ge2+2q_{\rm end}$.  It covers the
$O(K^2\mathfrak d_\lambda^2)$ scalar mode count and the fixed scaled chart
costs; it does not encode either boundary complexity.  Define the two
parameter-free recovery marks
\begin{equation}
\label{eq:section-recovery-mark}
 \mathcal R_\lambda^\pm:=C_R\left[1+\log\mathfrak d_\lambda
 +\log\operatorname{Mark}(\lambda)
 +\log\{1+Z(\mathcal G_\lambda^\pm)\}\right].
\end{equation}
For an integer $1\le R\le3$ allowed by the $C^3$ table regularity, define
\begin{equation}
\label{eq:event-physical-norm}
 \|\mathcal K_{\rm ev}\|_{\rm phys,R}:=
 \sum_\chi\|\zeta_\chi k_{\rm reg}\|_{W^{R,1}}
 +\sum_\lambda M_\lambda\|w_\lambda\|_{\mathcal A^\alpha}
 \mathfrak d_\lambda^{A_{\rm dom}}\operatorname{Mark}(\lambda)^2
 \{1+Z(\mathcal G_\lambda^-)+Z(\mathcal G_\lambda^+)\}.
\end{equation}
For bounded $F:M^q\to\mathbb C$ let
\begin{equation}
\label{eq:separate-dynamic-holder-norm}
 \|F\|_{\rm SH}:=\|F\|_\infty+
 \sum_{i=1}^q\sup_{x_j:j\ne i}\sup_{x_i\ne y_i}
 \frac{|F(\ldots,x_i,\ldots)-F(\ldots,y_i,\ldots)|}
      {\theta_0^{s_i(x_i,y_i)}}.
\end{equation}
The stable or unstable separation time appropriate to the $i$th slot is used,
as in \cite[Definition~2.2]{LS}.
\end{definition}

We now define the low-frequency object on the scalar quotient on which it
actually lives.  Fix an even cutoff $\chi\in C_c^\infty((-2,2))$ with
$\chi=1$ on $[-1,1]$.  After the branch-section reduction, the assembled
graph piece induces the finite signed scalar measure
\begin{equation}
\label{eq:dominant-pushforward-measure}
 \nu_\lambda=(X_\lambda,Y_\lambda)_*
       (M_\lambda w_\lambda\,d\bar s),\qquad
 \widehat\nu_\lambda(k,l)=
 \int e^{-i(kX_\lambda+lY_\lambda)}M_\lambda w_\lambda\,d\bar s .
\end{equation}
Define the anisotropic multiplier and its polynomial density by
\begin{align}
 m_{\lambda,K}(k,l)
 &=\chi(a_{X,\lambda}k/K)\chi(a_{Y,\lambda}l/K),
 \label{eq:dominant-physical-cutoff}\\
 A_{\lambda,K}^{\rm dom}(x,y)
 &=(2\pi)^{-2}\sum_{(k,l)\in\mathbb Z^2}
 m_{\lambda,K}(k,l)\widehat\nu_\lambda(k,l)e^{i(kx+ly)}.
 \label{eq:canonical-dominant-polynomial}
\end{align}
This is a density on the scalar quotient and nowhere else.  Define the two
mode currents on the branch sections by
\begin{align}
 \eta^-_{\lambda,k}(h)
 &:=\int h(\iota^-_\lambda(X))\chi^-_\lambda(X)e^{ikX}\,dX,
 \label{eq:past-section-mode-current}\\
 \eta^+_{\lambda,l}(f)
 &:=\int f(\iota^+_\lambda(Y))\chi^+_\lambda(Y)e^{ilY}\,dY,
 \label{eq:future-section-mode-current}
\end{align}
where the fixed scalar cutoffs are part of the localized quotient chart.
They equal one on a neighborhood of
$\operatorname{supp}\nu_\lambda$ and their supports stay inside the
homogeneous section domains.
The branch-section lift $\mathfrak L^{\rm sec}_{\lambda,K}$ is
\begin{equation}
\label{eq:section-lifted-low-block}
 \int h(x)f(y)\,d\mathfrak L^{\rm sec}_{\lambda,K}(x,y)
 :=\int h(\iota^-_\lambda(X))f(\iota^+_\lambda(Y))
       \chi^-_\lambda(X)\chi^+_\lambda(Y)
       A_{\lambda,K}^{\rm dom}(X,Y)\,dX\,dY.
\end{equation}
If $\mathsf P_{\lambda,K}$ denotes the even quotient multiplier in
\eqref{eq:dominant-physical-cutoff} and
$H(X,Y)=h(\iota^-_\lambda(X))f(\iota^+_\lambda(Y))$, then self-adjointness of
the multiplier and the support clause give the exact localized identity
\begin{equation}
\label{eq:section-quotient-low-high}
 \int H\,d\nu_\lambda-
 \int \chi^-_\lambda\chi^+_\lambda H A_{\lambda,K}^{\rm dom}\,dX\,dY
 =\int (I-\mathsf P_{\lambda,K})
       (\chi^-_\lambda\chi^+_\lambda H)\,d\nu_\lambda.
\end{equation}
Thus the only localization error is the same quotient high-pass term treated
by the balanced/unbalanced phase estimates; it is not an unrecorded leakage
of the Fourier polynomial outside the branch chart.
Fourier inversion on the quotient gives the exact
finite-rank identity
\begin{equation}
\label{eq:section-current-finite-rank}
 \mathfrak L^{\rm sec}_{\lambda,K}
 =(2\pi)^{-2}\sum_{k,l}
 m_{\lambda,K}(k,l)\widehat\nu_\lambda(k,l)
 \eta^-_{\lambda,k}\otimes\eta^+_{\lambda,l}.
\end{equation}
Thus finite rank is obtained in the tensor product of two concrete measured
standard-curve currents.  The original graph current is not claimed to be
nuclear or absolutely continuous.

For a finite signed kernel $\sigma$ on $M\times M$ put
\begin{equation}
\label{eq:kernel-operator-identification}
 \langle\operatorname{Op}(\sigma)h,f\rangle
 :=\int_{M\times M}h(x)f(y)\,d\sigma(x,y).
\end{equation}

Fix once and for all a finite product atlas, a subordinate smooth product
partition $\{\zeta_\chi:=\zeta_\chi^-(x)\zeta_\chi^+(y)\}$, product super-cutoffs
$\widetilde\zeta_\chi=1$ on $\operatorname{supp}\zeta_\chi$, and extension
maps $E_\chi$ supported where the corresponding super-cutoff is one and
bounded simultaneously on $L^1$ and on $W^{r,1}$ for every
$1\le r\le R\le3$.  We use the fixed periodic box
$\mathbb T^4=(\mathbb R/2\pi\mathbb Z)^4$ with normalized Haar density
$(2\pi)^{-4}dz$ and Fourier convention
$\widehat h(q)=(2\pi)^{-4}\int_{\mathbb T^4}h(z)e^{-iq\cdot z}dz$.
Choose a real even
$\vartheta\in C_c^\infty(\mathbb R^4)$ with
$\vartheta=1$ on $[-1,1]^4$ and $\vartheta=0$ outside $[-2,2]^4$, and set
\begin{equation}
\label{eq:flat-bandlimited-regular-smoother}
 \mathcal P_K^{(4)}k
 :=\sum_\chi\widetilde\zeta_\chi\,
   \kappa_\chi^*\!\left[
    \sum_{q\in\mathbb Z^4}\vartheta(q/K)
      \widehat{E_\chi(\zeta_\chi k)}(q)e^{iq\cdot z}
   \right].
\end{equation}
This is a signed flat spectral smoother, not a positive ordinary mollifier.
On each chart it contains at most $CK^4$ characters and satisfies
\begin{equation}
\label{eq:flat-smoother-approximation}
 \|(I-\mathcal P_K^{(4)})k\|_{L^1}
 \le C_RK^{-R}\|k\|_{W^{R,1}},\qquad 1\le R\le3.
\end{equation}
Let $\check\vartheta$ be the Euclidean inverse Fourier transform.  Flatness
at the origin gives
$\int_{\mathbb R^4}z^\alpha\check\vartheta(z)\,dz=0$ for
$1\le|\alpha|\le R-1$, while its Schwartz decay gives a finite absolute
$R$th moment.  The Fourier series in
\eqref{eq:flat-bandlimited-regular-smoother} is convolution by the
periodization of $K^4\check\vartheta(Kz)$, equivalently Euclidean convolution
of the periodic extension.  Taylor's integral remainder on $\mathbb R^4$
therefore gives the $K^{-R}W^{R,1}$ bound without asserting polynomial
moment cancellation on a chosen torus fundamental domain.  The simultaneous
$L^1/W^{R,1}$ extension bounds and the fixed partition give
\eqref{eq:flat-smoother-approximation}.  The super-cutoffs preserve the
localized input and only multiply the finitely many characters by fixed
product chart functions.  Hence every four-dimensional character remains a
product of two endpoint characters, and the super-cutoffs create neither a
$K$-dependent mode count nor an extra derivative loss.

\begin{lemma}[Oscillatory section currents are recoverable]
\label{lem:oscillatory-section-recovery}
For each mode current in
\eqref{eq:past-section-mode-current}--
\eqref{eq:future-section-mode-current}, its real and imaginary parts are
finite signed differences of positive measured standard families.  Their
total masses are bounded by the section mark, their scaled density marks are
at most
$CK\mathfrak d_\lambda\operatorname{Mark}(\lambda)$, and they are proper after
\begin{equation}
\label{eq:section-mode-recovery-time}
 R_{\lambda,K}^{\pm}\le C_0+C_1\log K+\mathcal R_\lambda^\pm.
\end{equation}
The same assertion holds for the reversed collision map.
If $\tau_\pm\in(0,1)$ are the two proper-family coupling rates, then the
centered mode currents obey the delayed one-sided bounds
\begin{equation}
\label{eq:section-mode-delayed-coupling}
 |(\eta^\pm_{\lambda,k})^\circ(Q_\pm^jh)|
 \le C\operatorname{Mark}(\lambda)
 \tau_\pm^{(j-R_{\lambda,K}^\pm)_+}\|h\|_{C^1}.
\end{equation}
\end{lemma}

\begin{proof}
Relative to the positive coordinate/arclength density on the section, the
multiplier is a fixed cutoff times $e^{ikX}$.  Split real and imaginary parts
and add twice their supremum to obtain positive densities, exactly as in
Lemma~\ref{lem:signed-observable-multiplier}.  One scalar derivative costs
$C|k|\le CK/a_{X,\lambda}$ (and similarly in $Y$), which gives the displayed
scaled mark.  The density recursion contracts this mark in
$O(\log(K\mathfrak d_\lambda\operatorname{Mark}(\lambda)))$ iterates.  The
boundary growth recursion has the affine form
$Z(F^{n_*}\mathcal G)\le\gamma Z(\mathcal G)+Z_0$ and therefore costs
$O(\log(1+Z))$ iterates.  Forward/reverse proper-family growth and coupling
are uniform on the compact table class, proving
\eqref{eq:section-mode-recovery-time}.  Apply exponential coupling only after
that time and total variation before it.  This proves
\eqref{eq:section-mode-delayed-coupling} without converting the two recovery
windows into a product of initial boundary complexities.
\end{proof}

There are two different centering operations, and we keep their types
separate.  If $\sigma$ is a finite signed measure on $M\times M$, write
$\sigma_\pm=(\operatorname{pr}_\pm)_*\sigma$ and
$|\sigma|_0=\sigma(M\times M)$, and define
\begin{equation}
\label{eq:current-centering}
 \mathsf D_{\rm cur}\sigma
 :=\sigma-\sigma_-\otimes\nu^+-\nu^-\otimes\sigma_+
       +|\sigma|_0\nu^-\otimes\nu^+.
\end{equation}
Its test-side adjoint is
\begin{align}
 \mathsf D_{\rm test}F(x,y)
 :=F(x,y)&-\int F(x,y')\,d\nu^+(y')
          -\int F(x',y)\,d\nu^-(x') \notag\\
 &+\int F(x',y')\,d\nu^-(x')d\nu^+(y'),
 \label{eq:test-centering}
\end{align}
so that
\begin{equation}
 \label{eq:centering-adjoint-identity}
 \int F\,d(\mathsf D_{\rm cur}\sigma)
 =\int \mathsf D_{\rm test}F\,d\sigma.
\end{equation}
In particular a current is never silently fed to an operator defined on
functions.  On the regular density use the flat band-limited smoother
\eqref{eq:flat-bandlimited-regular-smoother} and define
\begin{align}
\label{eq:dominant-centered-smoother}
\mathsf S_K^{\rm dom}\mathcal K_{\rm ev}:=
 \mathsf D_{\rm cur}\left(
 [\mathcal P_K^{(4)}k_{\rm reg}]dm_{M\times M}
 +\sum_\lambda\mathfrak L^{\rm sec}_{\lambda,K}\right),
 \notag\\
 \mathsf H_K^{\rm dom}(\mathsf D_{\rm cur}\mathcal K_{\rm ev}):=
 \mathsf D_{\rm cur}\mathcal K_{\rm ev}-
 \mathsf S_K^{\rm dom}\mathcal K_{\rm ev}.
\end{align}
All artificial endpoints are assembled before this definition.  Expanding
$\mathsf D_{\rm cur}$ centers each pair of section-mode currents and includes
the two marginal and constant product blocks explicitly.
The definition of $\mathsf H_K^{\rm dom}$ gives the exact operator decomposition
\begin{equation}
\label{eq:current-operator-low-high-identity}
 \operatorname{Op}(\mathsf D_{\rm cur}\mathcal K_{\rm ev})
 =\operatorname{Op}(\mathsf S_K^{\rm dom}\mathcal K_{\rm ev})
  +\operatorname{Op}\!\left(
       \mathsf H_K^{\rm dom}(\mathsf D_{\rm cur}\mathcal K_{\rm ev})\right).
\end{equation}

\begin{theorem}[Branch-section low projection and exact finite rank]
\label{thm:canonical-dominant-smoothing}
For $K\ge2$ the centered current in
\eqref{eq:dominant-centered-smoother} has the exact finite-rank expansion
\begin{equation}
\label{eq:finite-rank-low-kernel}
 \mathsf S_K^{\rm dom}\mathcal K_{\rm ev}
 =\sum_{j\in\mathcal J_K}c_j
       (\eta_j^-)^\circ\otimes(\eta_j^+)^\circ,
\end{equation}
where $\eta^\circ=\eta-\eta(1)\nu^\pm$.  Face factors are the concrete
section currents \eqref{eq:past-section-mode-current}--
\eqref{eq:future-section-mode-current}; regular-density factors are ordinary
absolutely continuous currents.  Before either recovery estimate is used,
the face modes have the total-variation mass budget
\begin{equation}
\label{eq:finite-rank-mode-mass}
 \sum_{\lambda}\sum_{k,l}|c_{\lambda,k,l}|
 \|(\eta^-_{\lambda,k})^\circ\|_{\rm TV}
 \|(\eta^+_{\lambda,l})^\circ\|_{\rm TV}
 \le CK^2\sum_\lambda M_\lambda\|w_\lambda\|_{L^1}
 \mathfrak d_\lambda^2\operatorname{Mark}(\lambda)^2.
\end{equation}
The constant is independent of return and grazing truncations, and
\eqref{eq:current-operator-low-high-identity} holds before any absolute value
is taken.  This is a finite-rank statement only after scalar cutoff; no
nuclear or absolutely continuous representation of the original graph current
is asserted.  In particular, no product
$(1+Z(\mathcal G^-_\lambda))(1+Z(\mathcal G^+_\lambda))$ is hidden in
\eqref{eq:finite-rank-mode-mass}.
\end{theorem}

\begin{proof}
The coefficient bound
$|\widehat\nu_\lambda(k,l)|\le
M_\lambda\|w_\lambda\|_{L^1}$ follows directly from
\eqref{eq:dominant-pushforward-measure}.  The multiplier has at most
$CK^2\mathfrak d_\lambda^2$ nonzero pairs.  Formula
\eqref{eq:section-current-finite-rank} gives their exact product-current
expansion.  Each centered section current has total variation at most a fixed
multiple of $\operatorname{Mark}(\lambda)$, independently of its recovery
time.  Absolute summation of the coefficients and modes proves
\eqref{eq:finite-rank-mode-mass}.  Formula
\eqref{eq:current-centering} replaces each factor by its centered version,
and \eqref{eq:centering-adjoint-identity} proves the four marginal identities.
The definition of the remainder proves the exact low/high operator identity;
Fourier inversion is used only on the scalar quotient.  The preliminary
replacement of the physical graph trace by the branch sections is charged
separately to \eqref{eq:contracted-coordinate-error} before the cutoff.
\end{proof}

\begin{definition}[Separate recovery-tail input $\mathrm{REC}(q_{\rm rec})$]
\label{def:separate-recovery-tail}
For $q_{\rm rec}>0$, the terminal face-label law satisfies
$\mathrm{REC}(q_{\rm rec})$ if, separately for the two orientations,
\begin{equation}
\label{eq:separate-recovery-moments}
 \sum_\lambda M_\lambda\|w_\lambda\|_{\mathcal A^\alpha}
 \mathfrak d_\lambda^{A_{\rm dom}}\operatorname{Mark}(\lambda)^2
 e^{q_{\rm rec}\mathcal R_\lambda^\pm}<\infty.
\end{equation}
The two displayed sums are taken under the same terminal label law, but their
integrands are not multiplied.  Thus this input is a pair of marginal
exponential recovery tails, not a pointwise product--$Z$ moment.
\end{definition}

\begin{lemma}[Regular high-frequency tail]
\label{lem:regular-high-tail}
For the regular $M\times M$ density,
\begin{equation}
\label{eq:regular-high-tail}
 \|(I-\mathcal P_K^{(4)})k_{\rm reg}\|_{L^1(M\times M)}
 \le CK^{-R}\|k_{\rm reg}\|_{W^{R,1}(M\times M)}.
\end{equation}
Consequently its pairing with $Q^mg$ and $Q^nf$ is bounded by the right-hand
side times $\|g\|_\infty\|f\|_\infty$, uniformly in $m,n$.
\end{lemma}

\begin{proof}
This is \eqref{eq:flat-smoother-approximation}; Markov normalization gives
$\|Q^mh\|_\infty\le2\|h\|_\infty$.
\end{proof}

\begin{lemma}[Regular low projector]
\label{lem:regular-low-projector}
Let $c_-,c_+>0$ be the two one-state mixing exponents.  The centered regular
low block satisfies
\begin{equation}
\label{eq:regular-low-projector}
 \left|\left\langle Q^n\operatorname{Op}\!\left(
 \mathsf D_{\rm cur}[(\mathcal P_K^{(4)}k_{\rm reg})dm_{M\times M}]
 \right)Q^mg,f\right\rangle\right|
 \le CK^6e^{-c_-m-c_+n}\|k_{\rm reg}\|_{L^1}
       \|g\|_{C^1}\|f\|_{C^1}.
\end{equation}
\end{lemma}

\begin{proof}
In a fixed product atlas the four-dimensional cutoff has $O(K^4)$ Fourier
modes and every coefficient is bounded by $C\|k_{\rm reg}\|_{L^1}$.
Each two-dimensional endpoint character has $C^1$ cost at most $CK$.
After current centering, Fubini makes every mode a product of two centered
one-state correlations.  Their two coupling bounds contribute
$e^{-c_-m-c_+n}$, while the mode count and endpoint costs contribute
$K^4K^2$.  Summing the fixed atlas proves
\eqref{eq:regular-low-projector}.
\end{proof}

\begin{proposition}[Central low frequency by delayed bi-recovery]
\label{prop:functional-correlation-low-frequency}
Let $c_-,c_+>0$ be the uniform past and future two-state mixing rates supplied
by proper-family coupling, let $D\ge1$ define the diagonal band
$D^{-1}m\le n\le Dm$, and assume $\mathrm{REC}(q_{\rm rec})$.  Fix
$0<\epsilon<1/2$, put $N=\max\{m,n\}$ and
$K=e^{\eta_{\rm lo}N}$, and choose
\begin{equation}
\label{eq:delayed-recovery-eta-window}
 6\eta_{\rm lo}<\tfrac12\min\{c_-,c_+\},\qquad
 2\eta_{\rm lo}<(1-\epsilon)\min\{c_-,c_+\},\qquad
 (q_{\rm rec}C_1+2)\eta_{\rm lo}<q_{\rm rec}\epsilon/D.
\end{equation}
Then, for some $c_{\rm low}>0$ independent of $m,n$,
\begin{equation}
\label{eq:functional-correlation-low-frequency}
 \left|\left\langle Q^n
 \operatorname{Op}(\mathsf S_K^{\rm dom}\mathcal K_{\rm ev})Q^mg,
 f\right\rangle\right|
 \le C e^{-c_{\rm low}(m+n)}\|g\|_{C^r}\|f\|_{C^r}.
\end{equation}
\end{proposition}

\begin{proof}
Insert \eqref{eq:finite-rank-low-kernel} into the product-current pairing.
After the face modes of one label are summed, Fubini,
\eqref{eq:section-mode-delayed-coupling}, and
\eqref{eq:finite-rank-mode-mass} give
\[
 C M_\lambda\|w_\lambda\|_{L^1}K^2\mathfrak d_\lambda^2
 \operatorname{Mark}(\lambda)^2
 e^{-c_-(m-R^-_{\lambda,K})_+}
 e^{-c_+(n-R^+_{\lambda,K})_+}
 \|g\|_{C^1}\|f\|_{C^1}.
\]
There is no condition $y=Tx$: the two factors are independent one-time
correlations.  Split the labels into the good set
$R^-_{\lambda,K}\le\epsilon m$, $R^+_{\lambda,K}\le\epsilon n$ and the two
bad sets where one of these inequalities fails.  On the good set the first
two inequalities in \eqref{eq:delayed-recovery-eta-window} absorb the
$K^2$ mode count.  On the past-bad set,
\eqref{eq:section-mode-recovery-time} implies
\[
 \mathcal R^-_\lambda>\epsilon m-C_0-C_1\eta N.
\]
Use total variation for both factors and the minus instance of
\eqref{eq:separate-recovery-moments}.  Its contribution is at most a constant
times
\[
 \exp\{-[q_{\rm rec}\epsilon/D-(q_{\rm rec}C_1+2)\eta]N\}.
\]
The plus-bad set is identical and uses only the plus recovery moment.  A
union bound, rather than a product of the two initial complexities, completes
the face estimate.  The flat band-limited four-dimensional low projector of the
regular density is bounded by Lemma~\ref{lem:regular-low-projector} and is absorbed by the
first inequality in \eqref{eq:delayed-recovery-eta-window}.  Since
$N\asymp_Dm+n$, decreasing the exponent proves
\eqref{eq:functional-correlation-low-frequency}.
\end{proof}

\begin{corollary}[Regular defect via functional correlation]
\label{cor:regular-functional-correlation}
Assume all hypotheses of
Proposition~\ref{prop:functional-correlation-low-frequency} and
Lemmas~\ref{lem:regular-low-projector} and
\ref{lem:regular-high-tail}, including the proper-family coupling,
recovery moments, common event/test norms, and a uniform
$W^{R,1}$ regular-density ledger.  Then
The regular part of $K_V$ has central-band exponential decay: its finite-rank
low part is covered by Proposition~\ref{prop:functional-correlation-low-frequency}, and
Lemma~\ref{lem:regular-high-tail} gives
$K^{-R}=e^{-\eta_{\rm lo}RN}$ for its high
part.  No tensor, wavelet, or face-phase estimate is used.
\end{corollary}

\subsection{Radial faces: direct flux boundary control}

We retain every connected component.  Let
$\mathcal W_\lambda(j,k)=\{W_{\lambda,j,k,r}\}_r$ be the maximal components
of one coarea face with incoming strip $j$ and outgoing strip $k$.  Put
\[
 p_{\lambda,j,k,r}=\int_{W_{\lambda,j,k,r}}\rho^{\rm raw}_\lambda\,ds,
 \quad M_\lambda=\sum_{j,k,r}p_{\lambda,j,k,r},\quad
 \bar p_{\lambda,j,k,r}=p_{\lambda,j,k,r}/M_\lambda.
\]
Thus $\bar p$, not $p$, is the weight of the normalized standard family.

\begin{lemma}[Oriented refinement bookkeeping]
\label{lem:cone-monotone-components}
Assume a face has been cut into $N_{\rm or}(\lambda)$ signed charts on each
of which either the forward or the reversed input and output angle
coordinates are monotone.  Then for fixed incoming/outgoing signed
homogeneity ranks $(j,k)$ there are at most $4N_{\rm or}(\lambda)$ connected
components and
\begin{equation}
\label{eq:derived-weighted-component-complexity}
 \sum_{j,k,r}\frac1{(1+j)^2(1+k)^2}
 \le C N_{\rm or}(\lambda).
\end{equation}
The existence of such a cover and every estimate for $N_{\rm or}$ are not
deduced from cone transversality or a qualitative jet condition; their
weighted form is an explicit part of $\mathrm{FZ}_{\rm geom}$ below.
\end{lemma}

\begin{proof}
On an unstable graph the cone bounds give $dr/dt$ and $d\varphi/dt$ one fixed
sign after the orientation of $t$ is chosen.  Invariance of the unstable cone
gives the same statement after one collision; the reversed cone gives the
backward version.  A signed homogeneity strip is an interval in the angle
coordinate.  The preimages of an incoming and an outgoing signed strip are
therefore intervals on the oriented chart, and their intersection is an
interval.  The two grazing signs in each time orientation account for the
factor four.  Summing over charts and then over $j,k$ proves
\eqref{eq:derived-weighted-component-complexity}.  This proves only the
bookkeeping implication from an already constructed oriented cover; it makes
no analytic zero-count assertion.
\end{proof}

\begin{lemma}[Full tangent coarea Jacobian]
\label{lem:two-flux-coarea}
Let $\gamma(t)=(r_-(t),\varphi_-(t))$ be a unit-speed parametrization of an
input coarea face and let $W=T\gamma$.  On every oriented component the
unnormalized density with respect to output arclength is
\begin{equation}
\label{eq:radial-two-flux-density}
 \rho^{\rm raw}_{W,\lambda}(T\gamma(t))=
 \frac{|a_\lambda|J_{{\rm chart},\lambda}J_{{\rm hol},\lambda}
       \cos\varphi_-(t)}
 {|\nabla G_\lambda(\gamma(t))|\,
       |D T(\gamma(t))\gamma'(t)|}.
\end{equation}
If output impact coordinate $r_+$ is used instead of arclength, the last
factor in the denominator is $|d(r_+\circ\gamma)/dt|$.  For the vertical
family $\gamma(t)=(r_0,t)$ this becomes
\begin{equation}
\label{eq:vertical-two-flux-special-case}
 \rho^{\rm raw}_{r_+}=
 \frac{|a_\lambda|J_{{\rm chart},\lambda}J_{{\rm hol},\lambda}}
 {|\nabla G_\lambda|}\frac{\cos\varphi_-\cos\varphi_+}{\tau},
\end{equation}
which is precisely the calculation of Lemma~5.1 of
\cite{CanestrariHoles}.  Formula \eqref{eq:radial-two-flux-density}, not the
vertical special case, is used for a general physical face.  The reversed
orientation has the same change-of-variables form.
\end{lemma}

\begin{proof}
Input arclength coarea and the collision--SRB density give the measure
\[
 b_\lambda(t)\,dt,
 \qquad b_\lambda(t)=
 \frac{|a_\lambda|J_{{\rm chart},\lambda}J_{{\rm hol},\lambda}
       \cos\varphi_-(t)}{|\nabla G_\lambda(\gamma(t))|}.
\]
The one-dimensional area formula for the injective oriented branch
$t\mapsto T\gamma(t)$ divides this density by the curve Jacobian
$J_\gamma T(t)=|DT(\gamma(t))\gamma'(t)|$, proving
\eqref{eq:radial-two-flux-density}.  Parametrizing the image by $r_+$ instead
divides by $|d(r_+\circ\gamma)/dt|$.  When $r_-$ is fixed, the billiard
differential identity
$dr_+/d\varphi_-=-\tau/\cos\varphi_+$ yields
\eqref{eq:vertical-two-flux-special-case}.  Applying the same area formula to
$T^{-1}$ proves the reversed statement without asserting that a general
tangent is vertical.
\end{proof}

\begin{definition}[Geometric flux and component condition
$\mathrm{FZ}_{\rm geom}$]
\label{def:flux-complexity}
A face family satisfies $\mathrm{FZ}_{\rm geom}(\theta_Z)$ if the density of
Lemma~\ref{lem:two-flux-coarea} has marks
$\mathfrak M_\lambda,\mathfrak C_\lambda\ge1$ such that
\begin{align}
 p_{\lambda,j,k,r}
 &\le C\mathfrak M_\lambda
 \frac{|W_{\lambda,j,k,r}|}{(1+j)^2(1+k)^2},
 \label{eq:radial-component-flux}\\
 \sum_{j,k,r}\frac1{(1+j)^2(1+k)^2}
 &\le\mathfrak C_\lambda:=CN_{\rm or}(\lambda),
 \label{eq:weighted-component-complexity}
\end{align}
and, for some $q_F>1$,
\begin{equation}
\label{eq:primitive-flux-ledger}
 \sum_\lambda M_\lambda
 \left[\operatorname{Mark}(\lambda)^2
 \left\{1+\frac{\mathfrak M_\lambda\mathfrak C_\lambda}
 {M_\lambda}\right\}\right]^{q_F}<\infty.
\end{equation}
Pieces whose tangent makes angle at most $\epsilon$ with both usable cones
have total marked mass at most $C\epsilon^{\theta_Z}$.  This is the required
near-stable angle-sublevel input.  On the complementary good-angle pieces,
the product of all inverse-angle costs in the graph chart, density,
nonstationary phase, and recovery maps is at most
$C\epsilon^{-A_{\rm ang}}$ for a fixed $A_{\rm ang}>0$.  Both constants are
part of $\mathrm{FZ}_{\rm geom}$.  This condition contains no multi-time
mixing or split-law assertion.
\end{definition}

\begin{definition}[Split-law standard-family and positive-envelope interface]
\label{def:split-law-envelope}
Let $F_{s,L}^{\rm col}$ be the sharp four-state functional of the assembled
finite-parameter collar, after coincident artificial faces have cancelled.
For every $(m,n)$, expansion of the global test-side double projection gives
the exact signed law
\begin{equation}
\label{eq:exact-split-law-expansion}
 \Sigma_{m,n}^{\rm law}
 =\rho_{m,n}^{0}-\sum_{j=1}^{J}a_j\rho_{m,n}^{j},
 \qquad \rho_{m,n}^{0}=\rho_{m,n}^{\rm orig},
\end{equation}
where the $\rho_{m,n}^{j}$ are the block-product laws in the functional-
correlation split; all $\rho_{m,n}^{j}$, including $\rho_{m,n}^{0}$, are
probability laws, $J\le J_0$, and
$Z_{\rm law}:=1+\sum_j|a_j|\le Z_0$ uniformly.  Define the positive envelope
\begin{equation}
\label{eq:positive-envelope-law}
 \bar\rho_{m,n}:=Z_{\rm law}^{-1}
 \left(\rho_{m,n}^{0}+\sum_{j=1}^{J}|a_j|\rho_{m,n}^{j}\right).
\end{equation}
The interface $\mathrm{SF}_{\rm split}$ requires that the restriction of every
constituent law in \eqref{eq:exact-split-law-expansion}, in either usable time
orientation, disintegrates into bounded-mark proper standard families with the
same conditional mark--$Z$ ledger as the original collar.  This is a primitive
joint-law statement, not a consequence of the one-coordinate marginals.

Let $\mathcal P_r^{m,n}$ be the common rank-$r$ product-cylinder partition,
refined into connected event-chart components.  For a bounded $F$ put
\begin{equation}
\label{eq:envelope-cylinder-oscillation}
 \Omega_r^{\rm env}(F):=
 \sum_{C\in\mathcal P_r^{m,n}}\bar\rho_{m,n}(C)
 \operatorname*{ess\,osc}_{C,\bar\rho_{m,n}}F.
\end{equation}
The connected-cylinder pullback condition $\mathrm{CP}_{\rm env}$ says that,
outside pieces of total envelope mass $Ce^{-c_{\rm bad}r}$, every such
component pulls back to a good-angle homogeneous event cell of physical
diameter at most $Ce^{-\chi r}$, with the conditional density and component
marks in $\mathrm{FZ}_{\rm geom}$ and $\mathrm{SF}_{\rm split}$.  Grazing,
near-stable, and homogeneity-boundary cells are included in the displayed bad
mass, with every component counted.
There are fixed $A_{\rm amp}>0$ and $\alpha_{\rm amp}\in(0,1]$ such that, on
a good cell, the smooth coefficient multiplying each sharp collar indicator
has oscillation at most
$C|s|^{-1}e^{A_{\rm amp}L}e^{-\chi\alpha_{\rm amp}r}$.
The connected refinement is required to be separation-time admissible: two
distinct retained atoms either receive the same cylinder value or differ in
some coordinate before rank $r$.  Hence a function constant on its atoms has
separate dynamic-H\"older cost at most $Ce^{q_{\rm cyl}r}$ times its
supremum.  In particular no artificial chart boundary may create an
uncharged jump between symbolically indistinguishable points.
\end{definition}

\begin{proposition}[Geometric envelope dynamical tube]
\label{prop:envelope-dynamical-tube}
Assume $\mathrm{FZ}_{\rm geom}(\theta_Z)$,
$\mathrm{SF}_{\rm split}$, and $\mathrm{CP}_{\rm env}$.  Then there are
$A_{\rm col}>0$ and $\alpha_{\rm col}\in(0,1]$ such that
\begin{equation}
\label{eq:FZ-dynamical-tube}
 \sup_{m,n}\Omega_r^{\rm env}(F_{s,L}^{\rm col})
 \le C|s|^{-1}e^{A_{\rm col}L}e^{-\alpha_{\rm col}r},
 \qquad r\ge1.
\end{equation}
\end{proposition}

\begin{proof}
On a good connected component, split the oscillation into variation of the
smooth coefficient and variation of a sharp indicator.  The first part is
bounded by the scaled H\"older clause in $\mathrm{CP}_{\rm env}$.  The second
produces two points on opposite sides of a genuine face.  Connectedness and
the diameter bound imply that the whole component lies in the
$Ce^{-\chi r}$ tube of that face.  By
\eqref{eq:positive-envelope-law}, its envelope mass is the fixed positive
combination of the actual and split-law masses.  Apply the full tangent area
formula and the standard-family conditional density bound to every
constituent.  The good-angle part costs
$C|s|^{-1}e^{A_{\rm col}L}e^{-\chi\alpha_0r}$ for a geometric tube exponent
$\alpha_0>0$, after taking
$A_{\rm col}\ge A_{\rm amp}$ and decreasing $\alpha_0$ to at most
$\alpha_{\rm amp}$.  The complement costs
$C|s|^{-1}e^{A_{\rm col}L}e^{-c_{\rm bad}r}$ by
$\mathrm{CP}_{\rm env}$.  Balancing the angle cutoff with the
$\epsilon^{\theta_Z}$ sublevel and decreasing the exponent once gives
\eqref{eq:FZ-dynamical-tube}.  Artificial boundaries were cancelled before
\eqref{eq:exact-split-law-expansion}, and connected components were indexed
before the diameter argument, so neither can create an uncharged crossing.
The standard-pair strip coverage of \cite{CanestrariPerturb} is a template
for verifying $\mathrm{CP}_{\rm env}$; it is not used as a black box.
\end{proof}

\begin{proposition}[Flux-weighted radial boundary lemma]
\label{prop:radial-direct-Z}
Assume $\mathrm{FZ}_{\rm geom}(\theta_Z)$ and decompose the signed amplitude into
positive pieces.  Then the normalized standard family obeys
\begin{equation}
\label{eq:radial-direct-Z-bound}
 M_\lambda Z(\mathcal G_\lambda)
 =\sum_{j,k,r}\frac{p_{\lambda,j,k,r}}
 {|W_{\lambda,j,k,r}|}
 \le C\mathfrak M_\lambda\mathfrak C_\lambda,
\end{equation}
and, for the actual physical labels,
\begin{equation}
\label{eq:joint-mark-Z-ledger}
 \sum_\lambda M_\lambda
 \operatorname{Mark}(\lambda)^2
 \{1+Z(\mathcal G_\lambda^-)+Z(\mathcal G_\lambda^+)\}<\infty.
\end{equation}
and the near-stable pieces have the marked
$\epsilon^{\theta_Z}$ tail in Definition~\ref{def:flux-complexity}.
\end{proposition}

\begin{proof}
Divide \eqref{eq:radial-component-flux} by the component length and sum over
the full triple $(j,k,r)$.  Equation
\eqref{eq:weighted-component-complexity} gives
\eqref{eq:radial-direct-Z-bound}; no component is suppressed.  Multiply by
the squared density/amplitude mark and use
\eqref{eq:primitive-flux-ledger} (first with power one) to obtain
\eqref{eq:joint-mark-Z-ledger}.  The full tangent Jacobian, weighted
component complexity, and angle sublevel are the three falsifiable geometric
inputs; none is deduced merely from the homogeneity labels.
\end{proof}

\begin{lemma}[Explicit angle balance]
\label{lem:FZ-angle-balance}
Suppose the good-angle core of a central or off-diagonal estimate is
$Ce^{-c_{\rm core}N}$ before inverse-angle loss.  Set
\[
 \epsilon_N=e^{-\zeta N},\qquad
 \zeta=\frac{c_{\rm core}}{A_{\rm ang}+\theta_Z}.
\]
Then the sum of the bad-angle and good-angle contributions is at most
$Ce^{-c_ZN}$, where
\begin{equation}
\label{eq:explicit-cZ}
 c_Z=\frac{c_{\rm core}\theta_Z}
 {A_{\rm ang}+\theta_Z}>0.
\end{equation}
\end{lemma}

\begin{proof}
The bad-angle mass is $e^{-\zeta\theta_ZN}$, while the good-angle estimate is
$e^{-(c_{\rm core}-A_{\rm ang}\zeta)N}$.  The displayed choice makes the two
exponents equal.
\end{proof}

\begin{corollary}[Uniform radial recovery]
\label{cor:radial-bi-recovery}
For each time orientation, the radial family is regular after at most
\begin{equation}
\label{eq:radial-log-Z-recovery}
 R^\pm\le \alpha_1\log\{1+Z(\mathcal G_{\rm rad}^\pm)\}+\alpha_2
\end{equation}
iterates.  After multiplication by physical mass, the recovery-window sum in
all four time quadrants is finite and is controlled by
\eqref{eq:joint-mark-Z-ledger}.
\end{corollary}

\begin{proof}
Apply \cite[Corollary~6.15(c)]{CanestrariHoles} to the positive standard
families in each orientation.  Use
$\log(1+Z)\le C_\varepsilon(1+Z)^\varepsilon$ and multiply by
$M_\lambda$.  The joint mark--$Z$ ledger then sums both orientations and all
pre-recovery iterates.  This is the weighted hypothesis of
Proposition~\ref{prop:recovery-window-summability}.
\end{proof}

\begin{lemma}[Observable multiplier to positive proper families]
\label{lem:signed-observable-multiplier}
Let $(W,\rho)$ be a positive standard pair and let
$q_m=(g\circ T^{-m})\rho$ on the time orientation for which $W$ is unstable.
There is $C_m\le C\|g\|_{C^\alpha}$ such that
\[
 q_m=(q_m+C_m\rho)-C_m\rho
\]
is a difference of positive densities.  Their total masses are bounded by
$C\|g\|_{C^\alpha}$ and their scaled log--H\"older marks are at most
$Ce^{\chi m}\operatorname{Mark}(W,\rho)$.  For every admissible slowly
moving sequence in the compact table class, after
\begin{equation}
\label{eq:multiplier-recovery-time}
 R_m\le C_0+C_1m+C_2\log\{1+Z(\mathcal G)\}
\end{equation}
iterates, both positive families are proper.  The reversed collision map has
the analogous statement.
\end{lemma}

\begin{proof}
Take $C_m=2\|g\|_\infty$.  Positivity and the mass bound are immediate.
Hyperbolicity and bounded distortion give
$[g\circ T^{-m}]_{\alpha,\rm sc}\le Ce^{\chi m}\|g\|_{C^\alpha}$.
Since the added density is at least $C_m\rho/2$, division by it gives the same
scaled log--H\"older bound.  More explicitly, the density recursion in
Lemma~13 of \cite{SYZ} has the form
$H_{j+1}\le\theta H_j+C$ on every homogeneous descendant.  Starting from
$H_0\le Ce^{\chi m}\operatorname{Mark}(W,\rho)$, it reaches the fixed regular
cone after $C_0+C_1m$ iterates.  The normalized boundary recursion in
Lemma~16 of \cite{SYZ} is
$Z_j/M\le C\{1+\theta_p^j Z_0/M\}$, so a further
$C_2\log(1+Z_0/M)$ iterates make the family proper.  Theorems~1'--2' of
\cite{SYZ} make all constants uniform for the compact slowly moving table
class.  The fixed-map statement follows by taking a constant sequence, and
time reversal gives the last assertion.
\end{proof}

\begin{proposition}[Signed proper-family mapping and off-diagonal bounds]
\label{prop:signed-off-diagonal-mapping}
Assume the forward and reversed billiard growth/coupling theorem on positive
proper standard families, $\mathrm{FZ}_{\rm geom}(\theta_Z)$, and the scaled density
bound of Proposition~\ref{prop:scaled-face-density}.  Decompose every signed
face amplitude as the difference of two positive amplitudes before
normalization.  Then there are positive exponents
$\alpha_\pm,\beta_\pm$ such that the complete face current satisfies
\begin{align}
 |B_{m,n}|&\le Ce^{\beta_-m-\alpha_+n},
 \label{eq:verified-forward-cone}\\
 |B_{m,n}|&\le Ce^{-\alpha_-m+\beta_+n}.
 \label{eq:verified-backward-cone}
\end{align}
The constants include all recovery windows and are uniform under the
physical and flux ledgers.
\end{proposition}

\begin{proof}
Apply Lemma~\ref{lem:signed-observable-multiplier} on each positive face
piece.  Its recovery time is bounded by
\eqref{eq:multiplier-recovery-time}.  Equations
\eqref{eq:radial-direct-Z-bound} and
\eqref{eq:radial-log-Z-recovery} control the boundary part of that recovery.
Coupling for the remaining
$n-R^+_\lambda$ iterates gives
\[
 CM_\lambda\operatorname{Mark}(\lambda)^2(1+Z_\lambda^+)
 e^{\beta_-m-\alpha_+(n-R^+_\lambda)_+}.
\]
Before recovery use total variation.  The inequality
 $e^{\alpha_+R^+_\lambda}\le
 Ce^{\alpha_+C_1m}(1+Z_\lambda^+)^{\alpha_+C_2}$ and
\eqref{eq:joint-mark-Z-ledger}, with a weakened coupling exponent chosen below
the available mark--$Z$ moment, sum all labels and recovery windows.  Absorb
$e^{\alpha_+C_1m}$ into $e^{\beta_-m}$.  This is
\eqref{eq:verified-forward-cone}.  Apply the same argument to the reversed
collision map for \eqref{eq:verified-backward-cone}; positive/negative pieces
are recombined only after the two bounds are established.
\end{proof}

\section{Slope-energy and fractional-phase proof architecture}

The proof now has six noncircular links.
\begin{enumerate}[label=\textup{(F\arabic*)}]
\item Quotient the exponential Young return on the regular magnet to the
      one-state countable $C^{1+\alpha}$ IFS of Proposition
      \ref{prop:marked-full-branch-quotient}.  Physical side labels remain
      in the signed amplitude fiber.
\item Lemma~\ref{lem:homogeneity-grazing-moment} combines the exponential
      return tail with the maximal grazing rank and proves the joint moment
      needed for the countable quotient and the marked coarea amplitudes.
\item The normal-cycle value $\mathfrak A=-1$ motivates the clean template;
      $\mathrm{SS}$ and the primitive physical escape
      $\mathrm{PPE}_{N}$, $\mathrm{CEM\_DOM}$, and
      $\mathrm{TENERGY}_{W}+\mathrm{CLOCK\_PROB}_{W}
      +\mathrm{CLOCK\_SRC}_{W}$, together with the
      mass-preserving complete stopped-source tree, give the weighted energy
      theorem and Jensen.
\item Proposition~\ref{prop:exact-reduced-slope} and smooth value-space
      cutoffs and Lemma~\ref{lem:fractional-nonstationary-phase} give the
      balanced Fourier estimate without curvature or componentwise phase
      partitions.
\item Global centering is performed before the physical chart sum.
      Lemma~\ref{lem:dominant-coordinate-reduction} removes contracted modes;
      Proposition~\ref{prop:functional-correlation-low-frequency} treats the
      remaining low modes, while Proposition
      \ref{prop:central-phase-verification} treats the high modes.
\item At each cutoff coarea currents converge on a fixed core in the explicit
      bounded--Lipschitz current topology
      \eqref{eq:face-DQ-weak-convergence}, with the boundary tightness
      \eqref{eq:face-DQ-boundary-tightness}.  The no-$|s|^{-1}$ middle tail
      in $\mathrm{FACE}_{\rm 2CUT}$ then upgrades this to growing depth.  The
      component-indexed $\mathrm{FZ}_{\rm geom}$ estimate, the
      $\mathrm{SF}_{\rm split}$ and $\mathrm{CP}_{\rm env}$ interfaces,
      $\mathrm{QT}_{\rm path}$, and the strict scalar-tail margin remove
      $L,H$ before recovery summation.
\end{enumerate}

This order makes clear which external results are used.  Leclerc's
Lemma~3.15 supplies the cohomological bracket and his
Propositions~4.31--4.33 supply an equilibrium/Cantor cutting template with
all-scale polynomial escape on that factor.  They do not supply the
actual-collision-SRB prescribed-depth estimate $\mathrm{PPE}_{N}$ used here;
that estimate remains a primitive input, with the adaptive full-mass alphabet
recorded only as a nonvacuity strategy.
The final oscillatory estimate is the elementary fractional phase lemma
proved above; LPS is retained only as background for the quotient route.
Bertozzi's triple inducing supplies a useful model for the one-state gate
architecture, not the slope conclusion.  The positive entrance/exit comparison is the symbolic local
product argument of \cite{ClimenhagaLPS}.  The
one-step expansion, linear singularity complexity, and uniform billiard
constants used by the master atlas are summarized in \cite{DLreview}.

The route deliberately does not use the recent curvilinear two-ends theorem
or $L^2$ flattening: their nonlinear applications require $C^2$ or higher
geometry, whereas the billiard quotient available here is only
$C^{1+\alpha}$.  It also does not infer a product tower tail from two
marginal Chernoff bounds: the diagonal product exponent comes from time
comparability and functional correlation, while the off-diagonal regions are
handled before the three-region assembly.

\subsection{Execution audit}

The central analytic chain is now represented by explicit propositions.
The remaining primitive geometric inputs are
\[
 \begin{gathered}
 \mathrm{SS},\ \mathrm{ACTUAL\_SLOPE\_BRIDGE},\\
 \mathrm{PAIR\_SAFE\_REFERENCE},\ \mathrm{PAIR\_TENERGY},\
 \mathrm{PPE}_{N},\\
 \mathrm{BR}_{\rm sec},\ \mathrm{CEM\_DOM},\
 \mathrm{TENERGY}_{W},\ \mathrm{DUAL\_MASS\_KERNEL},\\
 \mathrm{MPD},\ \mathrm{POST\_CARRIER\_DOM},\\
 \mathrm{PHYS\_TEST\_LIFT},\ \mathrm{PHYS\_TREE\_MATCH},\\
 \mathrm{PROP\_Q\_MATCH},\\
 \mathrm{MAIN\_CLOCK\_COVERAGE},\
 \mathrm{SEL\_JSL}_{p_{\rm joint}},\\
 \mathrm{CLOCK\_PROB}_{W}
 \end{gathered}
\]
\[
 \begin{gathered}
 \mathrm{CLOCK\_SRC}_{W},\ \mathrm{KRH}_{p_W},\
 \mathrm{FZ}_{\rm geom},\ \mathrm{SF}_{\rm split},\\
 \mathrm{CP}_{\rm env},\ \mathrm{FACE}_{\rm 2CUT},\
 \mathrm{FACE\_TIME}_{\rm CM2},\
 \mathrm{QT}_{\rm JET},\ \mathrm{QT}_{\rm LEDGER}
 \end{gathered}.
\]
The exact gate interface is
Lemma~\ref{lem:cohomological-bracket-identification}, Proposition
\ref{prop:context-polynomial-blowup}, and Lemma~\ref{lem:gate-buffer}; the
slope-energy theorem is Theorem~\ref{thm:direct-slope-small-ball}; the
balanced phase is Proposition~\ref{prop:central-phase-verification}; and the
low-frequency collision-time estimate is Proposition
\ref{prop:functional-correlation-low-frequency}.  Radial recovery is applied
only after these estimates through Corollary~\ref{cor:radial-bi-recovery} and
Proposition~\ref{prop:recovery-window-summability}.

The finite-parameter collar is treated below by the actual/split-law cylinder
approximation theorem; no fixed-slice smoothing estimate is used.

\begin{lemma}[Parameter derivatives of contracting return words]
\label{lem:parametric-inverse-word}
Let $h_{j,s}$ be the inverse one-step branches in a fixed compact regular
atlas and assume, uniformly in $j$ and $|s|<s_0$,
\[
 \|Dh_{j,s}\|_\infty\le\lambda<1,
 \qquad
 \|\partial_sh_{j,s}\|_{C^1}
 +\|\partial_s\log|h_{j,s}'|\|_\infty\le A.
\]
For a clean word $a=(a_1,\ldots,a_R)$ let
$h_{a,s}=h_{a_1,s}\circ\cdots\circ h_{a_R,s}$.  Then
\begin{equation}
\label{eq:parametric-word-bounds}
 \|\partial_s h_{a,s}\|_\infty\le \frac{A}{1-\lambda},
 \qquad
 \|\partial_s\log|h_{a,s}'|\|_\infty\le C(1+R).
\end{equation}
The same linear-in-$R$ bound holds for products of uniformly $C^1$
Jacobian, stable-holonomy, and coarea factors.  Hence an exponential return
moment dominates the difference quotients of all clean marked amplitudes and
their return tail.
\end{lemma}

\begin{proof}
Differentiate the composition.  The contribution created at position $k$ is
multiplied on the left by at least $k-1$ contracting derivatives, so its
$C^0$ norm is bounded by $A\lambda^{k-1}$; summation gives the first bound.
The logarithmic derivative of a composition is a sum of $R$ one-step
logarithmic derivatives.  Differentiating it gives a uniformly bounded direct
term at every position and orbit-displacement terms controlled by the first
bound and uniform $C^2$ distortion.  This gives $C(1+R)$.  The product
statement follows after taking logarithms.  Finally
$R^p e^{-\epsilon R}$ is summable for every fixed $p$, so the exponential
return moment absorbs all these factors.
\end{proof}

This lemma shows that long return words do not create an exponential
parameter-derivative loss on the inverse-branch side.  The remaining issue is
not the clean word calculus but the common refinement of words that approach
or cross a moving singularity; those pieces must be charged to the already
constructed face/recovery tree rather than differentiated as clean branches.

That singular refinement is most naturally made before fixing the parameter.

\subsection{Parameter-state geometry and cutoff removal}
\label{sec:parameter-state}

\begin{definition}[Pathwise quantitative jet condition
$\mathrm{QT}_{\rm JET}$]
\label{def:quantitative-circular-jets}
Let $a(s)$ be the chosen analytic or circular table deformation with
$a(0)=a_0$.  For every face pair of depth at most $L$, let
$\Delta_{\lambda\lambda'}(s)$ be a finite collection of jet minors covering
the compact common-zero set after coincident germs are grouped.  Add the
individual face--cone boundary determinants and their critical-point jets;
these are the functions used to form the oriented tangent cuts in Lemma
\ref{lem:cone-monotone-components}.  The path satisfies
$\mathrm{QT}_{\rm JET}$ if there are $A,B,c_{\rm lab}>0$ and a
base-dependent $c_{\rm base}>0$ such that
\begin{align}
 |\Delta_{\lambda\lambda'}(0)|&\ge c_{\rm base}e^{-AL},
 &\sup_{|s|<\bar s}|\partial_s\Delta_{\lambda\lambda'}(s)|&\le Ce^{BL},
 \label{eq:pathwise-minor-bounds}\\
 \#\{(\lambda,\lambda'):|\lambda|\le L\}&\le Ce^{c_{\rm lab}L}.
 \label{eq:pathwise-label-count}
\end{align}
The constant $c_{\rm base}$ absorbs the finitely many depths not covered by
an asymptotic Borel--Cantelli statement.  No return-tail or geometric-loss
claim is part of $\mathrm{QT}_{\rm JET}$.
\end{definition}

\begin{definition}[Pathwise geometric-loss and return-tail ledger
$\mathrm{QT}_{\rm LEDGER}$]
\label{def:quantitative-path-ledger}
Condition $\mathrm{QT}_{\rm LEDGER}$ requires constants
$A_{\rm loss},c_{\rm tail}>0$ such that all inverse-angle, coarea, chart,
density, parameter-derivative, component, and Mark--$Z$ losses on the
depth-$L$ good restriction are at most $Ce^{A_{\rm loss}L}$.  Separately,
the actual unsmoothed positive-envelope mass of the complementary
return/grazing failure owner whose first required depth exceeds $L$ is at
most $Ce^{-c_{\rm tail}L}$.  No decay of the good mass itself is asserted.
Relative to the $A,B$ furnished by
$\mathrm{QT}_{\rm JET}$, the strict path margin is
\begin{equation}
\label{eq:strict-jet-tail-margin}
 c_{\rm tail}-A_{\rm loss}>A+B.
\end{equation}
Parameter randomness for the jet minors does not prove either ledger bound.
\end{definition}

\begin{definition}[Complete quantitative path condition]
\label{def:complete-QT-path}
Write
\[
 \mathrm{QT}_{\rm path}:=\mathrm{QT}_{\rm JET}
 \wedge\mathrm{QT}_{\rm LEDGER}.
\]
For a finite-dimensional ambient family, the sublevel estimate
\[
 \operatorname{Leb}\{a:|\Delta_{\lambda\lambda'}(a)|<e^{-AL}\}
 \le Ce^{-A\tau L},\qquad A\tau>c_{\rm lab},
\]
is a sufficient Borel--Cantelli mechanism for the first inequality at almost
every base parameter, but it is not a substitute for the derivative bound or
for $\mathrm{QT}_{\rm LEDGER}$.
\end{definition}

\begin{theorem}[Product-a.e. compact analytic realization of the pathwise jet bound]
\label{thm:analytic-QT-realization}
Let each scatterer boundary vary in a fixed complex analytic strip of width
$\sigma_a>0$.  Fix a strip slack $\delta_a>0$, put
$\sigma_p=\sigma_a+\delta_a$, and parametrize a neighborhood by
\begin{equation}
\label{eq:compact-analytic-product-probe}
 a_q=e^{-\sigma_p|q|}b_q,\qquad b_q\in[-1,1].
\end{equation}
Equip the $b_q$ with an independent compactly supported product law whose
one-dimensional densities are uniformly bounded.  The support of
\eqref{eq:compact-analytic-product-probe} is compact in the width-$\sigma_a$
analytic topology and hence in the chosen compact $C^3$ table chart.  When
this theorem is used with the fixed-gauge CM2 criterion, this product chart
is additionally required to be an analytic coordinate chart on the
constraint manifold \eqref{eq:constant-component-mass-vector}; independent
unconstrained modes on different components do not qualify.  Suppose
the depth-$L$ face atlas has at most $Ce^{c_{\rm lab}L}$ noncoincident jet
charts and that
\begin{equation}
\label{eq:localized-analytic-jet-direction}
 \begin{gathered}
 \text{each compact jet chart has a cover by at most }
 Ce^{c_{\rm zero}L}\text{ fiber boxes;}\\
 \text{on every box one index }q\le C_0L\text{ is fixed for the }b_q
 \text{ restriction;}\\
 \text{that restriction has at most }Ce^{c_{\rm zero}L}
 \text{ monotonicity intervals, and}\\
 |\partial_{b_q}\Delta|\ge
 e^{-(B_0+\sigma_pC_0)L}
 \quad\text{whenever }|\Delta|\le2e^{-AL}.
 \end{gathered}
\end{equation}
The fiber condition may equivalently be supplied by a uniform
Remez--Cartan/valency estimate with the same exponent $c_{\rm zero}$.  Then,
for every
\begin{equation}
\label{eq:analytic-QT-window}
 B_0+\sigma_pC_0+c_{\rm lab}+2c_{\rm zero}
 <A,
\end{equation}
product-almost every analytic base table satisfies
$|\Delta_{\lambda\lambda'}(0)|\ge e^{-AL}$ for all sufficiently large $L$.
After decreasing a base-dependent constant $c_{\rm base}>0$, the bound
$|\Delta_{\lambda\lambda'}(0)|\ge c_{\rm base}e^{-AL}$ holds at every
finite depth as well.  (The finitely many noncoincident minors have zero sets
of product measure zero under the fiber hypothesis.)
For every analytic direction whose depth-$L$ derivative cost is at most
$e^{BL}$, the resulting path satisfies $\mathrm{QT}_{\rm JET}$.  The theorem
does not supply $\mathrm{QT}_{\rm LEDGER}$, the return/grazing tail, or the
strict combined margin.  Small
analytic perturbations remain in the finite-horizon dispersing class.

For a fixed finite branch family, the adapted boundary-perturbation
construction of \cite{BFLS2026} supplies the relevant finite-jet
surjectivity template: perturb a boundary arc at the earliest differing
collision and then realize the required finite jet by an analytic
trigonometric polynomial.  Condition
\eqref{eq:localized-analytic-jet-direction} is the quantitative, simultaneous
version needed here.  Finite-branch jet surjectivity alone does not control
fiber valency.  For analytic minors one may verify the stated condition by a
quantitative propagation-of-smallness/Remez estimate such as
\cite{HWWRemez}; if the branch equations form a Nash family, the
Bernstein--Remez theorem \cite{BarbieriNash} is an alternative.  The theorem
applies only when the corresponding strip, valency, and normalized-coordinate
costs are verified for every face and cone-angle jet in the master atlas.  It
makes no assertion for the finite-dimensional circular subfamily.
\end{theorem}

\begin{proof}
Condition on all normalized Fourier coefficients except $b_q$ furnished
by \eqref{eq:localized-analytic-jet-direction}.  The one-dimensional change
of variables on each monotonicity interval and the bounded conditional density
give
\[
 \mathbb P\{|\Delta|<e^{-AL}\}
 \le C e^{2c_{\rm zero}L}
 e^{-\{A-B_0-\sigma_pC_0\}L}.
\]
The first $e^{c_{\rm zero}L}$ counts fiber boxes and the second counts
monotonicity intervals; the exponential Fourier weight is already present in
$\partial_{b_q}$.  Union over at most $Ce^{c_{\rm lab}L}$ labels.  The first inequality
in \eqref{eq:analytic-QT-window} makes the resulting series in $L$
summable, so Borel--Cantelli gives the base-minor bound.  Analytic strip norms
control the derivative along a fixed analytic direction by $Ce^{BL}$.
The finitely many remaining noncoincident minors give $c_{\rm base}$ after
excluding their null zero sets.  Finally, strict convexity, disjointness, and finite horizon are
open in the chosen compact $C^3$ neighborhood.
\end{proof}

\begin{lemma}[Persistence up to a depth depending on the parameter]
\label{lem:finite-depth-face-separation}
Assume $\mathrm{QT}_{\rm path}$ and choose
\begin{equation}
\label{eq:kappa-path-window}
 \frac1{c_{\rm tail}-A_{\rm loss}}<\kappa_L<\frac1{A+B},
 \qquad L(s)=\lfloor\kappa_L\log(1/|s|)\rfloor.
\end{equation}
For all sufficiently small $s$ and all $L\le L(s)$,
\[
 |\Delta_{\lambda\lambda'}(\theta s)|\ge\tfrac12c_{\rm base}e^{-AL},
 \qquad 0\le\theta\le1.
\]
Thus every noncoincident depth-$L(s)$ face family remains transverse along the
whole coarea segment from $0$ to $s$.
\end{lemma}

\begin{proof}
The mean-value theorem and
\eqref{eq:pathwise-minor-bounds} give
\[
 |\Delta(\theta s)-\Delta(0)|\le C|s|e^{BL}.
\]
For $L\le L(s)$ the ratio of this error to
$c_{\rm base}e^{-AL}$ is at most
$C(c_{\rm base})|s|^{1-\kappa_L(A+B)}=o(1)$.  The interval for $\kappa_L$ is nonempty precisely
because of \eqref{eq:strict-jet-tail-margin}.
\end{proof}

\begin{definition}[Shallow spacetime common refinement
$\mathrm{UCR}_{\rm sh}$]
\label{def:spacetime-common-refinement}
For every sufficiently small $s$ and every $L\le L(s)$, cut
$[0,1]_\theta\times M$ by all transported singularity, homogeneity, gate,
and coarea hypersurfaces used by words of depth at most $L$.  Condition
$\mathrm{UCR}_{\rm sh}$ requires a measurable master refinement with the
following properties.
\begin{enumerate}[label=\textup{(UCR\arabic*)}]
\item On each shallow spacetime cell the word, orientation,
      connected-component index, and actual/split-law labels are independent
      of $\theta$; the branch and endpoint maps extend to the cell and obey
      the same cone, distortion, scaled-density, and parameter-derivative
      bounds uniformly for $0\le\theta\le1$.
\item The actual native stopping antichains, branch sections,
      $\mathrm{PPE}_{N}$ labels, full-tangent
      coarea charts, positive envelopes, and connected cylinder partitions
      used at depth at most $L$ are restrictions of this one refinement.
      Adjacent artificial boundaries cancel before absolute values, and every
      genuine moving boundary keeps its component index and Mark--$Z$ charge.
      No assertion about all deeper words is included in this clause.
\item The union of shallow cells on which such a continuation fails has, for
      every actual or split constituent law, total marked mass at most
      $Ce^{-c_{\rm ucr}L}$ after the loss already recorded in
      $A_{\rm loss}$.  We replace $c_{\rm tail}$ everywhere by
      $\min\{c_{\rm tail},c_{\rm ucr}\}$ and require the same strict margin
      \eqref{eq:strict-jet-tail-margin}.  The shallow refinements are
      compatible when $L$ increases.
\end{enumerate}
The constants in
\[
 \begin{gathered}
 \mathrm{NST}_{\rm phys},\ \mathrm{BR}_{\rm sec},\ \mathrm{PPE}_{N},\\
 \mathrm{ACTUAL\_SLOPE\_BRIDGE},\ \mathrm{PAIR\_SAFE\_REFERENCE},\\
 \mathrm{PAIR\_TENERGY},\ \mathrm{FZ}_{\rm geom},\
 \mathrm{SF}_{\rm split},\ \mathrm{CP}_{\rm env},\
 \mathrm{REC}(q_{\rm rec})
 \end{gathered}
\]
\[
 \begin{gathered}
 \mathrm{CEM\_DOM},\ \mathrm{TENERGY}_{W},\
 \mathrm{DUAL\_MASS\_KERNEL},\ \mathrm{MPD},\\
 \mathrm{POST\_CARRIER\_DOM},\ \mathrm{PHYS\_TEST\_LIFT},\\
 \mathrm{PHYS\_TREE\_MATCH},\ \mathrm{PROP\_Q\_MATCH},\\
 \mathrm{MAIN\_CLOCK\_COVERAGE},\
 \mathrm{SEL\_JSL}_{p_{\rm joint}},\\
 \mathrm{CLOCK\_PROB}_{W},\\
 \mathrm{CLOCK\_SRC}_{W},\ \mathrm{KRH}_{p_W},\
 \mathrm{USC}_{\rm end}^{W}(q_g^W)
 \end{gathered}
\]
on the complete actual/split tree laws are required to be uniform on these
cells, including \eqref{eq:main-clock-coverage}, the selection-safe
likelihood bounds \eqref{eq:selection-safe-likelihood-rates}, and the
main-source base bounds \eqref{eq:main-source-base-mass}, together with
\eqref{eq:unselected-time-reference-energy} whenever the main time--Jensen
class is nonempty.
This is a primitive moving-singularity interface; compactness of the table
chart alone does not imply it.
\end{definition}

\begin{definition}[Normalized shallow long-time source
$\mathrm{SHCOL}_{\rm env}$]
\label{def:shallow-collar-envelope}
Let $\mathcal K_s^{\le L}$ be the assembled difference quotient carried by
face words of depth at most $L$.  Condition
$\mathrm{SHCOL}_{\rm env}$ is the conjunction of the following three typed
interfaces.
\begin{enumerate}[label=\textup{(SH\arabic*)}]
\item $\mathrm{SHSRC}^{+}$: before signed recombination, the exact
      actual/split expansion is dominated by propagated positive-envelope
      pieces.  If $C^{\rm sh}_{\lambda,s,L}$ are their shallow owner cells,
      put
      $M^{\rm sh}_{\lambda,s,L}:=
      \overline M_{s,\mathfrak d}^{q}
      (C^{\rm sh}_{\lambda,s,L})$, uniformly in every suppressed master
      coordinate, and
      \begin{equation}
      \label{eq:shallow-positive-source-mass}
       \Delta_{\rm sh}(s,L):=\sum_\lambda
       M^{\rm sh}_{\lambda,s,L}
       \le C|s|^{-1}e^{A_{\rm sh}L},\qquad
       \|\mathcal K_s^{\le L}\|_{\rm TV}\le\Delta_{\rm sh}(s,L).
      \end{equation}
      Artificial-face cancellation is used only after this positive size is
      recorded.  Thus $A_{\rm sh}$ already includes the operator, source
      norm, recovery, and opposite-amplitude growth absorbed in $q$; no
      unweighted physical mass is substituted for this envelope.
\item $\mathrm{SHCP}_{\rm env}$: the corresponding actual/split laws have
      one positive envelope and one connected, separation-time-admissible
      cylinder approximant.  For the bounded shallow collar functional
      $F^{\rm sh}_{s,L}$,
\begin{equation}
\label{eq:shallow-collar-envelope}
 \|F^{\rm sh}_{s,L}\|_\infty
 \le C|s|^{-1}e^{A_{\rm sh}L},\qquad
 \Omega_r^{\rm env}(F^{\rm sh}_{s,L})
 \le C|s|^{-1}e^{A_{\rm sh}L}e^{-\alpha_{\rm sh}r}.
\end{equation}
      The constants $c_{\rm FC},q_{\rm cyl}$ are those in
      Theorem~\ref{thm:law-adapted-cylinder-FCB}.
\item $\mathrm{SHREC}(q_{\rm sh})$: when
      $\Delta_{\rm sh}(s,L)>0$, put
      $\pi^{\rm sh}_{s,L}(\lambda)=
      M^{\rm sh}_{\lambda,s,L}/\Delta_{\rm sh}(s,L)$ and attach the two
      recovery times $R^{\pm,\rm sh}_{\lambda,s,L}$ and a coupling mark
      $\mathcal W^{\rm sh}_{\lambda,s,L}\ge1$.  Separately in the two
      orientations,
      \begin{equation}
      \label{eq:normalized-shallow-recovery}
       \sup_{s,L}\sum_\lambda\pi^{\rm sh}_{s,L}(\lambda)
       \mathcal W^{\rm sh}_{\lambda,s,L}
       e^{q_{\rm sh}R^{\pm,\rm sh}_{\lambda,s,L}}\le C,
       \qquad q_{\rm sh}>c_{\rm cpl}^{\rm sh}.
      \end{equation}
      After recovery the forward and reverse standard families couple at
      rate $c_{\rm cpl}^{\rm sh}$.
\end{enumerate}
This is an all-time positive-source condition; it is not inferred from a
depth-$L$ common partition.  The small-hole estimate $Ct\gamma^n$ and the
affine boundary recursion in \cite{CanestrariHoles} are its structural
model: normalize the small source, recover, and only then multiply back its
mass.
\end{definition}

\begin{lemma}[Shallow recovery keeps the positive source prefactor]
\label{lem:normalized-shallow-recovery}
Fix the central-band constant $\rho_{\rm sh}>0$ and choose
$0<\varepsilon<\rho_{\rm sh}/2$.  Under
$\mathrm{SHREC}(q_{\rm sh})$ put
\begin{equation}
\label{eq:explicit-shallow-recovery-rate}
 \begin{aligned}
 c_{\rm cen}^{\rm sh}&:=
 \min\{c_{\rm cpl}^{\rm sh}(1-2\varepsilon),q_{\rm sh}\varepsilon\},\\
 c_{\rm one}^{\rm sh}&:=
 \min\{c_{\rm cpl}^{\rm sh}(1-\varepsilon),q_{\rm sh}\varepsilon\},\\
 c_{\rm off}^{\rm sh}&:=\min\{c_{\rm cen}^{\rm sh},c_{\rm one}^{\rm sh}\}.
 \end{aligned}
\end{equation}
In the central band $\min\{m,n\}\ge\rho_{\rm sh}(m+n)$,
\begin{multline}
\label{eq:weighted-shallow-delayed-coupling}
 \sum_\lambda M^{\rm sh}_{\lambda,s,L}
 \mathcal W^{\rm sh}_{\lambda,s,L}
 \exp\{-c_{\rm cpl}^{\rm sh}[(m-R^-_\lambda)_+
 +(n-R^+_\lambda)_+]\}\\
 \le C\Delta_{\rm sh}(s,L)e^{-c_{\rm off}^{\rm sh}(m+n)} .
\end{multline}
The one-sided versions with rate $c_{\rm one}^{\rm sh}$ hold in the two
off-diagonal cones, measured in their respective long time coordinate.
\end{lemma}

\begin{proof}
Split each orientation into $R^\pm\le\varepsilon(m+n)$ and its complement.
On the central good set the recovered time is at least
$(1-2\varepsilon)(m+n)$; the two bad sets each cost
$e^{-q_{\rm sh}\varepsilon(m+n)}$.  Markov's inequality under
$\pi^{\rm sh}_{s,L}$ and
\eqref{eq:normalized-shallow-recovery} charges the two bad sets separately.
This proves the first minimum in
\eqref{eq:explicit-shallow-recovery-rate}.  In an off-diagonal cone only the
long orientation is coupled; splitting at $R^\pm\le\varepsilon$ times that
long coordinate gives the second minimum.  Multiply the normalized estimate
by $\Delta_{\rm sh}(s,L)$.  No product of the two initial recovery marks is
formed.
\end{proof}

\begin{definition}[Positive deep envelope tail
$\mathrm{TAIL}_{\rm env}^{+}$]
\label{def:deep-envelope-tail}
Let $\mathcal R_s^{>L}$ be the complete double-centered remainder formed by
words deeper than $L$.  Before signed recombination, dominate every
actual/split constituent by the restriction of the propagated positive
envelope to its deep owner cell $C_{\lambda,s,L}^{\rm deep}$ and put
$M_{\lambda,s,L}:=
\overline M_{s,\mathfrak d}^{q}(C_{\lambda,s,L}^{\rm deep})$, uniformly in
every suppressed master coordinate.  Define
\begin{equation}
\label{eq:deep-positive-envelope-mass}
 \Delta_{s,L}:=\sum_\lambda M_{\lambda,s,L}
 \le C|s|^{-1}e^{-(c_{\rm tail}-A_{\rm loss})L},
 \qquad
 \|\mathcal R_s^{>L}\|_{\rm TV}\le\Delta_{s,L}.
\end{equation}
The finite split coefficients and all operator/source/recovery costs are
included in this same $q$-mass, so $A_{\rm loss}$ is the post-absorption
loss and the gap $c_{\rm tail}-A_{\rm loss}$ must be recomputed after any
operator growth.  Connected-pullback
and boundary-oscillation constants are uniform in $\theta,s,L$.
Artificial faces cancel only in the final signed recombination; no reverse
comparison of positive envelope mass by the assembled signed TV norm is
assumed.
\end{definition}

\begin{definition}[Normalized positive-envelope deep recovery
$\mathrm{NREC}_{\rm deep}(q_{\rm deep})$]
\label{def:normalized-deep-recovery}
For $\Delta_{s,L}>0$ put
$\pi_{s,L}(\lambda)=M_{\lambda,s,L}/\Delta_{s,L}$ and attach recovery
times $R^\pm_{\lambda,s,L}$ and coupling mark
$\mathcal W_{\lambda,s,L}\ge1$.  Condition
$\mathrm{NREC}_{\rm deep}(q_{\rm deep})$ requires, separately in the two
orientations,
\begin{equation}
\label{eq:normalized-deep-recovery}
 \sup_{s,L}\sum_\lambda\pi_{s,L}(\lambda)
 \mathcal W_{\lambda,s,L}e^{q_{\rm deep}R^\pm_{\lambda,s,L}}\le C,
 \qquad q_{\rm deep}>\max\{\alpha_+,\alpha_-\}.
\end{equation}
Equivalently,
\begin{equation}
\label{eq:weighted-deep-recovery}
 \sum_\lambda M_{\lambda,s,L}\mathcal W_{\lambda,s,L}
 e^{q_{\rm deep}R^\pm_{\lambda,s,L}}\le C\Delta_{s,L}.
\end{equation}
The small positive source is normalized first, recovered second, and
multiplied back last.
\end{definition}

\begin{lemma}[Normalized standard-family criterion for deep recovery]
\label{lem:normalized-deep-recovery-criterion}
Assume that, after division by $\Delta_{s,L}$, every actual/split deep
constituent is a probability standard family and
\begin{equation}
\label{eq:normalized-initial-Z-moment}
 \sup_{s,L}\sum_\lambda\pi_{s,L}(\lambda)
 \mathcal W_{\lambda,s,L}
 (1+Z_{\lambda,s,L})^{q_{\rm deep}C_1}<\infty .
\end{equation}
If the growth lemma gives
$R^\pm_\lambda\le C_0+C_1\log(1+Z^\pm_\lambda)$ after a refinement with
the same conditional tail, then
$\mathrm{NREC}_{\rm deep}(q_{\rm deep})$ holds.
\end{lemma}

\begin{proof}
Exponentiate the recovery inequality and integrate against $\pi_{s,L}$.
Equation~\eqref{eq:normalized-initial-Z-moment} gives
\eqref{eq:normalized-deep-recovery}; multiplying by $\Delta_{s,L}$ gives
\eqref{eq:weighted-deep-recovery}.  The affine boundary recursion and leaky
standard-family construction of \cite{CanestrariHoles} are the model for
the normalized estimate.  No cancellation-stability inequality is used.
\end{proof}

\begin{proposition}[Depth-dependent coarea and pathwise scalar tail]
\label{prop:parameter-state-master-refinement}
Assume $|\nabla_xG_\lambda|\ge c_G>0$ on every active face and
$\mathrm{QT}_{\rm path}$, $\mathrm{UCR}_{\rm sh}$,
$\mathrm{TAIL}_{\rm env}^{+}$, the joint Mark--$Z$ ledger
\eqref{eq:joint-mark-Z-ledger}, and the physical/source identification
\eqref{eq:propagated-q-physical-match}.  In particular the return/grazing
tail and $A_{\rm loss}$ in the conclusion are the actual positive-envelope
quantities furnished by $\mathrm{QT}_{\rm LEDGER}$, not consequences of
$\mathrm{UCR}_{\rm sh}$.  For all depths
$L\le L(s)$,
\begin{equation}
\label{eq:master-coarea-difference}
 \frac{\one_{\{G_\lambda(s,\cdot)>0\}}
             -\one_{\{G_\lambda(0,\cdot)>0\}}}{s}
 =\int_0^1\delta(G_\lambda(\theta s,\cdot))
   \partial_sG_\lambda(\theta s,\cdot)\,d\theta
\end{equation}
is a uniformly controlled coarea current.  The omitted deeper words satisfy
\begin{equation}
\label{eq:finite-depth-tail}
 |s|^{-1}\,Ce^{-(c_{\rm tail}-A_{\rm loss})L(s)}
 \le C|s|^{\kappa_L(c_{\rm tail}-A_{\rm loss})-1}=o(1).
\end{equation}
This is a scalar total-variation tail.  Its conversion to an $o(1)$ tail in
the infinite CM2 double sum is a separate time--depth argument below; no such
conversion is inferred from \eqref{eq:finite-depth-tail} alone.  No fixed
neighborhood is asserted to be uniformly transverse at every depth.
\end{proposition}

\begin{proof}
Lemma~\ref{lem:finite-depth-face-separation} preserves noncoincident faces,
while $\mathrm{UCR}_{\rm sh}$ supplies the common labels, components, branches, and
uniform geometric bounds along $\theta s$ for shallow words.  The return/grazing tail,
the joint mark--$Z$ ledger, and the allowed geometric loss give
$Ce^{-(c_{\rm tail}-A_{\rm loss})L}$ before the difference quotient.  Apply
the total-variation bound to both endpoints for words deeper than $L(s)$ and
divide by $|s|$.  The first inequality in
\eqref{eq:kappa-path-window} gives \eqref{eq:finite-depth-tail}.
\end{proof}

\begin{lemma}[Finite-parameter collar as a law-adapted functional]
\label{lem:finite-s-collar-admissibility}
Assume $\mathrm{FZ}_{\rm geom}(\theta_Z)$,
$\mathrm{SF}_{\rm split}$, $\mathrm{CP}_{\rm env}$,
$\mathrm{UCR}_{\rm sh}$, $\mathrm{TAIL}_{\rm env}^{+}$, and
$\mathrm{NREC}_{\rm deep}(q_{\rm deep})$.
Let $\mathcal R_s^{>L}$ be the complete double-centered part of
$s^{-1}(P_s-P_0)$ carried by return/grazing words deeper than $L$.  Before
the coarea limit it has a full-dimensional density in every local event chart,
but its CM2 pairing is
\begin{equation}
\label{eq:finite-s-collar-law-pairing}
 B^{\rm deep}_{m,n}(s)=
 \int F_{s,L}^{\rm col}\,d\Sigma_{m,n}^{\rm law},
\end{equation}
where $\Sigma_{m,n}^{\rm law}$ is the exact finite expansion
\eqref{eq:exact-split-law-expansion}.  It is therefore not an integral against
$\mu^4$.
There are fixed $A_{\rm col}>0$ and $\alpha_{\rm col}\in(0,1]$ such that
\begin{align}
 \|\mathcal R_s^{>L}\|_{\rm TV}
 &\le \Delta_{s,L}\le C|s|^{-1}e^{-(c_{\rm tail}-A_{\rm loss})L},
 \label{eq:deep-collar-TV}\\
 \|F_{s,L}^{\rm col}\|_\infty
 &\le C|s|^{-1}e^{A_{\rm col}L},\qquad
 \Omega_r^{\rm env}(F_{s,L}^{\rm col})
 \le C|s|^{-1}e^{A_{\rm col}L}e^{-\alpha_{\rm col}r}.
 \label{eq:deep-collar-dynamical-tube}
\end{align}
The same pieces, disintegrated in either time orientation, are signed
differences of regular measured unstable families after their recovery
times, and their recovery-weighted mass satisfies
\eqref{eq:weighted-deep-recovery} with
$\Delta_{s,L}$ bounded by the right side of
\eqref{eq:deep-collar-TV}.
\end{lemma}

\begin{proof}
On the spacetime cells supplied by $\mathrm{UCR}_{\rm sh}$, the difference of the two moving domains is
a collar of width $O(|s|)$ and division by $s$ costs exactly one power of
$|s|^{-1}$.  Expanding the test-side projection before taking an absolute
value gives \eqref{eq:finite-s-collar-law-pairing} and
\eqref{eq:exact-split-law-expansion}.  The sharp indicator is not dynamically
H\"older.  Apply instead Proposition~\ref{prop:envelope-dynamical-tube}; this gives
\eqref{eq:deep-collar-dynamical-tube}.  Absolute summation with the
return/grazing mass gives \eqref{eq:deep-collar-TV}.  Disintegrating the collar
along the usable unstable orientation preserves positive envelope mass.
Normalize this disintegration by $\Delta_{s,L}$ before applying the affine $Z$ recursion,
the short-curve tail, and Lemma~\ref{lem:signed-observable-multiplier}.
The primitive normalized estimate
$\mathrm{NREC}_{\rm deep}(q_{\rm deep})$ then gives
\eqref{eq:weighted-deep-recovery}; only afterwards is the factor
$\Delta_{s,L}$ restored.  Eventual recovery of each piece alone would not
imply this weighted estimate.
\end{proof}

\begin{theorem}[Positive-envelope cylinder approximation]
\label{thm:law-adapted-cylinder-FCB}
Let $F$ be a bounded four-state functional and let the signed law and positive
envelope be \eqref{eq:exact-split-law-expansion}--
\eqref{eq:positive-envelope-law}.  Suppose
$\Omega_r^{\rm env}(F)\le Ve^{-\alpha r}$ and the common connected partition
is separation-time admissible as in $\mathrm{CP}_{\rm env}$.  On every rank-$r$ connected
product cylinder choose a value in the
$\bar\rho_{m,n}$-essential range of $F$ and call the resulting common cylinder
function $F_r$.  There is a fixed $q_{\rm cyl}>0$ such that
\begin{align}
 \|F_r\|_{\rm SH}&\le C\|F\|_\infty e^{q_{\rm cyl}r},
 \label{eq:cylinder-approximant-SH}\\
 \left|\int(F-F_r)\,d\Sigma_{m,n}^{\rm law}\right|
 &\le Z_0Ve^{-\alpha r}.
 \label{eq:cylinder-approximant-error}
\end{align}
Consequently a functional-correlation split across a gap $N$ obeys
\begin{equation}
\label{eq:cylinder-functional-correlation}
 |\operatorname{FC}_N(F)|
 \le C\left\{Ve^{-\alpha r}
       +\|F\|_\infty e^{-c_{\rm FC}N+q_{\rm cyl}r}\right\}.
\end{equation}
If two mutually singular constituent laws see the values zero and one on the
same cylinder, their positive envelope sees both values and the envelope
oscillation is one.  For the diagonal example, a value constant on every
constituent support is represented exactly; if a split constituent sees the
opposite side, the envelope charges that discrepancy.  Thus the construction
never replaces a law-adapted zero error by coordinatewise mollification across
the diagonal.
\end{theorem}

\begin{proof}
On each cylinder, the envelope-essential difference from the chosen value is
bounded by the envelope-essential oscillation.  Using
\eqref{eq:positive-envelope-law},
\[
 \int|F-F_r|\,d\rho^0+
 \sum_j|a_j|\int|F-F_r|\,d\rho^j
 =Z_{\rm law}\int|F-F_r|\,d\bar\rho_{m,n},
\]
which proves \eqref{eq:cylinder-approximant-error}.  Separation-time
admissibility in $\mathrm{CP}_{\rm env}$ says that a first possible jump
occurs only after some separation time drops below $r$; it costs at most
$C\theta_0^{-r}=Ce^{q_{\rm cyl}r}$ in that slot.  This proves
\eqref{eq:cylinder-approximant-SH}.  Apply \cite[Theorems~2.3--2.4]{LS} to
$F_r$ for the exact signed split and use
\eqref{eq:cylinder-approximant-error}.  The same $F_r$ serves all constituent
laws because it was constructed from their common positive envelope.  No
assertion about conditional coordinate marginals is used.
\end{proof}

\begin{theorem}[Two-threshold removal of shallow long times]
\label{thm:shallow-two-threshold}
Let
\[
 L(s)=\lfloor\kappa_L\log(1/|s|)\rfloor,
 \qquad T_{\rm sh}(s)=\lfloor\kappa_T^{\rm sh}\log(1/|s|)\rfloor,
\]
and put
\begin{equation}
\label{eq:shallow-collar-rate}
 c_{\rm sh}:=\min\left\{c_{\rm off}^{\rm sh},
 \frac{\alpha_{\rm sh}c_{\rm FC}}
 {q_{\rm cyl}+\alpha_{\rm sh}}\right\}.
\end{equation}
Let
\begin{equation}
\label{eq:intermediate-resolving-constant}
 C_{\rm mid}(\eta_{\rm hi})
 :=C_{\rm ep}+C_{\rm hom}
   +C_{\rm buf}\delta(C_{\rm scl}+\eta_{\rm hi}),
\end{equation}
where $C_{\rm hom}$ contains the fixed homogeneous and component buffers.
Assume the complete labelled ledger \eqref{eq:global-ledger-sum} and
\[
 \begin{gathered}
 \mathrm{UCR}_{\rm sh},\ \mathrm{NST}_{\rm phys},\ \mathrm{SS},\
 \mathrm{ACTUAL\_SLOPE\_BRIDGE},\\
 \mathrm{PAIR\_SAFE\_REFERENCE},\ \mathrm{PAIR\_TENERGY},\\
 \mathrm{PPE}_{N},\ \mathrm{CEM\_DOM},\ \mathrm{TENERGY}_{W},\\
 \mathrm{DUAL\_MASS\_KERNEL},\ \mathrm{MPD},\\
 \mathrm{POST\_CARRIER\_DOM},\ \mathrm{PHYS\_TEST\_LIFT},\\
 \mathrm{PHYS\_TREE\_MATCH},\ \mathrm{PROP\_Q\_MATCH},\\
 \mathrm{MAIN\_CLOCK\_COVERAGE},\
 \mathrm{SEL\_JSL}_{p_{\rm joint}}
 \end{gathered}
\]
\[
\begin{gathered}
 \mathrm{CLOCK\_PROB}_{W},\ \mathrm{CLOCK\_SRC}_{W},\
 \mathrm{KRH}_{p_W},\ \mathrm{USC}_{\rm end}^{W},\
 \mathrm{BR}_{\rm sec},\ \mathrm{REC},\ \mathrm{MT}_{\rm DQ},\\
 \mathrm{FACE\_TIME}_{\rm CM2},\ \mathrm{SHCOL}_{\rm env}.
\end{gathered}
\]
Assume also $\mathrm{MAIN\_CLOCK\_COVERAGE}$ and the main-source likelihood
and base bounds \eqref{eq:selection-safe-likelihood-rates} and
\eqref{eq:main-source-base-mass}, uniformly on the same shallow refinement.
For every nonempty main time--Jensen cell also assume
\eqref{eq:unselected-time-reference-energy} on the same frozen parents.
Require also the strict window
\begin{equation}
\label{eq:shallow-time-window}
 C_{\rm mid}(\eta_{\rm hi})\kappa_T^{\rm sh}<\kappa_L,
 \qquad
 \kappa_T^{\rm sh}c_{\rm sh}>1+\kappa_LA_{\rm sh},
 \qquad \eta_{\rm hi}\rho_\infty>a_{\rm ultra}.
\end{equation}
For $N=\max\{m,n\}\le T_{\rm sh}(s)$ the $r/$PPE/time--Jensen and TV
owners have the product-CM2 majorant from the three-frequency proof, while
every main and nonmain face/word owner has the independent majorant
\eqref{eq:face-time-CM2-majorant}.  Thus the complete shallow difference
quotient has the same product-CM2 majorant as the limiting shallow current.
For
$N>T_{\rm sh}(s)$,
\begin{equation}
\label{eq:shallow-long-time-bound}
 |B^{\rm sh}_{m,n}(s)|
 \le C|s|^{-1-A_{\rm sh}\kappa_L}e^{-c_{\rm sh}N},
\end{equation}
and consequently
\begin{equation}
\label{eq:shallow-long-time-sum}
 \sum_{\max\{m,n\}>T_{\rm sh}(s)}
 |B^{\rm sh}_{m,n}(s)|=o(1).
\end{equation}
For fixed parameters satisfying the third inequality, the two-sided
$\kappa_T^{\rm sh}$ interval is nonempty precisely when
$C_{\rm mid}(\eta_{\rm hi})(1+\kappa_LA_{\rm sh})
 <\kappa_Lc_{\rm sh}$.
\end{theorem}

\begin{proof}
If $N\le T_{\rm sh}(s)$, the low block uses only the endpoint/section
buffers.  In the intermediate block, \eqref{eq:intermediate-max-resolving-depth}
and the fixed homogeneous buffers give total depth at most
$C_{\rm mid}(\eta_{\rm hi})N<L(s)$.  The ultra-high block uses
\eqref{eq:ultra-high-observable-tail} and no escape or resolving refinement.
Scale-bad labels are removed by Lemma~\ref{lem:endpoint-upper-scale}; on
scale-good labels the $C_{\rm scl}$ term in $C_{\rm mid}$ is exactly the
upper endpoint-scale cost.  Thus the common shallow refinement supports the
complete three-frequency proof for its Jensen/TV handlers without an
infinite-depth call.  The face/word handlers removed from that proof are
estimated, including their fixed core, by
Lemma~\ref{lem:face-time-majorant}; $\mathrm{FACE}_{\rm 2CUT}$ is used only
for their word-depth convergence.  For $N>T_{\rm sh}(s)$ do not continue
that Fourier proof.  In the central band apply
Theorem~\ref{thm:law-adapted-cylinder-FCB} to the finite-$s$ collar in
Definition~\ref{def:shallow-collar-envelope} with
\[
 r_N=\left\lfloor\frac{c_{\rm FC}N}
 {q_{\rm cyl}+\alpha_{\rm sh}}\right\rfloor .
\]
The two off-diagonal cones use
Lemma~\ref{lem:normalized-shallow-recovery}; its retained factor
$\Delta_{\rm sh}(s,L)$ is bounded by
\eqref{eq:shallow-positive-source-mass}.  The definition of $c_{\rm sh}$ makes
\eqref{eq:shallow-long-time-bound} valid there as well.  Since there are $O(N)$ pairs
with maximum $N$, its sum is bounded by
\[
 C|s|^{-1-A_{\rm sh}\kappa_L}
 \sum_{N>T_{\rm sh}(s)}Ne^{-c_{\rm sh}N}=o(1)
\]
by the second inequality in \eqref{eq:shallow-time-window}.  The first
inequality is an upper bound for $\kappa_T^{\rm sh}$ and the second a lower
bound, which gives the last assertion.
\end{proof}

\begin{remark}[Noncircular order of exponent choices]
\label{rem:parameter-selection-order}
First prove $\mathrm{CEM\_DOM}$, fix the complete event weight
$\mathscr W_{\lambda,*}$, and fix the five refinement-moment exponents and
the genuine complete-tree cemetery exponent $q_C$.  Compute $\chi_W$ in
\eqref{eq:full-product-reciprocal-cost}; if $\chi_W<1$, choose the actual
$p_W$ from \eqref{eq:available-full-product-pW}.  For a registry, first
subtract aggregate complexity and then degrade the resulting partition rate
by $1/p_W'$ only under the appropriate global probability or normalized
incidence-law moment.  Context, stopping, PPE, and grazing rates are degraded
on their corresponding parent kernels only after
Lemma~\ref{lem:kernelwise-product-KRH}, or an independent
$\mathrm{KRH}_{p_W}$ primitive, has been verified.  Otherwise
$\mathrm{TENERGY}_{W}$,
$\mathrm{CLOCK\_PROB}_{W}$, $\mathrm{CLOCK\_SRC}_{W}$,
$\mathrm{KRH}_{p_W}$, and $\mathrm{USC}_{\rm end}^{W}$ remain independent
primitive inputs.  Complete both the probability partition and the
source-incidence clock matrix before choosing any phase parameter; no bare
atomic event rate survives this step.  Next fix the coefficient
of $\log S_\lambda^\pm$ in
$\widetilde G_N^\pm$ and choose $\zeta_g$; this simultaneously sets
$c_{\rm scl}=q_g^W\zeta_g$ and enlarges $C_{\rm scl}$.  Then choose $\delta$
so that $\alpha-A_{\rm frac}\delta>0$.  Next choose
the smallest admissible $C_{\rm buf}$ that gives
$\kappa^{N_{\rm PPE}(r)}\le r/100$ and the named tail bounds in
\eqref{eq:complete-phase-exponent}.  Then choose $\eta_{\rm hi}$ so that
\eqref{eq:ultra-high-exponent-window} holds, and verify that
\[
 \frac{1+\kappa_LA_{\rm sh}}{c_{\rm sh}}
 <\kappa_T^{\rm sh}<
 \frac{\kappa_L}{C_{\rm mid}(\eta_{\rm hi})}
\]
is genuinely nonempty.  Only afterwards are $\eta_{\rm lo}$ and the recovery
splitting parameters chosen.  Increasing $C_{\rm buf}$ is not monotone for
the proof: it improves the derived Jensen value rate
$\beta_{\rm PPE,J}^{\rm val}$ and the direct TV rate
$\beta_{\rm PPE,TV}^{\rm src}$ but enlarges $C_{\rm mid}$.
Likewise increasing $\zeta_g$ improves the weighted scale-bad exponent while
enlarging $C_{\rm scl}$ and hence $C_{\rm mid}$; neither choice may be made
after the shallow window has been declared nonempty.
If the new cemetery moment makes $\chi_W\ge1$, or if any nonempty clock
class has zero degraded rate, this route is a genuine no-go unless the
corresponding weighted interface is assumed directly.  In particular one
may not conceal the loss by choosing $\delta$, $C_{\rm buf}$, or
$\eta_{\rm lo}$ after the fact.
\end{remark}

\begin{theorem}[$(L,T,r)$ removal of the deep tail]
\label{thm:time-depth-diagonal}
Assume $\mathrm{FZ}_{\rm geom}(\theta_Z)$,
$\mathrm{SF}_{\rm split}$, $\mathrm{CP}_{\rm env}$,
$\mathrm{UCR}_{\rm sh}$, $\mathrm{TAIL}_{\rm env}^{+}$, and
$\mathrm{NREC}_{\rm deep}(q_{\rm deep})$, together with the finite-horizon
functional-correlation estimate and the forward/reverse moving-family
coupling bounds used in Lemma~\ref{lem:three-region}.  Assume also
$\mathrm{QT}_{\rm path}$, including the constants $A,B$ and the path window
\eqref{eq:kappa-path-window}.  Thus
Lemma~\ref{lem:finite-s-collar-admissibility} and
Theorem~\ref{thm:law-adapted-cylinder-FCB} apply uniformly to the same
master cells.  No conclusion of this theorem is available from a scalar TV
tail alone.
Put $c_d=c_{\rm tail}-A_{\rm loss}$ and choose
$L(s)=\lfloor\kappa_L\log(1/|s|)\rfloor$ with
$\kappa_Lc_d>1$ and $\kappa_L(A+B)<1$, as in
\eqref{eq:kappa-path-window}.  Set
\begin{equation}
\label{eq:time-depth-threshold}
 c_{\rm col}:=
 \frac{\alpha_{\rm col}c_{\rm FC}}
      {q_{\rm cyl}+\alpha_{\rm col}},\qquad
 T(s)=\lfloor\kappa_T\log(1/|s|)\rfloor,\qquad
 \kappa_Tc_{\rm col}>1+\kappa_LA_{\rm col}.
\end{equation}
Let
\[
 B^{\rm deep}_{m,n}(s)=
 \langle Q_s^n\mathcal R_s^{>L(s)}Q_s^mg,f\rangle
\]
with the fixed-gauge centered operators; replacing one side by $Q_0$ gives
the same conclusion.  Then
\begin{equation}
\label{eq:deep-tail-CM2-sum}
 \sum_{m,n\ge0}|B^{\rm deep}_{m,n}(s)|=o(1),
 \qquad s\to0.
\end{equation}
\end{theorem}

\begin{proof}
Write
$\Delta_s=C|s|^{-1}e^{-c_dL(s)}
=O(|s|^{\kappa_Lc_d-1})=o(1)$.
For $N=\max\{m,n\}\le T(s)$, total variation and Markov normalization give
\[
 \sum_{m,n\le T(s)}|B^{\rm deep}_{m,n}(s)|
 \le C\Delta_sT(s)^2=o(1).
\]
In the two off-diagonal regions, disintegrate the collar as in Lemma
\ref{lem:finite-s-collar-admissibility}.  Forward and reverse moving-table
coupling for one piece gives, in the two orientations,
\begin{align*}
 C M_\lambda\mathcal W_\lambda
 e^{\beta_-m-\alpha_+n+\alpha_+R^+_\lambda},\qquad
 C M_\lambda\mathcal W_\lambda
 e^{-\alpha_-m+\beta_+n+\alpha_-R^-_\lambda}.
\end{align*}
Because $q_{\rm deep}>\max\{\alpha_+,\alpha_-\}$,
\eqref{eq:weighted-deep-recovery} bounds each sum over $\lambda$ by
$C\Delta_s$ times the corresponding cone exponential.  Choosing the cone
slopes as in Lemma~\ref{lem:three-region} therefore makes the complete
off-diagonal sums $O(\Delta_s)$.  This is the only step where the normalized
deep-recovery hypothesis is used; total mass plus eventual recovery would
not suffice.

It remains to consider the central band with $N>T(s)$.  The collar itself,
not its singular coarea limit, is the sharp four-state function
$F_{s,L(s)}^{\rm col}$.  It is not dynamically H\"older.  Apply
Theorem~\ref{thm:law-adapted-cylinder-FCB} with the $N$-dependent rank
\[
 r_N=\left\lfloor
 \frac{c_{\rm FC}N}{q_{\rm cyl}+\alpha_{\rm col}}
 \right\rfloor.
\]
Global double centering kills both product terms, and
\eqref{eq:deep-collar-dynamical-tube}--
\eqref{eq:cylinder-functional-correlation} give
\[
 |B^{\rm deep}_{m,n}(s)|
 \le C|s|^{-1-A_{\rm col}\kappa_L}e^{-c_{\rm col}N}.
\]
There are $O(N)$ pairs with maximum $N$, so
\[
 \sum_{\substack{m\asymp n\\\max(m,n)>T(s)}}
 |B^{\rm deep}_{m,n}(s)|
 \le C|s|^{-1-A_{\rm col}\kappa_L}
       \sum_{N>T(s)}Ne^{-c_{\rm col}N}=o(1)
\]
by \eqref{eq:time-depth-threshold}.  The three regions exhaust the quadrant,
proving \eqref{eq:deep-tail-CM2-sum}.
\end{proof}

This construction is why a differentiable parameter family of Young towers
is unnecessary.  The dynamical return alphabet is fixed on clean master
cells, while every atom which changes its itinerary is represented by the
physical coarea current that the response formula already contains.

The clean cohomological bracket and the master refinement can now be put on one ledger.  The
point of the next proposition is not to create a second symbolic dynamics.
It records, on the amplitude side, which already existing full branches and
recovery descendants contain a physical face piece.

\begin{lemma}[Global complete-tree refinement and full-product $L^{p_W}$]
\label{lem:measurable-refinement-moments}
Let $\mathcal F_0\subset\cdots\subset\mathcal F_5$ be the finite or countable
label filtration.  Its successive increments record the physical
chart/return word, the two endpoint contexts, the native stopping antichain,
the homogeneity cuts with component index, and the time
orientation/path-depth label.  At every stage let
$R_i(\lambda_{i-1},d\lambda_i)$ be the normalized
mass kernel of the refinement.  Assume
\begin{equation}
\label{eq:refinement-kernel-mass}
 R_i(\lambda_{i-1},\Lambda_i)=1
 \quad\text{for every parent, including cemetery continuations,}
\end{equation}
and assume the complete-tree tilt normalization already required in
\eqref{eq:complete-energy-tree-tilt}:
\begin{equation}
\label{eq:full-product-weight-normalization}
 1\le\mathbb E_{\Pi_{s,N}}\mathscr W_*\le C_W.
\end{equation}
and, with $\mathfrak F$ denoting a positive mark--$Z$ envelope and
$\rho_{\rm PT}$ the predictable factor from
\eqref{eq:physical-tree-RN-identity}, require the
pointwise domination
\begin{equation}
\label{eq:full-mark-cemetery-domination}
 (1+\rho_{\rm PT})\operatorname{Mark}^2
 \{1+Z(\mathcal G^-)+Z(\mathcal G^+)\}\le C\mathfrak F
\end{equation}
on every actual, split, survivor, and cemetery continuation.  Put
$\mathcal C=1$ on survivors and $\mathcal C=\mathcal C_i$ on $B_i$.
With unobserved pre-failure factors replaced by the harmless upper-bound
value one only after $\mathcal C_i$ has been inserted, one has on the entire
tree
\begin{equation}
\label{eq:full-cemetery-product-envelope}
 (1+\rho_{\rm PT})\mathscr W_*\le
 e^{\epsilon(|c^-|+|c^+|)}(1+H)^{a_H}
 \mathfrak d^{A_{\rm dom}}\mathfrak F
 S^{a_{\rm scl}}\mathcal C.
\end{equation}
Assume
\begin{align}
 \mathbb E[(1+H)^{q_H}\mid\mathcal F_1]&\le C,&
 \mathbb E[\mathfrak F^{q_F}\mid\mathcal F_2]&\le C,
 \label{eq:conditional-refinement-moments}\\
 \mathbb E[\mathfrak d^{A_{\rm dom}q_D}\mid\mathcal F_1]&\le C,&
 \mathbb E[S_\lambda^{a_{\rm scl}q_S}\mid\mathcal F_3]&\le C,
 \notag\\
 \mathbb E[e^{q_R\epsilon(|c^-|+|c^+|)}\mid\mathcal F_0]&\le C.
 \notag
\end{align}
For the piecewise cemetery factor require instead the genuine complete-tree
joint moment
\begin{equation}
\label{eq:complete-tree-cemetery-moment}
 \sup_{s,N,\,\mathrm{actual/split}}
 \mathbb E_{\Pi_{s,N}}[\mathcal C^{q_C}]\le C.
\end{equation}
Equivalently one may specify a strictly earlier parent sigma-field
$\mathcal G_i$ and a normalized child kernel $K_i(\omega,d\xi)$ which has not
yet revealed $\mathcal C_i$, and require
\begin{equation}
\label{eq:parent-kernel-cemetery-moment}
 \int \mathcal C_i(\omega,\xi)^{q_C}K_i(\omega,d\xi)\le C.
\end{equation}
Conditioning on $\mathcal F_{\tau_i}$ itself is not an alternative: because
$\mathcal C_i$ is already $\mathcal F_{\tau_i}$-measurable, such a condition
would say $\mathcal C_i^{q_C}\le C$ pointwise and would not express an
unbounded moment.  These estimates are imposed on each complete
actual/split tree law with its cemetery continuations; they are not asserted
after conditioning on every rare frozen parent.  The conditional
$q_F$ estimate is a primitive stronger than the unconditional flux sum
\eqref{eq:primitive-flux-ledger}, unless a separate transfer theorem to all
these laws has been proved.  If $\mathcal C$ is pointwise absorbed by an
already listed factor, take $q_C=\infty$ and omit the last reciprocal cost;
otherwise $q_C$ is a genuinely new primitive moment.
Define the full reciprocal cost
\begin{equation}
\label{eq:full-product-reciprocal-cost}
 \chi_W:=\frac{a_H}{q_H}+\frac1{q_F}+\frac1{q_D}
       +\frac1{q_S}+\frac1{q_R}+\frac1{q_C}.
\end{equation}
If $\chi_W<1$, then every
\begin{equation}
\label{eq:available-full-product-pW}
 1<p_W<\chi_W^{-1}
\end{equation}
is available uniformly on the complete actual/split tree as a global law,
and
\begin{equation}
\label{eq:proved-full-product-Lp}
 \|(1+\rho_{\rm PT})\mathscr W_*\|_{L^{p_W}(\Pi_{s,N})}
 \le C,\qquad
 1\le\mathbb E_{\Pi_{s,N}}\mathscr W_*\le C_W.
\end{equation}
Thus both the bare and the physical-Radon--Nikodym-weighted versions have
the same $L^{p_W}$ margin.  In particular
\eqref{eq:full-product-Lp-energy-route} follows from the
displayed moment assumptions with this same $p_W$.  This is a statement on
the complete global law: it does not imply \(\mathrm{KRH}_{p_W}\) on rare
frozen-parent cells and it does not control source-incidence mass unless an
incidence-space product moment is separately assumed.  All coarser global
marginals, including the positive cemetery masses, are unchanged by the
refinement.
\end{lemma}

\begin{proof}
Equation \eqref{eq:refinement-kernel-mass} and Ionescu--Tulcea construct one
complete probability law and show by iterated conditional expectation that
each earlier marginal is preserved.  The displayed conditional estimates
imply their global moments on this same complete law.  Raise the piecewise weight
$(1+\rho_{\rm PT})\mathscr W_*$ to $p_W$.  On a survivor use
\eqref{eq:global-event-weight}; on $B_i$ use
\eqref{eq:cemetery-kernel-domination}.  Generalized H\"older on the complete tree
uses reciprocal costs
$p_W/q_F$, $p_W/q_D$, $p_W/q_S$, $p_W/q_R$, and
$p_Wa_H/q_H$, while the cemetery step uses $p_W/q_C$.  Their sum is
$p_W\chi_W<1$, leaving one strict conjugate slot.  The genuine joint moment
\eqref{eq:complete-tree-cemetery-moment}, or its earlier-parent kernel
realization \eqref{eq:parent-kernel-cemetery-moment}, therefore proves
\eqref{eq:proved-full-product-Lp}.  The normalization
\eqref{eq:full-product-weight-normalization}, not a false pointwise floor on
every cemetery envelope, gives the second part of that display.  This proves an actual
$L^{p_W}$ margin rather than mere integrability and does not multiply moments
taken under different label laws.
Coincident artificial faces are cancelled before $R_1$ is defined, so the
construction also preserves the signed physical marginal.
\end{proof}

\begin{lemma}[Kernelwise product criterion for \(\mathrm{KRH}_p\)]
\label{lem:kernelwise-product-KRH}
For every frozen parent $\eta$ of
Definition~\ref{def:kernel-reverse-holder}, suppose the normalized descendant
kernel $K_\eta$ carries the same pointwise product envelope as
\eqref{eq:full-cemetery-product-envelope}, with all component moments in
\eqref{eq:conditional-refinement-moments}--
\eqref{eq:parent-kernel-cemetery-moment} valid under $K_\eta$, uniformly in
$\eta$.  The assertion includes actual, split, survivor, cemetery, and
source-incidence parents.  If the reciprocal cost is bounded uniformly by
$\chi_W<1$, also require the genuinely relative normalization
\begin{equation}
\label{eq:kernelwise-weight-normalization}
 0<c_K\le\int W_\eta\,dK_\eta\le C_K<\infty
 \qquad\text{uniformly in }\eta.
\end{equation}
Then every $1<p_W<\chi_W^{-1}$ satisfies
\begin{equation}
\label{eq:kernelwise-product-KRH}
 \sup_\eta\int
 \left(\frac{W_\eta}{\int W_\eta\,dK_\eta}\right)^{p_W}dK_\eta\le C.
\end{equation}
Hence \(\mathrm{KRH}_{p_W}\) holds.
\end{lemma}

\begin{proof}
Generalized H\"older is applied separately inside each normalized kernel
$K_\eta$, using the same strict reciprocal margin.  Its numerator has a
uniform $L^{p_W}$ bound, while
\eqref{eq:kernelwise-weight-normalization} controls the denominator.  On a
source-incidence parent one first normalizes the positive incidence mass;
the same displayed lower normalization is required for the resulting
likelihood.  The
constants are independent of $m,n,s$, Fourier mode, polynomial, and final
observable.  No global-to-conditional passage is used.
\end{proof}

\begin{corollary}[Global full-product H\"older transfer]
\label{cor:full-product-holder-energy}
Assume the complete hypothesis package of
Lemma~\ref{lem:measurable-refinement-moments}: normalized refinement
kernels, pointwise Mark--$Z$ domination, the full product envelope, and
$\chi_W<1$.  Assume also global unweighted aggregate
probability-clock bounds on that same complete tree.  Choose $p_W$ by
\eqref{eq:available-full-product-pW}.  Then the corresponding global
\(\mathrm{CLOCK\_PROB}_{W}\) bounds hold with every already aggregated
polynomial and exponential rate degraded by
\[
 \Theta_{\rm agg}^{W}=\Theta_{\rm agg}/p_W',\qquad
 c_{\rm agg}^{W}=c_{\rm agg}/p_W',
\]
provided registry complexity has already been subtracted before this single
H\"older step.  A separate global incidence-space $L^{p_W}$ bound gives the
analogous global \(\mathrm{CLOCK\_SRC}_{W}\) transfer.  Neither conclusion
implies parentwise \(\mathrm{TENERGY}_{W}\), \(\mathrm{USC}_{\rm end}^{W}\),
or PPE tails.
\end{corollary}

\begin{proof}
For any tree event $B$, H\"older and
\eqref{eq:proved-full-product-Lp} give
\[
 \Pi_{s,N}^{W}(B)
 =\frac{\mathbb E_{\Pi_{s,N}}[\mathscr W_*\one_B]}
        {\mathbb E_{\Pi_{s,N}}\mathscr W_*}
 \le C\Pi_{s,N}(B)^{1/p_W'}.
\]
Apply this once to each already aggregated global probability event.  For
source occurrences the same calculation is valid only after integrating
the positive envelope measure $M_{\rm src}^{q}$ of
Definition~\ref{def:source-incidence-space}.  No branchwise Jensen,
conditional transfer, or signed cancellation occurs here.
\end{proof}

\begin{corollary}[Parent-kernel H\"older transfer]
\label{cor:parent-kernel-holder-energy}
Assume the raw context, stopping, buffer, gate, continuation, retained PPE,
and upper-scale tails uniformly under every frozen parent kernel $K_\eta$.
If \(\mathrm{KRH}_{p_W}\) holds, then the corresponding parentwise tilted
tails in \(\mathrm{TENERGY}_{W}\), retained PPE, and
\(\mathrm{USC}_{\rm end}^{W}\) hold, with each raw exponent divided by
$p_W'$.  If Lemma~\ref{lem:kernelwise-product-KRH} is used as the sufficient
route, all of its component moments and the uniform lower normalization
\eqref{eq:kernelwise-weight-normalization} must be verified under every such
kernel; the global moment Lemma~\ref{lem:measurable-refinement-moments} is
not a substitute.
\end{corollary}

\begin{proof}
Apply Lemma~\ref{lem:parentwise-holder-transfer} inside each normalized
kernel, retaining the survivor indicator on PPE and upper-scale events.
Uniformity in $\eta$ gives exactly the quantifiers required by the three
parentwise interfaces.
\end{proof}

\begin{definition}[Physical/tree-law identification
$\mathrm{PHYS\_TREE\_MATCH}$]
\label{def:physical-tree-match}
Let $\mathfrak M_s^{\rm phys}:=\sum_\lambda M_\lambda\delta_\lambda$
be the positive physical label measure used by the face, shallow, and deep
ledgers, and let $p_{\rm phys}$ be the Borel projection from the complete
refinement tree to its physical label.  The interface requires a predictable
Radon--Nikodym factor $\rho_{\rm PT}\ge0$ and a normalization
$0<Z_-\le Z_{s,\mathfrak d}\le Z_+<\infty$ such that, for every nonnegative
Borel $F$,
\begin{equation}
\label{eq:physical-tree-RN-identity}
 \sum_\lambda M_\lambda F(\lambda)
 =Z_{s,\mathfrak d}\int
   F(p_{\rm phys}\omega)\rho_{\rm PT}(\omega)
       \,d\Pi_{s,\mathfrak d}(\omega).
\end{equation}
Because $p_{\rm phys}$ may be many-to-one, the tree weight is not inserted
into this formula as though it were a label function.  Disintegrate the
finite measure $\rho_{\rm PT}\Pi_{s,\mathfrak d}$ over $p_{\rm phys}$:
\begin{equation}
\label{eq:physical-weight-fibre-disintegration}
 \rho_{\rm PT}(\omega)d\Pi_{s,\mathfrak d}(\omega)
 =Z_{s,\mathfrak d}^{-1}
   d\mathfrak M_s^{\rm phys}(\lambda)
   K_{s,\mathfrak d,\lambda}^{\rm PT}(d\omega),
 \qquad p_{\rm phys}\omega=\lambda,
\end{equation}
where $K_{s,\mathfrak d,\lambda}^{\rm PT}$ is a probability kernel.  Define
the physical label weight, rather than assuming descent, by
\begin{equation}
\label{eq:physical-weight-conditional-definition}
 \mathscr W_{\lambda,*}^{\rm phys}
 :=\int\mathscr W_*(\omega)
      K_{s,\mathfrak d,\lambda}^{\rm PT}(d\omega).
\end{equation}
Then the exact commuting identity is
\begin{equation}
\label{eq:physical-weight-RN-identity}
 \sum_\lambda M_\lambda\mathscr W_{\lambda,*}^{\rm phys}
 =Z_{s,\mathfrak d}\int\rho_{\rm PT}(\omega)\mathscr W_*(\omega)
       d\Pi_{s,\mathfrak d}(\omega).
\end{equation}
Below, $\mathscr W_{\lambda,*}$ in a physical-label sum means precisely
$\mathscr W_{\lambda,*}^{\rm phys}$ from this display.  If the complete
weight is already $p_{\rm phys}$-measurable, this reduces to its ordinary
descent.
The factor $\rho_{\rm PT}$ is fixed before $(m,n,\xi,P)$ and the final test.
It does not alter the exact definition of $\mathscr W_*$.  Instead the
single domination \eqref{eq:full-mark-cemetery-domination} controls both
$\mathscr W_*$ and $\rho_{\rm PT}\mathscr W_*$ in the declared
$\mathfrak F$ moment.  No undeclared separate reciprocal H\"older slot is
allowed.  Thus no proof may pass from a
$\Pi$-moment to a physical $M_\lambda$-sum without using
\eqref{eq:physical-tree-RN-identity}.  In the common special case
$\rho_{\rm PT}=1$, the physical positive measure is exactly the bounded
normalization of the physical marginal of the complete tree.
\end{definition}

\begin{definition}[Exact propagated physical/source bridge
$\mathrm{PROP\_Q\_MATCH}$]
\label{def:propagated-q-match}
Let $\Omega_{\rm tree}$ be the complete refinement tree and let
$\Omega_{\rm occ}^{\rm phys}$ be a standard-Borel exact
occurrence/component refinement with maps
$p_{\rm tree}:\Omega_{\rm occ}^{\rm phys}\to\Omega_{\rm tree}$ and
$p_{\rm lab}=p_{\rm phys}\circ p_{\rm tree}$.  The interface supplies the
actual finite positive MPD coefficient kernel
$D_\omega^m(d\upsilon)$, supported on
$p_{\rm tree}\upsilon=\omega$, and defines, rather than freely chooses,
\begin{equation}
\label{eq:physical-occurrence-m-measure}
 \int F(\upsilon)d\mathfrak M_{s,\mathfrak d}^{m,\rm phys}(\upsilon)
 :=Z_{s,\mathfrak d}\int \rho_{\rm PT}(\omega)
       \int F(\upsilon)D_\omega^m(d\upsilon)
       d\Pi_{s,\mathfrak d}(\omega).
\end{equation}
Writing $c_{\rm MPD}(\omega):=D_\omega^m(\Omega_{\rm occ}^{\rm phys})$,
the forgetful marginal is the exact identity
\begin{equation}
\label{eq:physical-occurrence-forgetful-marginal}
 (p_{\rm lab})_*\mathfrak M^{m,\rm phys}(A)
 =Z_{s,\mathfrak d}\int\one_{\{p_{\rm phys}\omega\in A\}}
   \rho_{\rm PT}(\omega)c_{\rm MPD}(\omega)d\Pi_{s,\mathfrak d}(\omega).
\end{equation}
The coefficient kernel is normalized by the physical signed-recombination
identity
\begin{equation}
\label{eq:physical-occurrence-signed-recombination}
 J_\omega^{\rm phys}
 =\int s_{\rm pol}(\upsilon)
       U_\upsilon\nu_\upsilon\,D_\omega^m(d\upsilon)
\end{equation}
in the fixed physical test dual.  Hence multiplying
$D_\omega^m$ by an arbitrary scalar is not allowed.

There is also one predictable Borel component envelope density
$\mathscr C_{\rm occ}^{\rm phys}(\upsilon)$, fixed before the forbidden
choices and containing the exact component-specific multiplicity,
source/operator, recovery, and response costs.  It is not an extra factor
appended after the tree ledger: the defining disintegration bound is
\begin{equation}
\label{eq:physical-occurrence-envelope-disintegration}
 \int\mathscr C_{\rm occ}^{\rm phys}(\upsilon)
       D_\omega^m(d\upsilon)\le\mathscr W_*(\omega)
 \quad\text{for }\Pi_{s,\mathfrak d}\text{-a.e. }\omega.
\end{equation}
The same inequality is required after inserting every record-measurable
owner/handler restriction used to transfer a physical rate; otherwise that
rate is assumed directly on the propagated $q$-restriction.  Define
\begin{equation}
\label{eq:physical-occurrence-q-measure}
 d\mathfrak M_{s,\mathfrak d}^{q,\rm phys}
 :=\mathscr C_{\rm occ}^{\rm phys}\,
   d\mathfrak M_{s,\mathfrak d}^{m,\rm phys}.
\end{equation}
Consequently
\begin{equation}
\label{eq:physical-occurrence-total-ledger}
 \mathfrak M_{s,\mathfrak d}^{q,\rm phys}
       (\Omega_{\rm occ}^{\rm phys})
 \le Z_{s,\mathfrak d}\int\rho_{\rm PT}\mathscr W_*\,d\Pi_{s,\mathfrak d}
 =\sum_\lambda M_\lambda\mathscr W_{\lambda,*}^{\rm phys}.
\end{equation}
Thus no component cost can be introduced after the proved label ledger.  The two measures have
Borel record map
$\operatorname{Rec}_{\rm phys}$ whose fields are
\begin{equation}
\label{eq:physical-occurrence-record}
 (j,o,\operatorname{type},\varepsilon_{\rm or},s_{\rm pol},
   \ell_{\rm word},\kappa,\pi,\nu,U,K,R,\gamma,C_U,C_{\rm car},
   \mathscr A^{\rm opp}).
\end{equation}
Let $\operatorname{Rec}_{\rm src}$ be the identical record extracted from
\eqref{eq:source-incidence-record}.  There is one jointly Borel Markov
occurrence kernel $\mathcal K^{\rm occ}$, projectively fixed before
$(m,n,\xi,P)$ and the final test, whose support lies in the exact
compatibility graph
\begin{equation}
\label{eq:prop-q-compatibility-graph}
 \operatorname{supp}\mathcal K^{\rm occ}
 \subset\{(\upsilon;\lambda,a):
  \operatorname{Rec}_{\rm phys}(\upsilon)
   =\operatorname{Rec}_{\rm src}(\lambda,a)\}.
\end{equation}
and it is mass preserving:
\begin{equation}
\label{eq:prop-q-occurrence-kernel-markov}
 \mathcal K^{\rm occ}(\upsilon;I)=1
 \quad\text{for }\mathfrak M^{m,\rm phys}\text{-a.e. }\upsilon.
\end{equation}
In particular the kernel cannot swap two labels having the same mass but
different owner, type, orientation, polarity, word depth, clock, passage,
or occurrence index.  It satisfies the exact positive coefficient identity
and propagated envelope domination
\begin{align}
 \int \!\mathcal K^{\rm occ}(\upsilon;B)
       \,d\mathfrak M_{s,\mathfrak d}^{m,\rm phys}(\upsilon)
 &=\int_Bd\Pi_{s,\mathfrak d}(\lambda)m_\lambda(da),
 \label{eq:prop-q-exact-m-match}\\
 \overline M_{s,\mathfrak d}^{q}(B)
 &\le C\int\!\mathcal K^{\rm occ}(\upsilon;B)
       \,d\mathfrak M_{s,\mathfrak d}^{q,\rm phys}(\upsilon)
 \label{eq:propagated-q-physical-match}
\end{align}
for every Borel $B\subset I$.  Finally, for every
$\Phi\in\mathfrak T_{\rm src}$ the physical occurrence current equals the
grouped source current obtained from
\eqref{eq:marked-source-assembled-pairing}; equivalently, after inserting
$\one_B$, both sides of \eqref{eq:prop-q-exact-m-match} integrate the same
recorded scalar
$s_{\rm pol}\langle U\nu,\Phi\rangle$.  This exact current identity is part
of the interface and is stronger than arbitrary measure domination.
The physical $q$-measure in \eqref{eq:propagated-q-physical-match} is exactly
\eqref{eq:physical-occurrence-q-measure}.  Equations
\eqref{eq:physical-occurrence-m-measure}--
\eqref{eq:physical-occurrence-signed-recombination}, the Markov identity,
and the two pushforward relations form one commuting physical--tree--source
diagram.  No parallel physical mass or independently rescaled occurrence
decomposition may be substituted.
\end{definition}

\begin{proposition}[Global physical ledger]
\label{prop:global-labelled-ledger}
Assume all hypotheses of
Lemma~\ref{lem:measurable-refinement-moments}, including its normalized
refinement kernels, pointwise domination, full product envelope,
return/context moments, and complete-tree cemetery moment, and assume
\[
 \begin{gathered}
 \mathrm{REC}(q_{\rm rec}),\ \mathrm{NST}_{\rm phys},\ \mathrm{SS}(d_Z),\\
 \mathrm{ACTUAL\_SLOPE\_BRIDGE},\ \mathrm{PAIR\_SAFE\_REFERENCE},\
 \mathrm{PAIR\_TENERGY},\\
 \mathrm{PPE}_{N}(d_Z),\ \mathrm{CEM\_DOM},\ \mathrm{TENERGY}_{W}
 \end{gathered}
\]
\[
 \begin{gathered}
 \mathrm{DUAL\_MASS\_KERNEL},\ \mathrm{MPD},\
 \mathrm{POST\_CARRIER\_DOM},\ \mathrm{PHYS\_TEST\_LIFT},\\
 \mathrm{PHYS\_TREE\_MATCH},\ \mathrm{PROP\_Q\_MATCH},\\
 \mathrm{MAIN\_CLOCK\_COVERAGE},\
 \mathrm{SEL\_JSL}_{p_{\rm joint}},\\
 \mathrm{CLOCK\_PROB}_{W},\ \mathrm{CLOCK\_SRC}_{W},\\
 \mathrm{KRH}_{p_W},\ \mathrm{USC}_{\rm end}^{W}(q_g^W),\
 \mathrm{BR}_{\rm sec}
 \end{gathered}
\]
\[
 \mathrm{FZ}_{\rm geom}(\theta_Z),\ \mathrm{SF}_{\rm split},\
 \mathrm{CP}_{\rm env},\ \mathrm{UCR}_{\rm sh}.
\]
Assume in addition the exact physical/tree and record-preserving
physical/source bridges of
Definitions~\ref{def:physical-tree-match} and
\ref{def:propagated-q-match}.  In particular,
\eqref{eq:physical-tree-RN-identity} is used before any tree moment is
converted to an $M_\lambda$-sum, while
\eqref{eq:prop-q-exact-m-match} identifies the exact occurrence current and
\eqref{eq:propagated-q-physical-match} controls its propagated envelope.
Neither statement follows merely from
$w^q\ge w^mK\gamma^{-R}\overline{\mathscr A}^{\rm opp}C_{\rm prop}$.
Assume $\mathrm{MAIN\_CLOCK\_COVERAGE}$ and the selection-safe likelihood
and base bounds \eqref{eq:selection-safe-likelihood-rates} and
\eqref{eq:main-source-base-mass} with the same constants.  For every
nonempty main time--Jensen cell also assume the unselected parentwise
reference energy \eqref{eq:unselected-time-reference-energy}; it is not
supplied by $\mathrm{CLOCK\_PROB}_{W}$.
Assume also the complete path condition $\mathrm{QT}_{\rm path}$ of
Definition~\ref{def:complete-QT-path}, including both the jet condition and
the independent loss/return-tail ledger.  Let
$q_H,q_F,q_D,q_S,q_R,q_C>1$ be the available context, flux, inverse-scale,
positive-scale, return/context, and complete-tree cemetery moment exponents,
and assume
\begin{equation}
\label{eq:holder-conjugate-margin}
 \chi_W=\frac{a_H}{q_H}+q_F^{-1}+q_D^{-1}+q_S^{-1}+q_R^{-1}
        +q_C^{-1}<1.
\end{equation}
After artificial-face cancellation, the complete first defect has a disjoint
signed refinement
\begin{equation}
\label{eq:global-label}
 \lambda=(\chi,a,k,c^-,g,c^+;Q;j,k',r,\sigma),
\end{equation}
where $\chi$ is a physical chart, $a$ a return word, $k$ the insertion time,
$c^-,g,c^+$ the single-source gate data, $Q$ a native gate ball,
$(j,k',r)$ the incoming/outgoing strips and connected-component index, and
$\sigma$ the chosen time orientation.  The refinement does not promote
amplitude cuts to new inverse branches.  For $a_H>0$ and $a_{\rm scl}>0$
below their respective moment thresholds,
\begin{equation}
\label{eq:global-ledger-sum}
\sum_\lambda M_\lambda\mathscr W_{\lambda,*}<\infty,
\end{equation}
where $\mathscr W_{\lambda,*}$ means the physical conditional weight
\eqref{eq:physical-weight-conditional-definition} of the exact tree weight
\eqref{eq:global-stopped-source-weight}; no descent through a many-to-one
physical projection is assumed, and no smaller unweighted ledger is
substituted when scale-bad labels are removed.  The occurrence-level bound
\eqref{eq:physical-occurrence-envelope-disintegration} and total identity
\eqref{eq:physical-occurrence-total-ledger} first convert this label ledger
to the physical occurrence $q$-ledger.  Then $\mathrm{PROP\_Q\_MATCH}$ makes
its record-preserving occurrence-kernel pushforward dominate the
corresponding restriction of
$d\overline M_{s,\mathfrak d}^{q}
=d\Pi_{s,\mathfrak d}q_\lambda$.  Conversely, every
source-rate assertion below is made directly for that propagated
$q$-restriction, not inferred from a parallel bare mass.  This identification
is the one used by the FACE, shallow, and deep clauses above.  The exact
$m$-pushforward and grouped current identity are simultaneously those of
\eqref{eq:prop-q-exact-m-match}; domination alone is not used to identify a
physical occurrence.
All constants in the functional-correlation, native-energy,
fractional-phase, return, grazing, and direct-$Z$ estimates are uniform on
the spacetime master cells of the quantitative table family.
\end{proposition}

\begin{proof}
Use the normalized refinement kernels of
Lemma~\ref{lem:measurable-refinement-moments}.  The context and $H$ moments
are transported through the normalized stopped kernel in
$\mathrm{NST}_{\rm phys}$; the underlying bounds come from the full-branch
quotient and exponential gate hitting;
Proposition~\ref{prop:radial-direct-Z} gives the component-indexed
mark--boundary conditional moment before the boundary-oscillation theorem is
applied; and the dominant chart hypothesis gives the conditional
$\mathfrak d^{A_{\rm dom}q_D}$ moment.  The two separate recovery moments
\eqref{eq:separate-recovery-moments} are retained on the same complete law
but are not multiplied with one another.  Return/grazing tails and Proposition
\ref{prop:parameter-state-master-refinement} provide the last normalized
kernel; the strict margin \eqref{eq:strict-jet-tail-margin} absorbs its
finite-depth loss.  Apply Lemma~\ref{lem:measurable-refinement-moments} and
the strict full reciprocal margin \eqref{eq:holder-conjugate-margin}.  The
resulting $p_W$ in \eqref{eq:available-full-product-pW} proves the required
complete-tree moment.  Insert the charged factor $\rho_{\rm PT}$ and use the
exact identity \eqref{eq:physical-weight-RN-identity} to obtain
\eqref{eq:global-ledger-sum}; without this step the tree moment does not
control the physical $M_\lambda$-measure.  This transfers only global
probability aggregates.  Parentwise energy and upper-scale tails use
$\mathrm{KRH}_{p_W}$, while additive remainders use the independent
incidence-space ledger.  No unrecorded $p>2$ moment and no intersection of
unrelated label laws is used.  The ledger is not assumed again in the main
theorem.
The positive initial endpoint scale is included through
$S_\lambda^{a_{\rm scl}}$.  Its excess is then split off by the
\emph{tilted} estimate \eqref{eq:upper-scale-excess-tail}, on that same
weighted complete law, before the phase estimate is applied.  Context,
stopping, buffer, gate, continuation, and PPE losses are separately charged
by $\mathrm{TENERGY}_{W}$ on this identical tilt, while each additional
bad-support class is charged through $\mathrm{CLOCK\_PROB}_{W}$ and every
additive occurrence through $\mathrm{CLOCK\_SRC}_{W}$.  For cemetery pieces the
passage from weighted probability to actual discarded current is exactly
\eqref{eq:cemetery-current-weighted-bound}; it is not inferred from the
flux sum.
The occurrence-level disintegration
\eqref{eq:physical-occurrence-envelope-disintegration} prevents a component
cost from being appended after the tree moment.  The exact
physical-to-source identification is \eqref{eq:prop-q-exact-m-match}
together with its compatibility graph, and its envelope domination is precisely
\eqref{eq:propagated-q-physical-match}; none of the moment arguments in this
proof is claimed to derive either statement.
\end{proof}

\begin{theorem}[Conditional pathwise quantitative CM2 criterion]
\label{thm:global-cm2-assembly}
Let an analytic-boundary or circular deformation $a(s)$ satisfy
$\mathrm{QT}_{\rm path}$ and the gauge-compatibility restriction
\eqref{eq:constant-component-mass-vector}.
Assume a compact finite-horizon chart and, along every segment
$a(\theta s)$ and every $L\le L(s)$, one spacetime refinement satisfying
$\mathrm{UCR}_{\rm sh}$, with:
\begin{enumerate}[label=\textup{(H\arabic*)}]
\item all hypotheses of
      Lemma~\ref{lem:measurable-refinement-moments} hold on the same complete
      actual/split tree.  This explicitly includes the normalized refinement
      kernels \eqref{eq:refinement-kernel-mass}, the weight normalization
      \eqref{eq:full-product-weight-normalization}, the pointwise Mark--$Z$
      domination \eqref{eq:full-mark-cemetery-domination}, the full cemetery
      product envelope \eqref{eq:full-cemetery-product-envelope}, all
      conditional moments, and the genuine complete-tree cemetery moment;
      no collection of separate marginal moments is substituted for this
      package.  In addition assume the finite-horizon functional-correlation bound of \cite{LS}, the
      standard one-state marked quotient, and exponential return/context
      moments; a dominant inverse-scale moment
      $\sum_\lambda M_\lambda
      \mathfrak d_\lambda^{A_{\rm dom}q_D}<\infty$ and a positive initial-scale
      moment $\sum_\lambda M_\lambda
      S_\lambda^{a_{\rm scl}q_S}<\infty$; the separate terminal-law
      recovery input $\mathrm{REC}(q_{\rm rec})$, with
      $q_{\rm rec},\epsilon,\eta_{\rm lo}$ chosen to satisfy
      \eqref{eq:delayed-recovery-eta-window}; and
      exponents $q_H,q_F,q_D,q_S,q_R,q_C$ satisfying
      \eqref{eq:holder-conjugate-margin}, with $p_W$ chosen in
      \eqref{eq:available-full-product-pW}; the refinement moments hold
      conditionally along the common refinement as in
      \eqref{eq:conditional-refinement-moments}, whereas the possibly
      unbounded cemetery factor satisfies the genuine joint law
      \eqref{eq:complete-tree-cemetery-moment} (or the earlier-parent kernel
      version \eqref{eq:parent-kernel-cemetery-moment});
\item forward/reverse proper-family growth and coupling uniformly for the
      compact slowly moving table class, in the form of Theorems~1'--2' of
      \cite{SYZ}; Lemma~\ref{lem:signed-observable-multiplier} then supplies
      the signed multiplier bridge;
\item On the persistent clean gate assume
      \[
      \begin{gathered}
       \mathrm{NST}_{\rm phys},\ \mathrm{SS}(d_Z),\
       \mathrm{ACTUAL\_SLOPE\_BRIDGE},\\
       \mathrm{PAIR\_SAFE\_REFERENCE},\ \mathrm{PAIR\_TENERGY},\
       \mathrm{PPE}_{N}(d_Z),
      \end{gathered}
      \]
      together with
      \[
       \begin{gathered}
       \mathrm{CEM\_DOM},\ \mathrm{TENERGY}_{W},\\
       \mathrm{DUAL\_MASS\_KERNEL},\ \mathrm{MPD},\
       \mathrm{POST\_CARRIER\_DOM},\ \mathrm{PHYS\_TEST\_LIFT},\\
       \mathrm{PHYS\_TREE\_MATCH},\ \mathrm{PROP\_Q\_MATCH},\\
       \mathrm{MAIN\_CLOCK\_COVERAGE},\
       \mathrm{SEL\_JSL}_{p_{\rm joint}},\\
       \mathrm{CLOCK\_PROB}_{W},\ \mathrm{CLOCK\_SRC}_{W},\
       \mathrm{KRH}_{p_W}
       \end{gathered}
      \]
      \[
       \mathrm{USC}_{\rm end}^{W}(q_g^W),\ \mathrm{BR}_{\rm sec}.
      \]
      The marked kernel satisfies the global finite-envelope closure
      \eqref{eq:global-finite-marked-envelope}, the exhaustive owner identity
      \eqref{eq:exhaustive-source-owner-partition}, the resolution coverage
      \eqref{eq:nonmain-clock-coverage}--
      \eqref{eq:main-clock-coverage}, the exact product reference
      \eqref{eq:unselected-reference-law}--\eqref{eq:pair-tenergy-bound},
      the bounded bilinear physical-test lift and its exact commutation
      \eqref{eq:bounded-physical-test-lift}--
      \eqref{eq:physical-test-lift-commutation}, the physical/tree identity
      \eqref{eq:physical-tree-RN-identity}, and the record-preserving
      physical/source bridge
      \eqref{eq:physical-occurrence-m-measure}--
      \eqref{eq:propagated-q-physical-match}.  It also satisfies the
      projective consistency
      \eqref{eq:source-kernel-projective-consistency} and the main-source
      likelihood/base
      bounds \eqref{eq:selection-safe-likelihood-rates} and
      \eqref{eq:main-source-base-mass} uniformly on the same parent cells.
      For every nonempty main time--Jensen cell the unselected reference
      energy \eqref{eq:unselected-time-reference-energy} holds uniformly on
      those same frozen parents.
      The selection-safe joint source likelihood is verified either by the
      same-carrier clause of Lemma~\ref{lem:joint-source-factor-criterion} or
      by its explicit $\mathrm{CARRIER\_MATCH}$ alternative.
      Here \(\mathrm{KRH}_{p_W}\) is uniform on every frozen parent kernel.
      These interfaces and the dominant chart bounds
      \eqref{eq:primitive-dominant-chart-bounds} have the
      parameter-uniform constants required by $\mathrm{UCR}_{\rm sh}$ for
      $C_{\rm mid}(\eta_{\rm hi})N\le L(s)$.  The cutoffs
      \eqref{eq:three-frequency-cutoffs} satisfy
      $0<\eta_{\rm lo}<\eta_{\rm hi}$ and
      $r_{\rm obs}>2+\beta_{\rm phase}^{W}$,
      \eqref{eq:ultra-high-exponent-window}, and the complete phase exponent
      \eqref{eq:complete-phase-exponent}; longer shallow times are not estimated
      by this interface;
\item the two-chart nondegeneracy $|\nabla G|\ge c_G$ and
      $\mathrm{FZ}_{\rm geom}(\theta_Z)$, including the full tangent coarea identity,
      weighted component complexity, angle sublevel, and good-angle exponent
      $A_{\rm ang}$, together with $\mathrm{SF}_{\rm split}$ and
      $\mathrm{CP}_{\rm env}$ on the same spacetime cells;
      Proposition~\ref{prop:envelope-dynamical-tube}
      then supplies \eqref{eq:FZ-dynamical-tube}; the shallow finite-$s$
      source satisfies $\mathrm{SHCOL}_{\rm env}$ and the strict nonempty
      window \eqref{eq:shallow-time-window}; the terminal tail satisfies
      $\mathrm{TAIL}_{\rm env}^{+}$ and
      $\mathrm{NREC}_{\rm deep}(q_{\rm deep})$ with
      $q_{\rm deep}>\max\{\alpha_+,\alpha_-\}$;
\item $\mathrm{MT}_{\rm DQ}$ holds for the fixed-gauge operator quotient,
      including its exact typed decomposition and
      \eqref{eq:uniform-gauge-DQ-regular}--
      \eqref{eq:fixed-time-evolution-convergence}.  The regular defect
      satisfies, for some
      $1\le R\le3$,
      \[
       \sum_\chi\|\zeta_\chi k_{\rm reg}\|_{W^{R,1}}<\infty.
      \]
      The local face amplitudes have the scaled
      $\mathcal A^\alpha$ regularity used in
      Proposition~\ref{prop:scaled-face-density}; observable regularity is
      above the dominant two-dimensional Fourier threshold.  The joint face
      part of \eqref{eq:event-physical-norm} is concluded from
      Proposition~\ref{prop:global-labelled-ledger}, not assumed here.
      The growing-depth face assembly satisfies
      $\mathrm{FACE}_{\rm 2CUT}$, in particular the no-$|s|^{-1}$ middle
      estimate \eqref{eq:face-middle-tail}.  Independently, all main and
      nonmain face/word owners satisfy
      $\mathrm{FACE\_TIME}_{\rm CM2}$ on the same propagated $q$-restriction.
\end{enumerate}
Then the complete double-centered fixed-gauge first operator defect satisfies
\[
 |\langle Q^nK_VQ^mg,f\rangle|
 \le C\tau^m\vartheta^n\|g\|_{C^r}\|f\|_{C^r}.
\]
The assembled difference quotients converge after the independent word-depth
and time-depth splits.  At every fixed $(m,n)$ the response,
product-current, regular, and face blocks converge respectively by
\eqref{eq:strong-DQ-convergence},
\eqref{eq:product-current-DQ-convergence},
\eqref{eq:regular-DQ-weak-convergence}, and
Theorem~\ref{thm:face-two-cutoff}.  The shallow long-time part converges to zero by
Theorem~\ref{thm:shallow-two-threshold}, and the deep remainder does so by
Theorem~\ref{thm:time-depth-diagonal}; hence the CM2 series converges term by
term as $s\to0$.  More
precisely, if $c_{\rm off}$ is the off-diagonal coupling exponent and $c_Z$
is the exponent obtained by balancing the $\mathrm{FZ}_{\rm geom}$ angle sublevel, one
may choose
\begin{equation}
\label{eq:final-CM2-exponent}
 c_*<\min\left\{
 \begin{gathered}
 c_{\rm off},c_{\rm low},-2\log\theta_-,-2\log\theta_+,\\
 \eta_{\rm lo}\beta_{\rm phase}^{W},
 \eta_{\rm hi}\rho_\infty-a_{\rm ultra},c_{\rm scl},\\
 c_{\rm time,J}^{\rm val},c_{\rm time,TV}^{\rm src},
 \eta_{\rm lo}R,c_{\rm face,time},c_{\rm tail}-A_{\rm loss},c_Z
 \end{gathered}
 \right\}
\end{equation}
and then take $\tau=\vartheta=e^{-c_*/2}$ after changing the constant.
\end{theorem}

\begin{proof}
Lemma~\ref{lem:fixed-gauge-centering} puts the limiting operator and every
finite-$s$ operator quotient under the same $\Pi_0$.  Coordinate
commutators remain in the response block, and their uniform product decay is
derived from \eqref{eq:strong-side-memory}--\eqref{eq:response-side-memory},
while \eqref{eq:strong-DQ-convergence}--
\eqref{eq:fixed-time-evolution-convergence} gives each fixed two-time limit; no
transported observable changes the target.  The mixed marginal/rank-one
terms remain in the fourth product-current type and form the exact centered
cluster \eqref{eq:typed-face-product-cluster} with their graph face.
They obey \eqref{eq:product-current-type} and vanish only against the two individually
centered endpoint tests; they are not absorbed into the regular or face
norm.  Every graph-supported
commutator boundary remains in the physical face atlas under \textup{(H5)}.
The regular fixed-term limit is
\eqref{eq:regular-DQ-weak-convergence}.  The face fixed-core limit follows
from \eqref{eq:face-DQ-weak-convergence} and the boundary-neighbourhood
estimate \eqref{eq:face-DQ-boundary-tightness};
Theorem~\ref{thm:face-two-cutoff} then removes the intermediate triangular
band before the terminal deep remainder is invoked.  Its uniform CM2
majorant is instead supplied, for every fixed core and every other
face/word owner, by Lemma~\ref{lem:face-time-majorant}; the event/section
ledgers alone are not used to infer time decay.
Thus no
uniform norm bound is used as a substitute for DQ convergence in either
typed block.
Proposition~\ref{prop:global-labelled-ledger} first supplies the complete face
part of \eqref{eq:event-physical-norm} from the primitive moment hypotheses.
Conditions $\mathrm{NST}_{\rm phys}$, $\mathrm{SS}$,
$\mathrm{ACTUAL\_SLOPE\_BRIDGE}$, $\mathrm{PAIR\_SAFE\_REFERENCE}$,
$\mathrm{PAIR\_TENERGY}$,
$\mathrm{PPE}_{N}$, $\mathrm{CEM\_DOM}$, and
$\mathrm{TENERGY}_{W}+\mathrm{CLOCK\_PROB}_{W}
+\mathrm{CLOCK\_SRC}_{W}$, together with the complete
positive stopped-source tree and its cemetery labels, give
Theorem~\ref{thm:direct-slope-small-ball}.  The high-frequency balanced
estimate is Proposition~\ref{prop:central-phase-verification}.  The regular
four-dimensional $M\times M$ defect is included in the low-frequency Proposition
\ref{prop:functional-correlation-low-frequency}.  Choosing
$K=e^{\eta_{\rm lo}\max\{m,n\}}$ and
$H=e^{\eta_{\rm hi}\max\{m,n\}}$ produces the low, intermediate-phase,
ultra-high observable, scale-bad, and regular-tail exponents in
\eqref{eq:final-CM2-exponent}.

Outside the diagonal band use Proposition
\ref{prop:signed-off-diagonal-mapping};
Lemma~\ref{lem:three-region} converts the three estimates to product decay.
For radial pieces Proposition~\ref{prop:radial-direct-Z}, Lemma
\ref{lem:FZ-angle-balance}, and Corollary
\ref{cor:radial-bi-recovery} give the common physical majorant without a
component-length estimate and with every component explicitly summed.
Proposition~\ref{prop:global-labelled-ledger} gives the global physical
ledger and its $L^{p_W}$ margin from the primitive hypotheses.  Lemma
\ref{lem:finite-depth-face-separation} and Proposition
\ref{prop:parameter-state-master-refinement} give the common coarea atlas at
each fixed depth.  Theorem~\ref{thm:face-two-cutoff} alone upgrades those
fixed cores through the middle band to $L(s)$.  Theorem
\ref{thm:shallow-two-threshold} removes shallow times beyond
$T_{\rm sh}(s)$.  Lemma~\ref{lem:finite-s-collar-admissibility} and Theorem
\ref{thm:time-depth-diagonal} give the required $o(1)$ CM2-norm tail beyond
$L(s)$.  No fixed neighborhood is
claimed to be transverse at every depth, and no qualitative residual
assertion is used in this step.  The four fixed-term topologies,
$\mathrm{FACE}_{\rm 2CUT}$, and the fixed-term uniform ledgers meet in
\eqref{eq:typed-common-dominated-convergence}.  The shallow exponent window
converts its triangular bound into a common summable exponential array.
The deep block is different: Theorem~\ref{thm:time-depth-diagonal} gives
direct $\ell^1(\mathbb N^2)$-norm convergence to zero.  Thus the final
limit--double-sum exchange is dominated convergence for the fixed/shallow
blocks plus this separate uniform-integrability estimate for the deep
remainder; it is not claimed to be one literal application of
\eqref{eq:typed-common-dominated-convergence} to every triangular block.
\end{proof}

\begin{corollary}[Product-a.e. analytic paths and the circular boundary]
\label{cor:analytic-pathwise-CM2}
Assume all hypotheses of Theorem~\ref{thm:global-cm2-assembly} other than
$\mathrm{QT}_{\rm JET}$, uniformly on the product-a.e. bases and directions
to which the analytic realization below applies.  In
particular this includes $\mathrm{QT}_{\rm LEDGER}$ and
\eqref{eq:constant-component-mass-vector}, as well as
\[
\begin{gathered}
 \mathrm{MT}_{\rm DQ},\ \mathrm{NST}_{\rm phys},\\
 \mathrm{ACTUAL\_SLOPE\_BRIDGE},\ \mathrm{PAIR\_SAFE\_REFERENCE},\\
 \mathrm{PAIR\_TENERGY},\\
 \mathrm{PPE}_{N},\ \mathrm{CEM\_DOM},\ \mathrm{TENERGY}_{W},\\
 \mathrm{DUAL\_MASS\_KERNEL},\ \mathrm{MPD},\
 \mathrm{POST\_CARRIER\_DOM},\ \mathrm{PHYS\_TEST\_LIFT},\\
 \mathrm{PHYS\_TREE\_MATCH},\ \mathrm{PROP\_Q\_MATCH},\\
 \mathrm{MAIN\_CLOCK\_COVERAGE},\
 \mathrm{SEL\_JSL}_{p_{\rm joint}},\\
 \mathrm{CLOCK\_PROB}_{W},\ \mathrm{CLOCK\_SRC}_{W},\
 \mathrm{KRH}_{p_W},\ \mathrm{USC}_{\rm end}^{W},\ \mathrm{BR}_{\rm sec},\\
 \mathrm{REC},\ \mathrm{FZ}_{\rm geom},\ \mathrm{SF}_{\rm split},\
 \mathrm{CP}_{\rm env},\ \mathrm{UCR}_{\rm sh},\\
 \mathrm{FACE}_{\rm 2CUT},\ \mathrm{FACE\_TIME}_{\rm CM2},\
 \mathrm{SHCOL}_{\rm env},\
 \mathrm{TAIL}_{\rm env}^{+},\ \mathrm{NREC}_{\rm deep}.
\end{gathered}
\]
If the analytic boundary family
satisfies the localized jet-direction condition
\eqref{eq:localized-analytic-jet-direction} and the exponent window
\eqref{eq:analytic-QT-window}, then for almost every base table in the
analytic product model, restricted to the gauge-compatible analytic chart,
and every admissible tangent direction with the
recorded derivative exponent, CM2 holds along the path.
The product randomness in this statement realizes only
$\mathrm{QT}_{\rm JET}$; it supplies neither
$\mathrm{QT}_{\rm LEDGER}$ nor any dynamical/source interface in the
preceding display.
For the finite-dimensional circular family, assume instead all hypotheses
of the theorem including the explicitly stated $\mathrm{QT}_{\rm path}$;
then the same conclusion holds only for center motions with fixed radii or
for constrained radius/shape paths satisfying
\eqref{eq:constant-component-mass-vector}.  No automatic circular
transversality or independent-radius gauge is asserted.
\end{corollary}

\begin{proof}
Apply Theorem~\ref{thm:analytic-QT-realization} to obtain
$\mathrm{QT}_{\rm JET}$.  Combine it with the separately assumed
$\mathrm{QT}_{\rm LEDGER}$, and then apply Theorem
\ref{thm:global-cm2-assembly}.  The last sentence is the constrained
circular branch of the latter theorem.
\end{proof}

\begin{proposition}[Pathwise coarea-current convergence]
\label{prop:common-refinement-majorant}
In a compact table chart satisfying \textup{(H1)--(H5)}, assume
$\mathrm{QT}_{\rm path}$ and
\[
\begin{gathered}
 \mathrm{UCR}_{\rm sh},\ \mathrm{CEM\_DOM},\ \mathrm{TENERGY}_{W},\\
 \mathrm{ACTUAL\_SLOPE\_BRIDGE},\ \mathrm{PAIR\_SAFE\_REFERENCE},\
 \mathrm{PAIR\_TENERGY},\\
 \mathrm{DUAL\_MASS\_KERNEL},\ \mathrm{MPD},\
 \mathrm{POST\_CARRIER\_DOM},\ \mathrm{PHYS\_TEST\_LIFT},\\
 \mathrm{PHYS\_TREE\_MATCH},\ \mathrm{PROP\_Q\_MATCH},\\
 \mathrm{MAIN\_CLOCK\_COVERAGE},\
 \mathrm{SEL\_JSL}_{p_{\rm joint}},\\
 \mathrm{CLOCK\_PROB}_{W},\ \mathrm{CLOCK\_SRC}_{W},\
 \mathrm{KRH}_{p_W},\ \mathrm{FACE}_{\rm 2CUT},\
 \mathrm{FACE\_TIME}_{\rm CM2},\\
 \mathrm{SHCOL}_{\rm env},\ \mathrm{TAIL}_{\rm env}^{+},\
 \mathrm{NREC}_{\rm deep}.
\end{gathered}
\]
Then all constants in the
functional-correlation, native-energy, fractional-phase, direct-$Z$, and
parameter-tail estimates are uniform up to $L(s)$.  The shallow coarea current
converges at every fixed depth in
\eqref{eq:face-DQ-weak-convergence} and satisfies
\eqref{eq:face-DQ-boundary-tightness}; the middle band is uniformly
integrable by \eqref{eq:face-middle-tail}, hence the growing-depth current
converges by Theorem~\ref{thm:face-two-cutoff}.  Independently, it has the
CM2 majorant \eqref{eq:face-time-CM2-majorant}, while the deep
remainder is $o(1)$ in the CM2 sum.
\end{proposition}

\begin{proof}
The spacetime cells in $\mathrm{UCR}_{\rm sh}$, rather than compactness alone, make
hyperbolicity, distortion, the clean gate, moving branch sections, split
laws, connected cylinders, and single-face nondegeneracy uniform along the
whole coarea segment.  The
upper-scale excess tail and every energy-tree failure are uniform on the
same complete tilted labels by $\mathrm{USC}_{\rm end}^{W}$ and
$\mathrm{TENERGY}_{W}+\mathrm{CLOCK\_PROB}_{W}
+\mathrm{CLOCK\_SRC}_{W}$; no raw lattice cutoff or unweighted tail is used to
bound the complete marked phase mass.  The
fixed-depth coarea formula gives chartwise convergence of orientations and
amplitudes, hence bounded--Lipschitz convergence of the assembled current.
The uniform component-indexed boundary-$Z$ ledger bounds the mass of a
$\rho$-neighbourhood of the finite singularity set by
$C_{L,m,n}\rho^{\theta_{L,m,n}}$, which is exactly
\eqref{eq:face-DQ-boundary-tightness}.  Thus the current convergence may be
tested against the fixed-time branchwise observables; this is the same
weak-* plus boundary-neighbourhood mechanism as
\cite[Lemma~5.4]{CanestrariHoles}, now required on the common moving-face
atlas.  The per-depth positive coarea/return estimate
\eqref{eq:face-middle-per-depth} is then summed before any signed
recombination to obtain the middle tail with no $|s|^{-1}$ loss.  The
functional-correlation constants are uniform for a compact finite-horizon
family, and the return/context moments and both one-sided coupling estimates
are uniform by \textup{(H1)--(H2)}.  Proposition
\ref{prop:radial-direct-Z} fixes the component-indexed density/boundary
majorant.  Use the strict path margin
\eqref{eq:strict-jet-tail-margin}.  Proposition
\ref{prop:parameter-state-master-refinement} supplies the
depth-dependent scalar quotient tail on the same labels.  Apply the low/high
central estimates and off-diagonal couplings only through $T_{\rm sh}(s)$,
sum \eqref{eq:global-ledger-sum}, use Theorem
\ref{thm:shallow-two-threshold} for the remaining shallow times, and use Theorem
\ref{thm:time-depth-diagonal} for the deep part.
\end{proof}

The construction does not require a differentiable parameter family of
countable towers.  The coarea depth grows like $L(s)$, while the remaining
collision/grazing words are removed in the CM2 norm by the time--depth split
of Theorem~\ref{thm:time-depth-diagonal}, not by scalar total variation alone.

\section{Global ledger and quantitative falsification checks}

For reference, the logical dependencies of the active proof are
\begin{enumerate}[label=\textup{(D\arabic*)}]
\item \textup{(H1)--(H5)}, $\mathrm{UCR}_{\rm sh}$, the global physical
      ledger (including the joint Mark--$Z$ consequence of
      $\mathrm{FZ}_{\rm geom}$), and the clock/source interfaces give the
      prescribed-depth Jensen/TV low--high rows.
\item $\mathrm{FACE\_TIME}_{\rm CM2}$ and the bounded test lift give every
      main/nonmain face/word product-time bound.
\item \textup{(H1),(H2),(H4),(H5)}, the scaled face-density ledger,
      $\mathrm{SHCOL}_{\rm env}$, and $\mathrm{NREC}_{\rm deep}$ give
      shallow long-time and the two off-diagonal removals.
\item The functional-correlation and forward/reverse coupling inputs in
      \textup{(H1),(H2),(H4)}, together with
      $\mathrm{FZ}_{\rm geom}$, $\mathrm{SF}_{\rm split}$,
      $\mathrm{CP}_{\rm env}$, $\mathrm{UCR}_{\rm sh}$,
      $\mathrm{QT}_{\rm JET}$, $\mathrm{QT}_{\rm LEDGER}$,
      $\mathrm{TAIL}_{\rm env}^{+}$, $\mathrm{NREC}_{\rm deep}$, and the
      Mark--$Z$ ledger give envelope crossing and the $(L,T,r)$ deep-tail
      removal.
\item $\mathrm{MT}_{\rm DQ}$ and $\mathrm{FACE}_{\rm 2CUT}$ give the four
      typed fixed limits and the independent word-depth face limit.
\end{enumerate}
These outputs are combined only after their types and time ranges have been
fixed.

The proof has seven independent exponent ledgers.
\begin{enumerate}[label=\textup{(A\arabic*)}]
\item For $N\le T_{\rm sh}(s)$, the low-frequency ledger is
\begin{align*}
 &K^2e^{-c_-(1-\epsilon)m-c_+(1-\epsilon)n}\\
 &\quad+K^2e^{-q_{\rm rec}\epsilon m+q_{\rm rec}C_1\log K}
 +K^2e^{-q_{\rm rec}\epsilon n+q_{\rm rec}C_1\log K},
 \qquad K=e^{\eta_{\rm lo}N},
\end{align*}
plus the regular-density term $K^6e^{-c_-m-c_+n}$.
\item The prescribed-depth intermediate face ledger is
$K^{-\beta_{\rm phase}^{W}}$; the ultra-high original-observable tail is
$e^{a_{\rm ultra}N}H^{-\rho_\infty}$ with
$H=e^{\eta_{\rm hi}N}$, the scale-bad tail is $e^{-c_{\rm scl}N}$, and the
regular high tail is $K^{-R}$.  The definition
\eqref{eq:complete-phase-exponent} charges every $r$- or PPE-depth energy
tail with its registered Jensen status; genuine-time events enter
\eqref{eq:complete-central-exponent} separately and no event is counted
twice.
\item The first inequality in \eqref{eq:shallow-time-window} guarantees
that every resolution used in items 1--2 lies inside
$\mathrm{UCR}_{\rm sh}$.
The face/word owner is not contained in those two ledgers; it has the
independent positive-envelope rate
$e^{-c_{\rm face,time}(m+n)}$ from
\eqref{eq:face-time-CM2-majorant}.
\item The shallow long-time collar has ledger
\[
 |s|^{-1-A_{\rm sh}\kappa_L}
 e^{-c_{\rm sh}T_{\rm sh}(s)},
\]
which vanishes by the second inequality in
\eqref{eq:shallow-time-window}.
\item Lemma~\ref{lem:three-region} supplies the two off-diagonal coupling
ledgers, while the pathwise return/grazing quotient tail is multiplied by
the joint mark--$Z$ sum \eqref{eq:joint-mark-Z-ledger}.
\item The deep positive source is removed by
\[
 |s|^{\kappa_L(c_{\rm tail}-A_{\rm loss})-1}\log^2(1/|s|)
 \quad\hbox{and}\quad
 |s|^{-1-A_{\rm col}\kappa_L}e^{-c_{\rm col}T(s)}.
\]
The off-diagonal recovery prefactor retains $\Delta_{s,L}$ by
\eqref{eq:weighted-deep-recovery}.
\item The fixed-gauge operator quotient is typed before absolute values.  Its
regular and face terms use their strong corresponding ledgers for the
majorant and the weaker topologies
\eqref{eq:regular-DQ-weak-convergence} and
\eqref{eq:face-DQ-weak-convergence} for the regular limit and fixed face
core.  The middle face band is charged by
\eqref{eq:face-middle-tail}, while the
product-current terms use \eqref{eq:product-current-type}--
\eqref{eq:product-current-DQ-convergence} and vanish only against the two
centered endpoint tests.  The response term is charged to the derived estimate
\eqref{eq:derived-strong-weak-CM2}; its difference quotient converges in the
weaker topology \eqref{eq:strong-DQ-convergence}, rather than through an
assumed CM2 norm.
\end{enumerate}
All exponents are required strictly positive; in particular the
$\kappa_T^{\rm sh}$ interval must be nonempty.

The primitive inventory is falsifiable term by term.  $\mathrm{MT}_{\rm DQ}$ must
retain the fixed-gauge commutator, double-project regular and response blocks,
keep every graph face with its exact marginal product correction, and
separate response-space mixing from difference-quotient convergence.  Its
fourth type uses the concrete strong one-coordinate TV space
$\mathfrak R_R=\mathcal M(M)$ and the contractive marginal maps
\eqref{eq:bounded-face-marginal-maps}, while fixed-term factor convergence
is taken in the weaker $\mathfrak R_{\rm BL}$ topology.  No unnamed stronger
marginal norm is imported from the face ledger.
$\mathrm{NST}_{\rm phys}$ must provide a
mass-preserving actual stopping law with
$\sigma_Q\asymp\sigma(r)$.  $\mathrm{SS}$ must contain exactly one rough source, and
$\mathrm{ACTUAL\_SLOPE\_BRIDGE}$ must identify that source pointwise with
the logarithmic derivative ratio of the two actual $\mathrm{BR}_{\rm sec}$
endpoints.  $\mathrm{PPE}_{N}$ must give actual-SRB escape at the prescribed
depth, while $\mathrm{PAIR\_SAFE\_REFERENCE}$ fixes the exact unselected
branch-mixture product law and slope and $\mathrm{PAIR\_TENERGY}$ controls
their energy at the same PPE window.
$\mathrm{CEM\_DOM}$ must dominate every discarded current by its predictable
stopped-source envelope, while $\mathrm{TENERGY}_{W}$ must put that escape
and every failure cemetery label under the one full event tilt.
$\mathrm{DUAL\_MASS\_KERNEL}$ must measurably retain every additive
occurrence even on overlapping supports, separate exact positive-component
mass from norm
envelope, absorb operator growth into that same propagated $q$, and exhaust
all owners before scalar weak assembly.  $\mathrm{MPD}$ must resolve every
regular signed source into measurable positive components with
occurrence-constant polarity.  $\mathrm{PROP\_Q\_MATCH}$ must identify every
occurrence with a record-preserving exact-$m$ pushforward and a uniformly
dominated $q$ pushforward of the physical ledger; finiteness of the latter
alone does not bound an independently enlarged $q$.
$\mathrm{POST\_CARRIER\_DOM}$ separately controls the variation created by
propagation before a $J$ row is restricted in its carrier coordinate.
$\mathrm{PHYS\_TEST\_LIFT}$ supplies one uniformly bounded lift of the two
physical observables into the heterogeneous product test space, and
$\mathrm{PHYS\_TREE\_MATCH}$ gives the exact Radon--Nikodym identity between
the physical $M_\lambda$ ledger and the complete-tree law.  The master
source kernel and all these records are projectively fixed before time,
mode, polynomial, and final test.  $\mathrm{MAIN\_CLOCK\_COVERAGE}$ supplies the
separate resolution-dependent partition and handler identity inside the
main owner.  $\mathrm{SEL\_JSL}_{p_{\rm joint}}$ must compare each
unnormalized selected source measure with one unselected probability proposal
law; the selector, reciprocal proposal, and carrier changes all belong to
the unique density $Z_A$.  Theorem~\ref{thm:main-source-energy-lift} then
derives the main marked value mass from unselected reference energy,
selection likelihood, declared base-mass bounds, and one global incidence
Cauchy step; it is not a primitive source-clock row.  $\mathrm{CLOCK\_PROB}_{W}$
controls only bad-support
unions and survivor mass, while $\mathrm{CLOCK\_SRC}_{W}$ controls the full
clock--passage matrix of positive current masses after source complexity.
$\mathrm{KRH}_{p_W}$ supplies the uniform frozen-parent event-weight bound;
a global $L^{p_W}$ moment is not substituted for it.
$\mathrm{USC}_{\rm end}^{W}$ must charge positive endpoint-scale growth to the
homogeneity indices and the initial scale mark $S_\lambda$ under the complete
tilted event weight on that same complete actual/split law.
$\mathrm{BR}_{\rm sec}$ and $\mathrm{REC}$ control the same-branch
sections and their two separate recovery tails.
$\mathrm{FZ}_{\rm geom}$ counts every component;
$\mathrm{SF}_{\rm split}$ represents every constituent law;
$\mathrm{CP}_{\rm env}$ gives connected cells;
$\mathrm{UCR}_{\rm sh}$ continues only finite shallow labels;
$\mathrm{FACE}_{\rm 2CUT}$ supplies uniform integrability between a fixed
face core and $L(s)$, while $\mathrm{FACE\_TIME}_{\rm CM2}$ separately
charges the propagated positive mass of every main/nonmain face/word owner
in product time;
$\mathrm{SHCOL}_{\rm env}$ gives all-time shallow positive-envelope
mixing; $\mathrm{TAIL}_{\rm env}^{+}$ gives the deep positive source mass;
$\mathrm{NREC}_{\rm deep}$ preserves that source under recovery weighting;
$\mathrm{QT}_{\rm JET}$ supplies only minor separation and path derivatives,
and $\mathrm{QT}_{\rm LEDGER}$ separately supplies the actual loss/tail
margin.  Their conjunction is $\mathrm{QT}_{\rm path}$.
They are primitive quantitative conditions, not consequences hidden in
compactness or a Banach norm.

\begin{proposition}[Eighty-four obstruction tests]
\label{prop:v44-obstruction-tests}
The active interfaces reject the following countermodels.
\begin{enumerate}[label=\textup{(T\arabic*)}]
\item If $\nu\ne\mu_0$ and $J=\delta_{(x,y)}$, fixed and moving
      centerings disagree.  The measure-trivializing gauge excludes it.
\item If the physical dynamics is fixed but the auxiliary gauge is changed,
      the transported-observable physical derivative can vanish while the
      fixed-gauge commutator is nonzero.  $\mathrm{MT}_{\rm DQ}$ retains the
      commutator and therefore does not identify the two targets.
\item A clean Bernoulli gate with $p<1$ has depth-$J$ mass $p^J$;
      Lemma~\ref{lem:negative-pressure-trial-obstruction} forbids an
      $O(J)$-trial exponential physical cover.
\item If $X^{[J]}$ has good sublevel bounds but $X^{[N]}$ collapses to $P$
      off a zero-volume clean core, an existential-depth PPE would pass while
      the phase energy fails.  $\mathrm{PPE}_{N}$ tests $X^{[N(r)]}$
      directly.
\item If a purported energy cell has $\sigma_Q=1$ while
      $\sigma(r)\to0$, the remainder
      $H_\beta\sigma_Q^{1+\alpha_B}$ need not be small even when the
      marginal $H_\beta$ gate passes.  $\mathrm{NST}_{\rm phys}$ enforces
      \eqref{eq:native-stopping-comparability} and the joint gate
      \eqref{eq:joint-native-good-gate}.
\item A fixed shallow word with $N>L(s)$ lies outside the finite-depth
      refinement.  Theorem~\ref{thm:shallow-two-threshold} sends it to the
      all-time collar estimate instead of the Fourier proof.
\item For a fixed shallow word and $N\le T_{\rm sh}(s)$, a Fourier mode with
      $\Lambda=e^{L(s)^2}$ would require a PPE depth larger than $L(s)$.
      The cutoff $H_N$ sends it to
      \eqref{eq:ultra-high-observable-tail}, which uses no PPE.
\item The family $k_s=s h_s$ may have a limiting weak derivative while
      $\|h_s\|_{\rm strong}\to\infty$.  $\mathrm{MT}_{\rm DQ}$
      requires the finite-DQ uniform bound rather than only a limiting
      derivative.
\item Two equal positive pieces with opposite signs have assembled TV zero
      but positive envelope mass two.  $\mathrm{TAIL}_{\rm env}^{+}$ uses
      the latter and never infers it from signed TV.
\item A deep piece of mass $\Delta$ and recovery time
      $q^{-1}\log(1/\Delta)$ has bounded unnormalized moment but divergent
      normalized moment.  $\mathrm{NREC}_{\rm deep}$ rejects it.
\item The same example for a shallow collar would make its off-diagonal
      prefactor diverge despite a uniform post-recovery rate.
      $\mathrm{SHREC}$ rejects it through
      \eqref{eq:normalized-shallow-recovery}.
\item If one symbolic cylinder has two disconnected event components,
      $\mathrm{CP}_{\rm env}$ assigns distinct component indices before
      oscillations are taken.
\item A graph current may have nonzero marginals even though its complete
      current-side projection is centered.  Applying the response-side
      zero-mass estimate to the graph and marginal corrections separately is
      illegal.  The exact joint identity
      \eqref{eq:typed-face-product-cluster}, together with the individual
      projection rule \eqref{eq:typed-block-double-projection} only for the
      regular and response rows, rejects this model.
\item Two zero-mass operators $K^A,K^B:\cB_2\to\cB_1$ may have the same
      uniform response bound while
      $\widehat K_{s,\rm sw}^{\rm DQ}$ alternates between them as $s\to0$.
      Equation~\eqref{eq:strong-DQ-convergence} rejects the resulting
      nonconvergent difference quotient.
\item A raw mode with $\Lambda=1$ may have $a_X=e^{\chi N}$ and hence
      physical frequency $\Xi=e^{\chi N}$.  The raw cutoff alone misses it;
      $\mathrm{USC}_{\rm end}^{W}$ either puts it in the scale-good cost
      $C_{\rm scl}$ or charges it to $e^{-c_{\rm scl}N}$.
\item If the derived main or exceptional Jensen value rate is smaller than
      the direct TV or fractional-integration rate, a phase minimum formed
      from probability energy alone overstates the saving.  The complete
      value rates \eqref{eq:complete-Jensen-value-rates} and the typed
      minimum \eqref{eq:complete-phase-exponent} return the correct smaller
      exponent without applying a second square root.
\item Let $\mathbb P(j)\asymp e^{-qj}$,
      $\widetilde G=j$, and $\mathscr W_j=e^{qj}/j^2$.  The unweighted
      exponential tail and $L^1$ integrability both hold, but
      $\mathbb E[\mathscr W\one_{\{\widetilde G>u\}}]\asymp u^{-1}$.
      The tilted estimate \eqref{eq:upper-scale-excess-tail}, or its explicit
      parent-kernel $\mathrm{KRH}_{p_W}$ sufficient condition, rejects this
      correlated-mark model.
\item Let an auxiliary failure have raw probability $e^{-qN}$ and let the
      full event weight grow like $e^{aN}$ on precisely those branches, with
      an available $L^{p_W}$ moment.  Its marked tail is degraded and cannot
      retain exponent $q$.  Equations
      \eqref{eq:tilted-auxiliary-tails} and
      Corollary~\ref{cor:parent-kernel-holder-energy} use only the degraded
      parentwise rate.
\item A weight may concentrate on context- or PPE-bad branches while every
      unweighted energy estimate and the unweighted Jensen step remain
      valid.  The complete law \eqref{eq:complete-energy-tree-tilt}, including
      cemetery labels, and $\mathrm{TENERGY}_{W}$ force Jensen to be taken in
      the actual positive phase law.
\item A regular DQ kernel may alternate between two
      $W^{R,1}$-bounded densities, and a face DQ current may alternate between
      two amplitudes on the same common atlas.  Their uniform majorant ledgers
      hold but neither difference quotient converges.  Equations
      \eqref{eq:regular-DQ-weak-convergence} and
      \eqref{eq:face-DQ-weak-convergence} reject both versions.
\item At $m=0$ one can have $G_0^-=0$ but
      $S_\lambda^-=e^J$ from an initial coordinate change, and hence
      unbounded $a_X$.  Equation \eqref{eq:endpoint-upper-scale-growth} records
      $\log S_\lambda^-$, and $\mathrm{USC}_{\rm end}^{W}$ charges its excess
      under the same event weight.
\item A triangular family may converge on every fixed face depth and have
      no mass beyond $L(s)$ while carrying unit oscillating mass exactly at
      depth $L(s)$.  Fixed-depth BL$^*$ convergence does not reject it;
      \eqref{eq:face-middle-tail} does.
\item A proposed middle-tail estimate of the form
      $|s|^{-1}e^{-cL_0}$ cannot be made small by first fixing $L_0$ and then
      taking $s\to0$.  The positive coarea estimate
      \eqref{eq:face-middle-per-depth} and
      $\mathrm{FACE}_{\rm 2CUT}$ explicitly forbid this factor.
\item Let a cemetery label set every unknown future factor to one, while the
      actual discarded current carries density $e^{aN}$ on a failure set of
      mass $e^{-qN}$.  Its probability tail has rate $q$, but its current
      has only rate $q-a$.  Equation \eqref{eq:cemetery-kernel-domination}
      rejects the missing envelope.
\item If a remaining propagator $U_i$ expands total variation, replacing it
      by a unit future factor underestimates the discarded source.  The
      predictable $\mathcal C_i$ and
      Lemma~\ref{lem:stopped-source-propagation} charge that expansion before
      the cemetery transition.
\item A tail $e^{-cN_{\rm PPE}(r)}$ with $r=\Xi^{-\delta}$ and
      $\Xi=e^{\eta_{\rm lo}N_{\rm time}}$ contributes only
      $e^{-\eta_{\rm lo}\delta C_{\rm buf}cN_{\rm time}}$ before Jensen; it
      cannot be entered as $e^{-cN_{\rm time}}$.  The clock registry rejects
      this substitution.
\item If the same PPE-depth event passes through the slope-energy Jensen
      step and its base occurrence mass is uniformly bounded, its saving is
      $\Xi^{-\delta C_{\rm buf}c/2}$, not
      $\Xi^{-\delta C_{\rm buf}c}$.  The class
      $\mathcal J_{\rm PPE}^{J}$ records the factor $1/2$.
\item Counting native stopping both as a named failure and again in an
      additional probability registry falsely improves a minimum.  Its
      named row and \(\mathrm{CLOCK\_PROB}_{W}\) forbid this double charge.
\item At depth $N$, let there be $e^{aN}$ disjoint events, each of mass
      $e^{-cN}$.  Every atomic rate is positive, but their union has only
      rate $c-a$ and is order one when $a=c$.  The aggregate probability
      partition bounds \eqref{eq:clock-prob-tails} reject the latter case.
\item In the preceding model, summing the raw partition and applying
      H\"older once gives $(c-a)/p_W'$, whereas tilting every atom first gives
      only $c/p_W'-a$.  Corollary~\ref{cor:full-product-holder-energy}
      fixes the global order and prevents the stronger rate from being
      borrowed by the weaker route.
\item If an unbounded $\mathcal C_i$ is already
      $\mathcal F_{\tau_i}$-measurable, then
      $\mathbb E[\mathcal C_i^{q_C}\mid\mathcal F_{\tau_i}]
      =\mathcal C_i^{q_C}$ and no uniform conditional bound exists.
      Equation \eqref{eq:complete-tree-cemetery-moment}, or the genuinely
      earlier kernel \eqref{eq:parent-kernel-cemetery-moment}, rejects this
      filtration tautology.
\item A branch entering cemetery before depth $N$ has no future strip indices,
      endpoint scale, or $\widetilde G_N^\pm$.  The survivor indicator in
      \eqref{eq:upper-scale-excess-tail} prevents the upper-scale ledger from
      evaluating a nonexistent random variable; cemetery mass is charged by
      $\mathrm{CEM\_DOM}$ and $\mathrm{CLOCK\_PROB}_{W}$.
\item A hierarchy may have exponentially decaying atomic activations but a
      partition function with nonpositive pressure gap.  Then neither
      survivor normalization nor a response majorant follows.  The positive
      gaps in \eqref{eq:clock-prob-tails},
      \eqref{eq:clock-source-matrix-tails},
      \eqref{eq:clock-source-TV-tails}, and
      \eqref{eq:exceptional-Jensen-value-rates} declare this an explicit
      no-go rather than a tunable parameter loss.
\item Let one rare support have weighted probability $e^{-cN}$ but carry
      $e^{aN}$ distinct positive remainder currents of unit envelope.  A
      first-hit probability partition sees only $e^{-cN}$, whereas
      \eqref{eq:source-incidence-TV-bound} and
      \eqref{eq:clock-source-partition} record the true
      $e^{-(c-a)N}$ source mass.
\item It is possible that every bad-support union in
      $\mathrm{CLOCK\_PROB}_{W}$ is exponentially small while the sum of
      additive currents is order one.  The separate incidence partition
      \eqref{eq:clock-source-partition} makes
      $\mathrm{CLOCK\_SRC}_{W}$, rather than survivor probability, the
      required phase input.
\item A genuine $N_{\rm time}$ source with bounded base occurrence mass and
      energy $e^{-cN}$ contributes only $e^{-cN/2}$ after the
      energy/Jensen passage.  The two time
      cells in \eqref{eq:clock-source-matrix-tails} and
      \eqref{eq:clock-source-TV-tails}, together with the final minimum
      \eqref{eq:final-CM2-exponent} reject use of the full rate $c$.
\item Let parent cells $A_N$ have mass $2^{-2N}$, put a raw bad subset of
      conditional mass $2^{-N}$ inside $A_N$, and set $W=2^N$ only there.
      A global $L^p$ bound may hold for $p<3$, while the tilted conditional
      bad probability is order one.  Definition~\ref{def:kernel-reverse-holder}
      rejects this model; Lemma~\ref{lem:measurable-refinement-moments} alone
      does not.
\item If two overlapping named events activate two different currents,
      assigning the overlap to the first event deletes the second current
      from a probability first-hit partition.  The canonical owner map in
      Definition~\ref{def:source-incidence-space} identifies physical
      occurrences, not event supports, and therefore retains both.
\item Let $\lambda\in[0,1]$, one occurrence be active for every label, and
      let its current be $\delta_\lambda$.  The naked expression
      $\sum_{\lambda\in[0,1]}\langle\delta_\lambda,1\rangle$ is undefined,
      while \eqref{eq:marked-source-assembled-pairing} is the Lebesgue
      integral and equals one.
\item A countable set $I(\lambda)$ chosen by a nonmeasurable subset of the
      label space has finite fibres but does not define a kernel.  The Borel
      graph and jointly measurable weighted counting kernel
      \eqref{eq:marked-weighted-counting-kernel} reject it before Tonelli is
      invoked.
\item Let $M=\epsilon^{-1}$ occurrences have energies
      $E_a=\epsilon/M$.  Then $\sum_aE_a=\epsilon$ but
      $\sum_a\sqrt{E_a}=1$.  The global inequality
      \eqref{eq:global-incidence-cauchy} takes one square root after the base
      and energy masses are integrated and therefore does not accept the
      occurrencewise argument.
\item If $e^{aN}$ occurrences each have raw energy $e^{-cN}$, then direct TV
      summability may hold for $c>a$, whereas the Jensen value gap is
      $(c-2a)/2$ and fails whenever $c\le2a$.  The BASE--ENERGY ledger
      \eqref{eq:clock-source-matrix-tails} declares exactly this no-go.
\item Parentwise PPE may hold uniformly while the main occurrence base mass
      grows faster than its energy saving.  Then the probability theorem is
      true but the marked value tail is false.  The selection-likelihood and
      base bounds \eqref{eq:selection-safe-likelihood-rates} and
      \eqref{eq:main-source-base-mass}, together with positivity of
      \eqref{eq:source-value-small-ball} reject this model; the main source
      is not silently reinserted into $\mathrm{CLOCK\_SRC}_{W}$.
\item There may be one source occurrence, so the occurrence-level
      $\mathrm{KRH}_p$ is trivial, while the collision--SRB slope law is
      uniform and the source-section slope law is $\delta_0$.  Probability
      PPE then has energy $O(r)$ but source energy is one.
      $\mathrm{SEL\_JSL}_{p_{\rm joint}}$ in
      \eqref{eq:joint-source-likelihood} rejects this internal-law mismatch.
\item Let $W_{\rm prob}=W_{\rm src}=[0,1]$, $H={\rm id}$,
      $\kappa={\rm Leb}$, $\nu=\delta_0$, and $S(x)=x$.  The Jacobian is one
      and the slopes match exactly, but
      $\nu\not\ll H_*\kappa$; reference energy is $O(h)$ while source energy
      is one.  The pushforward law and RN condition
      \eqref{eq:carrier-match-pushforward-law}--
      \eqref{eq:carrier-match-density-ratio}, not a bounded Jacobian alone,
      reject it.
\item If the joint likelihood is only $L^p$ and the reference energy is
      $r^\Theta$, transferring it as $r^\Theta$ misses the H\"older loss.
      Equations \eqref{eq:joint-source-energy-transfer} and
      \eqref{eq:source-value-small-ball} use
      $\Theta/p_{\rm joint}'$ before the global square root.
\item Let the exact source mass be one while its boundary/atlas complexity
      tends to infinity.  Setting the envelope mass equal to one loses the
      source norm although the probability section remains normalized.
      Equations \eqref{eq:normalized-source-shape-cost} and
      \eqref{eq:source-incidence-envelope} force this complexity into $q$.
\item Every fibre can satisfy $q_\lambda(Y_{\rm src})<\infty$ while
      $\int q_\lambda(Y_{\rm src})d\Pi(\lambda)=\infty$.  Then pointwise
      scalar integrals do not define a bounded current on the test unit ball.
      The single propagated global requirement in
      \eqref{eq:global-finite-marked-envelope} reject this model.
\item A unit occurrence may have owner main, so
      \eqref{eq:exhaustive-source-owner-partition} holds, while at one
      resolution every main clock filter is empty.  All local rate hypotheses
      are then vacuous although an order-one resonant current remains.
      Identity \eqref{eq:main-clock-coverage}, not the owner partition,
      rejects this model.
\item Let the unselected reference law satisfy
      $\mathsf R_0(D_h^0)=\delta_h\to0$, and let a main selector isolate the
      $D_h^0$ branch on which the source energy is one.  Conditioning a
      reference law on that cell would inherit a false order-$\delta_h$
      bound.  The fixed-before-selection proposal
      \eqref{eq:proposal-lift} and the unnormalized density $Z_A$ in
      \eqref{eq:joint-source-likelihood} force the selector rarity into
      $a_A^L$.
\item If $\mathsf R_0(A)=\delta$, then
      $\|d\mathsf R_0(\cdot\mid A)/d\mathsf R_0\|_p
       =\delta^{-1/p'}$.  A scale-free normalized selected likelihood would
      hide exactly this blow-up.  Equation
      \eqref{eq:selection-safe-likelihood-rates} records it, and
      \eqref{eq:source-value-small-ball} rejects a nonpositive resulting
      gap.
\item Suppose $\Sigma_A(B)>0$ but the proposal gives
      $\overline{\mathsf R}(B)=0$.  Every separate occurrence and carrier
      moment may look finite, yet no source likelihood exists.  The absolute
      continuity and support clause in \eqref{eq:joint-source-likelihood}
      makes the main lift an explicit no-go.
\item On one fixed carrier let $\nu$ be normalized Lebesgue measure and let
      $\epsilon_N(x)=\operatorname{sgn}\sin(2\pi2^Nx)$.  The positive
      section has bounded source cost, while $\epsilon_N\nu$ has
      exponentially many cuts (or a smooth approximation has
      $C^{\alpha_*}$ cost of order $2^{\alpha_*N}$).  The constant-polarity
      decomposition \eqref{eq:measurable-positive-decomposition} forbids
      applying this uncharged carrier-dependent sign.
\item A naive Jordan split of a regular signed density may produce positive
      parts vanishing at their zero sets and hence outside a uniform
      log--H\"older cone.  It cannot be sent through
      $\mathrm{SOURCE\_PROP}$.  The baseline split
      \eqref{eq:baseline-cone-split} gives two strictly positive regular
      pieces, while the final clause of $\mathrm{MPD}$ sends an unavailable
      regular split to a TV/cemetery row.
\item At depth $N$, let bare physical mass be $e^{-cN}$ and let the operator
      norm be $C_U=e^{aN}$.  Every fixed-depth operator integral is finite,
      but the propagated current has only rate $c-a$ and does not decay when
      $a\ge c$.  The envelope inequality
      \eqref{eq:source-incidence-envelope} and
      Remark~\ref{rem:operator-absorption-rates} put this growth into $q$ and
      declare the zero gap a no-go.
\item A normalized proposal may attach a mark whose source slope is constant
      on the resonant set even though the $a$-free reference slope has small
      energy.  Pointwise normalization of $\zeta$ alone then preserves the
      wrong event.  The event compatibility in
      \eqref{eq:joint-source-energy-transfer}, or the transported inclusion
      \eqref{eq:carrier-match-slope-window}, rejects it.
\item Positive decompositions may exist separately on every fibre while the
      component choice is non-Borel.  Then neither the marked kernel nor
      Tonelli assembly is defined.  The standard-Borel measurable component
      extension in Definition~\ref{lem:measurable-positive-decomposition} rejects
      this model before any signed recombination.
\item Let the complete positive physical ledger have total mass one, but
      attach to its single source occurrence an independent propagated
      envelope of mass $e^{aN}$.  Fibrewise $q$ is finite and the physical
      moment ledger is finite, yet neither controls the other and every
      claimed source rate can lose $a$.  The occurrence-kernel domination
      \eqref{eq:propagated-q-physical-match} rejects this model; finiteness on
      the two spaces is not a substitute for their typed identification.
\item Let the only face/word current have word depth zero and
      $B_{m,n}^{\rm fw}=\one_{\{m=n\}}$.  Every fixed-time
      $\mathrm{FACE}_{\rm 2CUT}$ tail is then trivial, but the CM2 double sum
      diverges.  The propagated-envelope decay
      \eqref{eq:face-time-positive-envelope}, not word-depth convergence,
      rejects it.
\item For two collision components, vary one perimeter while fixing the
      other.  A near-identity component-preserving diffeomorphism cannot
      change either component's total mass, so no $A_s$ can satisfy
      \eqref{eq:measure-trivializing-gauge}.  Condition
      \eqref{eq:constant-component-mass-vector} rejects the path before any
      dynamical estimate.
\item A random analytic base may satisfy every jet-minor lower bound while
      its unsmoothed return mass is order one or its coarea loss is
      $e^{2c_{\rm tail}L}$.  Theorem
      \ref{thm:analytic-QT-realization} supplies only
      $\mathrm{QT}_{\rm JET}$; the separate
      $\mathrm{QT}_{\rm LEDGER}$ rejects this overclaim.
\item Double projection of a graph current produces
      $-m_J^-\otimes\mu_0$ and $-\mu_0\otimes m_J^+$, which are neither a
      four-dimensional density nor a one-dimensional graph face.  The fourth
      type \eqref{eq:product-current-type}, its exact joint cluster
      \eqref{eq:typed-face-product-cluster}, and its convergence
      \eqref{eq:product-current-DQ-convergence} reject either
      misclassification.
\item Let $J_s=\delta_{2s}\to\delta_0$ in $\mathrm{BL}^*$, and choose a
      continuous bump $\Phi_s$ supported in $(s,3s)$ with
      $\Phi_s(2s)=1$ but pointwise limit zero.  Uniform sup norm and
      branchwise pointwise convergence do not imply convergence of the
      pairing.  The common off-boundary BL modulus and convergence in
      \eqref{eq:moving-face-test-BL-convergence} reject this model.
\item Take the actual endpoints $\mathsf X(u)=\mathsf Y(u)=u$ and the
      balanced mode $(1,-1)$, so the physical phase derivative vanishes,
      but attach an unrelated abstract slope satisfying an excellent
      polynomial sublevel estimate.  Conditions $\mathrm{SS}$ and
      $\mathrm{PPE}_{N}$ alone do not see the mismatch.  The pointwise
      identities \eqref{eq:actual-derivative-ratio-slope}--
      \eqref{eq:actual-balanced-phase-derivative} reject it.
\item On the circle let a frozen-opposite-endpoint carrier law have density
      proportional to $d(x,x')^k$ relative to arclength.  Every conditional
      diagonal window can have size $O(h^{k+1})$, while the two endpoint
      marginals may both be uniform and their genuine iid diagonal energy is
      only $O(h)$.  The common-kernel requirement
      \eqref{eq:pair-tenergy-factorized-route}, or the direct pair bound
      \eqref{eq:pair-tenergy-bound}, forbids integrating the
      $x'$-dependent laws as $\kappa\otimes\kappa$.
\item Let $\mathsf R_0$ be uniform with slope $S(x)=x$, all likelihood and
      base costs equal one, but declare $h_{\rm PPE}(N)=1$.  The small-$r$
      reference energy holds while the PPE source energy is one for every
      $N$.  The scale identity \eqref{eq:ppe-clock-window-scale} and the
      event inclusion in $\mathrm{SEL\_JSL}$ reject this model.
\item For occurrence $j$ take $X_j=\{+,-\}$,
      $\nu_j=(1/2,1/2)$,
      $\|(x,y)\|_{\mathfrak S_j}=|x+y|+2^j|x-y|$,
      $U_j(x,y)=2^j(x-y)$, and $m_j=q_j=2^{-j}$.  The grouped operator
      pairing is bounded (indeed it cancels for the unit test), but
      $\sum_jq_j\int|U_j^*1|\,d\nu_j=\infty$.  The occurrencewise definition
      \eqref{eq:marked-source-assembled-pairing} makes the grouped scalar
      integral legitimate, while $\mathrm{POST\_CARRIER\_DOM}$ forbids any
      subsequent carrier restriction not controlled by positive variation.
\item Let the $j$th declared local test space be $\mathbb R$ with norm
      $2^j|t|$, and let the physical unit observable generate coordinate
      one in every row.  All coordinates exist separately but the tuple is
      not a bounded product test.  The uniform lift
      \eqref{eq:bounded-physical-test-lift} rejects this model.
\item For each mode datum $\mathfrak d$ let
      $q^{\mathfrak d}=\delta_{a_{\mathfrak d}}$ on a new mark.  Every
      fibrewise kernel is finite, yet no finite simultaneous kernel governs
      all modes.  The projective identity
      \eqref{eq:source-kernel-projective-consistency} rejects this
      post-selection construction.
\item Two physical occurrences may have equal masses but different owner,
      orientation, or word depth.  An arbitrary domination kernel can swap
      them and still preserve total mass.  The support condition
      \eqref{eq:prop-q-compatibility-graph} and the exact current identity
      following \eqref{eq:prop-q-exact-m-match} reject the swap.
\item Let $\Pi(j)=2^{-j}$, let the physical label mass be
      $M_j\asymp j^{-2}$, and attach a physical factor $\mathcal C_j=j$.
      Every $\Pi$-moment of $\mathcal C$ is finite, whereas
      $\sum_jM_j\mathcal C_j=\infty$.  The exact physical/tree
      Radon--Nikodym identity \eqref{eq:physical-tree-RN-identity}, with its
      density charged in the full product envelope, rejects the untyped
      passage between these measures.
\item A jointly linear map on
      $C^r(M)\times C^r(M)$ satisfying a product norm bound vanishes: its
      values on $(f,0)$ and $(0,g)$ are zero, hence so is its value on
      $(f,g)$.  The bounded \emph{bilinear} lift
      \eqref{eq:bounded-physical-test-lift} and the commuting identity
      \eqref{eq:physical-test-lift-commutation} reject this formally empty
      interface.
\item Let the physical Radon--Nikodym factor tend to zero while the bare
      Mark--$Z$ factor grows like its reciprocal.  Controlling only their
      product can leave every physical moment finite while the bare tree
      weight has no $L^{1+\epsilon}$ margin.  The factor
      $(1+\rho_{\rm PT})$ in
      \eqref{eq:full-mark-cemetery-domination}--
      \eqref{eq:proved-full-product-Lp} controls both laws and rejects the
      model.
\item An independently declared physical occurrence measure can be
      multiplied by any $\varepsilon>0$ and still satisfy a pair of
      homogeneous domination displays.  The defining formula
      \eqref{eq:physical-occurrence-m-measure}, its forgetful marginal, and
      the signed recombination identity
      \eqref{eq:physical-occurrence-signed-recombination} fix its scale
      before the Markov occurrence kernel is applied.
\item A conventional even mollifier has
      $1-\widehat\rho(1/K)\asymp K^{-2}$ and therefore cannot furnish a
      $K^{-3}W^{3,1}$ remainder, while a generic mollifier is not
      band-limited at all.  The explicitly flat finite-frequency multiplier
      \eqref{eq:flat-bandlimited-regular-smoother} and its Taylor estimate
      \eqref{eq:flat-smoother-approximation} reject both substitutions.
\item Take one TV occurrence with $K=m=q=1$, $U=100\,\mathrm{id}$,
      $C_U=100$, $R=0$, $\gamma=1/2$, and
      $\mathscr A^{\rm opp}=10^{-2}$.  An envelope using the unfloored
      opposite amplitude would falsely bound a current of size $100$ by a
      unit $q$-mass.  The global restrictions $0<\gamma<1$, $R\in\mathbb N_0$
      and the floor \eqref{eq:opposite-amplitude-floor} force the full
      operator cost into \eqref{eq:source-incidence-envelope}.
\item Let two complete-tree points have the same physical label but different
      charged weights.  Choosing either value as a putative descended label
      weight can undercount the physical sum while the Radon--Nikodym identity
      for label functions still holds.  The fibre kernel
      \eqref{eq:physical-weight-fibre-disintegration} and conditional
      definition \eqref{eq:physical-weight-conditional-definition} replace
      this invalid descent by the exact conditional average.
\item On one tree fibre let component coefficients be
      $D^m(j)=2^{-2j}$ and append costs
      $\mathscr C_{\rm occ}(j)=2^{2j}/j$.  The exact signed decomposition can
      have finite coefficient mass while its propagated occurrence envelope
      diverges.  The disintegration bound
      \eqref{eq:physical-occurrence-envelope-disintegration}, required again
      after every handler restriction, rejects this post-ledger component
      cost.
\item A map may be bounded on $W^{1,1}$ but not from $L^1$ to $L^1$: add to a
      valid extension a fixed exterior bump multiplied by an interior trace
      functional.  Narrow input bumps then have vanishing $L^1$ mass but a
      nonvanishing zero Fourier coefficient after extension.  The simultaneous
      $L^1/W^{r,1}$ bounds imposed before
      \eqref{eq:flat-bandlimited-regular-smoother} rule out this false low-mode
      estimate.
\item A positive graph ledger may have uniformly bounded total mass while an
      independently declared stronger marginal norm diverges on its projected
      marginals.  The concrete choice
      $\mathfrak R_R=\mathcal M(M)$, the contractive identities
      \eqref{eq:bounded-face-marginal-maps}, and the separate weak
      $\mathfrak R_{\rm BL}$ convergence close the product-current type
      without importing such an uncharged strong norm or demanding TV
      convergence of a moving collar.
\item For a checkerboard or fat-Cantor carrier event $E$, the restricted
      measure $\one_E\nu$ need not belong to the regular source cone even
      when $\nu$ does.  The scalar post-propagation measure
      \eqref{eq:post-carrier-restriction-identity}--
      \eqref{eq:post-carrier-full-pairing}, rather than an application of
      $U$ to that ill-typed source, defines the $J$ restriction.
\item Let a complete law be concentrated on a cemetery row and let a
      positive weight concentrate on events of probability $\delta_N\to0$.
      Without a uniform normalization floor, its $L^p$ norm divided by its
      mean grows like $\delta_N^{-1/p'}$.  The explicit normalization
      \eqref{eq:full-product-weight-normalization} together with the full
      product envelope rejects this hidden rare conditioning; no pointwise
      cemetery floor is asserted in its proof.
\item Separate moments of $\rho_{\rm PT}$ and of a cemetery weight can be
      finite while their correlated product has infinite mean.  The main
      theorem now imports the complete package of
      Lemma~\ref{lem:measurable-refinement-moments}, in particular
      \eqref{eq:full-cemetery-product-envelope}; listing only the component
      moments is not an admissible substitute.
\item Inside normalized parent kernels let $W_\eta$ equal one on a set of
      mass $\delta_\eta\downarrow0$ and $\delta_\eta^2$ elsewhere.  Uniform
      upper moments do not prevent the normalized likelihood
      $W_\eta/\int W_\eta dK_\eta$ from having divergent $L^p$ norm.  The
      parentwise lower bound
      \eqref{eq:kernelwise-weight-normalization}, which is independent of the
      global tree normalization, rejects this rare-parent countermodel.
\end{enumerate}
\end{proposition}

\begin{proof}
Items 1--3, 8--11, and 13--84 are direct finite-measure, frequency, or norm
calculations.  Item 4 is the missing endpoint-depth countermodel, item 5 is
the native-scale countermodel, item 6 is the word/time-depth countermodel,
item 7 is the ultra-high-frequency countermodel, and item 12 is exactly the
connected-refinement clause.
\end{proof}

\begingroup\raggedright
\begin{thebibliography}{99}
\bibitem{DLcones}
M.~Demers and C.~Liverani,
\emph{Projective cones for sequential dispersing billiards},
Commun. Math. Phys. 401 (2023), arXiv:2104.06947.

\bibitem{DLreview}
M.~Demers and C.~Liverani,
\emph{Recent progress in the application of transfer operators to dispersing
billiards}, arXiv:2606.10155 (2026).

\bibitem{FroylandPhalempin}
G.~Froyland and M.~Phalempin,
\emph{Optimal linear response for Anosov diffeomorphisms},
arXiv:2504.16532 (2025).

\bibitem{BahsounGalatolo}
W.~Bahsoun and S.~Galatolo,
\emph{Linear response due to singularities}, arXiv:2301.02301 (2023).

\bibitem{DLsequential}
M.~Demers and C.~Liverani,
\emph{Central limit theorem for sequential dynamical systems},
arXiv:2502.07765v2 (2025).

\bibitem{SYZ}
M.~Stenlund, L.-S.~Young and H.-K.~Zhang,
\emph{Dispersing billiards with moving scatterers},
arXiv:1210.0011.

\bibitem{LS}
J.~Lepp\"anen and M.~Stenlund,
\emph{A note on the finite-dimensional distributions of dispersing billiard processes},
arXiv:1702.01461.

\bibitem{FlemingFCB}
N.~Fleming-V\'azquez,
\emph{Functional correlation bounds and optimal iterated moment bounds for
slowly-mixing nonuniformly hyperbolic maps}, arXiv:2106.06486v2.

\bibitem{BFLS2026}
I.~Baldom\'a, A.~Florio, M.~Leguil and T.~M.-Seara,
\emph{Chaoticity of generic analytic convex billiards},
arXiv:2605.19897 (2026).

\bibitem{Wormell}
C.~Wormell,
\emph{Conditional mixing in deterministic chaos},
arXiv:2206.09291; and
\emph{On convergence of linear response formulae in piecewise hyperbolic maps},
arXiv:2206.09292.

\bibitem{BB}
S.~Baker and A.~Banaji,
\emph{Polynomial Fourier decay for fractal measures and their pushforwards},
Math. Ann. 392 (2025), 209--261, arXiv:2401.01241.

\bibitem{BKS}
S.~Baker, O.~Khalil and T.~Sahlsten,
\emph{Fourier decay from $L^2$-flattening},
arXiv:2407.16699.

\bibitem{CD}
V.~Climenhaga and J.~Day,
\emph{Every finite horizon Sinai billiard map has a unique measure of maximal entropy},
arXiv:2604.25881.

\bibitem{BNST}
P.~B\'alint, P.~N\'andori, D.~Sz\'asz and I.~P.~T\'oth,
\emph{Equidistribution for standard pairs in planar dispersing billiard flows},
Ann. Henri Poincar\'e 19 (2018), arXiv:1707.03068.

\bibitem{HWWRemez}
S.~Huang, G.~Wang and M.~Wang,
\emph{Observability inequality, log-type Hausdorff content and heat equations},
arXiv:2411.11573v2 (2024).

\bibitem{BarbieriNash}
S.~Barbieri,
\emph{Bernstein--Remez inequality for Nash functions: a complex analytic
approach}, arXiv:2207.13922 (2022).

\bibitem{BCC}
O.~Butterley, G.~Canestrari and R.~Castorrini,
\emph{Discontinuities cause essential spectrum on surfaces},
Ann. Henri Poincar\'e (2024), arXiv:2306.00484.

\bibitem{CanestrariPerturb}
G.~Canestrari,
\emph{On linear response for discontinuous perturbations of smooth
endomorphisms}, Advances in Mathematics 497 (2026), 111008,
arXiv:2411.16628v3.

\bibitem{CanestrariHoles}
G.~Canestrari,
\emph{Linear response for Sinai billiards with small holes},
arXiv:2604.19671v2 (2026).

\bibitem{AlvesBahsoun}
J.~F.~Alves and W.~Bahsoun,
\emph{Linear response for skew-product maps with contracting fibres},
arXiv:2602.02317v2 (2026).

\bibitem{VakarOng}
M.~V\'ak\'ar and L.~M.~Ong,
\emph{On $s$-finite measures and kernels},
Mathematical Foundations of Programming Semantics (2018),
arXiv:1810.01837.

\bibitem{NiKernel}
A.~Ni,
\emph{Ergodic and foliated kernel-differentiation method for linear
responses of random systems}, arXiv:2410.10138 (2024).

\bibitem{GalatoloLucarini}
S.~Galatolo and V.~Lucarini,
\emph{A mathematical framework for linear response theory for nonautonomous
systems}, arXiv:2603.19509v3 (2026).

\bibitem{Leclerc}
G.~Leclerc,
\emph{Fourier decay of equilibrium states for bunched attractors},
arXiv:2301.10623.

\bibitem{Leclerc2026}
G.~Leclerc,
\emph{Fourier decay of equilibrium states and the Fibonacci Hamiltonian},
arXiv:2507.23731v2 (2026).

\bibitem{LPS}
G.~Leclerc, S.~Paukkonen and T.~Sahlsten,
\emph{Fourier decay in parabolic $C^{1+\alpha}$ systems with overlaps},
arXiv:2505.15468v3 (2025).

\bibitem{ARHW}
A.~Algom, F.~Rodriguez Hertz and Z.~Wang,
\emph{Polynomial Fourier decay and a cocycle version of Dolgopyat's method
for self conformal measures}, arXiv:2306.01275v2 (2024).

\bibitem{LOP}
Y.~Lima, D.~Obata and M.~Poletti,
\emph{Measures of maximal entropy for non-uniformly hyperbolic maps},
arXiv:2405.04676v2 (2025).

\bibitem{Chernov1999}
N.~Chernov,
\emph{Decay of correlations and dispersing billiards},
J. Stat. Phys. 94 (1999), 513--556.

\bibitem{Teixeira}
P.~Teixeira,
\emph{On the conservative pasting lemma},
Ergodic Theory Dynam. Systems 40 (2020), 1402--1440,
arXiv:1611.01694.

\bibitem{Ruehr}
R.~R\"uhr,
\emph{Pressure inequalities for Gibbs measures of countable Markov shifts},
arXiv:2012.13226v3.

\bibitem{Bertozzi}
N.~Bertozzi,
\emph{Exponential mixing for hyperbolic flows on non-compact spaces},
arXiv:2603.11852v2 (2026).

\bibitem{ClimenhagaLPS}
V.~Climenhaga,
\emph{Gibbs measures have local product structure},
arXiv:2310.17495v2 (2025).

\bibitem{GouezelRenewal}
S.~Gou\"ezel,
\emph{Berry--Esseen theorem and local limit theorem for non uniformly
expanding maps}, Ann. Inst. H. Poincar\'e Probab. Statist. 41 (2005),
997--1024, arXiv:math/0310391.

\bibitem{Galli}
D.~Galli,
\emph{A cohomological approach to Ruelle--Pollicott resonances and speed of
mixing of Anosov diffeomorphisms}, arXiv:2405.17045.
\end{thebibliography}
\endgroup

\end{document}
